% Generated mechanically from the frozen derived.tex target.
% Do not edit this generated file; edit the bound locale source instead.

% OK, start here.
%

\setcounter{chapter}{12}
\chapter{导出范畴}



\phantomsection
\label{derived-section-phantom}


\section{引言}
\label{derived-section-introduction}

\noindent
我们首先讨论三角范畴以及三角范畴中的局部化。随后证明：加性范畴中的复形所成的
同伦范畴是三角范畴。在此基础上，我们把阿贝尔范畴的导出范畴定义为其同伦范畴
关于拟同构的局部化。Verdier 的学位论文 \cite{Verdier} 是一份很好的参考文献。



\section{三角范畴}
\label{derived-section-triangulated-categories}

\noindent
三角范畴是描述阿贝尔范畴的导出范畴所固有结构的一种便利工具。可参见
\cite{Verdier}、\cite{KS} 和 \cite{Neeman}。




\section{三角范畴的定义}
\label{derived-section-triangulated-definitions}

\noindent
本节汇集有关三角范畴和预三角范畴的大部分定义。

\begin{definition}
\label{derived-definition-triangle}
设 $\mathcal{D}$ 为加性范畴。设
$[1] : \mathcal{D} \to \mathcal{D}$，$E \mapsto E[1]$
为 $\mathcal{D}$ 的一个加性自等价函子。
\begin{enumerate}
\item 一个{\it 三角}是六元组
$(X, Y, Z, f, g, h)$，其中 $X, Y, Z \in \Ob(\mathcal{D})$，且
$f : X \to Y$、$g : Y \to Z$ 和 $h : Z \to X[1]$ 是
$\mathcal{D}$ 中的态射。
\item 三角之间的一个{\it 态射}
$(X, Y, Z, f, g, h) \to (X', Y', Z', f', g', h')$
由 $\mathcal{D}$ 中的态射 $a : X \to X'$、$b : Y \to Y'$ 和
$c : Z \to Z'$ 给出，并满足
$b \circ f = f' \circ a$、$c  \circ g = g' \circ b$ 以及
$a[1] \circ h = h' \circ c$。
\end{enumerate}
\end{definition}

\noindent
三角的态射可由下列交换图表示：
$$
\xymatrix{
X \ar[r] \ar[d]^a &
Y \ar[r] \ar[d]^b &
Z \ar[r] \ar[d]^c &
X[1] \ar[d]^{a[1]} \\
X' \ar[r] &
Y' \ar[r] &
Z' \ar[r] &
X'[1]
}
$$
在定义 \ref{derived-definition-triangle} 的情形中，记
$[0] = \text{id}$；当 $n > 0$ 时，以 $[n]$ 表示 $[1]$ 的 $n$ 次复合；
选择 $[1]$ 的一个拟逆 $[-1]$；并令 $[-n]$ 等于 $[-1]$ 的 $n$ 次复合。
于是 $\{[n]\}_{n \in \mathbf{Z}}$ 是一族以 $n \in \mathbf{Z}$ 为指标的
$\mathcal{D}$ 的加性自等价，并且给定函子同构
$[n] \circ [m] \cong [n + m]$。

\medskip\noindent
下面是 Verdier 学位论文中给出的三角范畴定义。

\begin{definition}
\label{derived-definition-triangulated-category}
一个{\it 三角范畴}由三元组
$(\mathcal{D}, \{[n]\}_{n\in \mathbf{Z}}, \mathcal{T})$
构成，其中
\begin{enumerate}
\item $\mathcal{D}$ 是加性范畴；
\item $[1] : \mathcal{D} \to \mathcal{D}$，$E \mapsto E[1]$
是加性自等价，而 $n \in \mathbf{Z}$ 时的 $[n]$ 如上所述；
\item $\mathcal{T}$ 是一组三角（定义 \ref{derived-definition-triangle}），
称为{\it 特异三角}。
\end{enumerate}
它们须满足下列条件：
\begin{enumerate}
\item[TR1] 与特异三角同构的任何三角都是特异三角。形如
$(X, X, 0, \text{id}, 0, 0)$ 的任何三角都是特异三角。
对 $\mathcal{D}$ 的任一态射 $f : X \to Y$，存在形如
$(X, Y, Z, f, g, h)$ 的特异三角。
\item[TR2] 三角 $(X, Y, Z, f, g, h)$ 是特异三角，当且仅当三角
$(Y, Z, X[1], g, h, -f[1])$ 是特异三角。
\item[TR3] 给定实线图
$$
\xymatrix{
X \ar[r]^f \ar[d]^a &
Y \ar[r]^g \ar[d]^b &
Z \ar[r]^h \ar@{-->}[d] &
X[1] \ar[d]^{a[1]} \\
X' \ar[r]^{f'} &
Y' \ar[r]^{g'} &
Z' \ar[r]^{h'} &
X'[1]
}
$$
其两行均为特异三角，并且满足 $b \circ f = f' \circ a$，则存在态射
$c : Z \to Z'$，使得 $(a, b, c)$ 是三角的态射。
\item[TR4] 给定 $\mathcal{D}$ 的对象 $X$、$Y$、$Z$，态射
$f : X \to Y$、$g : Y \to Z$，以及特异三角
$(X, Y, Q_1, f, p_1, d_1)$、
$(X, Z, Q_2, g \circ f, p_2, d_2)$ 和
$(Y, Z, Q_3, g, p_3, d_3)$，则存在态射
$a : Q_1 \to Q_2$ 和 $b : Q_2 \to Q_3$，使得
\begin{enumerate}
\item $(Q_1, Q_2, Q_3, a, b, p_1[1] \circ d_3)$ 是特异三角；
\item 三元组 $(\text{id}_X, g, a)$ 是三角之间的态射
$(X, Y, Q_1, f, p_1, d_1) \to (X, Z, Q_2, g \circ f, p_2, d_2)$；并且
\item 三元组 $(f, \text{id}_Z, b)$ 是三角之间的态射
$(X, Z, Q_2, g \circ f, p_2, d_2) \to (Y, Z, Q_3, g, p_3, d_3)$。
\end{enumerate}
\end{enumerate}
若 TR1、TR2 和 TR3 成立，则称
$(\mathcal{D}, [\ ], \mathcal{T})$ 为{\it 预三角范畴}。\footnote{这里用
$[\ ]$ 作为族 $\{[n]\}_{n\in \mathbf{Z}}$ 的缩写。}
\end{definition}

\noindent
TR4 的含义如下：若把 $Q_1$ 看作 $Y/X$，把 $Q_2$ 看作 $Z/X$，
把 $Q_3$ 看作 $Z/Y$，则 TR4(a) 表达同构
$(Z/X)/(Y/X) \cong Z/Y$，而 TR4(b) 与 TR4(c) 表明我们可以用三角的态射
来比较三角 $X \to Y \to Q_1 \to X[1]$ 等。关于这一思想更精确的重述，
参见引理 \ref{derived-lemma-two-split-injections} 的证明。

\medskip\noindent
TR2 中的符号意味着：若 $(X, Y, Z, f, g, h)$ 是特异三角，则长序列
\begin{equation}
\label{derived-equation-rotate}
\ldots \to
Z[-1] \xrightarrow{-h[-1]}
X \xrightarrow{f}
Y \xrightarrow{g}
Z \xrightarrow{h}
X[1] \xrightarrow{-f[1]}
Y[1] \xrightarrow{-g[1]}
Z[1] \to \ldots
\end{equation}
中每个四项序列都给出一个特异三角。

\medskip\noindent
照例，我们略微滥用记号，只说一个（预）三角范畴 $\mathcal{D}$，而不显式引入
附加数据的记号。预三角范畴这一概念可用于寻找与 TR4 等价的命题。

\medskip\noindent
下面定义三角函子。

\begin{definition}
\label{derived-definition-exact-functor-triangulated-categories}
设 $\mathcal{D}$、$\mathcal{D}'$ 为预三角范畴。从 $\mathcal{D}$ 到
$\mathcal{D}'$ 的一个{\it 正合函子}，或称{\it 三角函子}，是函子
$F : \mathcal{D} \to \mathcal{D}'$ 连同给定的函子同构
$\xi_X : F(X[1]) \to F(X)[1]$，并且对 $\mathcal{D}$ 的每个特异三角
$(X, Y, Z, f, g, h)$，三角
$(F(X), F(Y), F(Z), F(f), F(g), \xi_X \circ F(h))$
是 $\mathcal{D}'$ 的特异三角。
\end{definition}

\noindent
正合函子是加性函子，见引理 \ref{derived-lemma-exact-functor-additive}。
所谓两个三角范畴等价，是指它们在三角范畴的 $2$-范畴中等价。
该 $2$-范畴中的一个 $2$-态射 $a : (F, \xi) \to (F', \xi')$，
就是与 $\xi$ 和 $\xi'$ 相容的函子变换 $a : F \to F'$，亦即下图交换：
$$
\xymatrix{
F \circ [1] \ar[r]_\xi \ar[d]_{a \star 1} & [1] \circ F \ar[d]^{1 \star a} \\
F' \circ [1] \ar[r]^{\xi'} & [1] \circ F'
}
$$

\begin{definition}
\label{derived-definition-triangulated-subcategory}
设 $(\mathcal{D}, [\ ], \mathcal{T})$ 为预三角范畴。一个
{\it 预三角子范畴}\footnote{这一定义可能并非标准定义。若
$\mathcal{D}'$ 是满子范畴，则 $\mathcal{T}'$ 是 $\mathcal{D}'$ 中三角之集
与 $\mathcal{T}$ 的交，见引理 \ref{derived-lemma-triangulated-subcategory}。
在这种情形下，我们在记号中省略 $\mathcal{T}'$。}
是一个偶 $(\mathcal{D}', \mathcal{T}')$，满足
\begin{enumerate}
\item $\mathcal{D}'$ 是 $\mathcal{D}$ 的加性子范畴，它在 $[1]$ 下保持，
且 $[1] : \mathcal{D}' \to \mathcal{D}'$ 是自等价；
\item $\mathcal{T}' \subset \mathcal{T}$，并且对每个
$(X, Y, Z, f, g, h) \in \mathcal{T}'$，有
$X, Y, Z \in \Ob(\mathcal{D}')$，且
$f, g, h \in \text{Arrows}(\mathcal{D}')$；
\item $(\mathcal{D}', [\ ], \mathcal{T}')$ 是预三角范畴。
\end{enumerate}
若 $\mathcal{D}$ 是三角范畴，则当
$(\mathcal{D}', \mathcal{T}')$ 是预三角子范畴且
$(\mathcal{D}', [\ ], \mathcal{T}')$ 是三角范畴时，称其为
{\it 三角子范畴}。
\end{definition}

\noindent
在这种情形下，嵌入函子 $\mathcal{D}' \to \mathcal{D}$ 是正合函子，
其中 $\xi_X : X[1] \to X[1]$ 取为 $X[1]$ 上的恒等态射。

\medskip\noindent
我们将在引理 \ref{derived-lemma-composition-zero} 中看到：在预三角范畴的特异三角
$(X, Y, Z, f, g, h)$ 中，复合 $g \circ f : X \to Z$ 为零。
因此序列 (\ref{derived-equation-rotate}) 是复形。同调函子就是把这一复形变为长正合序列的函子。

\begin{definition}
\label{derived-definition-homological}
设 $\mathcal{D}$ 为预三角范畴，$\mathcal{A}$ 为阿贝尔范畴。
加性函子 $H : \mathcal{D} \to \mathcal{A}$ 称为{\it 同调的}，
如果对每个特异三角 $(X, Y, Z, f, g, h)$，序列
$$
H(X) \to H(Y) \to H(Z)
$$
在阿贝尔范畴 $\mathcal{A}$ 中正合。加性函子
$H : \mathcal{D}^{opp} \to \mathcal{A}$ 称为{\it 上同调的}，
如果相应的函子 $\mathcal{D} \to \mathcal{A}^{opp}$ 是同调的。
\end{definition}

\noindent
若 $H : \mathcal{D} \to \mathcal{A}$ 是同调函子，我们常记
$H^n(X) = H(X[n])$，从而 $H(X) = H^0(X)$。上面对 TR2 的讨论表明，
特异三角 $(X, Y, Z, f, g, h)$ 决定长正合序列
\begin{equation}
\label{derived-equation-long-exact-cohomology-sequence}
\xymatrix@C=3pc{
H^{-1}(Z) \ar[r]^{H(h[-1])} &
H^0(X) \ar[r]^{H(f)} &
H^0(Y) \ar[r]^{H(g)} &
H^0(Z) \ar[r]^{H(h)} &
H^1(X)
}
\end{equation}
它称为与该特异三角及同调函子相伴的{\it 长正合序列}。
如记号所示，我们不在长正合序列的态射上使用任何符号。其副作用是：
与一个特异三角的旋转（TR2）相伴的长正合序列中的映射，与上述序列中的映射
会相差某些符号。

\begin{definition}
\label{derived-definition-delta-functor}
设 $\mathcal{A}$ 为阿贝尔范畴，$\mathcal{D}$ 为三角范畴。
一个{\it 从 $\mathcal{A}$ 到 $\mathcal{D}$ 的 $\delta$-函子}由函子
$G : \mathcal{A} \to \mathcal{D}$ 以及如下规则给出：对每个短正合序列
$$
0 \to A \xrightarrow{a} B \xrightarrow{b} C \to 0
$$
指定态射 $\delta = \delta_{A \to B \to C} : G(C) \to G(A)[1]$，
使得
\begin{enumerate}
\item 对上述每个短正合序列，三角
$(G(A), G(B), G(C), G(a), G(b), \delta_{A \to B \to C})$
都是 $\mathcal{D}$ 的特异三角；并且
\item 对短正合序列的每个态射
$(A \to B \to C) \to (A' \to B' \to C')$，下图交换：
$$
\xymatrix{
G(C) \ar[d] \ar[rr]_{\delta_{A \to B \to C}} & &
G(A)[1] \ar[d] \\
G(C') \ar[rr]^{\delta_{A' \to B' \to C'}} & &
G(A')[1]
}
$$
\end{enumerate}
在这种情形下，称
$(G(A), G(B), G(C), G(a), G(b), \delta_{A \to B \to C})$
为给定 $\delta$-函子下该短正合序列的{\it 像}。
\end{definition}

\noindent
注意，$\delta$-函子带有附加结构。严格说来，称某个给定函子
$\mathcal{A} \to \mathcal{D}$ 为 $\delta$-函子并无意义，不过我们仍会经常这样说。










\section{三角范畴的初步结果}
\label{derived-section-elementary-results}

\noindent
本节的大多数结果都是对预三角范畴证明的，因而自然也适用于任何三角范畴。

\begin{lemma}
\label{derived-lemma-composition-zero}
设 $\mathcal{D}$ 为预三角范畴，$(X, Y, Z, f, g, h)$ 为特异三角。则
$g \circ f = 0$、$h \circ g = 0$ 且 $f[1] \circ h = 0$。
\end{lemma}

\begin{proof}
由 TR1 可知 $(X, X, 0, 1, 0, 0)$ 是特异三角。对下图应用 TR3：
$$
\xymatrix{
X \ar[r] \ar[d]^1 &
X \ar[r] \ar[d]^f &
0 \ar[r] \ar@{-->}[d] &
X[1] \ar[d]^{1[1]} \\
X \ar[r]^f &
Y \ar[r]^g &
Z \ar[r]^h &
X[1]
}
$$
虚线箭头当然是零映射。因此图的交换性给出 $g \circ f = 0$。
其余情形只需旋转三角，即应用 TR2。
\end{proof}

\begin{lemma}
\label{derived-lemma-representable-homological}
设 $\mathcal{D}$ 为预三角范畴。对 $\mathcal{D}$ 的任一对象 $W$，
函子 $\Hom_\mathcal{D}(W, -)$ 是同调函子，而函子
$\Hom_\mathcal{D}(-, W)$ 是上同调函子。
\end{lemma}

\begin{proof}
考虑特异三角 $(X, Y, Z, f, g, h)$。我们已经看到
$g \circ f = 0$，见引理 \ref{derived-lemma-composition-zero}。
设 $a : W \to Y$ 为满足 $g \circ a = 0$ 的态射。于是得到交换图
$$
\xymatrix{
W \ar[r]_1 \ar@{..>}[d]^b &
W \ar[r] \ar[d]^a &
0 \ar[r] \ar[d]^0 &
W[1] \ar@{..>}[d]^{b[1]} \\
X \ar[r] & Y \ar[r] & Z \ar[r] & X[1]
}
$$
两行都是特异三角（上行使用 TR1）。因此可补出虚线箭头 $b$
（先用 TR2 旋转，再应用 TR3，最后旋转回来）。这就证明了引理。
\end{proof}

\begin{lemma}
\label{derived-lemma-third-isomorphism-triangle}
设 $\mathcal{D}$ 为预三角范畴。设
$$
(a, b, c) : (X, Y, Z, f, g, h) \to (X', Y', Z', f', g', h')
$$
为特异三角的态射。若 $a, b, c$ 中有两个是同构，则第三个也是同构。
\end{lemma}

\begin{proof}
假设 $a$ 和 $c$ 是同构。对 $\mathcal{D}$ 的任一对象 $W$，记
$H_W( - ) = \Hom_\mathcal{D}(W, -)$。于是得到阿贝尔群的交换图
$$
\xymatrix{
H_W(Z[-1]) \ar[r] \ar[d] &
H_W(X) \ar[r] \ar[d] &
H_W(Y) \ar[r] \ar[d] &
H_W(Z) \ar[r] \ar[d] &
H_W(X[1]) \ar[d] \\
H_W(Z'[-1]) \ar[r] &
H_W(X') \ar[r] &
H_W(Y') \ar[r] &
H_W(Z') \ar[r] &
H_W(X'[1])
}
$$
由假设，右边两个和左边两个竖直箭头都是双射。由引理
\ref{derived-lemma-representable-homological}，$H_W$ 是同调函子；再由五引理
（同调，引理 \ref{homology-lemma-five-lemma}），中间的竖直箭头是同构。
于是由 Yoneda 引理（见范畴，引理 \ref{categories-lemma-yoneda}），
$b$ 是同构。旋转三角（使用 TR2）即可得到其他情形。
\end{proof}

\begin{remark}
\label{derived-remark-special-triangles}
设 $\mathcal{D}$ 为带有定义 \ref{derived-definition-triangle} 中平移函子 $[n]$
的加性范畴。若对 $\mathcal{D}$ 的每个对象 $W$，阿贝尔群长序列
$$
\ldots \to
\Hom_\mathcal{D}(W, X) \to
\Hom_\mathcal{D}(W, Y) \to
\Hom_\mathcal{D}(W, Z) \to
\Hom_\mathcal{D}(W, X[1]) \to \ldots
$$
正合，则称三角 $(X, Y, Z, f, g, h)$ 为{\it 特殊的}\footnote{这不是标准术语。}。
引理 \ref{derived-lemma-third-isomorphism-triangle} 的证明表明：若
$$
(a, b, c) : (X, Y, Z, f, g, h) \to (X', Y', Z', f', g', h')
$$
是特殊三角的态射，且 $a, b, c$ 中有两个是同构，则第三个也是同构。
对{\it 余特殊}三角也有对偶命题；所谓余特殊三角，是指应用函子
$\Hom_\mathcal{D}(-, W)$ 后变成长正合序列的三角。因此特异三角既特殊又余特殊，
但一般而言，（余）特殊三角远多于特异三角。
\end{remark}

\begin{lemma}
\label{derived-lemma-third-map-square-zero}
设 $\mathcal{D}$ 为预三角范畴。设
$$
(0, b, 0), (0, b', 0) : (X, Y, Z, f, g, h) \to (X, Y, Z, f, g, h)
$$
为某特异三角的自同态。则 $bb' = 0$。
\end{lemma}

\begin{proof}
考虑图
$$
\xymatrix{
X \ar[r] \ar[d]^0 &
Y \ar[r] \ar[d]^{b, b'} \ar@{..>}[ld]^\alpha &
Z \ar[r] \ar[d]^0 \ar@{..>}[ld]^\beta &
X[1] \ar[d]^0 \\
X \ar[r] & Y \ar[r] & Z \ar[r] & X[1]
}
$$
应用引理 \ref{derived-lemma-representable-homological}，可找到虚线箭头
$\alpha$ 和 $\beta$，使得 $b' = f \circ \alpha$ 且
$b = \beta \circ g$。于是
$bb' = \beta \circ g \circ f \circ \alpha = 0$，因为由引理
\ref{derived-lemma-composition-zero} 有 $g \circ f = 0$。
\end{proof}

\begin{lemma}
\label{derived-lemma-third-map-idempotent}
设 $\mathcal{D}$ 为预三角范畴，$(X, Y, Z, f, g, h)$ 为特异三角。
若
$$
\xymatrix{
Z \ar[r]_h \ar[d]_c & X[1] \ar[d]^{a[1]} \\
Z \ar[r]^h & X[1]
}
$$
交换，且 $a^2 = a$、$c^2 = c$，则存在满足 $b^2 = b$ 的态射
$b : Y \to Y$，使得 $(a, b, c)$ 是三角 $(X, Y, Z, f, g, h)$ 的自同态。
\end{lemma}

\begin{proof}
由 TR3，存在态射 $b'$，使得 $(a, b', c)$ 是
$(X, Y, Z, f, g, h)$ 的自同态。于是 $(0, (b')^2 - b', 0)$ 也是自同态。
由引理 \ref{derived-lemma-third-map-square-zero}，$(b')^2 - b'$ 的平方为零。
令 $b = b' - (2b' - 1)((b')^2 - b') = 3(b')^2 - 2(b')^3$。
直接计算表明 $(a, b, c)$ 是自同态，且
$b^2 - b = (4(b')^2 - 4b' - 3)((b')^2 - b')^2 = 0$。
\end{proof}

\begin{lemma}
\label{derived-lemma-cone-triangle-unique-isomorphism}
设 $\mathcal{D}$ 为预三角范畴，$f : X \to Y$ 为 $\mathcal{D}$ 的态射。
存在特异三角 $(X, Y, Z, f, g, h)$，而且它在三角的（不唯一）同构意义下唯一。
更精确地，给定另一个这样的特异三角 $(X, Y, Z', f, g', h')$，存在同构
$$
(1, 1, c) : (X, Y, Z, f, g, h) \longrightarrow (X, Y, Z', f, g', h')
$$
\end{lemma}

\begin{proof}
存在性由 TR1 得到。在同构意义下的唯一性由 TR3 和引理
\ref{derived-lemma-third-isomorphism-triangle} 得到。
\end{proof}

\begin{lemma}
\label{derived-lemma-uniqueness-third-arrow}
设 $\mathcal{D}$ 为预三角范畴。设
$$
(a, b, c) : (X, Y, Z, f, g, h) \to (X', Y', Z', f', g', h')
$$
为特异三角的态射。若下列条件之一成立：
\begin{enumerate}
\item $\Hom(Y, X') = 0$；
\item $\Hom(Z, Y') = 0$；
\item $\Hom(X, X') = \Hom(Z, X') = 0$；
\item $\Hom(Z, X') = \Hom(Z, Z') = 0$；或
\item $\Hom(X[1], Z') = \Hom(Z, X') = 0$，
\end{enumerate}
则在使 $(a, b, c)$ 成为三角态射的所有 $Y \to Y'$ 态射中，$b$ 是唯一的。
\end{lemma}

\begin{proof}
若还有第二个三角态射 $(a, b', c)$，则 $(0, b - b', 0)$ 是三角态射。
因此只需证明：若态射 $b : Y \to Y'$ 使得
$X \to Y \to Y'$ 与 $Y \to Y' \to Z'$ 均为零，则其必为 $0$。
下文将不再说明地使用引理 \ref{derived-lemma-representable-homological}。
特别地，条件 (3) 蕴含条件 (1)。在条件 (1) 下，若复合
$g' \circ b : Y \to Y' \to Z'$ 为零，则 $b$ 可提升为态射
$Y \to X'$，而后者必为零。这证明了 (1)。

\medskip\noindent
(2) 与 (4) 的证明与此论证对偶。

\medskip\noindent
假设 (5)。考虑图
$$
\xymatrix{
X \ar[r]_f \ar[d]^0 &
Y \ar[r]_g \ar[d]^b &
Z \ar[r]_h \ar[d]^0 \ar@{..>}[ld]^\epsilon &
X[1] \ar[d]^0 \\
X' \ar[r]^{f'} &
Y' \ar[r]^{g'} &
Z' \ar[r]^{h'} &
X'[1]
}
$$
可选择 $\epsilon$ 使 $b = \epsilon \circ g$。于是
$g' \circ \epsilon \circ g = 0$，故对某个
$\delta \in \Hom(X[1], Z')$ 有 $g' \circ \epsilon = \delta \circ h$。
由于 $\Hom(X[1], Z') = 0$，得到 $g' \circ \epsilon = 0$。
因而对某个 $\gamma \in \Hom(Z, X')$ 有
$\epsilon = f' \circ \gamma$。又因 $\Hom(Z, X') = 0$，可知
$\epsilon = 0$，从而如愿得到 $b = 0$。
\end{proof}

\begin{lemma}
\label{derived-lemma-third-object-zero}
设 $\mathcal{D}$ 为预三角范畴，$f : X \to Y$ 为 $\mathcal{D}$ 的态射。
下列条件等价：
\begin{enumerate}
\item $f$ 是同构；
\item $(X, Y, 0, f, 0, 0)$ 是特异三角；
\item 对任一特异三角 $(X, Y, Z, f, g, h)$，都有 $Z = 0$。
\end{enumerate}
\end{lemma}

\begin{proof}
由 TR1，三角 $(X, X, 0, 1, 0, 0)$ 是特异三角。设
$(X, Y, Z, f, g, h)$ 为特异三角。由 TR3，存在特异三角的态射
$(1, f, 0) : (X, X, 0) \to (X, Y, Z)$。若 $f$ 是同构，则由引理
\ref{derived-lemma-third-isomorphism-triangle}，$(1, f, 0)$ 是三角的同构，
故 $Z = 0$。反过来，若 $Z = 0$，则 $(1, f, 0)$ 同样是三角的同构，
因而 $f$ 是同构。
\end{proof}

\begin{lemma}
\label{derived-lemma-direct-sum-triangles}
设 $\mathcal{D}$ 为预三角范畴。设 $(X, Y, Z, f, g, h)$ 和
$(X', Y', Z', f', g', h')$ 为三角。下列条件等价：
\begin{enumerate}
\item $(X \oplus X', Y \oplus Y', Z \oplus Z',
f \oplus f', g \oplus g', h \oplus h')$
是特异三角；
\item $(X, Y, Z, f, g, h)$ 和 $(X', Y', Z', f', g', h')$ 都是特异三角。
\end{enumerate}
\end{lemma}

\begin{proof}
假设 (2)。由 TR1，可选取特异三角
$(X \oplus X', Y \oplus Y', Q, f \oplus f', g'', h'')$。
由 TR3，可找到特异三角的态射
$(X, Y, Z, f, g, h) \to
(X \oplus X', Y \oplus Y', Q, f \oplus f', g'', h'')$
以及
$(X', Y', Z', f', g', h') \to
(X \oplus X', Y \oplus Y', Q, f \oplus f', g'', h'')$。
取这些态射的直和，得到三角的态射
$$
\xymatrix{
(X \oplus X', Y \oplus Y', Z \oplus Z',
f \oplus f', g \oplus g', h \oplus h')
\ar[d]^{(1, 1, c)} \\
(X \oplus X', Y \oplus Y', Q, f \oplus f', g'', h'').
}
$$
用注 \ref{derived-remark-special-triangles} 的术语说，这是特殊三角之间的映射
（因为特殊三角的直和仍然特殊），因此 $c$ 是同构。故 (1) 成立。

\medskip\noindent
假设 (1)。我们证明 $(X, Y, Z, f, g, h)$ 是特异三角。先注意到
$(X, Y, Z, f, g, h)$ 是特殊三角
（使用注 \ref{derived-remark-special-triangles} 的术语），因为它是特异、因而特殊的三角
$(X \oplus X', Y \oplus Y', Z \oplus Z',
f \oplus f', g \oplus g', h \oplus h')$ 的直和因子。
用 TR1 取特异三角 $(X, Y, Q, f, g'', h'')$。由 TR3，存在特异三角的态射
$(X \oplus X', Y \oplus Y', Z \oplus Z',
f \oplus f', g \oplus g', h \oplus h') \to (X, Y, Q, f, g'', h'')$。
将它与嵌入映射复合，得到三角的态射
$$
(1, 1, c) :
(X, Y, Z, f, g, h)
\longrightarrow
(X, Y, Q, f, g'', h'')
$$
由注 \ref{derived-remark-special-triangles}，$c$ 是同构，因此 (2) 成立。
\end{proof}

\begin{lemma}
\label{derived-lemma-split}
设 $\mathcal{D}$ 为预三角范畴，$(X, Y, Z, f, g, h)$ 为特异三角。
\begin{enumerate}
\item 若 $h = 0$，则 $g$ 存在右逆 $s : Z \to Y$。
\item 对 $g$ 的任一右逆 $s : Z \to Y$，映射
$f \oplus s : X \oplus Z \to Y$ 是同构。
\item 对 $\mathcal{D}$ 的任意对象 $X', Z'$，三角
$(X', X' \oplus Z', Z', (1, 0), (0, 1), 0)$ 是特异三角。
\end{enumerate}
\end{lemma}

\begin{proof}
为证明 (1)，使用引理 \ref{derived-lemma-representable-homological} 所给出的正合性
$\Hom_\mathcal{D}(Z, Y) \to \Hom_\mathcal{D}(Z, Z) \to
\Hom_\mathcal{D}(Z, X[1])$。
同理，若 $s$ 如 (2) 所述，则 $h = 0$，且序列
$$
0 \to \Hom_\mathcal{D}(W, X) \to \Hom_\mathcal{D}(W, Y)
\to \Hom_\mathcal{D}(W, Z) \to 0
$$
是分裂正合的（由 $s : Z \to Y$ 分裂）。因此由 Yoneda 引理，
$X \oplus Z \to Y$ 是同构。最后一个断言由 TR1 和引理
\ref{derived-lemma-direct-sum-triangles} 得出。
\end{proof}

\begin{lemma}
\label{derived-lemma-when-split}
设 $\mathcal{D}$ 为预三角范畴，$f : X \to Y$ 为 $\mathcal{D}$ 的态射。
下列条件等价：
\begin{enumerate}
\item $f$ 有核；
\item $f$ 有余核；
% P02-E0019: frozen source has the grammatical typo “is the isomorphic”.
\item 对 $\mathcal{D}$ 的某些对象 $K, Z, Q$，$f$ 同构于投影和余投影的复合
$K \oplus Z \to Z \to Z \oplus Q$。
\end{enumerate}
\end{lemma}

\begin{proof}
任何同构于形如 $X' \oplus Z \to Z \oplus Y'$ 之映射的态射都有核和余核。
故 (3) $\Rightarrow$ (1), (2)。下面证明 (1) $\Rightarrow$ (3)。
先设 $f : X \to Y$ 是单态射，即其核为零。由 TR1，存在特异三角
$(X, Y, Z, f, g, h)$。由引理 \ref{derived-lemma-composition-zero}，复合
$f \circ h[-1] = 0$。由于 $f$ 是单态射，得到 $h[-1] = 0$，从而
$h = 0$。于是引理 \ref{derived-lemma-split} 给出 $Y = X \oplus Z$，即 (3) 成立。
再假设 $f$ 有核 $K$。由于 $K \to X$ 是单态射，可得
$X = K \oplus X'$，并且 $f|_{X'} : X' \to Y$ 是单态射。
因此 $Y = X' \oplus Y'$，结论成立。
(2) $\Rightarrow$ (3) 与此对偶。
\end{proof}

\begin{lemma}
\label{derived-lemma-products-sums-shifts-triangles}
设 $\mathcal{D}$ 为预三角范畴，$I$ 为集合。
\begin{enumerate}
\item 设 $X_i$，$i \in I$，为 $\mathcal{D}$ 的一族对象。
\begin{enumerate}
\item 若 $\prod X_i$ 存在，则 $(\prod X_i)[1] = \prod X_i[1]$。
\item 若 $\bigoplus X_i$ 存在，则 $(\bigoplus X_i)[1] = \bigoplus X_i[1]$。
\end{enumerate}
\item 设 $X_i \to Y_i \to Z_i \to X_i[1]$ 为 $\mathcal{D}$ 的一族特异三角。
\begin{enumerate}
\item 若 $\prod X_i$、$\prod Y_i$、$\prod Z_i$ 存在，则
$\prod X_i \to \prod Y_i \to \prod Z_i \to \prod X_i[1]$
是特异三角。
\item 若 $\bigoplus X_i$、$\bigoplus Y_i$、$\bigoplus Z_i$ 存在，则
$\bigoplus X_i \to \bigoplus Y_i \to \bigoplus Z_i \to \bigoplus X_i[1]$
是特异三角。
\end{enumerate}
\end{enumerate}
\end{lemma}

\begin{proof}
(1) 成立，因为 $[1]$ 是 $\mathcal{D}$ 的自等价，而直和与积都由范畴结构定义。
我们来证明 (2)(a)。选取特异三角
$\prod X_i \to \prod Y_i \to Z \to \prod X_i[1]$。
对每个 $j$，可用 TR3 选取态射 $p_j : Z \to Z_j$，使其与投影映射
$\prod X_i \to X_j$ 和 $\prod Y_i \to Y_j$ 一起构成特异三角的态射。
由积的定义，得到映射 $\prod p_i : Z \to \prod Z_i$，它构成从上述
特异三角到由积组成之三角的三角态射。注意，用注
\ref{derived-remark-special-triangles} 的术语说，``积''三角
$\prod X_i \to \prod Y_i \to \prod Z_i \to \prod X_i[1]$
是特殊的，因为阿贝尔群正合序列的积仍正合。因此注
\ref{derived-remark-special-triangles} 表明该三角态射是同构，再由 TR1 得出结论。
(2)(b) 的证明与此对偶。
\end{proof}

\begin{lemma}
\label{derived-lemma-projectors-have-images-triangulated}
设 $\mathcal{D}$ 为预三角范畴。若 $\mathcal{D}$ 有可数积，则
$\mathcal{D}$ 是 Karoubi 范畴。若 $\mathcal{D}$ 有可数余积，则
$\mathcal{D}$ 是 Karoubi 范畴。
\end{lemma}

\begin{proof}
假设 $\mathcal{D}$ 有可数积。由同调，引理
\ref{homology-lemma-projectors-have-images}，只需验证有右逆的态射都有核。
任何有右逆的态射都是满态射，故由引理 \ref{derived-lemma-when-split} 有核。
第二个断言与第一个对偶。
\end{proof}

\noindent
下面的引理使证明一个预三角范畴是三角范畴略为容易。

\begin{lemma}
\label{derived-lemma-easier-axiom-four}
设 $\mathcal{D}$ 为预三角范畴。为证明 TR4，只需证明：给定任意一对可复合态射
$f : X \to Y$ 与 $g : Y \to Z$，存在
\begin{enumerate}
\item 同构 $i : X' \to X$、$j : Y' \to Y$ 和 $k : Z' \to Z$；进而令
$f' = j^{-1}fi : X' \to Y'$、$g' = k^{-1}gj : Y' \to Z'$ 后，存在
\item 特异三角
$(X', Y', Q_1, f', p_1, d_1)$、
$(X', Z', Q_2, g' \circ f', p_2, d_2)$ 和
$(Y', Z', Q_3, g', p_3, d_3)$，
使 TR4 的断言成立。
\end{enumerate}
\end{lemma}

\begin{proof}
按我们对特异三角及其同构的定义，用 $X', Y', Z'$ 替换 $X, Y, Z$ 无关紧要。
由引理 \ref{derived-lemma-cone-triangle-unique-isomorphism}，特异三角
$(X', Y', Q_1, f', p_1, d_1)$、
$(X', Z', Q_2, g' \circ f', p_2, d_2)$ 和
$(Y', Z', Q_3, g', p_3, d_3)$ 在同构意义下唯一，故引理成立。
\end{proof}

\begin{lemma}
\label{derived-lemma-triangulated-subcategory}
设 $\mathcal{D}$ 为预三角范畴。假设 $\mathcal{D}'$ 是
$\mathcal{D}$ 的加性满子范畴。下列条件等价：
\begin{enumerate}
\item 存在一组三角 $\mathcal{T}'$，使得
$(\mathcal{D}', \mathcal{T}')$ 是 $\mathcal{D}$ 的预三角子范畴；
\item $\mathcal{D}'$ 在 $[1]$ 下保持，
$[1] : \mathcal{D}' \to \mathcal{D}'$ 是自等价，并且给定
$\mathcal{D}'$ 中的任一态射 $f : X \to Y$，在 $\mathcal{D}$ 中存在特异三角
$(X, Y, Z, f, g, h)$，其中 $Z$ 同构于 $\mathcal{D}'$ 的某个对象。
\end{enumerate}
在这种情形下，(1) 中的 $\mathcal{T}'$ 正是 $\mathcal{D}$ 的所有满足
$X, Y, Z \in \Ob(\mathcal{D}')$ 的特异三角
$(X, Y, Z, f, g, h)$ 之集。最后，若 $\mathcal{D}$ 是三角范畴，
则 (1) 和 (2) 还等价于
\begin{enumerate}
\item[(3)] $\mathcal{D}'$ 是三角子范畴。
\end{enumerate}
\end{lemma}

\begin{proof}
略。
\end{proof}
\begin{lemma}
\label{derived-lemma-exact-functor-additive}
预三角范畴之间的正合函子是加性函子。
\end{lemma}

\begin{proof}
设 $F : \mathcal{D} \to \mathcal{D}'$ 为预三角范畴之间的正合函子。
由于 $(0, 0, 0, 1_0, 1_0, 0)$ 是 $\mathcal{D}$ 的特异三角，三角
$$
(F(0), F(0), F(0), 1_{F(0)}, 1_{F(0)}, F(0))
$$
在 $\mathcal{D}'$ 中是特异三角。由引理 \ref{derived-lemma-composition-zero}，
这意味着 $1_{F(0)} \circ 1_{F(0)}$ 为零。因此 $F(0)$ 是
$\mathcal{D}'$ 的零对象。这也意味着 $F$ 把任一零态射映为零
（因为加性范畴中的态射为零，当且仅当它经过零对象分解）。接着，由引理
\ref{derived-lemma-split} 第 (3) 部分，
$(X, X \oplus Y, Y, (1, 0), (0, 1), 0)$ 是特异三角；故
$(F(X), F(X \oplus Y), F(Y), F(1, 0), F(0, 1), 0)$
也是特异三角。由引理 \ref{derived-lemma-split} 第 (2) 部分，映射
$F(X) \oplus F(Y) \to F(X \oplus Y)$ 是同构。最后应用同调，引理
\ref{homology-lemma-additive-functor} 即可。
\end{proof}

\begin{lemma}
\label{derived-lemma-exact-equivalence}
设 $F : \mathcal{D} \to \mathcal{D}'$ 为预三角范畴之间的全忠实正合函子。
则 $\mathcal{D}$ 的三角 $(X, Y, Z, f, g, h)$ 是特异三角，当且仅当
$(F(X), F(Y), F(Z), F(f), F(g), F(h))$ 是 $\mathcal{D}'$ 中的特异三角。
\end{lemma}

\begin{proof}
``仅当''部分显然。假设 $(F(X), F(Y), F(Z))$ 是 $\mathcal{D}'$ 中的特异三角。
在 $\mathcal{D}$ 中选取特异三角 $(X, Y, Z', f, g', h')$。
由引理 \ref{derived-lemma-cone-triangle-unique-isomorphism}，存在三角同构
$$
(1, 1, c') : (F(X), F(Y), F(Z)) \longrightarrow (F(X), F(Y), F(Z')).
$$
由于 $F$ 全忠实，存在态射 $c : Z \to Z'$，使得 $F(c) = c'$。
于是 $(1, 1, c)$ 是 $(X, Y, Z)$ 与 $(X, Y, Z')$ 之间的同构。
因此由 TR1，$(X, Y, Z)$ 是特异三角。
\end{proof}

\begin{lemma}
\label{derived-lemma-composition-exact}
设 $\mathcal{D}, \mathcal{D}', \mathcal{D}''$ 为预三角范畴，
$F : \mathcal{D} \to \mathcal{D}'$ 和
$F' : \mathcal{D}' \to \mathcal{D}''$ 为正合函子。则
$F' \circ F$ 是正合函子。
\end{lemma}

\begin{proof}
略。
\end{proof}

\begin{lemma}
\label{derived-lemma-exact-compose-homological-functor}
设 $\mathcal{D}$ 为预三角范畴，$\mathcal{A}$ 为阿贝尔范畴，
$H : \mathcal{D} \to \mathcal{A}$ 为同调函子。
\begin{enumerate}
\item 设 $\mathcal{D}'$ 为预三角范畴，
$F : \mathcal{D}' \to \mathcal{D}$ 为正合函子。则复合
$H \circ F$ 也是同调函子。
\item 设 $\mathcal{A}'$ 为阿贝尔范畴，
$G : \mathcal{A} \to \mathcal{A}'$ 为正合函子。则
$G \circ H$ 也是同调函子。
\end{enumerate}
\end{lemma}

\begin{proof}
略。
\end{proof}

\begin{lemma}
\label{derived-lemma-exact-compose-delta-functor}
设 $\mathcal{D}$ 为三角范畴，$\mathcal{A}$ 为阿贝尔范畴，
$G : \mathcal{A} \to \mathcal{D}$ 为 $\delta$-函子。
\begin{enumerate}
\item 设 $\mathcal{D}'$ 为三角范畴，
$F : \mathcal{D} \to \mathcal{D}'$ 为正合函子。则复合
$F \circ G$ 也是 $\delta$-函子。
\item 设 $\mathcal{A}'$ 为阿贝尔范畴，
$H : \mathcal{A}' \to \mathcal{A}$ 为正合函子。则
$G \circ H$ 也是 $\delta$-函子。
\end{enumerate}
\end{lemma}

\begin{proof}
略。
\end{proof}

\begin{lemma}
\label{derived-lemma-compose-delta-functor-homological}
设 $\mathcal{D}$ 为三角范畴，$\mathcal{A}$ 和 $\mathcal{B}$ 为阿贝尔范畴。
设 $G : \mathcal{A} \to \mathcal{D}$ 为 $\delta$-函子，
$H : \mathcal{D} \to \mathcal{B}$ 为同调函子。假设对
$\mathcal{A}$ 中所有 $A$ 都有 $H^{-1}(G(A)) = 0$。则族
$$
\{H^n \circ G, H^n(\delta_{A \to B \to C})\}_{n \geq 0}
$$
是从 $\mathcal{A} \to \mathcal{B}$ 的 $\delta$-函子；见同调，定义
\ref{homology-definition-cohomological-delta-functor}。
\end{lemma}

\begin{proof}
该记号表示如下内容。若
$0 \to A \xrightarrow{a} B \xrightarrow{b} C \to 0$ 是
$\mathcal{A}$ 中的短正合序列，则
$$
\delta = \delta_{A \to B \to C} : G(C) \to G(A)[1]
$$
是 $\mathcal{D}$ 中的态射，使得
$(G(A), G(B), G(C), a, b, \delta)$ 是特异三角；见定义
\ref{derived-definition-delta-functor}。于是
$H^n(\delta) : H^n(G(C)) \to H^n(G(A)[1]) = H^{n + 1}(G(A))$
显然关于短正合序列具有函子性。最后，长正合上同调序列
(\ref{derived-equation-long-exact-cohomology-sequence}) 与
$H^{-1}(G(C))$ 的消失合在一起，给出 $\mathcal{B}$ 中的长正合序列
$$
0 \to H^0(G(A)) \to H^0(G(B)) \to H^0(G(C))
\xrightarrow{H^0(\delta)} H^1(G(A)) \to \ldots
$$
如所需。
\end{proof}

\noindent
下面结果的证明使用 TR4。

\begin{proposition}
\label{derived-proposition-9}
设 $\mathcal{D}$ 为三角范畴。任何交换图
$$
\xymatrix{
X \ar[r] \ar[d] & Y \ar[d] \\
X' \ar[r] & Y'
}
$$
都可扩张为图
$$
\xymatrix{
X \ar[r] \ar[d] & Y \ar[r] \ar[d] & Z \ar[r] \ar[d] & X[1] \ar[d] \\
X' \ar[r] \ar[d] & Y' \ar[r] \ar[d] & Z' \ar[r] \ar[d] & X'[1] \ar[d] \\
X'' \ar[r] \ar[d] & Y'' \ar[r] \ar[d] & Z'' \ar[r] \ar[d] & X''[1] \ar[d] \\
X[1] \ar[r] & Y[1] \ar[r] & Z[1] \ar[r] & X[2]
}
$$
其中除右下方形反交换外，所有方形都交换。此外，每一行和每一列都是特异三角。
最后，底行（相应地，右列）中的态射，是对顶行（相应地，左列）的态射应用
$[1]$ 得到的。
\end{proposition}

\begin{proof}
为使证明易读，本证明中不写出箭头。选取特异三角
$(X, Y, Z)$、$(X', Y', Z')$、$(X, X', X'')$、$(Y, Y', Y'')$ 和
$(X, Y', A)$。注意，态射 $X \to Y'$ 既等于复合
$X \to Y \to Y'$，也等于复合 $X \to X' \to Y'$。因此可按 TR4 找到态射
\begin{enumerate}
\item $a : Z \to A$ 和 $b : A \to Y''$；以及
\item $a' : X'' \to A$ 和 $b' : A \to Z'$。
\end{enumerate}
记 $c : Y'' \to Z[1]$ 为复合 $Y'' \to Y[1] \to Z[1]$，并记
$c' : Z' \to X''[1]$ 为复合 $Z' \to X'[1] \to X''[1]$。
此次应用 TR4 的结论是
\begin{enumerate}
\item $(Z, A, Y'', a, b, c)$、$(X'', A, Z', a', b', c')$
是特异三角；
\item $(X, Y, Z) \to (X, Y', A)$、
$(X, Y', A) \to (Y, Y', Y'')$、
$(X, X', X'') \to (X, Y', A)$、
$(X, Y', A) \to (X', Y', Z')$ 是三角的态射。
\end{enumerate}
首先使用 $(X, X', X'') \to (X, Y', A)$ 和
$(X, Y', A) \to (Y, Y', Y'')$ 是三角态射这一事实，可知下列两图中的第一个交换：
$$
\vcenter{
\xymatrix{
X' \ar[r] \ar[d] & Y' \ar[d] \\
X'' \ar[r]^{b \circ a'} \ar[d] & Y'' \ar[d] \\
X[1] \ar[r] & Y[1]
}
}
\quad\text{且}\quad
\vcenter{
\xymatrix{
Y \ar[r] \ar[d] & Z \ar[d]^{b' \circ a} \ar[r] & X[1] \ar[d] \\
Y' \ar[r] & Z' \ar[r] & X'[1]
}
}
$$
使用 $(X, Y, Z) \to (X, Y', A)$ 和
$(X, Y', A) \to (X', Y', Z')$ 是三角态射这一事实，第二个图也交换。
此时，选取以映射 $b \circ a' : X'' \to Y''$ 开始的特异三角
$(X'', Y'' , Z'')$。

\medskip\noindent
接着，对态射 $X'' \to A \to Y''$ 及三角
$(X'', A, Z', a', b', c')$、$(X'', Y'', Z'')$ 和
$(A, Y'', Z[1], b, c , -a[1])$ 再应用一次 TR4，得到态射
$a'' : Z' \to Z''$ 和 $b'' : Z'' \to Z[1]$。
于是 $(Z', Z'', Z[1], a'', b'', - b'[1] \circ a[1])$ 是特异三角，
因而 $(Z, Z', Z'', -b' \circ a, a'', -b'')$ 也是特异三角，进而
$(Z, Z', Z'', b' \circ a, a'', b'')$ 也是特异三角。
此外，$(X'', A, Z') \to (X'', Y'', Z'')$ 和
$(X'', Y'', Z'') \to (A, Y'', Z[1], b, c , -a[1])$
都是三角的态射。至此，所有特异三角和所有态射都已定义，只需验证若干交换关系。

\medskip\noindent
为证明图的中间方形交换，注意箭头 $Y' \to Z'$ 分解为
$Y' \to A \to Z'$，因为 $(X, Y', A) \to (X', Y', Z')$
是三角的态射。同理，态射 $Y' \to Y''$ 分解为
$Y' \to A \to Y''$，因为 $(X, Y', A) \to (Y, Y', Y'')$
是三角的态射。因此，边为 $(A, Z', Z'', Y'')$ 的方形交换，因为
$(X'', A, Z') \to (X'', Y'', Z'')$ 是三角的态射（由 TR4）。
边为 $(Y'', Z'', Y[1], Z[1])$ 的方形交换，因为
$(X'', Y'', Z'') \to (A, Y'', Z[1], b, c , -a[1])$
是三角的态射，且 $c : Y'' \to Z[1]$ 是复合
$Y'' \to Y[1] \to Z[1]$。边为 $(Z', X'[1], X''[1], Z'')$ 的方形交换，
因为 $(X'', A, Z') \to (X'', Y'', Z'')$ 是三角的态射，且
$c' : Z' \to X''[1]$ 是复合 $Z' \to X'[1] \to X''[1]$。
最后必须证明边为 $(Z'', X''[1], Z[1], X[2])$ 的方形反交换。
这是因为 $(X'', Y'', Z'') \to (A, Y'', Z[1], b, c , -a[1])$
是三角的态射，证明完毕。
\end{proof}







\section{三角范畴的局部化}
\label{derived-section-localization}

\noindent
为从复形的同伦范畴构造导出范畴，我们将使用局部化过程。

\begin{definition}
\label{derived-definition-localization}
设 $\mathcal{D}$ 为预三角范畴。若乘法系 $S$ 满足下列两个条件，
则称它{\it 与三角结构相容}：
\begin{enumerate}
\item[MS5] 对 $\mathcal{D}$ 的态射 $f$，有
$f \in S \Leftrightarrow f[1] \in S$\footnote{见注 \ref{derived-remark-MS5}。}。
\item[MS6] 给定实线交换图
$$
\xymatrix{
X \ar[r] \ar[d]^s &
Y \ar[r] \ar[d]^{s'} &
Z \ar[r] \ar@{-->}[d] &
X[1] \ar[d]^{s[1]} \\
X' \ar[r] &
Y' \ar[r] &
Z' \ar[r] &
X'[1]
}
$$
其两行均为特异三角且 $s, s' \in S$，则存在 $S$ 中的态射
$s'' : Z \to Z'$，使得 $(s, s', s'')$ 是三角的态射。
\end{enumerate}
\end{definition}

\noindent
事实证明，这些公理并不独立于定义乘法系的公理。

\begin{lemma}
\label{derived-lemma-localization-conditions}
设 $\mathcal{D}$ 为预三角范畴，$S \subset \text{Arrows}(\mathcal{D})$。
\begin{enumerate}
\item 若 $S$ 包含所有恒等态射并且 MS6 成立
（定义 \ref{derived-definition-localization}），则 $\mathcal{D}$ 的每个同构都属于 $S$。
\item 若 MS1、MS5（范畴，定义
\ref{categories-definition-multiplicative-system}）和 MS6 成立，则 MS2 成立。
\end{enumerate}
\end{lemma}

\begin{proof}
假设 $S$ 包含所有恒等态射且 MS6 成立。设 $f : X \to Y$ 为
$\mathcal{D}$ 的同构。考虑图
$$
\xymatrix{
0 \ar[r] \ar[d]^1 &
X \ar[r]_1 \ar[d]^1 &
X \ar[r] \ar@{-->}[d] &
0[1] \ar[d]^{1[1]} \\
0 \ar[r] &
X \ar[r]^f &
Y \ar[r] &
0[1]
}
$$
由引理 \ref{derived-lemma-third-object-zero}，两行都是特异三角。
由 MS6，虚线箭头存在且属于 $S$，故 $f$ 属于 $S$。

\medskip\noindent
假设 MS1、MS5、MS6 成立。设 $f : X \to Y$ 为 $\mathcal{D}$ 的态射，
$t : X \to X'$ 为 $S$ 的元素。选取特异三角
$(X, Y, Z, f, g, h)$。再选取特异三角
$(X', Y', Z, f', g', t[1] \circ h)$（这里使用 TR1 与 TR2）。
由 MS5、MS6（并用 TR2 旋转），可在交换图
$$
\xymatrix{
X \ar[r] \ar[d]^t &
Y \ar[r] \ar@{..>}[d]^{s'} &
Z \ar[r] \ar[d]^1 &
X[1] \ar[d]^{t[1]} \\
X' \ar[r] &
Y' \ar[r] &
Z \ar[r] &
X'[1]
}
$$
中找到虚线箭头，并且 $s' \in S$。这证明了 LMS2。RMS2 的证明与此对偶。
\end{proof}

\begin{remark}
\label{derived-remark-MS5}
在 MS1 与 MS6 成立时，条件 MS5 等价于要求：对所有
$s \in S$ 和 $n \in \mathbf{Z}$，都有 $s[n] \in S$。
例如，假设 MS5 成立，已知 $s \in S$，而我们想证明 $s[-1] \in S$。
这并非立即成立，因为 $s[-1][1]$ 并不等于 $s$，而只是在
$\mathcal{D}$ 的箭头意义下同构于 $s$。不过，这仍意味着对某些同构 $f$、$g$ 有
$s[-1][1] = f \circ s \circ g$。由引理
\ref{derived-lemma-localization-conditions} (1)，$f, g \in S$；故由 MS1，
$s[-1][1] \in S$；再由 MS5，$s[-1] \in S$。
完整证明留给读者作为练习。
\end{remark}

\begin{lemma}
\label{derived-lemma-triangle-functor-localize}
设 $F : \mathcal{D} \to \mathcal{D}'$ 为预三角范畴之间的正合函子。令
$$
S = \{f \in \text{Arrows}(\mathcal{D})
\mid F(f)\text{ 为同构}\}
$$
则 $S$ 是饱和的（见范畴，定义
\ref{categories-definition-saturated-multiplicative-system}）乘法系，
并且与 $\mathcal{D}$ 上的三角结构相容。
\end{lemma}

\begin{proof}
需要证明公理 MS1 -- MS6；见范畴，定义
\ref{categories-definition-multiplicative-system}、
\ref{categories-definition-saturated-multiplicative-system} 以及定义
\ref{derived-definition-localization}。MS1、MS4 和 MS5 直接来自定义。
MS6 由 TR3 和引理 \ref{derived-lemma-third-isomorphism-triangle} 得出。
由引理 \ref{derived-lemma-localization-conditions}，MS2 成立。为完成证明，
还需证明 MS3。设 $f, g : X \to Y$ 为 $\mathcal{D}$ 的态射，
$t : Z \to X$ 为 $S$ 的元素，并满足 $f \circ t = g \circ t$。
由于 $\mathcal{D}$ 是加性范畴，这只意味着 $a \circ t = 0$，其中
$a = f - g$。用 TR1 选取特异三角 $(Z, X, Q, t, d, h)$。
由于 $a \circ t = 0$，由引理 \ref{derived-lemma-representable-homological}，
存在态射 $i : Q \to Y$，使得 $i \circ d = a$。最后，再次使用 TR1，
可选取三角 $(Q, Y, W, i, j, k)$。图示如下：
$$
\xymatrix{
Z \ar[r]_t & X \ar[r]_d \ar[d]^1 & Q \ar[r] \ar[d]^i & Z[1] \\
& X \ar[r]_a & Y \ar[d]^j \\
& & W
}
$$
现在对该图应用函子 $F$。由于 $t \in S$，由引理
\ref{derived-lemma-third-object-zero} 可知 $F(Q) = 0$。
故由同一引理，$F(j)$ 是同构，即 $j \in S$。最后，
$j \circ a = j \circ i \circ d = 0$，因为 $j \circ i = 0$。
因此 $j \circ f = j \circ g$，从而 LMS3 成立。RMS3 的证明与此对偶。
\end{proof}

\begin{lemma}
\label{derived-lemma-homological-functor-localize}
设 $H : \mathcal{D} \to \mathcal{A}$ 为从预三角范畴到阿贝尔范畴的同调函子。令
$$
S = \{f \in \text{Arrows}(\mathcal{D})
\mid \text{对所有 }i \in \mathbf{Z},\ H^i(f)\text{ 均为同构}\}
$$
则 $S$ 是饱和的（见范畴，定义
\ref{categories-definition-saturated-multiplicative-system}）乘法系，
并且与 $\mathcal{D}$ 上的三角结构相容。
\end{lemma}

\begin{proof}
需要证明公理 MS1 -- MS6；见范畴，定义
\ref{categories-definition-multiplicative-system}、
\ref{categories-definition-saturated-multiplicative-system} 以及定义
\ref{derived-definition-localization}。MS1、MS4 和 MS5 直接来自定义。
MS6 由 TR3 和长正合上同调序列
(\ref{derived-equation-long-exact-cohomology-sequence}) 得出。
由引理 \ref{derived-lemma-localization-conditions}，MS2 成立。为完成证明，
还需证明 MS3。设 $f, g : X \to Y$ 为 $\mathcal{D}$ 的态射，
$t : Z \to X$ 为 $S$ 的元素，并满足 $f \circ t = g \circ t$。
由于 $\mathcal{D}$ 是加性范畴，这只意味着 $a \circ t = 0$，其中
$a = f - g$。用 TR1 和 TR2 选取特异三角 $(Z, X, Q, t, u, v)$。
% P02-E0020: frozen source types the factorization as i \circ g = a; later lines and the diagram require u.
由于 $a \circ t = 0$，由引理 \ref{derived-lemma-representable-homological}，
存在态射 $i : Q \to Y$，使得 $i \circ g = a$。最后，再次使用 TR1，
可选取三角 $(Q, Y, W, i, j, k)$。图示如下：
$$
\xymatrix{
Z \ar[r]_t & X \ar[r]_u \ar[d]^1 & Q \ar[r]_v \ar[d]^i & Z[1] \\
& X \ar[r]_a & Y \ar[d]^j \\
& & W
}
$$
现在对该图应用各函子 $H^i$。由于 $t \in S$，由长正合上同调序列
(\ref{derived-equation-long-exact-cohomology-sequence}) 可知 $H^i(Q) = 0$。
故由同一论证，对所有 $i$，$H^i(j)$ 都是同构，即 $j \in S$。
最后，$j \circ a = j \circ i \circ u = 0$，因为 $j \circ i = 0$。
因此 $j \circ f = j \circ g$，从而 LMS3 成立。RMS3 的证明与此对偶。
\end{proof}

\begin{proposition}
\label{derived-proposition-construct-localization}
设 $\mathcal{D}$ 为预三角范畴，$S$ 为与三角结构相容的乘法系。
则 $S^{-1}\mathcal{D}$ 上存在唯一的预三角范畴结构，使得
$[1] \circ Q = Q \circ [1]$，并且局部化函子
$Q : \mathcal{D} \to S^{-1}\mathcal{D}$ 是正合的。
此外，若 $\mathcal{D}$ 是三角范畴，则 $S^{-1}\mathcal{D}$ 也是三角范畴。
\end{proposition}

\begin{proof}
我们已经在同调，引理 \ref{homology-lemma-localization-additive} 中看到，
$S^{-1}\mathcal{D}$ 是加性范畴，且局部化函子 $Q$ 是加性函子。
由 MS5，存在唯一的加性自等价
$[1] : S^{-1}\mathcal{D} \to S^{-1}\mathcal{D}$，使得
$Q \circ [1] = [1] \circ Q$（函子的等式）；细节从略。
若 $S^{-1}\mathcal{D}$ 的一个三角同构于某个特异三角在局部化函子 $Q$
下的像，则称其为特异三角。

\medskip\noindent
证明 TR1。这里只需证明：若 $a : Q(X) \to Q(Y)$ 是
$S^{-1}\mathcal{D}$ 的态射，则 $a$ 可嵌入一个特异三角。
对某个 $S$ 中的 $s : Y \to Y'$ 和某个 $f : X \to Y'$，写成
$a = Q(s)^{-1} \circ Q(f)$。在 $\mathcal{D}$ 中选取特异三角
$(X, Y', Z, f, g, h)$。于是
$(Q(X), Q(Y), Q(Z), a, Q(g) \circ Q(s), Q(h))$
是 $S^{-1}\mathcal{D}$ 的特异三角。

\medskip\noindent
证明 TR2。由定义立即得到。

\medskip\noindent
证明 TR3。注意，在把两个三角替换为与之同构的三角之后，也可以证明所需虚线箭头的存在性。
因此可设给定 $\mathcal{D}$ 的特异三角
$(X, Y, Z, f, g, h)$ 和 $(X', Y', Z', f', g', h')$，以及交换图
$$
\xymatrix{
Q(X) \ar[r]_{Q(f)} \ar[d]_a & Q(Y) \ar[d]^b \\
Q(X') \ar[r]^{Q(f')} & Q(Y')
}
$$
它位于 $S^{-1}\mathcal{D}$ 中。现在应用范畴，引理
\ref{categories-lemma-left-localization-diagram}，在 $\mathcal{D}$ 中找到态射
$f'' : X'' \to Y''$ 以及交换图
$$
\xymatrix{
X \ar[d]_f \ar[r]_k & X'' \ar[d]^{f''} & X' \ar[d]^{f'} \ar[l]^s \\
Y \ar[r]^l & Y'' & Y' \ar[l]_t
}
$$
其中 $s, t \in S$，且 $a = s^{-1}k$、$b = t^{-1}l$。
此时可在 $\mathcal{D}$ 中使用 TR3 和 MS6，得到交换图
% P02-E0021: frozen source labels this vertical arrow g[1], although k[1] is type-correct.
$$
\xymatrix{
X \ar[r] \ar[d]^k &
Y \ar[r] \ar[d]^l &
Z \ar[r] \ar[d]^m &
X[1] \ar[d]^{g[1]} \\
X'' \ar[r] &
Y'' \ar[r] &
Z'' \ar[r] &
X''[1] \\
X' \ar[r] \ar[u]_s &
Y' \ar[r] \ar[u]_t &
Z' \ar[r] \ar[u]_r &
X'[1] \ar[u]_{s[1]}
}
$$
其中 $r \in S$。于是令 $c = Q(r)^{-1}Q(m)$，便得到所需的三角态射
$$
\xymatrix{
(Q(X), Q(Y), Q(Z), Q(f), Q(g), Q(h))
\ar[d]^{(a, b, c)} \\
(Q(X'), Q(Y'), Q(Z'), Q(f'), Q(g'), Q(h'))
}
$$

\medskip\noindent
这证明了命题的第一个断言。若 $\mathcal{D}$ 还是三角范畴，则为证明
$S^{-1}\mathcal{D}$ 也是三角范畴，还需证明 TR4。由引理
\ref{derived-lemma-easier-axiom-four}，可归结为下列断言：给定可复合态射
$a : Q(X) \to Q(Y)$ 与 $b : Q(Y) \to Q(Z)$，在可能用同构对象替换
$Q(X), Q(Y), Q(Z)$ 后，必须构造一个八面体。为此，可先把 $Y$ 替换为某个对象，
使得对 $\mathcal{D}$ 中某个态射 $f : X \to Y$ 有 $a = Q(f)$。
（更精确地，写成 $a = s^{-1}f$，其中 $s : Y \to Y'$ 属于 $S$，
$f : X \to Y'$，然后用 $Y'$ 替换 $Y$。）此后同样替换 $Z$，使得对
$\mathcal{D}$ 中某个态射 $g : Y \to Z$ 有 $b = Q(g)$。
现在可在 $\mathcal{D}$ 中找到特异三角
$(X, Y, Q_1, f, p_1, d_1)$、
$(X, Z, Q_2, g \circ f, p_2, d_2)$ 和
$(Y, Z, Q_3, g, p_3, d_3)$（由 TR1），以及 TR4 中的态射
$a : Q_1 \to Q_2$ 和 $b : Q_2 \to Q_3$。
立即可验证，对所有这些数据应用函子 $Q$，便在
$S^{-1}\mathcal{D}$ 中得到相应结构。
\end{proof}

\noindent
三角范畴局部化的泛性质如下（我们对预三角范畴表述，因此它自然也适用于三角范畴）。

\begin{lemma}
\label{derived-lemma-universal-property-localization}
设 $\mathcal{D}$ 为预三角范畴，$S$ 为与三角结构相容的乘法系。
设 $Q : \mathcal{D} \to S^{-1}\mathcal{D}$ 为局部化函子，见命题
\ref{derived-proposition-construct-localization}。
\begin{enumerate}
\item 若 $H : \mathcal{D} \to \mathcal{A}$ 是取值于阿贝尔范畴
$\mathcal{A}$ 的同调函子，并且对所有 $s \in S$，$H(s)$ 都是同构，
则满足 $H = H' \circ Q$ 的唯一分解
$H' : S^{-1}\mathcal{D} \to \mathcal{A}$（见范畴，引理
\ref{categories-lemma-properties-left-localization}）也是同调函子。
\item 若 $F : \mathcal{D} \to \mathcal{D}'$ 是取值于预三角范畴
$\mathcal{D}'$ 的正合函子，并且对所有 $s \in S$，$F(s)$ 都是同构，
则满足 $F = F' \circ Q$ 的唯一分解
$F' : S^{-1}\mathcal{D} \to \mathcal{D}'$（见范畴，引理
\ref{categories-lemma-properties-left-localization}）也是正合函子。
\end{enumerate}
\end{lemma}

\begin{proof}
本引理不证自明。细节从略。
\end{proof}

\begin{lemma}
\label{derived-lemma-localization-subcategory}
设 $\mathcal{D}$ 为预三角范畴，$\mathcal{D}' \subset \mathcal{D}$
为满的预三角子范畴。设 $S$ 为 $\mathcal{D}$ 的饱和乘法系，并与三角结构相容。
假设对 $\mathcal{D}$ 中每个 $X$，都存在 $S$ 中的 $s : X' \to X$，
使得 $X'$ 是 $\mathcal{D}'$ 的对象。则
$S' = S \cap \text{Arrows}(\mathcal{D}')$ 是与三角结构相容的饱和乘法系，且函子
$$
(S')^{-1}\mathcal{D}' \longrightarrow S^{-1}\mathcal{D}
$$
是预三角范畴的等价。
\end{lemma}

\begin{proof}
考虑命题 \ref{derived-proposition-construct-localization} 的商函子
$Q : \mathcal{D} \to S^{-1}\mathcal{D}$。由于 $S$ 饱和，态射
$f : X \to Y$ 属于 $S$，当且仅当 $Q(f)$ 可逆；见范畴，引理
\ref{categories-lemma-what-gets-inverted}。因此，$S'$ 正是被复合
$\mathcal{D}' \to \mathcal{D} \to S^{-1}\mathcal{D}$
变为同构的箭头之族。
% P02-E0022: frozen source says “S' is is a saturated”.
故由引理 \ref{derived-lemma-triangle-functor-localize}，$S'$ 是与三角结构相容的饱和乘法系。
由引理 \ref{derived-lemma-universal-property-localization}，得到预三角范畴的正合函子
$(S')^{-1}\mathcal{D}' \to S^{-1}\mathcal{D}$。由假设，此函子本质满。
设 $X', Y'$ 为 $\mathcal{D}'$ 的对象。由范畴，注
\ref{categories-remark-right-localization-morphisms-colimit}，有
$$
\Mor_{S^{-1}\mathcal{D}}(X', Y') =
\colim_{s : X \to X'\text{ 于 }S} \Mor_\mathcal{D}(X, Y')
$$
我们的假设蕴含：对 $S$ 中任意 $s : X \to X'$，都能找到 $S$ 中态射
$s' : X'' \to X$，其中 $X''$ 属于 $\mathcal{D}'$。于是
$s \circ s' : X'' \to X'$ 属于 $S'$，故上面的余极限等于
$$
\colim_{s'' : X'' \to X'\text{ 于 }S'} \Mor_{\mathcal{D}'}(X'', Y') =
\Mor_{(S')^{-1}\mathcal{D}'}(X', Y')
$$
这证明该函子也全忠实，从而证明完成。
\end{proof}

\noindent
下面的引理描述局部化函子的核（见定义 \ref{derived-definition-kernel-category}）。

\begin{lemma}
\label{derived-lemma-kernel-localization}
设 $\mathcal{D}$ 为预三角范畴，$S$ 为与三角结构相容的乘法系，
$Z$ 为 $\mathcal{D}$ 的对象。下列条件等价：
\begin{enumerate}
\item 在 $S^{-1}\mathcal{D}$ 中，$Q(Z) = 0$；
\item 存在 $Z' \in \Ob(\mathcal{D})$，使得 $0 : Z \to Z'$ 是 $S$ 的元素；
\item 存在 $Z' \in \Ob(\mathcal{D})$，使得 $0 : Z' \to Z$ 是 $S$ 的元素；并且
\item 存在对象 $Z'$ 和特异三角 $(X, Y, Z \oplus Z', f, g, h)$，
使得 $f \in S$。
\end{enumerate}
若 $S$ 饱和，则它们还等价于
\begin{enumerate}
\item[(5)] 态射 $0 \to Z$ 是 $S$ 的元素；
\item[(6)] 态射 $Z \to 0$ 是 $S$ 的元素；
\item[(7)] 存在特异三角 $(X, Y, Z, f, g, h)$，使得 $f \in S$。
\end{enumerate}
\end{lemma}

\begin{proof}
(1)、(2)、(3) 的等价性即同调，引理
\ref{homology-lemma-kernel-localization}。若 (2) 成立，则
$(Z'[-1], Z'[-1] \oplus Z, Z, (1, 0), (0, 1), 0)$ 是特异三角
（见引理 \ref{derived-lemma-split}），并带有“$0 \in S$”。旋转后得到 (4)。
若 $(X, Y, Z \oplus Z', f, g, h)$ 是满足 $f \in S$ 的特异三角，
则 $Q(f)$ 是同构，因而 $Q(Z \oplus Z') = 0$，进而 $Q(Z) = 0$。
所以 (1) -- (4) 全部等价。

\medskip\noindent
下面假设 $S$ 饱和。注意，(5)、(6)、(7) 各自都蕴含等价条件 (1) -- (4)
之一。设 $Q(Z) = 0$。于是 $0 \to Z$ 是 $\mathcal{D}$ 中的态射，
并在 $S^{-1}\mathcal{D}$ 中变为同构。根据范畴，引理
\ref{categories-lemma-what-gets-inverted}，$S$ 的饱和性蕴含
$0 \to Z$ 属于 $S$。故 (1) $\Rightarrow$ (5)。对偶地，
(1) $\Rightarrow$ (6)。最后，若 $0 \to Z$ 属于 $S$，则三角
$(0, Z, Z, 0, \text{id}_Z, 0)$ 由 TR1 与 TR2 是特异三角，且正是 (7) 中的三角。
\end{proof}

\begin{lemma}
\label{derived-lemma-limit-triangles}
设 $\mathcal{D}$ 为预三角范畴，$S$ 为 $\mathcal{D}$ 中与三角结构相容的
饱和乘法系。设 $(X, Y, Z, f, g, h)$ 为 $\mathcal{D}$ 中的特异三角。
考虑三角态射所成的范畴
$$
\mathcal{I} =
\{(s, s', s'') : (X, Y, Z, f, g, h) \to (X', Y', Z', f', g', h')
\mid s, s', s'' \in S\}
$$
则 $\mathcal{I}$ 是滤过范畴，并且函子
$\mathcal{I} \to X/S$、$\mathcal{I} \to Y/S$ 和 $\mathcal{I} \to Z/S$
都是余终的。
\end{lemma}

\begin{proof}
我们强烈建议读者跳过本引理的证明，改在餐巾纸上自行推导。

\medskip\noindent
第一个观察是：旋转特异三角（TR2）给出 $\mathcal{I}$ 与特异三角
$(Y, Z, X[1], g, h, -f[1])$ 所对应范畴之间的等价。例如，由此可见，
若我们证明函子 $\mathcal{I} \to X/S$ 余终，那么函子
$\mathcal{I} \to Y/S$ 和 $\mathcal{I} \to Z/S$ 也余终。

\medskip\noindent
注意，若 $s : X \to X'$ 是 $S$ 的态射，则由 MS2 可找到
$s' : Y \to Y'$ 和 $f' : X' \to Y'$，使得
$f' \circ s = s' \circ f$；随后可用 MS6 将其补成 $\mathcal{I}$ 的一个对象。
因此函子 $\mathcal{I} \to X/S$ 在对象上满。由上述旋转，对
$\mathcal{I} \to Y/S$ 和 $\mathcal{I} \to Z/S$ 也同样成立。

\medskip\noindent
设给定 $X/S$ 中的对象 $s_1 : X \to X_1 $、$s_2 : X \to X_2$，
以及 $X/S$ 中的态射 $a : X_1 \to X_2$。由于 $S$ 饱和，由范畴，引理
\ref{categories-lemma-what-gets-inverted} 可知 $a \in S$。
按照上一段的论证，可把 $s_1 : X \to X_1$ 补成 $\mathcal{I}$ 中的对象
$(s_1, s'_1, s''_1) : (X, Y, Z, f, g, h) \to (X_1, Y_1, Z_1, f_1, g_1, h_1)$。
然后重复此步骤，找到
$(a, b, c) : (X_1, Y_1, Z_1, f_1, g_1, h_1) \to (X_2, Y_2, Z_2, f_2, g_2, h_2)$，
其中 $a, b, c \in S$，并补全给定的 $a : X_1 \to X_2$。
于是 $(a, b, c)$ 是 $\mathcal{I}$ 中的态射。因此函子
$\mathcal{I} \to X/S$ 在箭头上也满。再由上述旋转，对函子
$\mathcal{I} \to Y/S$ 和 $\mathcal{I} \to Z/S$ 亦然。

\medskip\noindent
范畴 $\mathcal{I}$ 非空，因为恒等态射给出一个对象。这证明了滤过范畴定义的条件 (1)；
见范畴，定义 \ref{categories-definition-directed}。

\medskip\noindent
我们为范畴 $\mathcal{I}$ 验证范畴，定义
\ref{categories-definition-directed} 的条件 (2)。设给定 $\mathcal{I}$ 中的对象
$(s_1, s'_1, s''_1) : (X, Y, Z, f, g, h) \to (X_1, Y_1, Z_1, f_1, g_1, h_1)$
和
$(s_2, s'_2, s''_2) : (X, Y, Z, f, g, h) \to (X_2, Y_2, Z_2, f_2, g_2, h_2)$。
我们要找 $\mathcal{I}$ 中的一个对象，使它同时是从
$(X_1, Y_1, Z_1, f_1, g_1, h_1)$ 和 $(X_2, Y_2, Z_2, f_2, g_2, h_2)$
出发之箭头的靶。由范畴，注
\ref{categories-remark-left-localization-morphisms-colimit}，范畴
$X/S$、$Y/S$、$Z/S$ 都是滤过的。因此可在 $X/S$ 中找到
$X \to X_3$ 以及态射 $s : X_2 \to X_3$ 和 $a : X_1 \to X_3$。
由上文可找到三角态射
$(s, s', s'') : (X_2, Y_2, Z_2, f_2, g_2, h_2) \to
(X_3, Y_3, Z_3, f_3, g_3, h_3)$，其中 $s', s'' \in S$。
把 $(X_2, Y_2, Z_2)$ 替换为 $(X_3, Y_3, Z_3)$ 后，可假设
$X/S$ 中存在态射 $a : X_1 \to X_2$。对 $Y$ 和 $Z$ 重复此论证
（用上述旋转），可假设存在 $X/S$ 中的态射 $a : X_1 \to X_2$、
$Y/S$ 中的态射 $b : Y_1 \to Y_2$，以及 $Z/S$ 中的态射
$c : Z_1 \to Z_2$。然而，这些态射未必给出特异三角的态射。
另一方面，所需各图在 $S^{-1}\mathcal{D}$ 中确实交换。因此，例如可见，
存在 $S$ 中的态射 $s'_2 : Y_2 \to Y_3$，使得
$s'_2 \circ f_2 \circ a = s'_2 \circ b \circ f_1$。
再像上面那样替换一次 $(X_2, Y_2, Z_2)$，便可进入
$f_2 \circ a = b \circ f_1$ 的情形。旋转并再应用两次同一论证，
可假设 $(a, b, c)$ 是三角的态射。这证明了条件 (2)。

\medskip\noindent
接着验证范畴，定义 \ref{categories-definition-directed} 的条件 (3)。
设 $(s_1, s_1', s_1'') : (X, Y, Z) \to (X_1, Y_1, Z_1)$ 和
$(s_2, s_2', s_2'') : (X, Y, Z) \to (X_2, Y_2, Z_2)$
是 $\mathcal{I}$ 的对象，并设 $(a, b, c), (a', b', c')$ 是它们之间的两个态射。
由于 $a \circ s_1 = a' \circ s_1$，存在态射 $s_3 : X_2 \to X_3$，使得
$s_3 \circ a = s_3 \circ a'$。使用上述满性断言，可将它补成三角态射
$(s_3, s_3', s_3'') : (X_2, Y_2, Z_2) \to (X_3, Y_3, Z_3)$，
其中 $s_3, s_3', s_3'' \in S$。因此
$(s_3 \circ s_2, s_3' \circ s_2', s_3'' \circ s_2'') :
(X, Y, Z) \to (X_3, Y_3, Z_3)$ 也是 $\mathcal{I}$ 的对象；将映射
$(a, b, c), (a', b', c')$ 与 $(s_3, s_3', s_3'')$ 复合后，得到
$a = a'$。通过旋转可同样得到 $b = b'$ 和 $c = c'$。

\medskip\noindent
最后验证 $\mathcal{I} \to X/S$ 余终；见范畴，定义
\ref{categories-definition-cofinal}。第一个条件成立，因为该函子满。
假设在 $X/S$ 中有对象 $s : X \to X'$，并且在 $\mathcal{I}$ 中有两个对象
$(s_1, s'_1, s''_1) : (X, Y, Z, f, g, h) \to (X_1, Y_1, Z_1, f_1, g_1, h_1)$
和
$(s_2, s'_2, s''_2) : (X, Y, Z, f, g, h) \to (X_2, Y_2, Z_2, f_2, g_2, h_2)$，
并有 $X/S$ 中的态射 $t_1 : X' \to X_1$ 和 $t_2 : X' \to X_2$。
由上面已证的 $\mathcal{I}$ 的性质 (2)，可找到 $\mathcal{I}$ 中的态射
$(s_3, s'_3, s''_3) : (X_1, Y_1, Z_1, f_1, g_1, h_1) \to
(X_3, Y_3, Z_3, f_3, g_3, h_3)$
和
$(s_4, s'_4, s''_4) : (X_2, Y_2, Z_2, f_2, g_2, h_2) \to
(X_3, Y_3, Z_3, f_3, g_3, h_3)$。
若复合 $X' \to X_1 \to X_3$ 与 $X' \to X_2 \to X_3$ 相等，
则已经完成（见范畴，定义 \ref{categories-definition-cofinal} 中的显示等式）。
若不相等，则因 $X/S$ 滤过，可选取 $X/S$ 中的态射 $X_3 \to X_4$，
使复合 $X' \to X_1 \to X_3 \to X_4$ 与
$X' \to X_2 \to X_3 \to X_4$ 相等。最后，把 $X_3 \to X_4$
补成 $\mathcal{I}$ 中的态射
$(X_3, Y_3, Z_3) \to (X_4, Y_4, Z_4)$，再与它复合，便得到结论。
\end{proof}







\section{三角范畴的商}
\label{derived-section-quotients}

\noindent
给定一个三角范畴及其一个三角子范畴，可以取“商”来构造另一个三角范畴。
这一构造使用局部化。它类似于阿贝尔范畴对 Serre 子范畴的商；见同调，
第 \ref{homology-section-serre-subcategories} 节。在实际构造之前，我们先简要讨论
正合函子的核。

\begin{definition}
\label{derived-definition-saturated}
设 $\mathcal{D}$ 为预三角范畴。若每当 $X \oplus Y$ 同构于
$\mathcal{D}'$ 的一个对象时，$X$ 和 $Y$ 都分别同构于
$\mathcal{D}'$ 的对象，则称 $\mathcal{D}$ 的满预三角子范畴
$\mathcal{D}'$ 是{\it 饱和的}。
\end{definition}

\noindent
饱和三角子范畴有时称为{\it 厚三角子范畴}。在某些文献中，
这一术语只用于严格满三角子范畴（有时定义本身即蕴含严格满性）。
还有另一个概念，称为 {\it 'epaisse 三角子范畴}。其定义是：给定交换图
$$
\xymatrix{
& S \ar[rd] \\
X \ar[ru] \ar[rr] & & Y \ar[r] & T \ar[r] & X[1]
}
$$
若第二行是特异三角，且 $S$ 与 $T$ 均同构于 $\mathcal{D}'$ 的对象，
则 $X$ 与 $Y$ 也同构于 $\mathcal{D}'$ 的对象。事实证明，这等价于饱和性
（这是初等事实，可见 \cite{Rickard-derived}），而饱和范畴这一概念更便于使用。

\begin{lemma}
\label{derived-lemma-triangle-functor-kernel}
设 $F : \mathcal{D} \to \mathcal{D}'$ 为预三角范畴之间的正合函子。
设 $\mathcal{D}''$ 为 $\mathcal{D}$ 的满子范畴，其对象为
$$
\Ob(\mathcal{D}'') =
\{X \in \Ob(\mathcal{D}) \mid F(X) = 0\}
$$
则 $\mathcal{D}''$ 是 $\mathcal{D}$ 的严格满、饱和预三角子范畴。
若 $\mathcal{D}$ 是三角范畴，则 $\mathcal{D}''$ 是三角子范畴。
\end{lemma}

\begin{proof}
显然，$\mathcal{D}''$ 在 $[1]$ 与 $[-1]$ 下保持。
若 $(X, Y, Z, f, g, h)$ 是 $\mathcal{D}$ 的特异三角，并且
$F(X) = F(Y) = 0$，则由于
$(F(X), F(Y), F(Z), F(f), F(g), F(h))$ 是特异三角，也有
$F(Z) = 0$。因此可应用引理 \ref{derived-lemma-triangulated-subcategory}，
得到 $\mathcal{D}''$ 是预三角子范畴（相应地，若 $\mathcal{D}$ 是三角范畴，
则它是三角子范畴）。最后，饱和性由
$F(X) \oplus F(Y) = 0 \Rightarrow F(X) = F(Y) = 0$ 得出。
\end{proof}

\begin{lemma}
\label{derived-lemma-homological-functor-kernel}
设 $H : \mathcal{D} \to \mathcal{A}$ 为从预三角范畴到阿贝尔范畴的同调函子。
设 $\mathcal{D}'$ 为 $\mathcal{D}$ 的满子范畴，其对象为
$$
\Ob(\mathcal{D}') =
\{X \in \Ob(\mathcal{D}) \mid
H(X[n]) = 0\text{ 对所有 }n \in \mathbf{Z}\}
$$
则 $\mathcal{D}'$ 是 $\mathcal{D}$ 的严格满、饱和预三角子范畴。
若 $\mathcal{D}$ 是三角范畴，则 $\mathcal{D}'$ 是三角子范畴。
\end{lemma}

\begin{proof}
显然，$\mathcal{D}'$ 在 $[1]$ 与 $[-1]$ 下保持。若
$(X, Y, Z, f, g, h)$ 是 $\mathcal{D}$ 的特异三角，并且对所有 $n$ 都有
$H(X[n]) = H(Y[n]) = 0$，则由长正合序列
(\ref{derived-equation-long-exact-cohomology-sequence})，对所有 $n$ 也有
$H(Z[n]) = 0$。因此可应用引理 \ref{derived-lemma-triangulated-subcategory}，
得到 $\mathcal{D}'$ 是预三角子范畴（相应地，若 $\mathcal{D}$ 是三角范畴，
则它是三角子范畴）。饱和性由
\begin{align*}
H((X \oplus Y)[n]) = 0 & \Rightarrow H(X[n] \oplus Y[n]) = 0 \\
& \Rightarrow H(X[n]) \oplus H(Y[n]) = 0 \\
& \Rightarrow H(X[n]) = H(Y[n]) = 0
\end{align*}
对所有 $n \in \mathbf{Z}$ 成立而得。
\end{proof}

\begin{lemma}
\label{derived-lemma-homological-functor-bounded}
设 $H : \mathcal{D} \to \mathcal{A}$ 为从预三角范畴到阿贝尔范畴的同调函子。
设 $\mathcal{D}_H^{+}, \mathcal{D}_H^{-}, \mathcal{D}_H^b$
为 $\mathcal{D}$ 的满子范畴，其对象为
$$
\begin{matrix}
\Ob(\mathcal{D}_H^{+}) =
\{X \in \Ob(\mathcal{D}) \mid
H(X[n]) = 0\text{ 对所有 }n \ll 0\} \\
\Ob(\mathcal{D}_H^{-}) =
\{X \in \Ob(\mathcal{D}) \mid
H(X[n]) = 0\text{ 对所有 }n \gg 0\} \\
\Ob(\mathcal{D}_H^b) =
\{X \in \Ob(\mathcal{D}) \mid
H(X[n]) = 0\text{ 对所有 }|n| \gg 0\}
\end{matrix}
$$
它们每一个都是 $\mathcal{D}$ 的严格满、饱和预三角子范畴。
若 $\mathcal{D}$ 是三角范畴，则它们每一个都是三角子范畴。
\end{lemma}

\begin{proof}
我们对 $\mathcal{D}_H^{+}$ 证明此事。显然，它在 $[1]$ 与 $[-1]$ 下保持。
若 $(X, Y, Z, f, g, h)$ 是 $\mathcal{D}$ 的特异三角，并且对所有
$n \ll 0$ 有 $H(X[n]) = H(Y[n]) = 0$，则由长正合序列
(\ref{derived-equation-long-exact-cohomology-sequence})，对所有 $n \ll 0$
也有 $H(Z[n]) = 0$。因此可应用引理
\ref{derived-lemma-triangulated-subcategory}，得到 $\mathcal{D}_H^{+}$ 是预三角子范畴
（相应地，若 $\mathcal{D}$ 是三角范畴，则它是三角子范畴）。饱和性由
\begin{align*}
H((X \oplus Y)[n]) = 0 & \Rightarrow H(X[n] \oplus Y[n]) = 0 \\
& \Rightarrow H(X[n]) \oplus H(Y[n]) = 0 \\
& \Rightarrow H(X[n]) = H(Y[n]) = 0
\end{align*}
对所有 $n \in \mathbf{Z}$ 成立而得。
\end{proof}

\begin{definition}
\label{derived-definition-kernel-category}
设 $\mathcal{D}$ 为（预）三角范畴。
\begin{enumerate}
\item 设 $F : \mathcal{D} \to \mathcal{D}'$ 为正合函子。
{\it $F$ 的核}是引理 \ref{derived-lemma-triangle-functor-kernel} 所描述的严格满、
饱和（预）三角子范畴。
\item 设 $H : \mathcal{D} \to \mathcal{A}$ 为同调函子。
{\it $H$ 的核}是引理 \ref{derived-lemma-homological-functor-kernel} 所描述的严格满、
饱和（预）三角子范畴。
\end{enumerate}
它们有时记为 $\Ker(F)$ 或 $\Ker(H)$。
\end{definition}

\noindent
下面引理的证明使用 TR4。

\begin{lemma}
\label{derived-lemma-construct-multiplicative-system}
设 $\mathcal{D}$ 为三角范畴，$\mathcal{D}' \subset \mathcal{D}$
为满三角子范畴。令
\begin{equation}
\label{derived-equation-multiplicative-system}
S =
\left\{
\begin{matrix}
f \in \text{Arrows}(\mathcal{D})
\text{ 使得存在一个特异三角 }\\
(X, Y, Z, f, g, h) \text{ 属于 }\mathcal{D}\text{ 其中 }
Z\text{ 同构于下列范畴中的对象： }\mathcal{D}'
\end{matrix}
\right\}
\end{equation}
则 $S$ 是与 $\mathcal{D}$ 上三角结构相容的乘法系。在此情形下，下列条件等价：
\begin{enumerate}
\item $S$ 是饱和乘法系；
\item $\mathcal{D}'$ 是饱和三角子范畴。
\end{enumerate}
\end{lemma}

\begin{proof}
为证明第一个断言，需要证明 MS1、MS2、MS3、MS5、MS6。

\medskip\noindent
证明 MS1。显然，恒等态射属于 $S$，因为对 $\mathcal{D}$ 的每个对象 $X$，
$(X, X, 0, 1, 0, 0)$ 都是特异三角，而且 $0$ 是 $\mathcal{D}'$ 的对象。
设 $f : X \to Y$ 和 $g : Y \to Z$ 为 $S$ 中的可复合态射。
选取特异三角 $(X, Y, Q_1, f, p_1, d_1)$、
$(X, Z, Q_2, g \circ f, p_2, d_2)$ 和
$(Y, Z, Q_3, g, p_3, d_3)$。由假设，$Q_1$ 和 $Q_3$ 同构于
$\mathcal{D}'$ 的对象。由 TR4，存在特异三角
$(Q_1, Q_2, Q_3, a, b, c)$。由于 $\mathcal{D}'$ 是三角子范畴，
$Q_2$ 同构于 $\mathcal{D}'$ 的对象。因此 $g \circ f \in S$。

\medskip\noindent
证明 MS3。设 $a : X \to Y$ 为态射，$t : Z \to X$ 为 $S$ 中满足
$a \circ t = 0$ 的元素。为证明 LMS3，只需找到 $s \in S$ 使得
$s \circ a = 0$；可与引理 \ref{derived-lemma-triangle-functor-localize} 的证明比较。
用 TR1 与 TR2 选取特异三角 $(Z, X, Q, t, g, h)$。
由于 $a \circ t = 0$，由引理 \ref{derived-lemma-representable-homological}，
存在态射 $i : Q \to Y$ 使得 $i \circ g = a$。最后再次使用 TR1，
选取三角 $(Q, Y, W, i, s, k)$。图示如下：
$$
\xymatrix{
Z \ar[r]_t & X \ar[r]_g \ar[d]^1 & Q \ar[r] \ar[d]^i & Z[1] \\
& X \ar[r]_a & Y \ar[d]^s \\
& & W
}
$$
由于 $t \in S$，$Q$ 同构于 $\mathcal{D}'$ 的对象。因此 $s \in S$。
最后，由引理 \ref{derived-lemma-composition-zero}，$s \circ i = 0$，从而
$s \circ a = s \circ i \circ g = 0$。故 LMS3 成立。RMS3 的证明与此对偶。

\medskip\noindent
证明 MS5。这由特异三角及 $\mathcal{D}'$ 在平移下稳定得出。

\medskip\noindent
证明 MS6。设给定交换图
$$
\xymatrix{
X \ar[r] \ar[d]^s &
Y \ar[d]^{s'} \\
X' \ar[r] &
Y'
}
$$
其中 $s, s' \in S$。由命题 \ref{derived-proposition-9}，可将其扩张为九方图。
由于 $s, s'$ 是 $S$ 的元素，$X'', Y''$ 同构于 $\mathcal{D}'$ 的对象。
又因 $\mathcal{D}'$ 是满三角子范畴，$Z''$ 也同构于 $\mathcal{D}'$ 的对象。
因而态射 $Z \to Z'$ 是 $S$ 的元素。这证明了 MS6。

\medskip\noindent
MS2 是 MS1、MS5、MS6 的形式推论；见引理
\ref{derived-lemma-localization-conditions}。这完成了引理第一个断言的证明。

\medskip\noindent
假设 $S$ 饱和。（下文将不再说明地使用特异三角的旋转。）
设 $X \oplus Y$ 是同构于 $\mathcal{D}'$ 中某对象的对象。考虑态射
$f : 0 \to X$。复合 $0 \to X \to X \oplus Y$ 属于 $S$，因为
$(0, X \oplus Y, X \oplus Y, 0, 1, 0)$ 是特异三角。
复合 $Y[-1] \to 0 \to X$ 属于 $S$，因为
$(X, X \oplus Y, Y, (1, 0), (0, 1), 0)$ 是特异三角；见引理
\ref{derived-lemma-split}。由于 $S$ 饱和，$0 \to X$ 属于 $S$。
因此，如所需，$X$ 同构于 $\mathcal{D}'$ 的对象。

\medskip\noindent
最后，假设 $\mathcal{D}'$ 是饱和三角子范畴。设
$$
W \xrightarrow{h}
X \xrightarrow{g}
Y \xrightarrow{f} Z
$$
为 $\mathcal{D}$ 的可复合态射，且 $fg, gh \in S$。
我们将逐步构造下图中的对象：
$$
\xymatrix{
 & &
Q_{12} \ar[rd] & &
Q_{23} \ar[rd] \\
 &
Q_1 \ar[ld]_{\! + \! 1} \ar[ru] & &
Q_2 \ar[ld]_{\! + \! 1} \ar[ll]_{\! + \! 1} \ar[ru] & &
Q_3 \ar[ld]_{\! + \! 1} \ar[ll]_{\! + \! 1} \\
W \ar[rr] & &
X \ar[lu] \ar[rr] & &
Y \ar[lu] \ar[rr] & &
Z \ar[lu]
}
$$
先选取特异三角 $(W, X, Q_1)$、$(X, Y, Q_2)$、$(Y, Z, Q_3)$、
$(W, Y, Q_{12})$ 和 $(X, Z, Q_{23})$。记
$s : Q_2 \to Q_1[1]$ 为复合 $Q_2 \to X[1] \to Q_1[1]$；记
$t : Q_3 \to Q_2[1]$ 为复合 $Q_3 \to Y[1] \to Q_2[1]$。
把 TR4 应用于复合 $W \to X \to Y$ 与复合 $X \to Y \to Z$，
得到使用态射 $s$ 和 $t$ 的特异三角
$(Q_1, Q_{12}, Q_2)$ 和 $(Q_2, Q_{23}, Q_3)$。
由于已假设 $W \to Y$ 与 $X \to Z$ 属于 $S$，对象 $Q_{12}$ 和
$Q_{23}$ 同构于 $\mathcal{D}'$ 的对象。又因 $S$ 在复合和平移下封闭，
$s[1]t$ 也属于 $S$。注意 $s[1]t = 0$，因为
$Y[1] \to Q_2[1] \to X[2]$ 为零；见引理
\ref{derived-lemma-composition-zero}。因此由引理 \ref{derived-lemma-split}，
$Q_3[1] \oplus Q_1[2]$ 同构于 $\mathcal{D}'$ 的对象。
由对 $\mathcal{D}'$ 的假设，$Q_3$ 和 $Q_1$ 同构于
$\mathcal{D}'$ 的对象。考察特异三角 $(Q_1, Q_{12}, Q_2)$，可知
$Q_2$ 也同构于 $\mathcal{D}'$ 的对象。最后考察特异三角
$(X, Y, Q_2)$，得到 $g \in S$。
% P02-E0023: frozen source has “It is also follows”.
（还可推出 $h, f \in S$，但我们不需要这一点。）
\end{proof}

\begin{definition}
\label{derived-definition-quotient-category}
设 $\mathcal{D}$ 为三角范畴，$\mathcal{B}$ 为满三角子范畴。
定义{\it 商范畴 $\mathcal{D}/\mathcal{B}$} 为
$\mathcal{D}/\mathcal{B} = S^{-1}\mathcal{D}$，其中 $S$ 是通过引理
\ref{derived-lemma-construct-multiplicative-system} 与 $\mathcal{B}$ 相伴的
$\mathcal{D}$ 的乘法系。此时局部化函子
$Q : \mathcal{D} \to \mathcal{D}/\mathcal{B}$ 称为{\it 商函子}。
\end{definition}

\noindent
注意，商函子 $Q : \mathcal{D} \to \mathcal{D}/\mathcal{B}$
是三角范畴之间的正合函子；见命题 \ref{derived-proposition-construct-localization}。
这一构造的泛性质如下。

\begin{lemma}
\label{derived-lemma-universal-property-quotient}
\begin{slogan}
Verdier 商的泛性质。
\end{slogan}
设 $\mathcal{D}$ 为三角范畴，$\mathcal{B}$ 为 $\mathcal{D}$ 的满三角子范畴，
$Q : \mathcal{D} \to \mathcal{D}/\mathcal{B}$ 为商函子。
\begin{enumerate}
\item 若 $H : \mathcal{D} \to \mathcal{A}$ 是取值于阿贝尔范畴
$\mathcal{A}$ 的同调函子，并且 $\mathcal{B} \subset \Ker(H)$，
则存在唯一分解 $H' : \mathcal{D}/\mathcal{B} \to \mathcal{A}$，
使得 $H = H' \circ Q$，并且 $H'$ 也是同调函子。
\item 若 $F : \mathcal{D} \to \mathcal{D}'$ 是取值于预三角范畴
$\mathcal{D}'$ 的正合函子，并且 $\mathcal{B} \subset \Ker(F)$，
则存在唯一分解 $F' : \mathcal{D}/\mathcal{B} \to \mathcal{D}'$，
使得 $F = F' \circ Q$，并且 $F'$ 也是正合函子。
\end{enumerate}
\end{lemma}

\begin{proof}
本引理由引理 \ref{derived-lemma-universal-property-localization} 得出。
事实上，若 $f : X \to Y$ 是 $\mathcal{D}$ 的态射，并且对某个特异三角
$(X, Y, Z, f, g, h)$，对象 $Z$ 同构于 $\mathcal{B}$ 的对象，
则在 (1)（相应地，(2)）的假设下，$H(f)$（相应地，$F(f)$）是同构。
细节从略。
\end{proof}

\noindent
商函子的核可描述如下。

\begin{lemma}
\label{derived-lemma-kernel-quotient}
设 $\mathcal{D}$ 为三角范畴，$\mathcal{B}$ 为满三角子范畴。
商函子 $Q : \mathcal{D} \to \mathcal{D}/\mathcal{B}$ 的核是
$\mathcal{D}$ 的严格满子范畴，其对象为
$$
\Ob(\Ker(Q)) =
\left\{
\begin{matrix}
Z \in \Ob(\mathcal{D})
\text{ 使得存在一个 }Z' \in \Ob(\mathcal{D}) \\
\text{ 使得 }Z \oplus Z'\text{ 同构于下列范畴中的对象： }\mathcal{B}
\end{matrix}
\right\}
$$
换言之，它是包含 $\mathcal{B}$ 的 $\mathcal{D}$ 的最小严格满、饱和三角子范畴。
\end{lemma}

\begin{proof}
先注意，核自动是严格满三角子范畴，并且包含其任一对象的所有直和因子；
见引理 \ref{derived-lemma-triangle-functor-kernel}。其对象的上述描述由定义和引理
\ref{derived-lemma-kernel-localization} 第 (4) 部分得出。
\end{proof}

\noindent
设 $\mathcal{D}$ 为三角范畴。现在已有一些构造，它们在下列两个偏序集之间诱导保序映射：
\begin{enumerate}
\item $\mathcal{D}$ 中与三角结构相容的乘法系 $S$ 所成的偏序集；
\item 满三角子范畴 $\mathcal{B} \subset \mathcal{D}$ 所成的偏序集。
\end{enumerate}
具体地，这些构造由
$S \mapsto \mathcal{B}(S) = \Ker(Q : \mathcal{D} \to S^{-1}\mathcal{D})$
和 $\mathcal{B} \mapsto S(\mathcal{B})$ 给出，其中 $S(\mathcal{B})$
是 (\ref{derived-equation-multiplicative-system}) 中的乘法集，即
$$
S(\mathcal{B}) =
\left\{
\begin{matrix}
f \in \text{Arrows}(\mathcal{D})
\text{ 使得存在一个特异三角 }\\
(X, Y, Z, f, g, h) \text{ 属于 }\mathcal{D}\text{ 其中 }
Z\text{ 同构于下列范畴中的对象： }\mathcal{B}
\end{matrix}
\right\}
$$
注意，这两个操作并不互为逆操作。

\begin{lemma}
\label{derived-lemma-operations}
设 $\mathcal{D}$ 为三角范畴。上述操作具有下列性质：
\begin{enumerate}
\item $S(\mathcal{B}(S))$ 是 $S$ 的“饱和”，即包含 $S$ 的
$\mathcal{D}$ 中最小饱和乘法系；
\item $\mathcal{B}(S(\mathcal{B}))$ 是 $\mathcal{B}$ 的“饱和”，
即包含 $\mathcal{B}$ 的 $\mathcal{D}$ 中最小严格满、饱和三角子范畴。
\end{enumerate}
特别地，这些构造在以下两个（偏序）集合之间定义互逆映射：
$\mathcal{D}$ 中与 $\mathcal{D}$ 上三角结构相容的饱和乘法系所成的集合，
以及 $\mathcal{D}$ 的严格满、饱和三角子范畴所成的集合。
\end{lemma}

\begin{proof}
先从满三角子范畴 $\mathcal{B}$ 出发。此时
$\mathcal{B}(S(\mathcal{B})) =
\Ker(Q : \mathcal{D} \to \mathcal{D}/\mathcal{B})$，
所以 (2) 正是引理 \ref{derived-lemma-kernel-quotient} 的内容。

\medskip\noindent
再设 $S$ 是 $\mathcal{D}$ 中与 $\mathcal{D}$ 上三角结构相容的乘法系。
于是 $\mathcal{B}(S) = \Ker(Q : \mathcal{D} \to S^{-1}\mathcal{D})$。
因此（在局部化范畴中使用引理 \ref{derived-lemma-third-object-zero}）
\begin{align*}
S(\mathcal{B}(S))
& =
\left\{
\begin{matrix}
f \in \text{Arrows}(\mathcal{D})
\text{ 使得存在一个特异}\\
\text{特异三角 }(X, Y, Z, f, g, h) \text{ 属于 }\mathcal{D}\text{ 其中 }Q(Z) = 0
\end{matrix}
\right\}.
\\
& =
\{f \in \text{Arrows}(\mathcal{D}) \mid Q(f)\text{ 为同构}\} \\
& = \hat S = S'
\end{align*}
其中记号取自范畴，引理 \ref{categories-lemma-what-gets-inverted}。
该引理的最后一个断言完成了证明。
\end{proof}

\begin{lemma}
\label{derived-lemma-acyclic-general}
设 $H : \mathcal{D} \to \mathcal{A}$ 为从三角范畴 $\mathcal{D}$ 到
阿贝尔范畴 $\mathcal{A}$ 的同调函子；见定义 \ref{derived-definition-homological}。
$\mathcal{D}$ 的子范畴 $\Ker(H)$ 是 $\mathcal{D}$ 的严格满、饱和三角子范畴，
而与其相应的饱和乘法系（见引理 \ref{derived-lemma-operations}）是集合
$$
S = \{f \in \text{Arrows}(\mathcal{D}) \mid
\text{对所有 }i \in \mathbf{Z},\ H^i(f)\text{ 均为同构}\}.
$$
函子 $H$ 经过商函子 $Q : \mathcal{D} \to \mathcal{D}/\Ker(H)$ 分解。
\end{lemma}

\begin{proof}
由引理 \ref{derived-lemma-homological-functor-kernel}，范畴 $\Ker(H)$
是 $\mathcal{D}$ 的严格满、饱和三角子范畴。由引理
\ref{derived-lemma-homological-functor-localize}，集合 $S$ 是与三角结构相容的
饱和乘法系。回忆与 $\Ker(H)$ 相应的乘法系是集合
$$
\left\{
\begin{matrix}
f \in \text{Arrows}(\mathcal{D})
\text{ 使得存在一个特异三角 }\\
(X, Y, Z, f, g, h)\text{ 其中 } H^i(Z) = 0 \text{ 对所有 }i
\end{matrix}
\right\}.
$$
由长正合上同调序列（见 (\ref{derived-equation-long-exact-cohomology-sequence})），
显然，$f$ 属于该集合当且仅当 $f$ 属于 $S$。最后，$H$ 经过 $Q$ 的分解
是引理 \ref{derived-lemma-universal-property-quotient} 的推论。
\end{proof}









\section{正合函子的伴随}
\label{derived-section-adjoints}

\noindent
本节给出三角范畴之间伴随函子的有关结果。

\begin{lemma}
\label{derived-lemma-adjoint-is-exact}
设 $F : \mathcal{D} \to \mathcal{D}'$ 为三角范畴之间的正合函子。
若 $F$ 有右伴随 $G: \mathcal{D'} \to \mathcal{D}$，则 $G$ 也是正合函子。
\end{lemma}

\begin{proof}
设 $\xi_X : F(X[1]) \to F(X)[1]$ 如定义
\ref{derived-definition-exact-functor-triangulated-categories} 所述。
设 $\epsilon_A : F(G(A)) \to A$ 为伴随映射。考虑复合
$$
F(G(A)[1]) \xrightarrow{\xi_{G(A)}} F(G(A))[1] \xrightarrow{\epsilon_A[1]} A[1]
$$
此映射的伴随是一个映射 $G(A)[1] \to G(A[1])$，我们断言它是同构。
为此，由 Yoneda 引理，只需证明：对 $\mathcal{D}$ 的每个对象 $X$ 应用
$\Mor_\mathcal{D}(X[1], -)$ 后，得到双射。
现在，每个态射 $f : X[1] \to G(A)[1]$ 都唯一地形如 $g[1]$，其中
$g : X \to G(A)$；每个 $g$ 又唯一地伴随于某个 $h : F(X) \to A$；
进而 $h[1] \circ \xi_X$ 唯一地伴随于某个 $i : X[1] \to G(A[1])$。
把 $f$ 送到 $i$ 的规则给出
$\Mor_\mathcal{D}(X[1], G(A)[1])$ 与
$\Mor_\mathcal{D}(X[1], G(A[1]))$ 之间的双射。
该双射确由上述态射给出，这可由范畴，第
\ref{categories-section-adjoint} 节的讨论以及下列计算得出：
\begin{align*}
\epsilon_A[1] \circ \xi_{G(A)} \circ F(f)
& =
\epsilon_A[1] \circ \xi_{G(A)} \circ F(g[1]) \\
& =
\epsilon_A[1] \circ F(g)[1] \circ \xi_X \\
& = 
(\epsilon_A \circ F(g))[1] \circ \xi_X \\
& =
h[1] \circ \xi_X
\end{align*}
略去一些细节。我们将用同构 $G(A)[1] \to G(A[1])$ 的逆证明
$G$ 是正合函子。这些同构关于 $A$ 具有函子性（细节从略）。

\medskip\noindent
设 $A \to B \to C \to A[1]$ 为 $\mathcal{D}'$ 中的特异三角。
在 $\mathcal{D}$ 中选取特异三角
$$
G(A) \to G(B) \to X \to G(A)[1]
$$
于是 $F(G(A)) \to F(G(B)) \to F(X) \to F(G(A))[1]$
是 $\mathcal{D}'$ 中的特异三角。由 TR3，可选取特异三角的态射
$$
\xymatrix{
F(G(A)) \ar[r] \ar[d]^{\epsilon_A} &
F(G(B)) \ar[r] \ar[d]^{\epsilon_B} &
F(X) \ar[r] \ar[d] &
F(G(A))[1] \ar[d]^{\epsilon_A[1]} \\
A \ar[r] & B \ar[r] & C \ar[r] & A[1]
}
$$
回忆，箭头 $F(X) \to F(G(A))[1]$ 是
$F(X) \to F(G(A)[1])$ 与 $\xi_{G(A)}$ 的复合。
因此利用 $G$ 是 $F$ 的右伴随及上述讨论，得到态射 $X \to G(C)$，
使下图交换：
$$
\xymatrix{
G(A) \ar[r] \ar[d] & G(B) \ar[r] \ar[d] & X \ar[r] \ar[d] & G(A)[1] \ar[d] \\
G(A) \ar[r] & G(B) \ar[r] & G(C) \ar[r] & G(A[1])
}
$$
且右边竖直箭头正是上面构造的态射。
% P02-E0024: frozen source applies Hom_{D'} to X and G(C), although both lie in D.
对 $\mathcal{D}'$ 的对象 $W$ 应用同调函子
$\Hom_{\mathcal{D}'}(W, -)$，由 $5$ 引理推得
$$
\Hom_{\mathcal{D}'}(W, X) \to \Hom_{\mathcal{D}'}(W, G(C))
$$
是双射；再次使用 Yoneda 引理，得到 $X \to G(C)$ 是同构。
因此 $G(A) \to G(B) \to G(C) \to G(A)[1]$ 是特异三角，正如所需。
\end{proof}

\begin{lemma}
\label{derived-lemma-fully-faithful-adjoint-kernel-zero}
设 $\mathcal{D}$、$\mathcal{D}'$ 为三角范畴，
$F : \mathcal{D} \to \mathcal{D}'$ 和
$G : \mathcal{D}' \to \mathcal{D}$ 为函子。假设
\begin{enumerate}
\item $F$ 与 $G$ 都是正合函子；
\item $F$ 全忠实；
\item $G$ 是 $F$ 的右伴随；
\item $G$ 的核为零。
\end{enumerate}
则 $F$ 是范畴等价。
\end{lemma}

\begin{proof}
由于 $F$ 全忠实，伴随映射 $\text{id} \to G \circ F$ 是同构
（范畴，引理 \ref{categories-lemma-adjoint-fully-faithful}）。
设 $X$ 为 $\mathcal{D}'$ 的对象。选取 $\mathcal{D}'$ 中的特异三角
$$
F(G(X)) \to X \to Y \to F(G(X))[1]
$$
应用 $G$，并使用 $G(F(G(X))) = G(X)$，得到特异三角
$$
G(X) \to G(X) \to G(Y) \to G(X)[1]
$$
因此 $G(Y) = 0$，故 $Y = 0$。所以 $F(G(X)) \to X$ 是同构。
\end{proof}








\section{同伦范畴}
\label{derived-section-homotopy}

\noindent
设 $\mathcal{A}$ 为加性范畴。$\mathcal{A}$ 的同伦范畴
$K(\mathcal{A})$ 是由 $\mathcal{A}$ 的复形组成的范畴，其态射是复形态射模同伦。
下面给出形式定义。

\begin{definition}
\label{derived-definition-complexes-notation}
设 $\mathcal{A}$ 为加性范畴。
\begin{enumerate}
\item 令 $\text{Comp}(\mathcal{A}) = \text{CoCh}(\mathcal{A})$
为{\it （上链）复形范畴}。
\item 若对所有 $n \ll 0$ 都有 $K^n = 0$，则称复形 $K^\bullet$
{\it 下有界}。
\item 若对所有 $n \gg 0$ 都有 $K^n = 0$，则称复形 $K^\bullet$
{\it 上有界}。
\item 若对所有 $|n| \gg 0$ 都有 $K^n = 0$，则称复形 $K^\bullet$
{\it 有界}。
\item 令 $\text{Comp}^{+}(\mathcal{A})$、$\text{Comp}^{-}(\mathcal{A})$、
相应地 $\text{Comp}^b(\mathcal{A})$ 为 $\text{Comp}(\mathcal{A})$
的满子范畴，其对象分别为下有界、上有界、相应地有界的复形。
\item 令 $K(\mathcal{A})$ 为与 $\text{Comp}(\mathcal{A})$ 有相同对象的范畴，
而其态射为复形映射的同伦类（见同调，引理
\ref{homology-lemma-compose-homotopy-cochain}）。
\item 令 $K^{+}(\mathcal{A})$、$K^{-}(\mathcal{A})$、
相应地 $K^b(\mathcal{A})$ 为 $K(\mathcal{A})$ 的满子范畴，
其对象分别为 $\mathcal{A}$ 的下有界、上有界、相应地有界复形。
\end{enumerate}
\end{definition}

\noindent
我们将看到，范畴 $K(\mathcal{A})$、$K^{+}(\mathcal{A})$、
$K^{-}(\mathcal{A})$ 和 $K^b(\mathcal{A})$
都是三角范畴。为了证明这一点，我们先发展一些
与锥和分裂正合序列有关的工具。




\section{锥与逐项分裂序列}
\label{derived-section-cones}

\noindent
设 $\mathcal{A}$ 为加性范畴，并设 $K(\mathcal{A})$ 表示由
$\mathcal{A}$ 的复形组成的范畴，其态射是模同伦的复形态射。
注意，复形上的平移函子 $[n]$（见同调，定义
\ref{homology-definition-shift-cochain}）诱导出函子
$[n] : K(\mathcal{A}) \to K(\mathcal{A})$，满足
$[n] \circ [m] = [n + m]$ 且 $[0] = \text{id}$。

\begin{definition}
\label{derived-definition-cone}
设 $\mathcal{A}$ 为加性范畴。设
$f : K^\bullet \to L^\bullet$ 为 $\mathcal{A}$ 的复形的态射。
$f$ 的{\it 锥}是复形 $C(f)^\bullet$，其定义为
$C(f)^n = L^n \oplus K^{n + 1}$，微分为
$$
d_{C(f)}^n =
\left(
\begin{matrix}
d^n_L & f^{n + 1} \\
0 & -d_K^{n + 1}
\end{matrix}
\right)
$$
它带有典范复形态射
$i : L^\bullet \to C(f)^\bullet$ 和 $p : C(f)^\bullet \to K^\bullet[1]$，
它们由显然的映射 $L^n \to C(f)^n \to K^{n + 1}$ 诱导。
\end{definition}

\noindent
换言之，$(K, L, C(f), f, i, p)$ 构成一个三角：
$$
K^\bullet \to L^\bullet \to C(f)^\bullet \to K^\bullet[1]
$$
这个三角的构造在下述意义下具有函子性。

\begin{lemma}
\label{derived-lemma-functorial-cone}
假设
$$
\xymatrix{
K_1^\bullet \ar[r]_{f_1} \ar[d]_a & L_1^\bullet \ar[d]^b \\
K_2^\bullet \ar[r]^{f_2} & L_2^\bullet
}
$$
是一个复形态射图，并且它交换至同伦。那么存在态射
$c : C(f_1)^\bullet \to C(f_2)^\bullet$，它在 $K(\mathcal{A})$ 中给出三角的态射
$(a, b, c) : (K_1^\bullet, L_1^\bullet, C(f_1)^\bullet, f_1, i_1, p_1)
\to
(K_2^\bullet, L_2^\bullet, C(f_2)^\bullet, f_2, i_2, p_2)$。
\end{lemma}

\begin{proof}
设 $h^n : K_1^n \to L_2^{n - 1}$ 为一族态射，满足
$b \circ f_1 - f_2 \circ a= d \circ h + h \circ d$。
用下列矩阵定义 $c^n$：
$$
c^n =
\left(
\begin{matrix}
b^n & h^{n + 1} \\
0 & a^{n + 1}
\end{matrix}
\right) :
L_1^n \oplus K_1^{n + 1} \to L_2^n \oplus K_2^{n + 1}
$$
% P02-E0025: frozen source has the grammatical error "A matrix computation show".
矩阵计算表明 $c$ 是复形态射。显然有
$c \circ i_1 = i_2 \circ b$，同样不难验证
$p_2 \circ c = a \circ p_1$。
\end{proof}

\noindent
注意，引理 \ref{derived-lemma-functorial-cone} 的证明中构造的态射
$c : C(f_1)^\bullet \to C(f_2)^\bullet$ 一般依赖于
$f_2 \circ a$ 与 $b \circ f_1$ 之间所选的同伦 $h$。

\begin{lemma}
\label{derived-lemma-map-from-cone}
假设 $f: K^\bullet \to L^\bullet$ 和 $g : L^\bullet \to M^\bullet$
是复形态射，且 $g \circ f$ 零同伦。那么
\begin{enumerate}
\item $g$ 经由某个态射 $C(f)^\bullet \to M^\bullet$ 分解；
\item $f$ 经由某个态射 $K^\bullet \to C(g)^\bullet[-1]$ 分解。
\end{enumerate}
\end{lemma}

\begin{proof}
这些假设表明图
$$
\xymatrix{
K^\bullet \ar[r]_f \ar[d] & L^\bullet \ar[d]^g \\
0 \ar[r] & M^\bullet
}
$$
交换至同伦。由于 $0 \to M^\bullet$ 的锥是 $M^\bullet$，
引理 \ref{derived-lemma-functorial-cone} 中的映射
$C(f)^\bullet \to C(0 \to M^\bullet) = M^\bullet$
就是 (1) 中的映射。$K^\bullet \to 0$ 的锥是
$K^\bullet[1]$，而应用引理 \ref{derived-lemma-functorial-cone}
得到映射 $K^\bullet[1] \to C(g)^\bullet$。应用 $[-1]$
便得到 (2) 中的映射。
\end{proof}

\noindent
注意，引理 \ref{derived-lemma-map-from-cone} 的证明中构造的态射
$C(f)^\bullet \to M^\bullet$ 和 $K^\bullet \to C(g)^\bullet[-1]$
一般依赖于所选的同伦。

\begin{definition}
\label{derived-definition-termwise-split-map}
设 $\mathcal{A}$ 为加性范畴。
{\it 逐项分裂单射 $\alpha : A^\bullet \to B^\bullet$}
是一个复形态射，使得每个 $A^n \to B^n$
都同构于某个直和项的嵌入。
{\it 逐项分裂满射 $\beta : B^\bullet \to C^\bullet$}
是一个复形态射，使得每个 $B^n \to C^n$
都同构于到某个直和项的投影。
\end{definition}

\begin{lemma}
\label{derived-lemma-make-commute-map}
设 $\mathcal{A}$ 为加性范畴。设
$$
\xymatrix{
A^\bullet \ar[r]_f \ar[d]_a & B^\bullet \ar[d]^b \\
C^\bullet \ar[r]^g & D^\bullet
}
$$
是一个交换至同伦的复形态射图。
若 $f$ 是逐项分裂单射，则 $b$ 同伦于一个使该图交换的态射。
若 $g$ 是逐项分裂满射，则 $a$ 同伦于一个使该图交换的态射。
\end{lemma}

\begin{proof}
设 $h^n : A^n \to D^{n - 1}$ 为一族态射，满足
$bf - ga = dh + hd$。假设 $\pi^n : B^n \to A^n$
是分裂态射 $f^n$ 的态射。取
$b' = b - dh\pi - h\pi d$。假设
$s^n : D^n \to C^n$ 是分裂态射
$g^n : C^n \to D^n$ 的态射。取 $a' = a + dsh + shd$。
计算从略。
\end{proof}

\noindent
下一引理可用来将复形态射替换为一个复形态射，
使它在每一次数上都是直和项的嵌入。

\begin{lemma}
\label{derived-lemma-make-injective}
设 $\mathcal{A}$ 为加性范畴。设
$\alpha : K^\bullet \to L^\bullet$ 为 $\mathcal{A}$ 的复形的态射。
存在分解
$$
\xymatrix{
K^\bullet \ar[r]^{\tilde \alpha} \ar@/_1pc/[rr]_\alpha &
\tilde L^\bullet \ar[r]^\pi &
L^\bullet
}
$$
使得
\begin{enumerate}
\item $\tilde \alpha$ 是逐项分裂单射（见定义
\ref{derived-definition-termwise-split-map}）；
\item 存在复形映射 $s : L^\bullet \to \tilde L^\bullet$，使得
$\pi \circ s = \text{id}_{L^\bullet}$，并且
$s \circ \pi$ 与 $\text{id}_{\tilde L^\bullet}$ 同伦。
\end{enumerate}
此外，若 $K^\bullet$ 和 $L^\bullet$ 均属于
$K^{+}(\mathcal{A})$、$K^{-}(\mathcal{A})$ 或 $K^b(\mathcal{A})$，
则 $\tilde L^\bullet$ 也属于同一范畴。
\end{lemma}

\begin{proof}
我们令
$$
\tilde L^n = L^n \oplus K^n \oplus K^{n + 1}
$$
并定义
$$
d^n_{\tilde L} =
\left(
\begin{matrix}
d^n_L & 0 & 0 \\
0 & d^n_K & \text{id}_{K^{n + 1}} \\
0 & 0 & -d^{n + 1}_K
\end{matrix}
\right)
$$
换言之，$\tilde L^\bullet = L^\bullet \oplus C(1_{K^\bullet})$。
再令
$$
\tilde \alpha =
\left(
\begin{matrix}
\alpha \\
\text{id}_{K^n} \\
0
\end{matrix}
\right)
$$
这显然是分裂单射，而且显然定义了复形态射。定义
$$
\pi =
\left(
\begin{matrix}
\text{id}_{L^n} &
0 &
0
\end{matrix}
\right)
$$
于是显然有 $\pi \circ \tilde \alpha = \alpha$。我们令
$$
s =
\left(
\begin{matrix}
\text{id}_{L^n} \\
0 \\
0
\end{matrix}
\right)
$$
于是 $\pi \circ s = \text{id}_{L^\bullet}$。最后，设
$h^n : \tilde L^n \to \tilde L^{n - 1}$ 是如下映射：
它通过恒等态射将 $\tilde L^n$ 的直和项 $K^n$
映到 $\tilde L^{n - 1}$ 的直和项 $K^n$。那么很容易（见下面备注中的计算）
证明
$$
\text{id}_{\tilde L^\bullet} - s \circ \pi
=
d \circ h + h \circ d
$$
这就证完了引理。
\end{proof}

\begin{remark}
\label{derived-remark-compute-modules}
为了验证上述证明中最后的等式，可以如下逐元论证。我们有
$s\pi(l, k, k^{+}) = (l, 0, 0)$。
因此等式左边的态射将
$(l, k, k^{+})$ 映到 $(0, k, k^{+})$。
另一方面，$h(l, k, k^{+}) = (0, 0, k)$ 且
$d(l, k, k^{+}) = (dl, dk + k^{+}, -dk^{+})$。
因此 $(dh + hd)(l, k, k^{+}) =
d(0, 0, k) + h(dl, dk + k^{+}, -dk^{+}) =
(0, k, -dk) + (0, 0, dk + k^{+}) = (0, k, k^{+})$，
正如所需。
\end{remark}

\begin{lemma}
\label{derived-lemma-make-surjective}
设 $\mathcal{A}$ 为加性范畴。设
$\alpha : K^\bullet \to L^\bullet$ 为 $\mathcal{A}$ 的复形的态射。
存在分解
$$
\xymatrix{
K^\bullet \ar[r]^i \ar@/_1pc/[rr]_\alpha &
\tilde K^\bullet \ar[r]^{\tilde \alpha} &
L^\bullet
}
$$
使得
\begin{enumerate}
\item $\tilde \alpha$ 是逐项分裂满射（见定义
\ref{derived-definition-termwise-split-map}）；
\item 存在复形映射 $s : \tilde K^\bullet \to K^\bullet$，使得
$s \circ i = \text{id}_{K^\bullet}$，并且
$i \circ s$ 与 $\text{id}_{\tilde K^\bullet}$ 同伦。
\end{enumerate}
此外，若 $K^\bullet$ 和 $L^\bullet$ 均属于
$K^{+}(\mathcal{A})$、$K^{-}(\mathcal{A})$ 或 $K^b(\mathcal{A})$，
则 $\tilde K^\bullet$ 也属于同一范畴。
\end{lemma}

\begin{proof}
这是引理 \ref{derived-lemma-make-injective} 的对偶命题。取
$$
\tilde K^n = K^n \oplus L^{n - 1} \oplus L^n
$$
并定义
$$
d^n_{\tilde K} =
\left(
\begin{matrix}
d^n_K & 0 & 0 \\
0 & - d^{n - 1}_L & \text{id}_{L^n} \\
0 & 0 & d^n_L
\end{matrix}
\right)
$$
换言之，$\tilde K^\bullet = K^\bullet \oplus C(1_{L^\bullet[-1]})$。
再令
$$
\tilde \alpha =
\left(
\begin{matrix}
\alpha &
0 &
\text{id}_{L^n}
\end{matrix}
\right)
$$
这显然是分裂满射，而且显然定义了复形态射。定义
$$
i =
\left(
\begin{matrix}
\text{id}_{K^n} \\
0 \\
0
\end{matrix}
\right)
$$
于是显然有 $\tilde \alpha \circ i = \alpha$。我们令
$$
s =
\left(
\begin{matrix}
\text{id}_{K^n} &
0 &
0
\end{matrix}
\right)
$$
于是 $s \circ i = \text{id}_{K^\bullet}$。最后，设
$h^n : \tilde K^n \to \tilde K^{n - 1}$ 是如下映射：
它通过恒等态射将 $\tilde K^n$ 的直和项 $L^{n - 1}$
映到 $\tilde K^{n - 1}$ 的直和项 $L^{n - 1}$。那么很容易证明
$$
\text{id}_{\tilde K^\bullet} - i \circ s
=
d \circ h + h \circ d
$$
这就证完了引理。
\end{proof}

\begin{definition}
\label{derived-definition-split-ses}
设 $\mathcal{A}$ 为加性范畴。
$\mathcal{A}$ 的复形的{\it 逐项分裂正合序列}
是一个复形之复形
$$
0 \to
A^\bullet \xrightarrow{\alpha}
B^\bullet \xrightarrow{\beta}
C^\bullet \to 0
$$
连同给定的直和分解 $B^n = A^n \oplus C^n$，
并且这些分解与 $\alpha^n$ 和 $\beta^n$ 相容。
我们常将直和分解诱导的映射记为
$s^n : C^n \to B^n$ 和 $\pi^n : B^n \to A^n$。
根据同调，引理 \ref{homology-lemma-ses-termwise-split-cochain}，
我们得到相应的复形态射
$$
\delta : C^\bullet \longrightarrow A^\bullet[1]
$$
它在次数 $n$ 上的映射为
$\pi^{n + 1} \circ d_B^n \circ s^n$。换言之，
$(A^\bullet, B^\bullet, C^\bullet, \alpha, \beta, \delta)$
构成三角
$$
A^\bullet \to B^\bullet \to C^\bullet \to A^\bullet[1]
$$
它称为{\it 与该复形的逐项分裂序列相关联的三角}。
\end{definition}

\begin{lemma}
\label{derived-lemma-triangle-independent-splittings}
设 $\mathcal{A}$ 为加性范畴。设
$0 \to A^\bullet \to B^\bullet \to C^\bullet \to 0$
为定义 \ref{derived-definition-split-ses} 中的逐项分裂正合序列。
设 $(\pi')^n$、$(s')^n$ 为另一组分裂。以
$\delta' : C^\bullet \longrightarrow A^\bullet[1]$ 表示
与这第二组分裂相关联的态射。那么
$$
(1, 1, 1) :
(A^\bullet, B^\bullet, C^\bullet, \alpha, \beta, \delta)
\longrightarrow
(A^\bullet, B^\bullet, C^\bullet, \alpha, \beta, \delta')
$$
是 $K(\mathcal{A})$ 中的三角同构。
\end{lemma}

\begin{proof}
该命题无非是说 $\delta$ 和 $\delta'$ 是同伦的复形映射。
这就是同调，引理
\ref{homology-lemma-ses-termwise-split-homotopy-cochain}。
\end{proof}

\begin{remark}
\label{derived-remark-make-commute}
设 $\mathcal{A}$ 为加性范畴。设
$0 \to A_i^\bullet \to B_i^\bullet \to C_i^\bullet \to 0$、$i = 1, 2$
为逐项分裂正合序列。假设
$a : A_1^\bullet \to A_2^\bullet$、
$b : B_1^\bullet \to B_2^\bullet$ 和
$c : C_1^\bullet \to C_2^\bullet$ 是复形态射，使得
$$
\xymatrix{
A_1^\bullet \ar[d]_a \ar[r] &
B_1^\bullet \ar[r] \ar[d]_b &
C_1^\bullet \ar[d]_c \\
A_2^\bullet \ar[r] & B_2^\bullet \ar[r] & C_2^\bullet
}
$$
在 $K(\mathcal{A})$ 中交换。一般而言，{\bf 不}存在与 $b$
同伦的态射 $b' : B_1^\bullet \to B_2^\bullet$，
使上图在复形范畴中交换。事实上，考虑例，等式
(\ref{examples-equation-commutes-up-to-homotopy})。
若能将其中的中间映射替换为一个与它同伦且使该图交换的映射，
那么我们就会得到并不成立的迹的可加性。
\end{remark}

\begin{lemma}
\label{derived-lemma-nilpotent}
设 $\mathcal{A}$ 为加性范畴。设
$0 \to A_i^\bullet \to B_i^\bullet \to C_i^\bullet \to 0$、$i = 1, 2, 3$
为复形的逐项分裂正合序列。设
$b : B_1^\bullet \to B_2^\bullet$ 和 $b' : B_2^\bullet \to B_3^\bullet$
为复形态射，使得
$$
\vcenter{
\xymatrix{
A_1^\bullet \ar[d]_0 \ar[r] &
B_1^\bullet \ar[r] \ar[d]_b &
C_1^\bullet \ar[d]_0 \\
A_2^\bullet \ar[r] & B_2^\bullet \ar[r] & C_2^\bullet
}
}
\quad\text{且}\quad
\vcenter{
\xymatrix{
A_2^\bullet \ar[d]^0 \ar[r] &
B_2^\bullet \ar[r] \ar[d]^{b'} &
C_2^\bullet \ar[d]^0 \\
A_3^\bullet \ar[r] & B_3^\bullet \ar[r] & C_3^\bullet
}
}
$$
在 $K(\mathcal{A})$ 中交换。那么
$b' \circ b = 0$ 在 $K(\mathcal{A})$ 中成立。
\end{lemma}

\begin{proof}
根据引理 \ref{derived-lemma-make-commute-map}，我们可以用同伦的映射
替换 $b$ 和 $b'$，使左图的右方块交换、右图的左方块交换。
换言之，我们有
$\Im(b^n) \subset \Im(A_2^n \to B_2^n)$ 以及
$\Ker((b')^n) \supset \Im(A_2^n \to B_2^n)$。
因而 $b' \circ b = 0$ 作为复形映射成立。
\end{proof}

\begin{lemma}
\label{derived-lemma-third-isomorphism}
设 $\mathcal{A}$ 为加性范畴。设
$f_1 : K_1^\bullet \to L_1^\bullet$ 和
$f_2 : K_2^\bullet \to L_2^\bullet$ 为复形态射。设
$$
(a, b, c) :
(K_1^\bullet, L_1^\bullet, C(f_1)^\bullet, f_1, i_1, p_1)
\longrightarrow
(K_2^\bullet, L_2^\bullet, C(f_2)^\bullet, f_2, i_2, p_2)
$$
是 $K(\mathcal{A})$ 中任意的三角态射。
若 $a$ 和 $b$ 是同伦等价，则 $c$ 也是同伦等价。
\end{lemma}

\begin{proof}
设 $a^{-1} : K_2^\bullet \to K_1^\bullet$ 为复形态射，
它在 $K(\mathcal{A})$ 中是 $a$ 的逆。设
$b^{-1} : L_2^\bullet \to L_1^\bullet$ 为复形态射，
它在 $K(\mathcal{A})$ 中是 $b$ 的逆。设
$c' : C(f_2)^\bullet \to C(f_1)^\bullet$ 是把引理
\ref{derived-lemma-functorial-cone} 应用于
$f_1 \circ a^{-1} = b^{-1} \circ f_2$ 所得的态射。
若能证明 $c \circ c'$ 和 $c' \circ c$ 是 $K(\mathcal{A})$ 中的同构，
则即得结论。因而只需证明下述命题：对于 $K(\mathcal{A})$ 中任意的三角态射
$(1, 1, c) : (K^\bullet, L^\bullet, C(f)^\bullet, f, i, p)$，
态射 $c$ 是 $K(\mathcal{A})$ 中的同构。
由假设，图中的两个方块
$$
\xymatrix{
L^\bullet \ar[r] \ar[d]_1 &
C(f)^\bullet \ar[r] \ar[d]_c &
K^\bullet[1] \ar[d]_1 \\
L^\bullet \ar[r] &
C(f)^\bullet \ar[r] &
K^\bullet[1]
}
$$
交换至同伦。由 $C(f)^\bullet$ 的构造，两行均构成复形的逐项分裂序列。
因此，由引理 \ref{derived-lemma-nilpotent}，我们看到
$(c - 1)^2 = 0$ 在 $K(\mathcal{A})$ 中成立。
故 $c$ 是 $K(\mathcal{A})$ 中的同构，其逆为 $2 - c$。
\end{proof}

\noindent
因此，若 $a$ 和 $b$ 是同伦等价，则所得的三角态射
是 $K(\mathcal{A})$ 中的三角同构。
事实上，$K(\mathcal{A})$ 中由锥给出的三角之类，
与 $K(\mathcal{A})$ 中由复形的逐项分裂序列给出的三角之类，
在同构意义下是相同的，至少相差一个符号！

\begin{lemma}
\label{derived-lemma-the-same-up-to-isomorphisms}
设 $\mathcal{A}$ 为加性范畴。
\begin{enumerate}
\item 给定复形的逐项分裂序列
$(\alpha : A^\bullet \to B^\bullet,
\beta : B^\bullet \to C^\bullet, s^n, \pi^n)$，
存在同伦等价 $C(\alpha)^\bullet \to C^\bullet$，使图
$$
\xymatrix{
A^\bullet \ar[r] \ar[d] & B^\bullet \ar[d] \ar[r] &
C(\alpha)^\bullet \ar[r]_{-p} \ar[d] & A^\bullet[1] \ar[d] \\
A^\bullet \ar[r] & B^\bullet \ar[r] &
C^\bullet \ar[r]^\delta & A^\bullet[1]
}
$$
在 $K(\mathcal{A})$ 中定义一个三角同构。
\item 给定复形态射 $f : K^\bullet \to L^\bullet$，
存在三角同构
$$
\xymatrix{
K^\bullet \ar[r] \ar[d] & \tilde L^\bullet \ar[d] \ar[r] &
M^\bullet \ar[r]_{\delta} \ar[d] & K^\bullet[1] \ar[d] \\
K^\bullet \ar[r] & L^\bullet \ar[r] &
C(f)^\bullet \ar[r]^{-p} & K^\bullet[1]
}
$$
其中上方三角是与某个逐项分裂正合序列
$K^\bullet \to \tilde L^\bullet \to M^\bullet$ 相关联的三角。
\end{enumerate}
\end{lemma}

\begin{proof}
证明 (1)。我们有 $C(\alpha)^n = B^n \oplus A^{n + 1}$，
并将 $C(\alpha)^n \to C^n$ 简单定义为先投影到 $B^n$，
再应用 $\beta^n$。这定义了复形态射，因为复合
$A^{n + 1} \to B^{n + 1} \to C^{n + 1}$ 都为零。
为得到同伦逆，我们取 $C^\bullet \to C(\alpha)^\bullet$，
它在次数 $n$ 上由 $(s^n , -\delta^n)$ 给出。
这是复形态射，因为态射 $\delta^n$
可表征为使下式成立的唯一态射 $C^n \to A^{n + 1}$：
$d \circ s^n - s^{n + 1} \circ d = \alpha \circ \delta^n$，
见同调，引理 \ref{homology-lemma-ses-termwise-split-cochain} 的证明。
复合 $C^\bullet \to C(\alpha)^\bullet \to C^\bullet$ 是恒等态射。
复合 $C(\alpha)^\bullet \to C^\bullet \to C(\alpha)^\bullet$
等于态射
$$
\left(
\begin{matrix}
s^n \circ \beta^n & 0 \\
-\delta^n \circ \beta^n & 0
\end{matrix}
\right)
$$
为了说明它与恒等映射同伦，使用由矩阵
$$
\left(
\begin{matrix}
0 & 0 \\
\pi^n & 0
\end{matrix}
\right) : C(\alpha)^n = B^n \oplus A^{n + 1} \to
B^{n - 1} \oplus A^n = C(\alpha)^{n - 1}
$$
给出的同伦 $h^n : C(\alpha)^n \to C(\alpha)^{n - 1}$。
直接验证即得
$$
\left(
\begin{matrix}
1 & 0 \\
0 & 1
\end{matrix}
\right)
-
\left(
\begin{matrix}
s^n \\
-\delta^n
\end{matrix}
\right)
\left(
\begin{matrix}
\beta^n & 0
\end{matrix}
\right)
=
\left(
\begin{matrix}
d & \alpha^n \\
0 & -d
\end{matrix}
\right)
\left(
\begin{matrix}
0 & 0 \\
\pi^n & 0
\end{matrix}
\right)
+
\left(
\begin{matrix}
0 & 0 \\
\pi^{n + 1} & 0
\end{matrix}
\right)
\left(
\begin{matrix}
d & \alpha^{n + 1} \\
0 & -d
\end{matrix}
\right)
$$
为完成 (1) 的证明，还需说明态射
$-p : C(\alpha)^\bullet \to A^\bullet[1]$（见定义
\ref{derived-definition-cone}）与 $C(\alpha)^\bullet \to C^\bullet \to A^\bullet[1]$
在同伦意义下相等。这从上述讨论是清楚的。事实上，
可以使用同伦逆 $(s, -\delta) : C^\bullet \to C(\alpha)^\bullet$，
转而验证两个映射 $C^\bullet \to A^\bullet[1]$ 相等。又注意到
$p \circ (s, -\delta) = -\delta$，正如所需。

\medskip\noindent
证明 (2)。设 $\tilde f : K^\bullet \to \tilde L^\bullet$、
$s : L^\bullet \to \tilde L^\bullet$ 和
$\pi : \tilde L^\bullet \to L^\bullet$ 如引理
\ref{derived-lemma-make-injective} 中所定义。由引理
\ref{derived-lemma-functorial-cone} 和 \ref{derived-lemma-third-isomorphism}，三角
$(K^\bullet, L^\bullet, C(f), i, p)$ 和
$(K^\bullet, \tilde L^\bullet, C(\tilde f), \tilde i, \tilde p)$
同构。注意，三角同构可以复合。因此可以用
$\tilde L^\bullet$ 替换 $L^\bullet$，用 $\tilde f$ 替换 $f$。
换言之，可假设 $f$ 是逐项分裂单射。
在此情形下，结论由 (1) 得到。
\end{proof}


\begin{lemma}
\label{derived-lemma-sequence-maps-split}
设 $\mathcal{A}$ 为加性范畴。设
$A_1^\bullet \to A_2^\bullet \to \ldots \to A_n^\bullet$
为一列可复合的复形态射。存在交换图
$$
\xymatrix{
A_1^\bullet \ar[r] &
A_2^\bullet \ar[r] &
\ldots \ar[r] &
A_n^\bullet \\
B_1^\bullet \ar[r] \ar[u] &
B_2^\bullet \ar[r] \ar[u] &
\ldots \ar[r] &
B_n^\bullet \ar[u]
}
$$
使得每个态射 $B_i^\bullet \to B_{i + 1}^\bullet$
都是逐项分裂单射，而每个 $B_i^\bullet \to A_i^\bullet$
都是同伦等价。此外，若所有 $A_i^\bullet$ 都属于
$K^{+}(\mathcal{A})$、$K^{-}(\mathcal{A})$ 或 $K^b(\mathcal{A})$，
则所有 $B_i^\bullet$ 也属于同一范畴。
\end{lemma}

\begin{proof}
$n = 1$ 的情形没有实质内容。引理
\ref{derived-lemma-make-injective} 是 $n = 2$ 的情形。
假设除 $B_n^\bullet$ 之外的图已经构造出来。
将引理 \ref{derived-lemma-make-injective} 应用于复合
$B_{n - 1}^\bullet \to A_{n - 1}^\bullet \to A_n^\bullet$。
所得结果是所需的分解
$B_{n - 1}^\bullet \to B_n^\bullet \to A_n^\bullet$。
\end{proof}

\begin{lemma}
\label{derived-lemma-rotate-triangle}
设 $\mathcal{A}$ 为加性范畴。设
$(\alpha : A^\bullet \to B^\bullet, \beta : B^\bullet \to C^\bullet, s^n,
\pi^n)$ 为复形的逐项分裂序列。设
$(A^\bullet, B^\bullet, C^\bullet, \alpha, \beta, \delta)$
为相关联的三角。那么三角
$(C^\bullet[-1], A^\bullet, B^\bullet, \delta[-1], \alpha, \beta)$
同构于三角
$(C^\bullet[-1], A^\bullet, C(\delta[-1])^\bullet, \delta[-1], i, p)$。
\end{lemma}

\begin{proof}
我们写作 $B^n = A^n \oplus C^n$，并将 $\alpha^n$ 和 $\beta^n$
分别识别为自然嵌入与投影。由 $\delta$ 的构造，有
$$
d_B^n =
\left(
\begin{matrix}
d_A^n & \delta^n \\
0 & d_C^n
\end{matrix}
\right)
$$
另一方面，$\delta[-1] : C^\bullet[-1] \to A^\bullet$ 的锥满足
$C(\delta[-1])^n = A^n \oplus C^n$，其微分正是上述矩阵！
因此引理得证。
\end{proof}

\begin{lemma}
\label{derived-lemma-rotate-cone}
设 $\mathcal{A}$ 为加性范畴。设
$f : K^\bullet \to L^\bullet$ 为复形态射。三角
$(L^\bullet, C(f)^\bullet, K^\bullet[1], i, p, f[1])$ 是与
由 $f$ 的锥的定义得到的逐项分裂序列
$$
0 \to L^\bullet \to C(f)^\bullet \to K^\bullet[1] \to 0
$$
相关联的三角。
\end{lemma}

\begin{proof}
立即由定义得到。
\end{proof}






\section{同伦范畴中的特异三角}
\label{derived-section-homotopy-triangulated}

\noindent
由于我们希望上同调长正合序列中的边缘映射
由蛇引理中不带符号的映射给出，我们如下定义同伦范畴中的特异三角。

\begin{definition}
\label{derived-definition-distinguished-triangle}
设 $\mathcal{A}$ 为加性范畴。$K(\mathcal{A})$ 的三角
$(X, Y, Z, f, g, h)$ 称为 $K(\mathcal{A})$ 的{\it 特异三角}，
若它同构于与某个复形的逐项分裂正合序列相关联的三角；
见定义 \ref{derived-definition-split-ses}。对
$K^{+}(\mathcal{A})$、$K^{-}(\mathcal{A})$ 和 $K^b(\mathcal{A})$
作相同定义。
\end{definition}

\noindent
注意，根据引理 \ref{derived-lemma-the-same-up-to-isomorphisms}，
形如 $(K^\bullet, L^\bullet, C(f)^\bullet, f, i, -p)$ 的三角
是特异三角。这确实给出一个三角范畴，见命题
\ref{derived-proposition-homotopy-category-triangulated}。在证明该命题之前，
还需要一个引理才能证明 TR4。

\begin{lemma}
\label{derived-lemma-two-split-injections}
设 $\mathcal{A}$ 为加性范畴。假设
$\alpha : A^\bullet \to B^\bullet$ 和 $\beta : B^\bullet \to C^\bullet$
是复形的分裂单射。那么存在特异三角
$(A^\bullet, B^\bullet, Q_1^\bullet, \alpha, p_1, d_1)$、
$(A^\bullet, C^\bullet, Q_2^\bullet, \beta \circ \alpha, p_2, d_2)$
和
$(B^\bullet, C^\bullet, Q_3^\bullet, \beta, p_3, d_3)$，
对它们 TR4 成立。
\end{lemma}

\begin{proof}
设 $\pi_1^n : B^n \to A^n$ 和 $\pi_3^n : C^n \to B^n$ 是分裂。
那么 $A^\bullet \to C^\bullet$ 也是分裂单射，其分裂为
$\pi_2^n = \pi_1^n \circ \pi_3^n$。以 $Q_1^\bullet$、$Q_2^\bullet$ 和
$Q_3^\bullet$ 表示“商”复形。换言之，
$Q_1^n = \Ker(\pi_1^n)$、$Q_3^n = \Ker(\pi_3^n)$ 且
$Q_2^n = \Ker(\pi_2^n)$。注意这些核存在。于是
$B^n = A^n \oplus Q_1^n$ 且 $C^n = B^n \oplus Q_3^n$，
其中我们把 $A^n$ 看作 $B^n$ 的子对象，依此类推。这意味着
$C^n = A^n \oplus Q_1^n \oplus Q_3^n$。注意
$\pi_2^n = \pi_1^n \circ \pi_3^n$ 在 $Q_1^n$ 和 $Q_3^n$ 上均为零。因而
$Q_2^n = Q_1^n \oplus Q_3^n$。考虑交换图
$$
\begin{matrix}
0 & \to & A^\bullet & \to & B^\bullet & \to & Q_1^\bullet & \to & 0 \\
  &     & \downarrow &     & \downarrow &     & \downarrow  & \\
0 & \to & A^\bullet & \to & C^\bullet & \to & Q_2^\bullet & \to & 0 \\
  &     & \downarrow &     & \downarrow &     & \downarrow  & \\
0 & \to & B^\bullet & \to & C^\bullet & \to & Q_3^\bullet & \to & 0
\end{matrix}
$$
此图各行都是逐项分裂正合序列，因而按定义决定特异三角。
此外，上图中向下的箭头与所选分裂相容，故定义三角态射
$$
(A^\bullet \to B^\bullet \to Q_1^\bullet \to A^\bullet[1])
\longrightarrow
(A^\bullet \to C^\bullet \to Q_2^\bullet \to A^\bullet[1])
$$
和
$$
(A^\bullet \to C^\bullet \to Q_2^\bullet \to A^\bullet[1])
\longrightarrow
(B^\bullet \to C^\bullet \to Q_3^\bullet \to B^\bullet[1]).
$$
注意，上图底行分裂序列的分裂 $Q_3^n \to C^n$
与 $C^n \to Q_2^n$ 复合后，给出分裂序列
$0 \to Q_1^\bullet \to Q_2^\bullet \to Q_3^\bullet \to 0$ 的一个分裂。
由此容易看出，相应特异三角
$$
(Q_1^\bullet \to Q_2^\bullet \to Q_3^\bullet \to Q_1^\bullet[1])
$$
中的态射 $\delta : Q_3^\bullet \to Q_1^\bullet[1]$
等于复合 $Q_3^\bullet \to B^\bullet[1] \to Q_1^\bullet[1]$。
因而我们得到公理 TR4 的结论中所规定的结构。
\end{proof}

\begin{proposition}
\label{derived-proposition-homotopy-category-triangulated}
设 $\mathcal{A}$ 为加性范畴。复形的模同伦范畴
$K(\mathcal{A})$，连同它的自然平移函子和上面定义的特异三角，
是一个三角范畴。
\end{proposition}

\begin{proof}
证明 TR1。依定义，任意同构于特异三角的三角都是特异的。
另外，任意三角 $(A^\bullet, A^\bullet, 0, 1, 0, 0)$ 都是特异的，因为
$0 \to A^\bullet \to A^\bullet \to 0 \to 0$ 是复形的逐项分裂序列。
最后，给定任意复形态射 $f : K^\bullet \to L^\bullet$，三角
$(K, L, C(f), f, i, -p)$ 由引理
\ref{derived-lemma-the-same-up-to-isomorphisms} 是特异的。

\medskip\noindent
证明 TR2。设 $(X, Y, Z, f, g, h)$ 为三角。假设
$(Y, Z, X[1], g, h, -f[1])$ 是特异的。那么存在复形的逐项分裂序列
$A^\bullet \to B^\bullet \to C^\bullet$，使得相关联的三角
$(A^\bullet, B^\bullet, C^\bullet, \alpha, \beta, \delta)$
同构于 $(Y, Z, X[1], g, h, -f[1])$。反向旋转后看到
$(X, Y, Z, f, g, h)$ 同构于
$(C^\bullet[-1], A^\bullet, B^\bullet, -\delta[-1], \alpha, \beta)$。
由引理 \ref{derived-lemma-rotate-triangle}，三角
$(C^\bullet[-1], A^\bullet, B^\bullet, \delta[-1], \alpha, \beta)$
同构于
$(C^\bullet[-1], A^\bullet, C(\delta[-1])^\bullet, \delta[-1], i, p)$。
将前一个三角同构与 $Y$ 上的 $-1$ 预复合，得到
$(X, Y, Z, f, g, h)$ 同构于
$(C^\bullet[-1], A^\bullet, C(\delta[-1])^\bullet, \delta[-1], i, -p)$。
因此它由引理 \ref{derived-lemma-the-same-up-to-isomorphisms} 是特异的。
反过来，假设 $(X, Y, Z, f, g, h)$ 是特异的。
由引理 \ref{derived-lemma-the-same-up-to-isomorphisms}，这意味着它同构于形如
$(K^\bullet, L^\bullet, C(f), f, i, -p)$ 的三角，其中 $f$
是某个复形态射。那么旋转后的三角
$(Y, Z, X[1], g, h, -f[1])$ 同构于
$(L^\bullet, C(f), K^\bullet[1], i, -p, -f[1])$，后者又同构于三角
$(L^\bullet, C(f), K^\bullet[1], i, p, f[1])$。
由引理 \ref{derived-lemma-rotate-cone}，这个三角是特异的。
因此 $(Y, Z, X[1], g, h, -f[1])$ 如所需是特异的。

\medskip\noindent
证明 TR3。设
$(X, Y, Z, f, g, h)$ 和 $(X', Y', Z', f', g', h')$
是 $K(\mathcal{A})$ 的特异三角，并设
$a : X \to X'$ 和 $b : Y \to Y'$ 为满足
$f' \circ a = b \circ f$ 的态射。由引理
\ref{derived-lemma-the-same-up-to-isomorphisms}，可假设
$(X, Y, Z, f, g, h) = (X, Y, C(f), f, i, -p)$ 且
$(X', Y', Z', f', g', h') = (X', Y', C(f'), f', i', -p')$。
此时，只需将引理 \ref{derived-lemma-functorial-cone}
应用于由 $f, f', a, b$ 给出的交换图。

\medskip\noindent
证明 TR4。目前我们已知 $K(\mathcal{A})$ 是预三角范畴。
因此可使用引理 \ref{derived-lemma-easier-axiom-four}。设
$A^\bullet \to B^\bullet$ 和 $B^\bullet \to C^\bullet$
是 $K(\mathcal{A})$ 中可复合的态射。由引理
\ref{derived-lemma-sequence-maps-split}，可假设
$A^\bullet \to B^\bullet$ 和 $B^\bullet \to C^\bullet$
是分裂单射。在此情形下，结论由引理
\ref{derived-lemma-two-split-injections} 得到。
\end{proof}

\begin{remark}
\label{derived-remark-boundedness-conditions-triangulated}
设 $\mathcal{A}$ 为加性范畴。命题
\ref{derived-proposition-homotopy-category-triangulated} 的证明原封不动地表明，
范畴 $K^{+}(\mathcal{A})$、$K^{-}(\mathcal{A})$ 和 $K^b(\mathcal{A})$
都是三角范畴。实际上，有界（上有界、下有界）复形之间的态射的锥
仍然有界（上有界、下有界）。不过下面我们将证明它们是
$K(\mathcal{A})$ 的三角子范畴，这给出另一个证明。
\end{remark}

\begin{lemma}
\label{derived-lemma-bounded-triangulated-subcategories}
设 $\mathcal{A}$ 为加性范畴。范畴
$K^{+}(\mathcal{A})$、$K^{-}(\mathcal{A})$ 和 $K^b(\mathcal{A})$
都是 $K(\mathcal{A})$ 的满三角子范畴。
\end{lemma}

\begin{proof}
所提及的每个范畴都是满加性子范畴。我们使用引理
\ref{derived-lemma-triangulated-subcategory} 的判别准则说明它们是三角子范畴。
显然，每个范畴
$K^{+}(\mathcal{A})$、$K^{-}(\mathcal{A})$ 和 $K^b(\mathcal{A})$
都在平移函子 $[1], [-1]$ 下保持。最后，假设
$f : A^\bullet \to B^\bullet$ 是
$K^{+}(\mathcal{A})$、$K^{-}(\mathcal{A})$ 或 $K^b(\mathcal{A})$ 中的态射。
那么 $(A^\bullet, B^\bullet, C(f)^\bullet, f, i, -p)$
是 $K(\mathcal{A})$ 的特异三角，且
$C(f)^\bullet \in K^{+}(\mathcal{A})$、
$K^{-}(\mathcal{A})$ 或 $K^b(\mathcal{A})$，这从锥的构造是明显的。
因此引理得证。（或者，
% P02-E0026: frozen source says "an termwise split injection".
$K^\bullet \to L^\bullet$ 在
$K^{+}(\mathcal{A})$、$K^{-}(\mathcal{A})$ 或 $K^b(\mathcal{A})$ 中
同构于一个复形的逐项分裂单射，见引理
\ref{derived-lemma-make-injective}，然后可直接取相关联的特异三角。）
\end{proof}

\begin{lemma}
\label{derived-lemma-additive-exact-homotopy-category}
设 $\mathcal{A}$、$\mathcal{B}$ 为加性范畴，
$F : \mathcal{A} \to \mathcal{B}$ 为加性函子。所诱导的函子
$$
\begin{matrix}
F : K(\mathcal{A}) \longrightarrow K(\mathcal{B}) \\
F : K^{+}(\mathcal{A}) \longrightarrow K^{+}(\mathcal{B}) \\
F : K^{-}(\mathcal{A}) \longrightarrow K^{-}(\mathcal{B}) \\
F : K^b(\mathcal{A}) \longrightarrow K^b(\mathcal{B})
\end{matrix}
$$
都是三角范畴的正合函子。
\end{lemma}

\begin{proof}
假设 $A^\bullet \to B^\bullet \to C^\bullet$
是 $\mathcal{A}$ 的复形的逐项分裂序列，其分裂为
$(s^n, \pi^n)$，相关联的态射为
$\delta : C^\bullet \to A^\bullet[1]$；见定义
\ref{derived-definition-split-ses}。那么
$F(A^\bullet) \to F(B^\bullet) \to F(C^\bullet)$
是复形的逐项分裂序列，其分裂为
$(F(s^n), F(\pi^n))$，相关联的态射为
$F(\delta) : F(C^\bullet) \to F(A^\bullet)[1]$。
因而 $F$ 将特异三角变换为特异三角。
\end{proof}

\begin{lemma}
\label{derived-lemma-improve-distinguished-triangle-homotopy}
设 $\mathcal{A}$ 为加性范畴，并设
$(A^\bullet, B^\bullet, C^\bullet, a, b, c)$ 为 $K(\mathcal{A})$
中的特异三角。那么存在一个同构的特异三角
$(A^\bullet, (B')^\bullet, C^\bullet, a', b', c)$，使对所有 $n$，
$0 \to A^n \to (B')^n \to C^n \to 0$ 都是分裂短正合序列。
\end{lemma}

\begin{proof}
我们将使用命题 \ref{derived-proposition-homotopy-category-triangulated}
所证的 $K(\mathcal{A})$ 是三角范畴这一事实。
设 $W^\bullet$ 是 $c : C^\bullet \to A^\bullet[1]$ 的锥，
其映射为 $i : A^\bullet[1] \to W^\bullet$ 和
$p : W^\bullet \to C^\bullet[1]$。那么由引理
\ref{derived-lemma-the-same-up-to-isomorphisms}，
$(C^\bullet, A^\bullet[1], W^\bullet, c, i, -p)$ 是特异三角。
向后旋转两次，看到
$(A^\bullet, W^\bullet[-1], C^\bullet, -i[-1], p[-1], c)$
是特异三角。由 TR3，存在特异三角的态射
$(\text{id}, \beta, \text{id}) : (A^\bullet, B^\bullet, C^\bullet, a, b, c) \to
(A^\bullet, W^\bullet[-1], C^\bullet, -i[-1], p[-1], c)$，
由引理 \ref{derived-lemma-third-isomorphism-triangle}，它必为同构。
这就完成了证明，因为根据第 \ref{derived-section-cones} 节中锥的构造，
$0 \to A^\bullet \to W^\bullet[-1] \to C^\bullet \to 0$
是复形的逐项分裂短正合序列。
\end{proof}

\begin{remark}
\label{derived-remark-homotopy-double}
设 $\mathcal{A}$ 为具有可数直和的加性范畴。以
$\text{DoubleComp}(\mathcal{A})$ 表示 $\mathcal{A}$ 中的双复形范畴，
见同调，第 \ref{homology-section-double-complexes} 节。
我们可用这个范畴构造两个三角范畴。
\begin{enumerate}
\item 我们可将 $\text{DoubleComp}(\mathcal{A})$ 的对象
$A^{\bullet, \bullet}$ 如下看作复形之复形：
$$
\ldots \to A^{\bullet, -1} \to A^{\bullet, 0} \to A^{\bullet, 1} \to \ldots
$$
并取同伦范畴 $K_{first}(\text{DoubleComp}(\mathcal{A}))$，
其相应的三角结构由命题
\ref{derived-proposition-homotopy-category-triangulated} 给出。
由同调，备注
\ref{homology-remark-double-complex-complex-of-complexes-first}，函子
$$
\text{Tot} :
K_{first}(\text{DoubleComp}(\mathcal{A}))
\longrightarrow
K(\mathcal{A})
$$
是三角范畴的正合函子。
\item 我们可将 $\text{DoubleComp}(\mathcal{A})$ 的对象
$A^{\bullet, \bullet}$ 如下看作复形之复形：
$$
\ldots \to A^{-1, \bullet} \to A^{0, \bullet} \to A^{1, \bullet} \to \ldots
$$
并取同伦范畴 $K_{second}(\text{DoubleComp}(\mathcal{A}))$，
其相应的三角结构由命题
\ref{derived-proposition-homotopy-category-triangulated} 给出。
由同调，备注
\ref{homology-remark-double-complex-complex-of-complexes-second}，函子
$$
\text{Tot} :
K_{second}(\text{DoubleComp}(\mathcal{A}))
\longrightarrow
K(\mathcal{A})
$$
是三角范畴的正合函子。
\end{enumerate}
\end{remark}

\begin{remark}
\label{derived-remark-double-complex-as-tensor-product-of}
设 $\mathcal{A}$、$\mathcal{B}$、$\mathcal{C}$ 为加性范畴，
并假设 $\mathcal{C}$ 具有可数直和。假设
$$
\otimes : \mathcal{A} \times \mathcal{B} \longrightarrow \mathcal{C},
\quad
(X, Y) \longmapsto X \otimes Y
$$
是一个在态射上双线性的函子。这决定一个函子
$$
\text{Comp}(\mathcal{A}) \times \text{Comp}(\mathcal{B})
\longrightarrow
\text{DoubleComp}(\mathcal{C}), \quad
(X^\bullet, Y^\bullet)
\longmapsto
X^\bullet \otimes Y^\bullet
$$
见同调，例
\ref{homology-example-double-complex-as-tensor-product-of}。
\begin{enumerate}
\item 对于 $\text{Comp}(\mathcal{A})$ 的一个固定对象 $X^\bullet$，函子
$$
K(\mathcal{B}) \longrightarrow K(\mathcal{C}), \quad
Y^\bullet \longmapsto \text{Tot}(X^\bullet \otimes Y^\bullet)
$$
是三角范畴的正合函子。
\item 对于 $\text{Comp}(\mathcal{B})$ 的一个固定对象 $Y^\bullet$，函子
$$
K(\mathcal{A}) \longrightarrow K(\mathcal{C}), \quad
X^\bullet \longmapsto \text{Tot}(X^\bullet \otimes Y^\bullet)
$$
是三角范畴的正合函子。
\end{enumerate}
这由备注 \ref{derived-remark-homotopy-double} 得到，因为函子
$\text{Comp}(\mathcal{A}) \to \text{DoubleComp}(\mathcal{C})$、
$Y^\bullet \mapsto X^\bullet \otimes Y^\bullet$ 以及
$\text{Comp}(\mathcal{B}) \to \text{DoubleComp}(\mathcal{C})$、
$X^\bullet \mapsto X^\bullet \otimes Y^\bullet$
显然与同伦和逐项分裂短正合序列相容，因而诱导三角范畴的正合函子
$$
K(\mathcal{B}) \to K_{first}(\text{DoubleComp}(\mathcal{C}))
\quad\text{且}\quad
K(\mathcal{A}) \to K_{second}(\text{DoubleComp}(\mathcal{C}))
$$
注意，对于这两个函子中的第一个，同构
$$
\text{Tot}(X^\bullet \otimes Y^\bullet[1]) \cong
\text{Tot}(X^\bullet \otimes Y^\bullet)[1]
$$
涉及符号（这可追溯到同调，备注
\ref{homology-remark-shift-double-complex} 中所选的符号）。
\end{remark}






\section{导出范畴}
\label{derived-section-derived-categories}

\noindent
本节通过在 $K(\mathcal{A})$ 中取拟同构的逆，构造阿贝尔范畴
$\mathcal{A}$ 的导出范畴。为此先回忆，函子
$H^i : \text{Comp}(\mathcal{A}) \to \mathcal{A}$
经由 $K(\mathcal{A})$ 分解，见同调，引理
\ref{homology-lemma-map-cohomology-homotopy-cochain}。此外，在同调，定义
\ref{homology-definition-cohomology-shift} 中，我们已定义识别
$H^i(K^\bullet[n]) = H^{i + n}(K^\bullet)$。此时，为了避免混淆和可能的符号错误，
将其重新定义为
$$
H^i(K^\bullet) = H^0(K^\bullet[i])
$$
是合理的。

\begin{lemma}
\label{derived-lemma-cohomology-homological}
设 $\mathcal{A}$ 为阿贝尔范畴。函子
$$
H^0 : K(\mathcal{A}) \longrightarrow \mathcal{A}
$$
是同调函子。
\end{lemma}

\begin{proof}
由于 $H^0$ 是函子，按我们对特异三角的定义，只需证明：
对给定的复形的逐项分裂短正合序列
$0 \to A^\bullet \to B^\bullet \to C^\bullet \to 0$，序列
$H^0(A^\bullet) \to H^0(B^\bullet) \to H^0(C^\bullet)$
正合。这由同调，引理
\ref{homology-lemma-long-exact-sequence-cochain} 得到。
\end{proof}

\noindent
特别地，该引理意味着，$K(\mathcal{A})$ 中的特异三角
$(X, Y, Z, f, g, h)$ 诱导上同调长正合序列
\begin{equation}
\label{derived-equation-long-exact-cohomology-sequence-D}
\xymatrix{
\ldots \ar[r] &
H^i(X) \ar[r]^{H^i(f)} &
H^i(Y) \ar[r]^{H^i(g)} &
H^i(Z) \ar[r]^{H^i(h)} &
H^{i + 1}(X) \ar[r] & \ldots
}
\end{equation}
见 (\ref{derived-equation-long-exact-cohomology-sequence})。此外，它与复形短正合序列所相关联的
上同调长正合序列相容
% P02-E0027: frozen source contains the editorial placeholder "insert future reference here".
（此处插入未来的引用）。例如，若
$(A^\bullet, B^\bullet, C^\bullet, \alpha, \beta, \delta)$
是与复形的逐项分裂正合序列相关联的特异三角（见定义
\ref{derived-definition-split-ses}），则上述上同调序列与用蛇引理定义的序列一致；
见同调，引理 \ref{homology-lemma-long-exact-sequence-cochain}，
关于序列的一致性，见同调，引理
\ref{homology-lemma-ses-termwise-split-long-cochain}。

\medskip\noindent
回忆，若对所有 $i \in \mathbf{Z}$ 都有 $H^i(K^\bullet) = 0$，
则称复形 $K^\bullet$ {\it 无环}。又回忆，复形态射
$f : K^\bullet \to L^\bullet$ 是{\it 拟同构}，当且仅当
$H^i(f)$ 对所有 $i$ 都是同构。见同调，定义
\ref{homology-definition-quasi-isomorphism-cochain}。

\begin{lemma}
\label{derived-lemma-acyclic}
设 $\mathcal{A}$ 为阿贝尔范畴。$K(\mathcal{A})$ 中由无环复形组成的满子范畴
$\text{Ac}(\mathcal{A})$ 是 $K(\mathcal{A})$ 的严格满饱和三角子范畴。
$K(\mathcal{A})$ 中相应的饱和乘法系（见引理
\ref{derived-lemma-operations}）是拟同构集 $\text{Qis}(\mathcal{A})$。
特别地，局部化函子
$Q : K(\mathcal{A}) \to \text{Qis}(\mathcal{A})^{-1}K(\mathcal{A})$
的核是 $\text{Ac}(\mathcal{A})$，且函子 $H^0$ 经由 $Q$ 分解。
\end{lemma}

\begin{proof}
由引理 \ref{derived-lemma-cohomology-homological}，我们知道 $H^0$ 是同调函子。
因而本引理是引理 \ref{derived-lemma-acyclic-general} 的特殊情形。
\end{proof}

\begin{definition}
\label{derived-definition-unbounded-derived-category}
设 $\mathcal{A}$ 为阿贝尔范畴。设
$\text{Ac}(\mathcal{A})$ 和 $\text{Qis}(\mathcal{A})$ 如引理
\ref{derived-lemma-acyclic} 中所定义。$\mathcal{A}$ 的{\it 导出范畴}
是三角范畴
$$
D(\mathcal{A}) =
K(\mathcal{A})/\text{Ac}(\mathcal{A}) =
\text{Qis}(\mathcal{A})^{-1} K(\mathcal{A}).
$$
以 $H^0 : D(\mathcal{A}) \to \mathcal{A}$ 表示唯一的函子，
它与商函子的复合恢复为上面定义的函子 $H^0$。使用引理
\ref{derived-lemma-homological-functor-bounded}，我们引入严格满饱和三角子范畴
$D^{+}(\mathcal{A}), D^{-}(\mathcal{A}), D^b(\mathcal{A})$，它们的对象集为
$$
\begin{matrix}
\Ob(D^{+}(\mathcal{A})) =
\{X \in \Ob(D(\mathcal{A})) \mid
H^n(X) = 0\text{ 对所有 }n \ll 0\} \\
\Ob(D^{-}(\mathcal{A})) =
\{X \in \Ob(D(\mathcal{A})) \mid
H^n(X) = 0\text{ 对所有 }n \gg 0\} \\
\Ob(D^b(\mathcal{A})) =
\{X \in \Ob(D(\mathcal{A})) \mid
H^n(X) = 0\text{ 对所有 }|n| \gg 0\}
\end{matrix}
$$
范畴 $D^b(\mathcal{A})$ 称为 $\mathcal{A}$ 的{\it 有界导出范畴}。
\end{definition}

\noindent
若 $K^\bullet$ 和 $L^\bullet$ 是 $\mathcal{A}$ 的复形，那么我们有时说
“$K^\bullet$ 与 $L^\bullet$ {\it 拟同构}”，意指
$K^\bullet$ 和 $L^\bullet$ 是 $D(\mathcal{A})$ 中的同构对象。

\begin{remark}
\label{derived-remark-existence-derived}
本章一贯处理“小”阿贝尔范畴（这是 Stacks 项目中的约定）。
对于“大”阿贝尔范畴 $\mathcal{A}$，导出范畴
$D(\mathcal{A})$ 是否存在并不清楚，因为导出范畴中的态射是否构成集并不清楚。
实际上，它们一般不构成集，见例，引理
\ref{examples-lemma-big-abelian-category}。不过，若 $\mathcal{A}$ 是 Grothendieck 阿贝尔范畴，
且 $K^\bullet, L^\bullet$ 属于 $K(\mathcal{A})$，那么由内射，定理
\ref{injectives-theorem-K-injective-embedding-grothendieck}，存在到 K-内射复形
$I^\bullet$ 的拟同构 $L^\bullet \to I^\bullet$，而引理
\ref{derived-lemma-K-injective} 表明
$$
\Hom_{D(\mathcal{A})}(K^\bullet, L^\bullet) =
\Hom_{K(\mathcal{A})}(K^\bullet, I^\bullet)
$$
是一个集。Grothendieck 阿贝尔范畴的一些例子包括环上模的范畴，
或更一般地，带环址点上的模层范畴。
\end{remark}

\noindent
每个变体 $D^{+}(\mathcal{A}), D^{-}(\mathcal{A}), D^b(\mathcal{A})$
都可构造为相应同伦范畴的局部化。这依赖于下面的简单引理。

\begin{lemma}
\label{derived-lemma-complex-cohomology-bounded}
设 $\mathcal{A}$ 为阿贝尔范畴，$K^\bullet$ 为复形。
\begin{enumerate}
\item 若对所有 $n \ll 0$ 都有 $H^n(K^\bullet) = 0$，
则存在拟同构 $K^\bullet \to L^\bullet$，其中 $L^\bullet$ 下有界。
\item 若对所有 $n \gg 0$ 都有 $H^n(K^\bullet) = 0$，
则存在拟同构 $M^\bullet \to K^\bullet$，其中 $M^\bullet$ 上有界。
\item 若对所有 $|n| \gg 0$ 都有 $H^n(K^\bullet) = 0$，
则存在复形态射的交换图
$$
\xymatrix{
K^\bullet \ar[r] & L^\bullet \\
M^\bullet \ar[u] \ar[r] & N^\bullet \ar[u]
}
$$
其中所有箭头都是拟同构，$L^\bullet$ 下有界，
$M^\bullet$ 上有界，而 $N^\bullet$ 是有界复形。
\end{enumerate}
\end{lemma}

\begin{proof}
选取 $a \ll 0 \ll b$，并令
$M^\bullet = \tau_{\leq b}K^\bullet$、
$L^\bullet = \tau_{\geq a}K^\bullet$ 以及
$N^\bullet = \tau_{\leq b}L^\bullet = \tau_{\geq a}M^\bullet$。
关于截断函子，见同调，第
\ref{homology-section-truncations} 节。
\end{proof}

\noindent
为了陈述下一引理，以
$\text{Ac}^{+}(\mathcal{A})$、$\text{Ac}^{-}(\mathcal{A})$、
相应地 $\text{Ac}^b(\mathcal{A})$ 表示
$K^{+}(\mathcal{A})$、$K^{-}(\mathcal{A})$、相应地 $K^b(\mathcal{A})$
与 $\text{Ac}(\mathcal{A})$ 的交。以
$\text{Qis}^{+}(\mathcal{A})$、$\text{Qis}^{-}(\mathcal{A})$、
相应地 $\text{Qis}^b(\mathcal{A})$ 表示
$K^{+}(\mathcal{A})$、$K^{-}(\mathcal{A})$、相应地 $K^b(\mathcal{A})$
与 $\text{Qis}(\mathcal{A})$ 的交。

\begin{lemma}
\label{derived-lemma-bounded-derived}
设 $\mathcal{A}$ 为阿贝尔范畴。子范畴
$\text{Ac}^{+}(\mathcal{A})$、$\text{Ac}^{-}(\mathcal{A})$、
相应地 $\text{Ac}^b(\mathcal{A})$
分别是 $K^{+}(\mathcal{A})$、$K^{-}(\mathcal{A})$、
相应地 $K^b(\mathcal{A})$ 的严格满饱和三角子范畴。
相应的饱和乘法系（见引理 \ref{derived-lemma-operations}）
分别是集 $\text{Qis}^{+}(\mathcal{A})$、$\text{Qis}^{-}(\mathcal{A})$、
相应地 $\text{Qis}^b(\mathcal{A})$。
\begin{enumerate}
\item 函子 $K^{+}(\mathcal{A}) \to D^{+}(\mathcal{A})$ 的核是
$\text{Ac}^{+}(\mathcal{A})$，并诱导三角范畴等价
$$
K^{+}(\mathcal{A})/\text{Ac}^{+}(\mathcal{A}) =
\text{Qis}^{+}(\mathcal{A})^{-1}K^{+}(\mathcal{A})
\longrightarrow
D^{+}(\mathcal{A})
$$
\item 函子 $K^{-}(\mathcal{A}) \to D^{-}(\mathcal{A})$ 的核是
$\text{Ac}^{-}(\mathcal{A})$，并诱导三角范畴等价
$$
K^{-}(\mathcal{A})/\text{Ac}^{-}(\mathcal{A}) =
\text{Qis}^{-}(\mathcal{A})^{-1}K^{-}(\mathcal{A})
\longrightarrow
D^{-}(\mathcal{A})
$$
\item 函子 $K^b(\mathcal{A}) \to D^b(\mathcal{A})$ 的核是
$\text{Ac}^b(\mathcal{A})$，并诱导三角范畴等价
$$
K^b(\mathcal{A})/\text{Ac}^b(\mathcal{A}) =
\text{Qis}^b(\mathcal{A})^{-1}K^b(\mathcal{A})
\longrightarrow
D^b(\mathcal{A})
$$
\end{enumerate}
\end{lemma}

\begin{proof}
最初的各项陈述由引理 \ref{derived-lemma-acyclic-general}
应用于同调函子 $H^0$ 的限制得到。
(1)、(2)、(3) 中关于核的陈述是下述事实的结果：
每一种情形中的定义。由引理
\ref{derived-lemma-complex-cohomology-bounded}，每个函子都本质满。
为完成证明，还需说明这些函子全忠实。
我们先对下有界的版本证明这一点。

\medskip\noindent
假设 $K^\bullet, L^\bullet$ 是下有界复形。它们在
$D(\mathcal{A})$ 中的态射具有形式 $s^{-1}f$，其中的一对态射为
$f : K^\bullet \to (L')^\bullet$、$s : L^\bullet \to (L')^\bullet$，
且 $s$ 是拟同构。这意味着 $(L')^\bullet$ 的上同调下有界。
因而由引理 \ref{derived-lemma-complex-cohomology-bounded}，可选取拟同构
$s' : (L')^\bullet \to (L'')^\bullet$，使 $(L'')^\bullet$ 下有界。
那么这一对 $(s' \circ f, s' \circ s)$ 在
$\text{Qis}^{+}(\mathcal{A})^{-1}K^{+}(\mathcal{A})$ 中定义一个态射。
因此该函子是“满”的。最后，假设一对
$f : K^\bullet \to (L')^\bullet$、$s : L^\bullet \to (L')^\bullet$
在 $\text{Qis}^{+}(\mathcal{A})^{-1}K^{+}(\mathcal{A})$ 中定义一个态射，
而该态射在 $D(\mathcal{A})$ 中为零。这意味着存在拟同构
$s' : (L')^\bullet \to (L'')^\bullet$，使 $s' \circ f = 0$。
再次使用引理 \ref{derived-lemma-complex-cohomology-bounded}，得到拟同构
$s'' : (L'')^\bullet \to (L''')^\bullet$，其中 $(L''')^\bullet$ 下有界。
于是 $s'' \circ s' \circ f = 0$，这意味着 $s^{-1}f$ 在
$\text{Qis}^{+}(\mathcal{A})^{-1}K^{+}(\mathcal{A})$ 中为零。
这就证明了 (1) 中的函子是等价。

\medskip\noindent
(2) 的证明对偶于 (1) 的证明。为证明 (3)，可使用 (2) 的结果。
因此只需证明函子
$\text{Qis}^b(\mathcal{A})^{-1}K^b(\mathcal{A})
\to \text{Qis}^{-}(\mathcal{A})^{-1}K^{-}(\mathcal{A})$
全忠实。前一段的论证可直接用于证明此事，
只需全程使用已经上有界的复形。
\end{proof}







\section{典范 \zhdelta-函子}
\label{derived-section-canonical-delta-functor}

\noindent
导出范畴应是万有上同调函子的容纳对象。为陈述该结果，
我们使用从阿贝尔范畴到三角范畴的 $\delta$-函子概念，
见定义 \ref{derived-definition-delta-functor}。

\medskip\noindent
考虑函子 $\text{Comp}(\mathcal{A}) \to K(\mathcal{A})$。
一般而言，这个函子{\bf 不}是 $\delta$-函子。
看出这一点的最简单方法，是考虑 $\mathcal{A}$ 的对象的一个不分裂短正合序列
$0 \to A \to B \to C \to 0$。由于
$\Hom_{K(\mathcal{A})}(C[0], A[1]) = 0$，
我们看到，由这个短正合序列产生的任何特异三角必形如
$(A[0], B[0], C[0], a, b, 0)$。但 $K(\mathcal{A})$ 中存在这样的特异三角
意味着扩张分裂，矛盾。

\medskip\noindent
事实上，函子 $\text{Comp}(\mathcal{A}) \to D(\mathcal{A})$
是 $\delta$-函子。为看出这一点，必须定义与阿贝尔范畴
$\mathcal{A}$ 中复形的短正合序列
$$
0 \to A^\bullet \xrightarrow{a} B^\bullet \xrightarrow{b} C^\bullet \to 0
$$
相关联的态射 $\delta$。考虑态射 $a$ 的锥 $C(a)^\bullet$。
我们有 $C(a)^n = B^n \oplus A^{n + 1}$，并将
$q^n : C(a)^n \to C^n$ 定义为先投影到 $B^n$，再应用 $b^n$。
由此得到复形态射
$$
q : C(a)^\bullet \longrightarrow C^\bullet.
$$
显然 $q \circ i = b$，其中 $i$ 如定义 \ref{derived-definition-cone} 中所定义。
注意，由于 $a^\bullet$ 在每个次数上都是单射，$q$ 的核可识别为
$\text{id}_{A^\bullet}$ 的锥，而该锥无环。因而 $q$ 是拟同构。
根据引理 \ref{derived-lemma-the-same-up-to-isomorphisms}，三角
$$
(A, B, C(a), a, i, -p)
$$
是 $K(\mathcal{A})$ 中的特异三角。由于局部化函子
$K(\mathcal{A}) \to D(\mathcal{A})$ 正合，我们看到
$(A, B, C(a), a, i, -p)$ 是 $D(\mathcal{A})$ 中的特异三角。
由于 $q$ 是拟同构，$q$ 在 $D(\mathcal{A})$ 中是同构。
因而我们推得
$$
(A, B, C, a, b, -p \circ q^{-1})
$$
是 $D(\mathcal{A})$ 的特异三角。这引出下面的引理。

\begin{lemma}
\label{derived-lemma-derived-canonical-delta-functor}
设 $\mathcal{A}$ 为阿贝尔范畴。所定义的函子
$\text{Comp}(\mathcal{A}) \to D(\mathcal{A})$
具有一个自然的 $\delta$-函子结构，其中
$$
\delta_{A^\bullet \to B^\bullet \to C^\bullet} = - p \circ q^{-1}
$$
$p$ 和 $q$ 如上所述。同一构造将函子
$\text{Comp}^{+}(\mathcal{A}) \to D^{+}(\mathcal{A})$、
$\text{Comp}^{-}(\mathcal{A}) \to D^{-}(\mathcal{A})$ 和
$\text{Comp}^b(\mathcal{A}) \to D^b(\mathcal{A})$
都变为 $\delta$-函子。
\end{lemma}

\begin{proof}
我们已经看到，每当给定复形的短正合序列时，这个选择都产生特异三角。
还需证明，给定交换图
$$
\xymatrix{
0 \ar[r] &
A^\bullet \ar[r]_a \ar[d]_f &
B^\bullet \ar[r]_b \ar[d]_g &
C^\bullet \ar[r] \ar[d]_h &
0 \\
0 \ar[r] &
(A')^\bullet \ar[r]^{a'} &
(B')^\bullet \ar[r]^{b'} &
(C')^\bullet \ar[r] &
0
}
$$
后，我们得到定义 \ref{derived-definition-delta-functor} (2) 中所需的交换图。
由引理 \ref{derived-lemma-functorial-cone}，态射对 $(f, g)$ 诱导典范态射
$c : C(a)^\bullet \to C(a')^\bullet$。简单计算表明
$q' \circ c = h \circ q$ 且 $f[1] \circ p = p' \circ c$。
结论由此直接得到。
\end{proof}

\begin{lemma}
\label{derived-lemma-derived-compare-triangles-ses}
设 $\mathcal{A}$ 为阿贝尔范畴。设
$$
\xymatrix{
0 \ar[r] &
A^\bullet \ar[r] \ar[d] &
B^\bullet \ar[r] \ar[d] &
C^\bullet \ar[r] \ar[d] &
0 \\
0 \ar[r] &
D^\bullet \ar[r] &
E^\bullet \ar[r] &
F^\bullet \ar[r] &
0
}
$$
是一个复形态射的交换图，其各行是复形短正合序列，
而各竖直箭头都是拟同构。上述引理
\ref{derived-lemma-derived-canonical-delta-functor} 的 $\delta$-函子将短正合序列
$0 \to A^\bullet \to B^\bullet \to C^\bullet \to 0$ 和
$0 \to D^\bullet \to E^\bullet \to F^\bullet \to 0$
映到同构的特异三角。
\end{lemma}

\begin{proof}
这直接由下述事实得到：
$K(\mathcal{A}) \to D(\mathcal{A})$ 将拟同构变成同构，
且所相关联的特异三角具有函子性。
\end{proof}

\begin{lemma}
\label{derived-lemma-derived-compare-triangles-split-case}
设 $\mathcal{A}$ 为阿贝尔范畴。设
$$
\xymatrix{
0 \ar[r] &
A^\bullet \ar[r] &
B^\bullet \ar[r] &
C^\bullet \ar[r] &
0
}
$$
% P02-E0028: frozen source says "a short exact sequences".
为复形短正合序列。假设该短正合序列逐项分裂。设
$(A^\bullet, B^\bullet, C^\bullet, \alpha, \beta, \delta)$
为 $K(\mathcal{A})$ 中与该序列相关联的特异三角。上述引理
\ref{derived-lemma-derived-canonical-delta-functor} 的 $\delta$-函子
将短正合序列 $0 \to A^\bullet \to B^\bullet \to C^\bullet \to 0$
映到一个同构于特异三角
$$
(A^\bullet, B^\bullet, C^\bullet, \alpha, \beta, \delta).
$$
的三角。
\end{lemma}

\begin{proof}
由引理 \ref{derived-lemma-the-same-up-to-isomorphisms} 得到。
\end{proof}

\begin{remark}
\label{derived-remark-truncation-distinguished-triangle}
设 $\mathcal{A}$ 为阿贝尔范畴，$K^\bullet$ 为 $\mathcal{A}$ 的复形，
$a \in \mathbf{Z}$。我们声称，$D(\mathcal{A})$ 中存在典范特异三角
$$
\tau_{\leq a}K^\bullet \to K^\bullet \to \tau_{\geq a + 1}K^\bullet \to
(\tau_{\leq a}K^\bullet)[1]
$$
这里使用了同调，第 \ref{homology-section-truncations} 节中的典范截断函子 $\tau$。
事实上，我们先取与复形短正合序列
$$
0 \to \tau_{\leq a}K^\bullet \to K^\bullet \to
K^\bullet/\tau_{\leq a}K^\bullet \to 0
$$
由我们的 $\delta$-函子（引理
\ref{derived-lemma-derived-canonical-delta-functor}）相关联的特异三角。
然后，由同调，第 \ref{homology-section-truncations} 节中上同调群的描述，
我们使用映射 $K^\bullet \to \tau_{\geq a + 1}K^\bullet$
经由拟同构
$K^\bullet/\tau_{\leq a}K^\bullet \to \tau_{\geq a + 1}K^\bullet$ 分解这一事实。
以类似方式得到典范特异三角
$$
\tau_{\leq a}K^\bullet \to \tau_{\leq a + 1}K^\bullet \to
H^{a + 1}(K^\bullet)[-a-1] \to (\tau_{\leq a}K^\bullet)[1]
$$
和
$$
H^a(K^\bullet)[-a] \to \tau_{\geq a}K^\bullet \to \tau_{\geq a + 1}K^\bullet
\to H^a(K^\bullet)[-a + 1]
$$
\end{remark}

\begin{lemma}
\label{derived-lemma-trick-vanishing-composition}
设 $\mathcal{A}$ 为阿贝尔范畴。设
$$
K_0^\bullet \to K_1^\bullet \to \ldots \to K_n^\bullet
$$
为复形映射，且满足
\begin{enumerate}
\item 对 $i > 0$，$H^i(K_0^\bullet) = 0$；
\item $H^{-j}(K_j^\bullet) \to H^{-j}(K_{j + 1}^\bullet)$ 为零。
\end{enumerate}
那么在 $D(\mathcal{A})$ 中，复合 $K_0^\bullet \to K_n^\bullet$
经由 $\tau_{\leq -n}K_n^\bullet \to K_n^\bullet$ 分解。对偶地，
给定复形映射
$$
K_n^\bullet \to K_{n - 1}^\bullet \to \ldots \to K_0^\bullet
$$
且满足
\begin{enumerate}
\item 对 $i < 0$，$H^i(K_0^\bullet) = 0$；
\item $H^j(K_{j + 1}^\bullet) \to H^j(K_j^\bullet)$ 为零，
\end{enumerate}
则在 $D(\mathcal{A})$ 中，复合 $K_n^\bullet \to K_0^\bullet$
经由 $K_n^\bullet \to \tau_{\geq n}K_n^\bullet$ 分解。
\end{lemma}

\begin{proof}
考虑 $n = 1$ 的情形。由于在 $D(\mathcal{A})$ 中
$\tau_{\leq 0}K_0^\bullet = K_0^\bullet$，可以用
$\tau_{\leq 0}K_0^\bullet$ 替换 $K_0^\bullet$，用
$\tau_{\leq 0}K_1^\bullet$ 替换 $K_1^\bullet$。考虑特异三角
$$
\tau_{\leq -1}K_1^\bullet \to K_1^\bullet \to
H^0(K_1^\bullet)[0] \to (\tau_{\leq -1}K_1^\bullet)[1]
$$
（备注 \ref{derived-remark-truncation-distinguished-triangle}）。复合
$K_0^\bullet \to K_1^\bullet \to H^0(K_1^\bullet)[0]$ 为零，
因为它等于
$K_0^\bullet \to H^0(K_0^\bullet)[0] \to H^0(K_1^\bullet)[0]$，
而后者由假设为零。函子
$\Hom_{D(\mathcal{A})}(K_0^\bullet, -)$ 是同调函子（引理
\ref{derived-lemma-representable-homological}）这一事实允许我们找到所需的分解。
对 $n = 2$，由 $n = 1$ 的情形得到分解
$K_0^\bullet \to \tau_{\leq -1}K_1^\bullet$，然后将 $n = 1$ 的情形
应用于复形映射
$\tau_{\leq -1}K_1^\bullet \to \tau_{\leq -1}K_2^\bullet$，
便得到分解
$\tau_{\leq -1}K_1^\bullet \to \tau_{\leq -2}K_2^\bullet$。
一般情形完全以相同方式证明。
\end{proof}




\section{过滤导出范畴}
\label{derived-section-filtered-derived-category}

\noindent
本节的一个参考文献是 \cite[I, Chapter V]{cotangent}。设
$\mathcal{A}$ 为阿贝尔范畴。本节将定义 $\mathcal{A}$ 的过滤导出范畴
$DF(\mathcal{A})$。简言之，我们将它定义为 $\mathcal{A}$ 中带有有限过滤的对象所构成的
正合范畴的导出范畴。（因此，我们的构造是更一般的正合范畴之导出范畴
构造的特殊情形；例见 \cite{Buhler}、\cite{Keller}。）
Illusie 的过滤导出范畴是我们范畴中过滤有限的那些对象构成的满子范畴。
（在我们的范畴中，过滤在每个次数上仍有限，但可能不一致有界。）
我们这样选择的理由是：此构造并不更难，并且可将讨论应用于引理
\ref{derived-lemma-two-ss-complex-functor} 的谱序列，又见备注
\ref{derived-remark-functorial-ss}。

\medskip\noindent
我们将使用同调，第 \ref{homology-section-filtrations} 节中引入的过滤对象记号。
$\mathcal{A}$ 的过滤对象范畴记为 $\text{Fil}(\mathcal{A})$。
约定所有过滤都是递减的。

\begin{definition}
\label{derived-definition-finite-filtered}
设 $\mathcal{A}$ 为阿贝尔范畴。
$\mathcal{A}$ 的{\it 有限过滤对象范畴}
是 $\mathcal{A}$ 中过滤 $F$ 有限的过滤对象 $(A, F)$ 所构成的范畴。
我们将其记为 $\text{Fil}^f(\mathcal{A})$。
\end{definition}

\noindent
因而 $\text{Fil}^f(\mathcal{A})$ 是 $\text{Fil}(\mathcal{A})$ 的满子范畴。
对每个 $p \in \mathbf{Z}$，有函子
$\text{gr}^p : \text{Fil}^f(\mathcal{A}) \to \mathcal{A}$。
还有函子
$$
\text{gr} = \bigoplus\nolimits_{p \in \mathbf{Z}} \text{gr}^p :
\text{Fil}^f(\mathcal{A}) \to \text{Gr}(\mathcal{A})
$$
其中 $\text{Gr}(\mathcal{A})$ 是 $\mathcal{A}$ 的分次对象范畴，
见同调，定义 \ref{homology-definition-graded}。
最后，有函子
$$
(\text{遗忘 }F) : \text{Fil}^f(\mathcal{A}) \longrightarrow \mathcal{A}
$$
它将过滤对象 $(A, F)$ 映到其在 $\mathcal{A}$ 中的底对象。
范畴 $\text{Fil}^f(\mathcal{A})$ 是加性范畴，但一般不是阿贝尔范畴，
见同调，例 \ref{homology-example-not-abelian}。

\medskip\noindent
由于函子 $\text{gr}^p$、$\text{gr}$、$(\text{遗忘 }F)$ 都是加性函子，
由引理 \ref{derived-lemma-additive-exact-homotopy-category}，它们诱导三角范畴的正合函子
$$
\text{gr}^p, (\text{遗忘 }F) :
K(\text{Fil}^f(\mathcal{A}))
\to
K(\mathcal{A})
\quad\text{且}\quad
\text{gr} :
K(\text{Fil}^f(\mathcal{A}))
\to
K(\text{Gr}(\mathcal{A}))
$$
类比阿贝尔范畴的同伦范畴情形，我们作出下列定义。

\begin{definition}
\label{derived-definition-filtered-acyclic}
设 $\mathcal{A}$ 为阿贝尔范畴。
\begin{enumerate}
\item 设 $\alpha : K^\bullet \to L^\bullet$ 为
$K(\text{Fil}^f(\mathcal{A}))$ 的态射。若态射
$\text{gr}(\alpha)$ 是拟同构，则称 $\alpha$ 是{\it 过滤拟同构}。
\item 设 $K^\bullet$ 为 $K(\text{Fil}^f(\mathcal{A}))$ 的对象。
若复形 $\text{gr}(K^\bullet)$ 无环，则称 $K^\bullet$ {\it 过滤无环}。
\end{enumerate}
\end{definition}

\noindent
注意，$\alpha : K^\bullet \to L^\bullet$ 是过滤拟同构，当且仅当每个
$\text{gr}^p(\alpha)$ 都是拟同构。类似地，复形 $K^\bullet$ 过滤无环，
当且仅当每个 $\text{gr}^p(K^\bullet)$ 都无环。

\begin{lemma}
\label{derived-lemma-filtered-cohomology-homological}
设 $\mathcal{A}$ 为阿贝尔范畴。
\begin{enumerate}
\item 函子
$K(\text{Fil}^f(\mathcal{A})) \longrightarrow \text{Gr}(\mathcal{A})$、
$K^\bullet \longmapsto H^0(\text{gr}(K^\bullet))$
是同调函子。
\item 函子
$K(\text{Fil}^f(\mathcal{A})) \rightarrow \mathcal{A}$、
$K^\bullet \longmapsto H^0(\text{gr}^p(K^\bullet))$
是同调函子。
\item 函子
$K(\text{Fil}^f(\mathcal{A})) \longrightarrow \mathcal{A}$、
$K^\bullet \longmapsto H^0((\text{遗忘 }F)K^\bullet)$
是同调函子。
\end{enumerate}
\end{lemma}

\begin{proof}
这由下述事实得到：$H^0 : K(\mathcal{A}) \to \mathcal{A}$
是同调函子（见引理 \ref{derived-lemma-cohomology-homological}），而函子
$\text{gr}, \text{gr}^p, (\text{遗忘 }F)$ 是三角范畴的正合函子。
见引理 \ref{derived-lemma-exact-compose-homological-functor}。
\end{proof}

\begin{lemma}
\label{derived-lemma-filtered-acyclic}
设 $\mathcal{A}$ 为阿贝尔范畴。
$K(\text{Fil}^f(\mathcal{A}))$ 中由过滤无环复形组成的满子范畴
$\text{FAc}(\mathcal{A})$ 是 $K(\text{Fil}^f(\mathcal{A}))$
的严格满饱和三角子范畴。
$K(\text{Fil}^f(\mathcal{A}))$ 中相应的饱和乘法系（见引理
\ref{derived-lemma-operations}）是过滤拟同构集 $\text{FQis}(\mathcal{A})$。
特别地，局部化函子
$$
Q :
K(\text{Fil}^f(\mathcal{A}))
\longrightarrow
\text{FQis}(\mathcal{A})^{-1}K(\text{Fil}^f(\mathcal{A}))
$$
的核是 $\text{FAc}(\mathcal{A})$，且函子 $H^0 \circ \text{gr}$
经由 $Q$ 分解。
\end{lemma}

\begin{proof}
由引理 \ref{derived-lemma-filtered-cohomology-homological}，我们知道
$H^0 \circ \text{gr}$ 是同调函子。因而本引理是引理
\ref{derived-lemma-acyclic-general} 的特殊情形。
\end{proof}

\begin{definition}
\label{derived-definition-filtered-derived}
设 $\mathcal{A}$ 为阿贝尔范畴。设
$\text{FAc}(\mathcal{A})$ 和 $\text{FQis}(\mathcal{A})$ 如引理
\ref{derived-lemma-filtered-acyclic} 中所定义。
$\mathcal{A}$ 的{\it 过滤导出范畴}是三角范畴
$$
DF(\mathcal{A}) =
K(\text{Fil}^f(\mathcal{A}))/\text{FAc}(\mathcal{A}) =
\text{FQis}(\mathcal{A})^{-1} K(\text{Fil}^f(\mathcal{A})).
$$
\end{definition}

\begin{lemma}
\label{derived-lemma-filtered-derived-functors}
函子 $\text{gr}^p, \text{gr}, (\text{遗忘 }F)$ 诱导典范正合函子
$$
\text{gr}^p, (\text{遗忘 }F):
DF(\mathcal{A})
\longrightarrow
D(\mathcal{A})
$$
和
$$
\text{gr}:
DF(\mathcal{A})
\longrightarrow
D(\text{Gr}(\mathcal{A}))
$$
它们与局部化函子交换。
\end{lemma}

\begin{proof}
只要能证明每个函子
$\text{gr}^p, \text{gr}, (\text{遗忘 }F)$
都将过滤拟同构变为拟同构，这就由局部化的万有性质（见引理
\ref{derived-lemma-universal-property-localization}）得到。对前两个函子，这依定义成立。
对最后一个，需要做一点工作。设
$f : K^\bullet \to L^\bullet$ 是 $K(\text{Fil}^f(\mathcal{A}))$ 中的过滤拟同构。
选取包含 $f$ 的特异三角
$(K^\bullet, L^\bullet, M^\bullet, f, g, h)$。那么由引理
\ref{derived-lemma-filtered-acyclic}，$M^\bullet$ 过滤无环。
因此，由 $K(\mathcal{A})$ 的相应引理，只需证明过滤无环复形
在遗忘过滤后是无环复形。这由同调，引理
\ref{homology-lemma-filtered-acyclic} 得到。
\end{proof}

\begin{definition}
\label{derived-definition-filtered-derived-bounded}
设 $\mathcal{A}$ 为阿贝尔范畴。{\it 有界过滤导出范畴}
$DF^b(\mathcal{A})$ 是 $DF(\mathcal{A})$ 的满子范畴，其对象是满足
$\text{gr}(X) \in D^b(\mathcal{A})$ 的那些 $X$。
对下有界过滤导出范畴 $DF^{+}(\mathcal{A})$
和上有界过滤导出范畴 $DF^{-}(\mathcal{A})$ 作类似定义。
\end{definition}

\begin{lemma}
\label{derived-lemma-filtered-complex-cohomology-bounded}
设 $\mathcal{A}$ 为阿贝尔范畴，并设
$K^\bullet \in K(\text{Fil}^f(\mathcal{A}))$。
\begin{enumerate}
\item 若对所有 $n < a$ 都有 $H^n(\text{gr}(K^\bullet)) = 0$，
则存在过滤拟同构 $K^\bullet \to L^\bullet$，使得对所有 $n < a$
都有 $L^n = 0$。
\item 若对所有 $n > b$ 都有 $H^n(\text{gr}(K^\bullet)) = 0$，
则存在过滤拟同构 $M^\bullet \to K^\bullet$，使得对所有 $n > b$
都有 $M^n = 0$。
\item 若对所有 $|n| \gg 0$ 都有 $H^n(\text{gr}(K^\bullet)) = 0$，
则存在复形态射的交换图
$$
\xymatrix{
K^\bullet \ar[r] & L^\bullet \\
M^\bullet \ar[u] \ar[r] & N^\bullet \ar[u]
}
$$
其中所有箭头都是过滤拟同构，$L^\bullet$ 下有界，
$M^\bullet$ 上有界，而 $N^\bullet$ 是有界复形。
\end{enumerate}
\end{lemma}

\begin{proof}
假设对所有 $n < a$ 都有 $H^n(\text{gr}(K^\bullet)) = 0$。由同调，引理
\ref{homology-lemma-filtered-acyclic}，序列
$$
K^{a - 2} \xrightarrow{d^{a - 2}} K^{a - 1} \xrightarrow{d^{a - 1}} K^a
$$
是 $\mathcal{A}$ 的对象的正合序列，且态射
$d^{a - 2}$ 和 $d^{a - 1}$ 都是严格的。因而在
$\text{Fil}^f(\mathcal{A})$ 中有
$\Coim(d^{a - 1}) = \Im(d^{a - 1})$，而映射
$\text{gr}(\Im(d^{a - 1})) \to \text{gr}(K^a)$
是单射，其像等于映射
$\text{gr}(K^{a - 1}) \to \text{gr}(K^a)$ 的像，见同调，引理
\ref{homology-lemma-characterize-strict}。这意味着到截断
$$
\tau_{\geq a}K^\bullet =
(\ldots \to 0 \to K^a/\Im(d^{a - 1}) \to K^{a + 1} \to \ldots)
$$
的映射 $K^\bullet \to \tau_{\geq a}K^\bullet$ 是过滤拟同构。
这证明了 (1)。(2) 的证明对偶于 (1) 的证明。
(3) 形式地由 (1) 和 (2) 得到。
\end{proof}

\noindent
为陈述下一引理，以
$\text{FAc}^{+}(\mathcal{A})$、$\text{FAc}^{-}(\mathcal{A})$、
相应地 $\text{FAc}^b(\mathcal{A})$ 表示
$K^{+}(\text{Fil}^f\mathcal{A})$、$K^{-}(\text{Fil}^f\mathcal{A})$、
相应地 $K^b(\text{Fil}^f\mathcal{A})$ 与 $\text{FAc}(\mathcal{A})$ 的交。
以 $\text{FQis}^{+}(\mathcal{A})$、$\text{FQis}^{-}(\mathcal{A})$、
相应地 $\text{FQis}^b(\mathcal{A})$ 表示
$K^{+}(\text{Fil}^f\mathcal{A})$、$K^{-}(\text{Fil}^f\mathcal{A})$、
相应地 $K^b(\text{Fil}^f\mathcal{A})$ 与 $\text{FQis}(\mathcal{A})$ 的交。

\begin{lemma}
\label{derived-lemma-filtered-bounded-derived}
设 $\mathcal{A}$ 为阿贝尔范畴。子范畴
$\text{FAc}^{+}(\mathcal{A})$、$\text{FAc}^{-}(\mathcal{A})$、
相应地 $\text{FAc}^b(\mathcal{A})$
分别是 $K^{+}(\text{Fil}^f\mathcal{A})$、
$K^{-}(\text{Fil}^f\mathcal{A})$、相应地
$K^b(\text{Fil}^f\mathcal{A})$ 的严格满饱和三角子范畴。
相应的饱和乘法系（见引理 \ref{derived-lemma-operations}）
分别是集 $\text{FQis}^{+}(\mathcal{A})$、
$\text{FQis}^{-}(\mathcal{A})$、相应地 $\text{FQis}^b(\mathcal{A})$。
\begin{enumerate}
\item 函子
$K^{+}(\text{Fil}^f\mathcal{A}) \to DF^{+}(\mathcal{A})$
的核是 $\text{FAc}^{+}(\mathcal{A})$，并诱导三角范畴等价
$$
K^{+}(\text{Fil}^f\mathcal{A})/\text{FAc}^{+}(\mathcal{A}) =
\text{FQis}^{+}(\mathcal{A})^{-1}K^{+}(\text{Fil}^f\mathcal{A})
\longrightarrow
DF^{+}(\mathcal{A})
$$
\item 函子
$K^{-}(\text{Fil}^f\mathcal{A}) \to DF^{-}(\mathcal{A})$
的核是 $\text{FAc}^{-}(\mathcal{A})$，并诱导三角范畴等价
$$
K^{-}(\text{Fil}^f\mathcal{A})/\text{FAc}^{-}(\mathcal{A}) =
\text{FQis}^{-}(\mathcal{A})^{-1}K^{-}(\text{Fil}^f\mathcal{A})
\longrightarrow
DF^{-}(\mathcal{A})
$$
\item 函子
$K^b(\text{Fil}^f\mathcal{A}) \to DF^b(\mathcal{A})$
的核是 $\text{FAc}^b(\mathcal{A})$，并诱导三角范畴等价
$$
K^b(\text{Fil}^f\mathcal{A})/\text{FAc}^b(\mathcal{A}) =
\text{FQis}^b(\mathcal{A})^{-1}K^b(\text{Fil}^f\mathcal{A})
\longrightarrow
DF^b(\mathcal{A})
$$
\end{enumerate}
\end{lemma}

\begin{proof}
这由上述结果，特别是引理
\ref{derived-lemma-filtered-complex-cohomology-bounded}，通过与引理
\ref{derived-lemma-bounded-derived} 的证明中完全相同的论证得到。
\end{proof}











\section{一般导出函子}
\label{derived-section-derived-functors}

\noindent
本节的一个参考文献是 \cite{SGA4} 中 Deligne 的 expos\'e XVII。
存在一种非常一般的右导出函子和左导出函子概念：
当我们有三角范畴之间的正合函子、源范畴中的乘法系，
并希望找到该正合函子到局部化范畴的“正确”延拓时，便可使用它。

\begin{situation}
\label{derived-situation-derived-functor}
此处 $F : \mathcal{D} \to \mathcal{D}'$ 是三角范畴之间的正合函子，
$S$ 是 $\mathcal{D}$ 中的饱和乘法系，且与
$\mathcal{D}$ 上的三角范畴结构相容。
\end{situation}

\noindent
设 $X \in \Ob(\mathcal{D})$。回忆范畴，备注
\ref{categories-remark-left-localization-morphisms-colimit} 中源为 $X$ 的
$S$ 中箭头 $s : X \to X'$ 所组成的过滤范畴 $X/S$。
对偶地，在范畴，备注
\ref{categories-remark-right-localization-morphisms-colimit} 中，我们定义了目标为 $X$ 的
$S$ 中箭头 $s : X' \to X$ 所组成的余过滤范畴 $S/X$。

\begin{definition}
\label{derived-definition-right-derived-functor-defined}
假设和记号如情形 \ref{derived-situation-derived-functor} 中所述。设
$X \in \Ob(\mathcal{D})$。
\begin{enumerate}
\item 若 ind-对象
$$
(X/S) \longrightarrow \mathcal{D}', \quad
(s : X \to X') \longmapsto F(X')
$$
本质常值\footnote{关于范畴的 ind-对象或 pro-对象何时本质常值的讨论，
参见范畴，第 \ref{categories-section-essentially-constant} 节。}，
则称{\it 右导出函子 $RF$ 在} $X$ {\it 处有定义}；在此情形下，其值
$\mathcal{D}'$ 中的 $Y$ 称为{\it $RF$ 在 $X$ 处的值}。
\item 若 pro-对象
$$
(S/X) \longrightarrow \mathcal{D}', \quad
(s: X' \to X) \longmapsto F(X')
$$
本质常值，则称{\it 左导出函子 $LF$ 在} $X$
{\it 处有定义}；在此情形下，其值 $\mathcal{D}'$ 中的 $Y$
称为{\it $LF$ 在 $X$ 处的值}。
\end{enumerate}
我们常滥用记号，将这些值简记为 $RF(X)$ 或 $LF(X)$。
\end{definition}

\noindent
事实上，$\mathcal{D}$ 中由 $RF$ 有定义的对象组成的满子范畴
是三角子范畴，而 $RF$ 将在此子范畴上定义一个函子，
它将 $S$ 中的态射变为同构。

\begin{lemma}
\label{derived-lemma-derived-functor}
假设和记号如情形 \ref{derived-situation-derived-functor} 中所述。设
$f : X \to Y$ 为 $\mathcal{D}$ 的态射。
\begin{enumerate}
\item 若 $RF$ 在 $X$ 和 $Y$ 处都有定义，则值之间存在唯一态射
$RF(f) : RF(X) \to RF(Y)$，使对任意交换图
$$
\xymatrix{
X \ar[d]_f \ar[r]_s & X' \ar[d]^{f'} \\
Y \ar[r]^{s'} & Y'
}
$$
若 $s, s' \in S$，则图
$$
\xymatrix{
F(X) \ar[d] \ar[r] & F(X') \ar[d] \ar[r] & RF(X) \ar[d] \\
F(Y) \ar[r] & F(Y') \ar[r] & RF(Y)
}
$$
交换。
\item 若 $LF$ 在 $X$ 和 $Y$ 处都有定义，则值之间存在唯一态射
$LF(f) : LF(X) \to LF(Y)$，使对任意交换图
$$
\xymatrix{
X' \ar[d]_{f'} \ar[r]_s & X \ar[d]^f \\
Y' \ar[r]^{s'} & Y
}
$$
若 $s, s'$ 属于 $S$，则图
$$
\xymatrix{
LF(X) \ar[d] \ar[r] & F(X') \ar[d] \ar[r] & F(X) \ar[d] \\
LF(Y) \ar[r] & F(Y') \ar[r] & F(Y)
}
$$
交换。
\end{enumerate}
\end{lemma}

\begin{proof}
只要假设余极限
$$
RF(X) = \colim_{s : X \to X'} F(X')
\quad\text{且}\quad
RF(Y) = \colim_{s' : Y \to Y'} F(Y')
$$
存在，(1) 就成立。事实上，给出余极限之间的态射
$RF(X) \to RF(Y)$，等价于对 $\Ob(X/S)$ 中的每个
$s : X \to X'$ 给出一个态射 $F(X') \to RF(Y)$，
并要求它们与范畴 $X/S$ 中的态射相容。
为得到该态射，按 MS2 选取交换图
$$
\xymatrix{
X \ar[d]_f \ar[r]_s & X' \ar[d]^{f'} \\
Y \ar[r]^{s'} & Y'
}
$$
其中 $s, s'$ 属于 $S$，并将 $F(X') \to RF(Y)$ 定义为复合
$F(X') \to F(Y') \to RF(Y)$。使用 MS3 可看出，这与上述图的选择无关。
细节从略。(2) 的证明与之对偶。
\end{proof}

\begin{lemma}
\label{derived-lemma-derived-inverts}
假设和记号如情形 \ref{derived-situation-derived-functor} 中所述。设
$s : X \to Y$ 是 $S$ 的元素。
\begin{enumerate}
\item $RF$ 在 $X$ 处有定义，当且仅当它在 $Y$ 处有定义。
在此情形下，值之间的映射 $RF(s) : RF(X) \to RF(Y)$ 是同构。
\item $LF$ 在 $X$ 处有定义，当且仅当它在 $Y$ 处有定义。
在此情形下，值之间的映射 $LF(s) : LF(X) \to LF(Y)$ 是同构。
\end{enumerate}
\end{lemma}

\begin{proof}
从略。
\end{proof}

\begin{lemma}
\label{derived-lemma-derived-shift}
\begin{slogan}
导出函子与平移相容
\end{slogan}
假设和记号如情形 \ref{derived-situation-derived-functor} 中所述。
设 $X$ 为 $\mathcal{D}$ 的对象，$n \in \mathbf{Z}$。
\begin{enumerate}
\item $RF$ 在 $X$ 处有定义，当且仅当它在 $X[n]$ 处有定义。
在此情形下，值之间存在典范同构
$RF(X)[n]= RF(X[n])$。
\item $LF$ 在 $X$ 处有定义，当且仅当它在 $X[n]$ 处有定义。
在此情形下，值之间存在典范同构
$LF(X)[n] \to LF(X[n])$。
\end{enumerate}
\end{lemma}

\begin{proof}
从略。
\end{proof}

\begin{lemma}
\label{derived-lemma-2-out-of-3-defined}
假设和记号如情形 \ref{derived-situation-derived-functor} 中所述。
设 $(X, Y, Z, f, g, h)$ 为 $\mathcal{D}$ 的特异三角。
若 $RF$ 在 $X, Y, Z$ 中的任意两个处有定义，则它在第三个处也有定义。
此外，在这种情形下
$$
(RF(X), RF(Y), RF(Z), RF(f), RF(g), RF(h))
$$
是 $\mathcal{D}'$ 中的特异三角。对 $LF$ 也有类似结论。
\end{lemma}

\begin{proof}
假设 $RF$ 在 $X, Y$ 处有定义，值分别为 $A, B$。
设 $RF(f) : A \to B$ 为诱导态射，见引理
\ref{derived-lemma-derived-functor}。可在 $\mathcal{D}'$ 中选取特异三角
$(A, B, C, RF(f), b, c)$。我们声称 $C$ 是 $RF$ 在 $Z$ 处的一个值。

\medskip\noindent
为证明这一点，选取 $S$ 中的 $s : X \to X'$，使存在如范畴，定义
\ref{categories-definition-essentially-constant-diagram} 中的态射
$\alpha : A \to F(X')$。根据 MS2，可选取交换图
$$
\xymatrix{
X \ar[d]_f \ar[r]_s & X' \ar[d]^{f'} \\
Y \ar[r]^{s'} & Y'
}
$$
其中 $s' \in S$。利用 $Y/S$ 是过滤的，我们可以（把 $s'$ 替换为
$S$ 中某个 $s'' : Y \to Y''$ 之后）假设存在如范畴，定义
\ref{categories-definition-essentially-constant-diagram} 中的态射
$\beta : B \to F(Y')$。考虑图
$$
\xymatrix{
A \ar[d]_{RF(f)} \ar[r]_-\alpha &
F(X') \ar[r] \ar[d]^{F(f')} &
A \ar[d]^{RF(f)} \\
B \ar[r]^-\beta & F(Y') \ar[r] & B
}
$$
左方块未必交换，但外部方块和右方块交换。
% P02-E0029: frozen source indexes the ind-object by s':Y'->Y, opposite to its own definition of Y/S.
假设 ind-对象 $\{F(Y')\}_{s' : Y' \to Y}$ 本质常值，这意味着存在
$S$ 中的 $s'' : Y \to Y''$ 和态射 $h : Y' \to Y''$，
使 $s'' = h \circ s'$，并且 $F(h)$ 等于复合
$F(Y') \to B \to F(Y') \to F(Y'')$。因此，将 $Y'$ 替换为 $Y''$，
将 $\beta$ 替换为 $F(h) \circ \beta$ 之后，上图将交换
（这可对箭头作直接计算）。

\medskip\noindent
使用 MS6 选取三角态射
$$
(s, s', s'') : (X, Y, Z, f, g, h) \longrightarrow (X', Y', Z', f', g', h')
$$
其中 $s'' \in S$。由 TR3，选取三角态射
$$
(\alpha, \beta, \gamma) :
(A, B, C, RF(f), b, c)
\longrightarrow
(F(X'), F(Y'), F(Z'), F(f'), F(g'), F(h'))
$$

\medskip\noindent
由引理 \ref{derived-lemma-derived-inverts}，只需证明 $RF(Z')$ 有定义，
并且箭头 $\gamma : C \to F(Z')$ 诱导同构 $C \to RF(Z')$。
事实上，此时我们将得到三角同构
$$
(A, B, C, RF(f), b, c)
\longrightarrow
(RF(X'), RF(Y'), RF(Z'), RF(f'), RF(g'), RF(h'))
$$
并由 TR1 得出其目标是特异三角。考虑引理
\ref{derived-lemma-limit-triangles} 中由三角组成的范畴 $\mathcal{I}$：
$$
\mathcal{I} =
\{(t, t', t'') : (X', Y', Z', f', g', h') \to (X'', Y'', Z'', f'', g'', h'')
\mid (t, t', t'') \in S\}
$$
要说明系 $F(Z'')$ 在范畴 $Z'/S$ 上本质常值，
等价于说明 $F(Z'')$ 的系在 $\mathcal{I}$ 上本质常值，
因为 $\mathcal{I} \to Z'/S$ 是余终的；见范畴，引理
\ref{categories-lemma-cofinal-essentially-constant}（余终性在引理
\ref{derived-lemma-limit-triangles} 中证明）。对 $\mathcal{D}'$ 的任意对象 $W$，考虑图
$$
\xymatrix{
\colim_\mathcal{I} \Mor_{\mathcal{D}'}(W, F(X'')) \ar[d] &
\Mor_{\mathcal{D}'}(W, A) \ar[l] \ar[d] \\
\colim_\mathcal{I} \Mor_{\mathcal{D}'}(W, F(Y'')) \ar[d] &
\Mor_{\mathcal{D}'}(W, B) \ar[d] \ar[l] \\
\colim_\mathcal{I} \Mor_{\mathcal{D}'}(W, F(Z'')) \ar[d] &
\Mor_{\mathcal{D}'}(W, C) \ar[d] \ar[l] \\
\colim_\mathcal{I} \Mor_{\mathcal{D}'}(W, F(X''[1])) \ar[d] &
\Mor_{\mathcal{D}'}(W, A[1]) \ar[d] \ar[l] \\
\colim_\mathcal{I} \Mor_{\mathcal{D}'}(W, F(Y''[1])) &
\Mor_{\mathcal{D}'}(W, B[1]) \ar[l]
}
$$
其水平箭头由与 $(\alpha, \beta, \gamma)$ 复合给出。
由于过滤余极限正合（代数，引理
\ref{algebra-lemma-directed-colimit-exact}），左列是正合序列。
因而 $5$ 引理（同调，引理 \ref{homology-lemma-five-lemma}）表明映射
$$
\colim_\mathcal{I} \Mor_{\mathcal{D}'}(W, F(Z''))
\longrightarrow
\Mor_{\mathcal{D}'}(W, C)
$$
是双射。由范畴，引理
\ref{categories-lemma-characterize-essentially-constant-ind} 的第 (4) 部分，
我们得出 $F(Z'')$ 在 $\mathcal{I}$ 上以 $C$ 为值本质常值。
\end{proof}

\begin{lemma}
\label{derived-lemma-direct-sum-defined}
假设和记号如情形 \ref{derived-situation-derived-functor} 中所述。
设 $X, Y$ 为 $\mathcal{D}$ 的对象。
\begin{enumerate}
\item 若 $RF$ 在 $X$ 和 $Y$ 处有定义，则 $RF$ 在 $X \oplus Y$ 处有定义。
\item 若 $\mathcal{D}'$ 是 Karoubi 范畴，且 $RF$ 在 $X \oplus Y$ 处有定义，
则 $RF$ 在 $X$ 和 $Y$ 处都有定义。
\end{enumerate}
两种情形下都有 $RF(X \oplus Y) = RF(X) \oplus RF(Y)$。
对 $LF$ 也有类似结论。
\end{lemma}

\begin{proof}
若 $RF$ 在 $X$ 和 $Y$ 处有定义，则由特异三角
$X \to X \oplus Y \to Y \to X[1]$（引理 \ref{derived-lemma-split}）以及引理
\ref{derived-lemma-2-out-of-3-defined}，$RF$ 在 $X \oplus Y$ 处有定义，且有特异三角
$RF(X) \to RF(X \oplus Y) \to RF(Y) \to RF(X)[1]$。
再对其应用引理 \ref{derived-lemma-split}，得到
$RF(X \oplus Y) = RF(X) \oplus RF(Y)$。这证明了 (1) 和最后的断言。

\medskip\noindent
反过来，假设 $RF$ 在 $X \oplus Y$ 处有定义，且 $\mathcal{D}'$
是 Karoubi 范畴。由于 $S$ 是饱和系，$S$ 是在加性局部化函子
$Q : \mathcal{D} \to S^{-1}\mathcal{D}$ 下变为可逆的箭头之集，
见范畴，引理 \ref{categories-lemma-what-gets-inverted}。
因此，对 $S$ 中任意 $s : X \to X'$ 和 $s' : Y \to Y'$，态射
$s \oplus s' : X \oplus Y \to X' \oplus Y'$ 都是 $S$ 的元素。
由此得到函子
$$
X/S \times Y/S \longrightarrow (X \oplus Y)/S
$$
回忆，范畴 $X/S, Y/S, (X \oplus Y)/S$ 都是过滤范畴（范畴，备注
\ref{categories-remark-left-localization-morphisms-colimit}）。由范畴，引理
\ref{categories-lemma-essentially-constant-over-product}，$X/S \times Y/S$ 是过滤的，而
$F|_{X/S} : X/S \to \mathcal{D}'$
% P02-E0030: frozen source changes the only functor F to an undefined G on Y/S.
（相应地 $G|_{Y/S} : Y/S \to \mathcal{D}'$）本质常值，当且仅当
$F|_{X/S} \circ \text{pr}_1 : X/S \times Y/S \to \mathcal{D}'$
（相应地 $G|_{Y/S} \circ \text{pr}_2 : X/S \times Y/S \to \mathcal{D}'$）
本质常值。下面我们将证明上述函子余终，因此由范畴，引理
\ref{categories-lemma-cofinal-essentially-constant}，我们看到，
$F|_{(X \oplus Y)/S}$ 本质常值意味着
$F|_{X/S} \circ \text{pr}_1 \oplus F|_{Y/S} \circ \text{pr}_2 :
X/S \times Y/S \to \mathcal{D}'$
本质常值。由同调，引理
\ref{homology-lemma-direct-sum-from-product-essentially-constant}
（此处用到 $\mathcal{D}'$ 是 Karoubi 范畴），我们看到，
$F|_{X/S} \circ \text{pr}_1 \oplus F|_{Y/S} \circ \text{pr}_2$
本质常值意味着
$F|_{X/S} \circ \text{pr}_1$ 和 $F|_{Y/S} \circ \text{pr}_2$
都本质常值，这证明了 $RF$ 在 $X$ 和 $Y$ 处有定义。

\medskip\noindent
证明上述函子余终。为此，选取 $S$ 中任意的
$t : X \oplus Y \to Z$。使用 MS2，可找到态射
$Z \to X'$ 和 $Z \to Y'$，以及 $S$ 中的
$s : X \to X'$ 和 $s' : Y \to Y'$，使得
$$
\xymatrix{
X \ar[d]^s & X \oplus Y \ar[d] \ar[l] \ar[r] & Y \ar[d]_{s'} \\
X' & Z \ar[l] \ar[r] & Y'
}
$$
交换。这证明在 $(X \oplus Y)/S$ 中存在映射
$Z \to X' \oplus Y'$，也就得到范畴，定义
\ref{categories-definition-cofinal} 的第 (1) 部分。为证明第 (2) 部分，
只需证明：给定 $t : X \oplus Y \to Z$ 以及
$(X \oplus Y)/S$ 中的态射
$s_i \oplus s'_i : Z \to X'_i \oplus Y'_i$、$i = 1, 2$，
我们能找到 $S$ 中的态射
$a : X'_1 \to X'$、$b : X'_2 \to X'$、
$c : Y'_1 \to Y'$、$d : Y'_2 \to Y'$，使得
$a \circ s_1 = b \circ s_2$ 且 $c \circ s'_1 = d \circ s'_2$。
为此，先选取任意 $X'$ 和 $Y'$ 以及 $S$ 中的映射 $a, b, c, d$；
这是可能的，因为 $X/S$ 和 $Y/S$ 是过滤范畴。然后两个映射
$a \circ s_1, b \circ s_2 : Z \to X'$ 在 $S^{-1}\mathcal{D}$ 中变得相等。
因而可找到 $S$ 中的态射 $X' \to X''$ 使它们相等。
类似地，可找到 $S$ 中的 $Y' \to Y''$ 使
$c \circ s'_1$ 和 $d \circ s'_2$ 相等。用 $X''$ 替换 $X'$、
用 $Y''$ 替换 $Y'$，即得
$a \circ s_1 = b \circ s_2$ 且 $c \circ s'_1 = d \circ s'_2$。

\medskip\noindent
% P02-E0031: frozen source has subject-verb disagreement in "The proof ... are dual".
对 $LF$ 的相应陈述之证明是对偶的。
\end{proof}

\begin{proposition}
\label{derived-proposition-derived-functor}
假设和记号如情形 \ref{derived-situation-derived-functor} 中所述。
\begin{enumerate}
\item $\mathcal{D}$ 中由 $RF$ 有定义的对象组成的满子范畴
$\mathcal{E}$ 是 $\mathcal{D}$ 的严格满三角子范畴。
\item 我们得到三角范畴的正合函子
$RF : \mathcal{E} \longrightarrow \mathcal{D}'$。
\item $S$ 中源或目标属于 $\mathcal{E}$ 的元素都是 $\mathcal{E}$ 的态射。
\item $S_\mathcal{E} = \text{Arrows}(\mathcal{E}) \cap S$ 的任意元素都被
$RF$ 映到同构。
\item 集 $S_\mathcal{E}$ 是 $\mathcal{E}$ 中与三角结构相容的饱和乘法系。
\item 函子 $S_\mathcal{E}^{-1}\mathcal{E} \to S^{-1}\mathcal{D}$
是三角范畴的全忠实正合函子。
\item 我们得到正合函子
$$
RF : S_\mathcal{E}^{-1}\mathcal{E} \longrightarrow \mathcal{D}'.
$$
\item 若 $\mathcal{D}'$ 是 Karoubi 范畴，则 $\mathcal{E}$ 是 $\mathcal{D}$
的饱和三角子范畴。
\end{enumerate}
对 $LF$ 有类似结果。
\end{proposition}

\begin{proof}
由于 $S$ 饱和，它包含所有同构（见范畴，定义
\ref{categories-definition-saturated-multiplicative-system} 之后的备注）。
因而 (1) 由引理 \ref{derived-lemma-derived-inverts}、
\ref{derived-lemma-2-out-of-3-defined} 和 \ref{derived-lemma-derived-shift} 得到。
(2) 由引理 \ref{derived-lemma-derived-functor}、\ref{derived-lemma-derived-shift}
和 \ref{derived-lemma-2-out-of-3-defined} 得到。
(3) 由引理 \ref{derived-lemma-derived-inverts} 得到。
(4) 由引理 \ref{derived-lemma-derived-inverts} 得到。
(5) 由定义和 (3) 得到。(6) 中的全忠实性由 (3) 和定义得到。
$S_\mathcal{E}^{-1}\mathcal{E} \to S^{-1}\mathcal{D}$ 正合，是因为
$S_\mathcal{E}^{-1}\mathcal{E}$ 中的三角是特异的，当且仅当它同构于
$\mathcal{E}$ 中某个特异三角的像；见命题
\ref{derived-proposition-construct-localization} 的证明。
$RF : \mathcal{E} \to \mathcal{D}'$ 经由正合函子
$S_\mathcal{E}^{-1}\mathcal{E} \to \mathcal{D}'$ 分解，这由引理
\ref{derived-lemma-universal-property-localization} 得到。
最后，(8) 由引理 \ref{derived-lemma-direct-sum-defined} 得到。
\end{proof}

\noindent
命题 \ref{derived-proposition-derived-functor} 告诉我们，$RF$ 定义在
$S^{-1}\mathcal{D}$ 的一个极大严格满三角子范畴上，
并且在这个三角范畴上是正合函子。图示为：
$$
\xymatrix{
\mathcal{D} \ar[d]_Q \ar[rrr]_F & & & \mathcal{D}' \\
S^{-1}\mathcal{D} & &
S_\mathcal{E}^{-1}\mathcal{E}
\ar[ll]_{\text{全忠实}}^{\text{正合}} \ar[ur]_{RF}
}
$$

\begin{definition}
\label{derived-definition-everywhere-defined}
在情形 \ref{derived-situation-derived-functor} 中，
若 $RF$ 在 $\mathcal{D}$ 的每个对象处都有定义，
则称 $F$ {\it 右可导}，或称 $RF$ {\it 处处有定义}。
若 $LF$ 在 $\mathcal{D}$ 的每个对象处都有定义，
则称 $F$ {\it 左可导}，或称 $LF$ {\it 处处有定义}。
\end{definition}

\noindent
在这种情形下，我们得到右（相应地左）导出函子
\begin{equation}
\label{derived-equation-everywhere}
RF : S^{-1}\mathcal{D} \longrightarrow \mathcal{D}',
\quad\text{（分别 }
LF : S^{-1}\mathcal{D} \longrightarrow \mathcal{D}'),
\end{equation}
见命题 \ref{derived-proposition-derived-functor}。
在大多数有趣的情形中，$RF \circ Q$ 并不等于 $F$。
事实上，典范映射 $F(X) \to RF(X)$ 甚至可能从不是同构。
实际中这不会发生，因为我们实际上只会通过下述方法证明 $F$
右可导：证明可在三角范畴 $\mathcal{D}$ 中适当选择的对象上取 $F$ 的值，
以计算 $RF$。这促使我们作出一个定义。

\begin{definition}
\label{derived-definition-computes}
在情形 \ref{derived-situation-derived-functor} 中：
\begin{enumerate}
\item 若 $RF$ 在 $\mathcal{D}$ 的对象 $X$ 处有定义，且典范映射
$F(X) \to RF(X)$ 是同构，则称 $X$ {\it 计算} $RF$。
\item 若 $LF$ 在 $\mathcal{D}$ 的对象 $X$ 处有定义，且典范映射
$LF(X) \to F(X)$ 是同构，则称 $X$ {\it 计算} $LF$。
\end{enumerate}
\end{definition}

\begin{lemma}
\label{derived-lemma-computes-shift}
假设和记号如情形 \ref{derived-situation-derived-functor} 中所述。
设 $X$ 为 $\mathcal{D}$ 的对象，$n \in \mathbf{Z}$。
\begin{enumerate}
\item $X$ 计算 $RF$，当且仅当 $X[n]$ 计算 $RF$。
\item $X$ 计算 $LF$，当且仅当 $X[n]$ 计算 $LF$。
\end{enumerate}
\end{lemma}

\begin{proof}
从略。
\end{proof}

\begin{lemma}
\label{derived-lemma-2-out-of-3-computes}
假设和记号如情形 \ref{derived-situation-derived-functor} 中所述。
设 $(X, Y, Z, f, g, h)$ 为 $\mathcal{D}$ 的特异三角。
若 $X, Y$ 计算 $RF$，则 $Z$ 也如此。对 $LF$ 有类似结论。
\end{lemma}

\begin{proof}
由引理 \ref{derived-lemma-2-out-of-3-defined}，我们知道 $RF$ 在 $Z$ 处有定义，
且将 $RF$ 应用于该三角会产生特异三角。考虑特异三角的态射
$$
\xymatrix{
(F(X), F(Y), F(Z), F(f), F(g), F(h)) \ar[d] \\
(RF(X), RF(Y), RF(Z), RF(f), RF(g), RF(h))
}
$$
三个映射中的两个是同构，因而第三个也是同构。
\end{proof}

\begin{lemma}
\label{derived-lemma-direct-sum-computes}
假设和记号如情形 \ref{derived-situation-derived-functor} 中所述。
设 $X, Y$ 为 $\mathcal{D}$ 的对象。若 $X \oplus Y$ 计算 $RF$，则
$X$ 和 $Y$ 都计算 $RF$。对 $LF$ 也有类似结论。
\end{lemma}

\begin{proof}
若 $X \oplus Y$ 计算 $RF$，则 $RF(X \oplus Y) = F(X) \oplus F(Y)$。
在引理 \ref{derived-lemma-direct-sum-defined} 的证明中，我们已经看到函子
$X/S \times Y/S \to (X \oplus Y)/S$、$(s, s') \mapsto s \oplus s'$
是余终的。因而由范畴，引理
\ref{categories-lemma-cofinal-essentially-constant} 和范畴，引理
\ref{categories-lemma-characterize-essentially-constant-ind} 的表征 (4)，我们知道，
对 $\mathcal{D}'$ 的任意对象 $W$，映射
$$
\Hom_{\mathcal{D}'}(F(X \oplus Y), W)
\longrightarrow
\colim_{s : X \to X', s' : Y \to Y'} \Hom_{\mathcal{D}'}(F(X' \oplus Y'), W)
$$
是双射。由于这个箭头显然与两边的直和分解相容，我们得出映射
$$
\Hom_{\mathcal{D}'}(F(X), W)
\longrightarrow
\colim_{s : X \to X'} \Hom_{\mathcal{D}'}(F(X'), W)
$$
是双射（一个小细节从略）。因而由范畴，引理
\ref{categories-lemma-characterize-essentially-constant-ind}，我们得出
$RF$ 在 $X$ 处有定义，其值为 $F(X)$。对 $Y$ 也同样。
\end{proof}

\begin{lemma}
\label{derived-lemma-existence-computes}
假设和记号如情形 \ref{derived-situation-derived-functor} 中所述。
\begin{enumerate}
\item 若对每个对象 $X \in \Ob(\mathcal{D})$，都存在 $S$ 中的箭头
$s : X \to X'$，使 $X'$ 计算 $RF$，则 $RF$ 处处有定义。
\item 若对每个对象 $X \in \Ob(\mathcal{D})$，都存在 $S$ 中的箭头
$s : X' \to X$，使 $X'$ 计算 $LF$，则 $LF$ 处处有定义。
\end{enumerate}
\end{lemma}

\begin{proof}
这由定义是清楚的。
\end{proof}

\begin{lemma}
\label{derived-lemma-find-existence-computes}
假设和记号如情形 \ref{derived-situation-derived-functor} 中所述。
若存在子集 $\mathcal{I} \subset \Ob(\mathcal{D})$，使得
\begin{enumerate}
\item 对所有 $X \in \Ob(\mathcal{D})$，存在 $S$ 中的
$s : X \to X'$，且 $X' \in \mathcal{I}$；
\item 对 $S$ 中每个箭头 $s : X \to X'$，若 $X, X' \in \mathcal{I}$，
则映射 $F(s) : F(X) \to F(X')$ 是同构，
\end{enumerate}
那么 $RF$ 处处有定义，而每个 $X \in \mathcal{I}$ 都计算 $RF$。
对偶地，若存在子集 $\mathcal{P} \subset \Ob(\mathcal{D})$，使得
\begin{enumerate}
\item 对所有 $X \in \Ob(\mathcal{D})$，存在 $S$ 中的
$s : X' \to X$，且 $X' \in \mathcal{P}$；
\item 对 $S$ 中每个箭头 $s : X \to X'$，若 $X, X' \in \mathcal{P}$，
则映射 $F(s) : F(X) \to F(X')$ 是同构，
\end{enumerate}
那么 $LF$ 处处有定义，而每个 $X \in \mathcal{P}$ 都计算 $LF$。
\end{lemma}

\begin{proof}
设 $X$ 为 $\mathcal{D}$ 的对象。假设 (1) 意味着，$S$ 中满足
$X' \in \mathcal{I}$ 的箭头 $s : X \to X'$ 在范畴 $X/S$ 中余终。
假设 (2) 意味着 $F$ 在这个余终子范畴上常值。
这显然意味着 $F : (X/S) \to \mathcal{D}'$ 本质常值，
对 $S$ 中任意满足 $X' \in \mathcal{I}$ 的 $s : X \to X'$，其值为 $F(X')$。
\end{proof}

\begin{lemma}
\label{derived-lemma-compose-derived-functors-general}
设 $\mathcal{A}, \mathcal{B}, \mathcal{C}$ 为三角范畴。设 $S$、相应地 $S'$
分别是 $\mathcal{A}$、相应地 $\mathcal{B}$ 中与三角结构相容的饱和乘法系。
设 $F : \mathcal{A} \to \mathcal{B}$ 和 $G : \mathcal{B} \to \mathcal{C}$ 为正合函子。
以 $F' : \mathcal{A} \to (S')^{-1}\mathcal{B}$ 表示 $F$ 与局部化函子的复合。
\begin{enumerate}
\item 若 $RF'$、$RG$、$R(G \circ F)$ 都处处有定义，则存在函子的典范变换
$t : R(G \circ F) \longrightarrow RG \circ RF'$。
\item 若 $LF'$、$LG$、$L(G \circ F)$ 都处处有定义，则存在函子的典范变换
$t : LG \circ LF' \to L(G \circ F)$。
\end{enumerate}
\end{lemma}

\begin{proof}
在本证明中，我们力求审慎。因此把导出函子看作函子
$$
RF' : S^{-1}\mathcal{A} \to (S')^{-1}\mathcal{B}, \quad
R(G \circ F) : S^{-1}\mathcal{A} \to \mathcal{C}, \quad
RG : (S')^{-1}\mathcal{B} \to \mathcal{C}.
$$
以 $Q_A : \mathcal{A} \to S^{-1}\mathcal{A}$ 和
$Q_B : \mathcal{B} \to (S')^{-1}\mathcal{B}$ 表示局部化函子。
那么 $F' = Q_B \circ F$。注意，对 $\mathcal{B}$ 的每个对象 $Y$，
存在典范映射
$$
G(Y) \longrightarrow RG(Q_B(Y))
$$
换言之，存在函子变换 $t' : G \to RG \circ Q_B$。
设 $X$ 为 $\mathcal{A}$ 的对象。我们有
\begin{align*}
R(G \circ F)(Q_A(X))
& = \colim_{s : X \to X' \in S} G(F(X')) \\
& \xrightarrow{t'} \colim_{s : X \to X' \in S} RG(Q_B(F(X'))) \\
& = \colim_{s : X \to X' \in S} RG(F'(X')) \\
& = RG(\colim_{s : X \to X' \in S} F'(X')) \\
& = RG(RF'(X)).
\end{align*}
系 $F'(X')$ 在范畴 $(S')^{-1}\mathcal{B}$ 中本质常值。
因而在上图的第三个等号中，可将余极限拉入函子 $RG$ 内；
见范畴，引理 \ref{categories-lemma-image-essentially-constant} 及其证明。
我们略去这定义函子变换的证明。左导出函子的情形类似。
\end{proof}




\section{导出范畴上的导出函子}
\label{derived-section-derived-functors-classical}

\noindent
实际中，导出函子最常出现在给定阿贝尔范畴之间的加性函子时。

\begin{situation}
\label{derived-situation-classical}
此处 $F : \mathcal{A} \to \mathcal{B}$ 是阿贝尔范畴之间的加性函子。
它诱导正合函子
$$
F : K(\mathcal{A}) \to K(\mathcal{B}), \quad
K^{+}(\mathcal{A}) \to K^{+}(\mathcal{B}), \quad
K^{-}(\mathcal{A}) \to K^{-}(\mathcal{B}).
$$
见引理 \ref{derived-lemma-additive-exact-homotopy-category}。
我们还以 $F$ 表示 $F$ 与局部化函子
$K(\mathcal{B}) \to D(\mathcal{B})$ 等的复合
$K(\mathcal{A}) \to D(\mathcal{B})$、
$K^{+}(\mathcal{A}) \to D^{+}(\mathcal{B})$ 和
$K^{-}(\mathcal{A}) \to D^-(\mathcal{B})$。
这一情形导致下面要考虑的四个导出函子。
\begin{enumerate}
\item 相对于乘法系 $\text{Qis}(\mathcal{A})$ 的
$F : K(\mathcal{A}) \to D(\mathcal{B})$ 的右导出函子。
\item 相对于乘法系 $\text{Qis}^{+}(\mathcal{A})$ 的
$F : K^{+}(\mathcal{A}) \to D^{+}(\mathcal{B})$ 的右导出函子。
\item 相对于乘法系 $\text{Qis}(\mathcal{A})$ 的
$F : K(\mathcal{A}) \to D(\mathcal{B})$ 的左导出函子。
\item 相对于乘法系 $\text{Qis}^-(\mathcal{A})$ 的
$F : K^{-}(\mathcal{A}) \to D^{-}(\mathcal{B})$ 的左导出函子。
\end{enumerate}
每一种情形都是情形 \ref{derived-situation-derived-functor} 的一个例子。
\end{situation}

\noindent
下一引理消除了可能出现的一部分歧义。

\begin{lemma}
\label{derived-lemma-irrelevant}
在情形 \ref{derived-situation-classical} 中：
\begin{enumerate}
\item 设 $X$ 为 $K^{+}(\mathcal{A})$ 的对象。
$K(\mathcal{A}) \to D(\mathcal{B})$ 的右导出函子在 $X$ 处有定义，
当且仅当 $K^{+}(\mathcal{A}) \to D^{+}(\mathcal{B})$ 的右导出函子
在 $X$ 处有定义。此外，这些值典范同构。
\item 设 $X$ 为 $K^{+}(\mathcal{A})$ 的对象。
$X$ 计算 $K(\mathcal{A}) \to D(\mathcal{B})$ 的右导出函子，
当且仅当 $X$ 计算
$K^{+}(\mathcal{A}) \to D^{+}(\mathcal{B})$ 的右导出函子。
\item 设 $X$ 为 $K^{-}(\mathcal{A})$ 的对象。
$K(\mathcal{A}) \to D(\mathcal{B})$ 的左导出函子在 $X$ 处有定义，
当且仅当 $K^{-}(\mathcal{A}) \to D^{-}(\mathcal{B})$ 的左导出函子
在 $X$ 处有定义。此外，这些值典范同构。
\item 设 $X$ 为 $K^{-}(\mathcal{A})$ 的对象。
$X$ 计算 $K(\mathcal{A}) \to D(\mathcal{B})$ 的左导出函子，当且仅当 $X$ 计算
$K^{-}(\mathcal{A}) \to D^{-}(\mathcal{B})$ 的左导出函子。
\end{enumerate}
\end{lemma}

\begin{proof}
设 $X$ 为 $K^{+}(\mathcal{A})$ 的对象。
考虑 $K(\mathcal{A})$ 中的拟同构 $s : X \to X'$。
由引理 \ref{derived-lemma-complex-cohomology-bounded}，
存在拟同构 $X' \to X''$，其中 $X''$ 下有界。
因此 $X/\text{Qis}^+(\mathcal{A})$ 在
$X/\text{Qis}(\mathcal{A})$ 中余终。由此 (1) 显然成立。
(2) 直接由 (1) 推出。(3) 和 (4) 分别与 (1) 和 (2) 对偶。
\end{proof}

\noindent
给定阿贝尔范畴 $\mathcal{A}$ 的对象 $A$，可得到复形
$$
A[0] = ( \ldots \to 0 \to A \to 0 \to \ldots )
$$
其中 $A$ 位于零次。于是得到函子
$\mathcal{A} \to K(\mathcal{A})$，$A \mapsto A[0]$。
暂且把定义在某个子范畴上的函子称为部分函子。

\begin{definition}
\label{derived-definition-derived-functor}
在情形 \ref{derived-situation-classical} 中：
\begin{enumerate}
\item $F$ 的{\it 右导出函子}是与情形
\ref{derived-situation-classical} 的 (1) 和 (2) 相对应的部分函子 $RF$。
\item $F$ 的{\it 左导出函子}是与情形
\ref{derived-situation-classical} 的 (3) 和 (4) 相对应的部分函子 $LF$。
\item 若 $A[0]$ 计算 $RF$，则称 $\mathcal{A}$ 的对象 $A$
对 $F$ {\it 右无环}，或对 $RF$ {\it 无环}。
\item 若 $A[0]$ 计算 $LF$，则称 $\mathcal{A}$ 的对象 $A$
对 $F$ {\it 左无环}，或对 $LF$ {\it 无环}。
\end{enumerate}
\end{definition}

\noindent
下面几个引理给出存在足够多无环对象的一些判据。

\begin{lemma}
\label{derived-lemma-subcategory-left-resolution}
设 $\mathcal{A}$ 为阿贝尔范畴。设
$\mathcal{P} \subset \Ob(\mathcal{A})$ 为包含 $0$ 的子集，
且 $\mathcal{A}$ 的每个对象都是 $\mathcal{P}$ 中某个元素的商。
设 $a \in \mathbf{Z}$。
\begin{enumerate}
\item 给定满足 $n > a$ 时 $K^n = 0$ 的 $K^\bullet$，
存在拟同构 $P^\bullet \to K^\bullet$，使得对所有 $n$，
$P^n \in \mathcal{P}$ 且 $P^n \to K^n$ 为满射，并且
$n > a$ 时 $P^n = 0$。
\item 给定满足 $n > a$ 时 $H^n(K^\bullet) = 0$ 的 $K^\bullet$，
存在拟同构 $P^\bullet \to K^\bullet$，使得对所有 $n$，
$P^n \in \mathcal{P}$，并且 $n > a$ 时 $P^n = 0$。
\end{enumerate}
\end{lemma}

\begin{proof}
先证明 (1)。考虑如下归纳假设 $IH_n$：
存在 $P^j \in \mathcal{P}$（$j \geq n$），且 $j > a$ 时 $P^j = 0$；
存在映射 $d^j : P^j \to P^{j + 1}$（$j \geq n$）；
还存在满射 $\alpha^j : P^j \to K^j$（$j \geq n$），使得图
$$
\xymatrix{
& &
P^n \ar[d]^\alpha \ar[r] &
P^{n + 1} \ar[d]^\alpha \ar[r] &
P^{n + 2} \ar[d]^\alpha \ar[r] &
\ldots \\
\ldots \ar[r] &
K^{n - 1} \ar[r] &
K^n \ar[r] &
K^{n + 1} \ar[r] &
K^{n + 2} \ar[r] &
\ldots
}
$$
交换；对 $j \geq n$ 有 $d^{j + 1} \circ d^j = 0$；
对 $j > n$，$\alpha$ 诱导同构
$\Ker(d^j)/\Im(d^{j - 1}) \to H^j(K^\bullet)$；并且
$\alpha : \Ker(d^n) \to \Ker(d_K^n)$ 为满射。
于是选取满射
$$
P^{n - 1}
\longrightarrow
K^{n - 1} \times_{K^n} \Ker(d^n) =
K^{n - 1} \times_{\Ker(d_K^n)} \Ker(d^n)
$$
其中 $P^{n - 1}$ 属于 $\mathcal{P}$。这使我们能把上图扩张为
$$
\xymatrix{
& P^{n - 1} \ar[d]^\alpha \ar[r] &
P^n \ar[d]^\alpha \ar[r] &
P^{n + 1} \ar[d]^\alpha \ar[r] &
P^{n + 2} \ar[d]^\alpha \ar[r] &
\ldots \\
\ldots \ar[r] &
K^{n - 1} \ar[r] &
K^n \ar[r] &
K^{n + 1} \ar[r] &
K^{n + 2} \ar[r] &
\ldots
}
$$
读者不难验证，在这一选择下 $IH_{n - 1}$ 成立。

\medskip\noindent
如下完成 (1) 的证明。首先注意，$n = a + 1$ 时 $IH_n$ 成立，
因为可以直接令 $j > a$ 时 $P^j = 0$。
因此按下降归纳法可构造复形 $P^\bullet$，它满足
$n > a$ 时 $P^n = 0$，由 $\mathcal{P}$ 中的对象组成，
并带有逐项满射的拟同构
$\alpha : P^\bullet \to K^\bullet$，正合所需。

\medskip\noindent
再证明 (2)。假设意味着态射
$\tau_{\leq a}K^\bullet \to K^\bullet$
（同调，章节 \ref{homology-section-truncations}）是拟同构。
应用 (1) 得到 $P^\bullet \to \tau_{\leq a}K^\bullet$。
复合 $P^\bullet \to K^\bullet$ 即所需拟同构。
\end{proof}

\begin{lemma}
\label{derived-lemma-subcategory-right-resolution}
设 $\mathcal{A}$ 为阿贝尔范畴。设
$\mathcal{I} \subset \Ob(\mathcal{A})$ 为包含 $0$ 的子集，
且 $\mathcal{A}$ 的每个对象都是 $\mathcal{I}$ 中某个元素的子对象。
设 $a \in \mathbf{Z}$。
\begin{enumerate}
\item 给定满足 $n < a$ 时 $K^n = 0$ 的 $K^\bullet$，
存在拟同构 $K^\bullet \to I^\bullet$，使得对所有 $n$，
$K^n \to I^n$ 为单射且 $I^n \in \mathcal{I}$，并且
$n < a$ 时 $I^n = 0$；
\item 给定满足 $n < a$ 时 $H^n(K^\bullet) = 0$ 的 $K^\bullet$，
存在拟同构 $K^\bullet \to I^\bullet$，使得
$I^n \in \mathcal{I}$，并且 $n < a$ 时 $I^n = 0$。
\end{enumerate}
\end{lemma}

\begin{proof}
本引理与引理 \ref{derived-lemma-subcategory-left-resolution} 对偶。
\end{proof}

\begin{lemma}
\label{derived-lemma-subcategory-right-acyclics}
在情形 \ref{derived-situation-classical} 中，设
$\mathcal{I} \subset \Ob(\mathcal{A})$ 为具有下列性质的子集：
\begin{enumerate}
\item $\mathcal{A}$ 的每个对象都是 $\mathcal{I}$ 中某个元素的子对象；
\item 对 $\mathcal{A}$ 中任意短正合序列
$0 \to P \to Q \to R \to 0$，若
$P, Q \in \mathcal{I}$，则 $R \in \mathcal{I}$，且
$0 \to F(P) \to F(Q) \to F(R) \to 0$ 正合。
\end{enumerate}
则 $\mathcal{I}$ 的每个对象都对 $RF$ 无环。
\end{lemma}

\begin{proof}
取 $A \in \mathcal{I}$。设 $A[0] \to K^\bullet$ 为拟同构，
其中 $K^\bullet$ 下有界。于是可找到拟同构
$K^\bullet \to I^\bullet$，其中 $I^\bullet$ 下有界且每个
$I^n \in \mathcal{I}$；见引理
\ref{derived-lemma-subcategory-right-resolution}\footnote{由 (1) 可知
$\mathcal{I}$ 非空。取 $\mathcal{I}$ 中的 $P$。
短正合序列 $0 \to P \to P \to 0 \to 0$ 与假设 (2)
表明 $0$ 属于 $\mathcal{I}$，故可应用该引理。}。
因此这些消解在范畴 $A[0]/\text{Qis}^{+}(\mathcal{A})$ 中余终。
为完成证明，只须说明：对任意拟同构
$A[0] \to I^\bullet$，若 $I^\bullet$ 下有界且
$I^n \in \mathcal{I}$，则
$F(A)[0] \to F(I^\bullet)$ 是拟同构。
为此，设 $n < n_0$ 时 $I^n = 0$。当然可假设 $n_0 < 0$。
从 $n = n_0$ 开始，利用性质 (2) 和正合序列
$$
0 \to \Ker(d^n) \to I^n \to \Im(d^n) \to 0.
$$
归纳证明 $\Im(d^{n - 1}) = \Ker(d^n)$ 以及
$\Im(d^{-1})$ 都是 $\mathcal{I}$ 的元素。
此外，性质 (2) 还保证复形
$$
0 \to F(I^{n_0}) \to F(I^{n_0 + 1}) \to \ldots \to F(I^{-1}) \to
F(\Im(d^{-1})) \to 0
$$
正合。正合序列
$0 \to \Im(d^{-1}) \to I^0 \to I^0/\Im(d^{-1}) \to 0$
意味着 $I^0/\Im(d^{-1})$ 是 $\mathcal{I}$ 的元素。
% P02-E0034: frozen source has grammatical “is an elements of”.
正合序列 $0 \to A \to I^0/\Im(d^{-1}) \to \Im(d^0) \to 0$
继而意味着 $\Im(d^0) = \Ker(d^1)$ 是 $\mathcal{I}$ 的元素；
此后同前可继续证明，对所有 $n > 0$，
$\Im(d^{n - 1}) = \Ker(d^n)$ 都是 $\mathcal{I}$ 的元素。
对上述每个短正合序列应用 $F$ 并使用 (2)，可见
% P02-E0033: frozen source says “is an isomorphism”; preceding requirement and argument establish quasi-isomorphism.
$F(A)[0] \to F(I^\bullet)$ 是同构，正合所需。
\end{proof}

\begin{lemma}
\label{derived-lemma-subcategory-left-acyclics}
在情形 \ref{derived-situation-classical} 中，设
$\mathcal{P} \subset \Ob(\mathcal{A})$ 为具有下列性质的子集：
\begin{enumerate}
\item $\mathcal{A}$ 的每个对象都是 $\mathcal{P}$ 中某个元素的商；
\item 对 $\mathcal{A}$ 中任意短正合序列
$0 \to P \to Q \to R \to 0$，若
$Q, R \in \mathcal{P}$，则 $P \in \mathcal{P}$，且
$0 \to F(P) \to F(Q) \to F(R) \to 0$ 正合。
\end{enumerate}
则 $\mathcal{P}$ 的每个对象都对 $LF$ 无环。
\end{lemma}

\begin{proof}
与引理 \ref{derived-lemma-subcategory-right-acyclics} 的证明对偶。
\end{proof}







\section{高阶导出函子}
\label{derived-section-higher-derived}

\noindent
下面这个简单引理表明，右导出函子“向右移动”。

\begin{lemma}
\label{derived-lemma-negative-vanishing}
设 $F : \mathcal{A} \to \mathcal{B}$ 为阿贝尔范畴之间的加性函子。
设 $K^\bullet$ 为 $\mathcal{A}$ 的复形，且 $a \in \mathbf{Z}$。
\begin{enumerate}
\item 若对所有 $i < a$ 有 $H^i(K^\bullet) = 0$，且
$RF$ 在 $K^\bullet$ 处有定义，则对所有 $i < a$ 有
$H^i(RF(K^\bullet)) = 0$。
\item 若 $RF$ 在 $K^\bullet$ 和 $\tau_{\leq a}K^\bullet$ 处都有定义，
则对所有 $i \leq a$ 有
$H^i(RF(\tau_{\leq a}K^\bullet)) = H^i(RF(K^\bullet))$。
\end{enumerate}
\end{lemma}

\begin{proof}
假设 $K^\bullet$ 满足 (1) 的条件。设
$s : K^\bullet \to L^\bullet$ 为任意拟同构。
由对 $K^\bullet$ 的假设，$K^\bullet \to \tau_{\geq a}L^\bullet$
也是拟同构。因此，在范畴
$K^\bullet/\text{Qis}^{+}(\mathcal{A})$ 中，满足 $L^n = 0$
（$n < a$）的拟同构 $s : K^\bullet \to L^\bullet$ 是余终的。
由范畴，引理 \ref{categories-lemma-cofinal-essentially-constant}，
可知 $RF$ 是这些 $s$ 所对应的本质常值 ind-对象
$F(L^\bullet)$ 的值。这意味着
$\text{id} : RF(K^\bullet) \to RF(K^\bullet)$ 分解通过某个
$F(L^\bullet)$，其中复形 $L^\bullet$ 满足 $n < a$ 时 $L^n = 0$。
由此，对 $i < a$ 有 $H^i(RF(K^\bullet)) = 0$。

\medskip\noindent
为证明 (2)，使用注记 \ref{derived-remark-truncation-distinguished-triangle}
中的特异三角
$$
\tau_{\leq a}K^\bullet \to K^\bullet \to \tau_{\geq a + 1}K^\bullet
\to (\tau_{\leq a}K^\bullet)[1]
$$
并由引理 \ref{derived-lemma-2-out-of-3-defined} 得到：
$RF$ 在 $\tau_{\geq a + 1}K^\bullet$ 处也有定义，并且在
$D(\mathcal{B})$ 中有特异三角
$$
RF(\tau_{\leq a}K^\bullet) \to RF(K^\bullet) \to RF(\tau_{\geq a + 1}K^\bullet)
\to RF(\tau_{\leq a}K^\bullet)[1]
$$
由 (1)，$RF(\tau_{\geq a + 1}K^\bullet)$ 在次数 $< a + 1$ 上
上同调为零。该特异三角的长正合上同调序列于是给出所需结论。
\end{proof}

\begin{definition}
\label{derived-definition-higher-derived-functors}
设 $F : \mathcal{A} \to \mathcal{B}$ 为阿贝尔范畴之间的加性函子。
假设 $RF : D^{+}(\mathcal{A}) \to D^{+}(\mathcal{B})$ 处处有定义。
设 $i \in \mathbf{Z}$。$F$ 的{\it 第 $i$ 个右导出函子 $R^iF$}是函子
$$
R^iF = H^i \circ RF :
\mathcal{A}
\longrightarrow
\mathcal{B}
$$
\end{definition}

\noindent
下面的引理表明，除非函子左正合，否则取其右导出函子其实没有多大意义。

\begin{lemma}
\label{derived-lemma-left-exact-higher-derived}
设 $F : \mathcal{A} \to \mathcal{B}$ 为阿贝尔范畴之间的加性函子，
并假设 $RF : D^{+}(\mathcal{A}) \to D^{+}(\mathcal{B})$ 处处有定义。
\begin{enumerate}
\item 对 $i < 0$ 有 $R^iF = 0$；
\item $R^0F$ 左正合；
\item 映射 $F \to R^0F$ 为同构，当且仅当 $F$ 左正合。
\end{enumerate}
\end{lemma}

\begin{proof}
设 $A$ 为 $\mathcal{A}$ 的对象。
% P02-E0032: frozen source math has an unmatched opening parenthesis.
由引理 \ref{derived-lemma-negative-vanishing}，对 $i < 0$ 有
$H^i(RF(A[0]) = 0$。这证明了 (1)。

\medskip\noindent
设 $0 \to A \to B \to C \to 0$ 为 $\mathcal{A}$ 的短正合序列。
由引理 \ref{derived-lemma-derived-canonical-delta-functor}，可在
$D^{+}(\mathcal{A})$ 中得到特异三角
$(A[0], B[0], C[0], a, b, c)$。
由长正合上同调序列（以及上面证明的 $i < 0$ 时的消失性），
可知 $0 \to R^0F(A) \to R^0F(B) \to R^0F(C)$ 正合。
故 $R^0F$ 左正合。当然，这也证明了：若
$F \to R^0F$ 为同构，则 $F$ 左正合。

\medskip\noindent
假设 $F$ 左正合。回忆 $RF(A[0])$ 是本质常值系
$F(K^\bullet)$ 的值，其中 $s : A[0] \to K^\bullet$ 遍历拟同构。
由此 $R^0F(A)$ 是本质常值系 $H^0(F(K^\bullet))$ 的值，
其中 $s : A[0] \to K^\bullet$ 遍历拟同构；见范畴，引理
\ref{categories-lemma-image-essentially-constant}。
但若 $s : A[0] \to K^\bullet$ 是拟同构，则
$A[0] \to \tau_{\geq 0}K^\bullet$ 是拟同构。
因此在范畴 $A[0]/\text{Qis}^{+}(\mathcal{A})$ 中，满足
$n < 0$ 时 $K^n = 0$ 的拟同构 $s : A[0] \to K^\bullet$ 是余终的。
由范畴，引理 \ref{categories-lemma-cofinal-essentially-constant}，
可知只须考虑这样的 $s$。此外，对这样的 $s$，序列
$$
0 \to A \to K^0 \to K^1
$$
正合。由于 $F$ 左正合，
$0 \to F(A) \to F(K^0) \to F(K^1)$ 也正合。
于是 $F(A) \to H^0(F(K^\bullet))$ 为同构，
而且该系实际上是取值为 $F(A)$ 的常值系。
故 $R^0F = F$，正合所需。
\end{proof}

\begin{lemma}
\label{derived-lemma-F-acyclic}
设 $F : \mathcal{A} \to \mathcal{B}$ 为阿贝尔范畴之间的加性函子，
并假设 $RF : D^{+}(\mathcal{A}) \to D^{+}(\mathcal{B})$ 处处有定义。
设 $A$ 为 $\mathcal{A}$ 的对象。
\begin{enumerate}
\item $A$ 对 $F$ 右无环，当且仅当
$F(A) \to R^0F(A)$ 为同构，且对所有 $i > 0$ 有 $R^iF(A) = 0$；
\item 若 $F$ 左正合，则 $A$ 对 $F$ 右无环，当且仅当
对所有 $i > 0$ 有 $R^iF(A) = 0$。
\end{enumerate}
\end{lemma}

\begin{proof}
若 $A$ 对 $F$ 右无环，则 $RF(A[0]) = F(A)[0]$；特别地，
$F(A) \to R^0F(A)$ 为同构，且 $i \not = 0$ 时
$R^iF(A) = 0$。反之，若 $F(A) \to R^0F(A)$ 为同构，且对所有
$i > 0$ 有 $R^iF(A) = 0$，则由引理
\ref{derived-lemma-left-exact-higher-derived} 的 (1)，
$F(A[0]) \to RF(A[0])$ 为拟同构，故 $A$ 无环。
若 $F$ 左正合，则 $F = R^0F$；见引理
\ref{derived-lemma-left-exact-higher-derived}。
\end{proof}

\begin{lemma}
\label{derived-lemma-F-acyclic-ses}
设 $F : \mathcal{A} \to \mathcal{B}$ 为阿贝尔范畴之间的左正合函子，
并假设 $RF : D^{+}(\mathcal{A}) \to D^{+}(\mathcal{B})$ 处处有定义。
设 $0 \to A \to B \to C \to 0$ 为 $\mathcal{A}$ 的短正合序列。
\begin{enumerate}
\item 若 $A$ 和 $C$ 都对 $F$ 右无环，则 $B$ 亦然。
\item 若 $A$ 和 $B$ 都对 $F$ 右无环，则 $C$ 亦然。
\item 若 $B$ 和 $C$ 都对 $F$ 右无环，且
$F(B) \to F(C)$ 为满射，则 $A$ 对 $F$ 右无环。
\end{enumerate}
在上述三种情形中，
$$
0 \to F(A) \to F(B) \to F(C) \to 0
$$
都是 $\mathcal{B}$ 的短正合序列。
\end{lemma}

\begin{proof}
由引理 \ref{derived-lemma-derived-canonical-delta-functor}，
可在 $D^{+}(\mathcal{A})$ 中得到特异三角
$(A[0], B[0], C[0], a, b, c)$。
由于 $RF$ 是正合函子，并且 $i < 0$ 时 $R^iF = 0$、
$R^0F = F$（引理 \ref{derived-lemma-left-exact-higher-derived}），
可在阿贝尔范畴 $\mathcal{B}$ 中得到正合上同调序列
$$
0 \to F(A) \to F(B) \to F(C) \to R^1F(A) \to \ldots
$$
所以本引理由引理 \ref{derived-lemma-F-acyclic} 中对无环对象的刻画推出。
\end{proof}

\begin{lemma}
\label{derived-lemma-right-derived-delta-functor}
设 $F : \mathcal{A} \to \mathcal{B}$ 为阿贝尔范畴之间的加性函子，
并假设 $RF : D^{+}(\mathcal{A}) \to D^{+}(\mathcal{B})$ 处处有定义。
\begin{enumerate}
\item 函子 $R^iF$（$i \geq 0$）带有从
$\mathcal{A} \to \mathcal{B}$ 出发的典范 $\delta$-函子结构；见
同调，定义 \ref{homology-definition-cohomological-delta-functor}。
\item 若 $\mathcal{A}$ 的每个对象都是某个对 $F$ 右无环对象的子对象，
则 $\{R^iF, \delta\}_{i \geq 0}$ 是泛 $\delta$-函子；见
同调，定义 \ref{homology-definition-universal-delta-functor}。
\end{enumerate}
\end{lemma}

\begin{proof}
函子 $\mathcal{A} \to \text{Comp}^{+}(\mathcal{A})$，
$A \mapsto A[0]$ 是正合的。函子
$\text{Comp}^{+}(\mathcal{A}) \to D^{+}(\mathcal{A})$
是 $\delta$-函子；见引理 \ref{derived-lemma-derived-canonical-delta-functor}。
函子 $RF : D^{+}(\mathcal{A}) \to D^{+}(\mathcal{B})$ 是正合的。
最后，函子 $H^0 : D^{+}(\mathcal{B}) \to \mathcal{B}$
是同调函子；见定义 \ref{derived-definition-unbounded-derived-category}。
因此，由引理 \ref{derived-lemma-compose-delta-functor-homological} 和
引理 \ref{derived-lemma-exact-compose-delta-functor}，可得到一个
$\delta$-函子结构。(2) 由同调，引理
\ref{homology-lemma-efface-implies-universal} 以及引理
\ref{derived-lemma-F-acyclic} 中对无环对象的描述推出。
\end{proof}

\begin{lemma}[Leray 无环性引理]
\label{derived-lemma-leray-acyclicity}
设 $F : \mathcal{A} \to \mathcal{B}$ 为阿贝尔范畴之间的加性函子。
设 $A^\bullet$ 为由对 $F$ 右无环的对象组成的下有界复形，
且 $RF$ 在 $A^\bullet$ 处有定义\footnote{例如，当
$RF : D^{+}(\mathcal{A}) \to D^{+}(\mathcal{B})$ 处处有定义时成立。}。
典范映射
$$
F(A^\bullet) \longrightarrow RF(A^\bullet)
$$
是 $D^{+}(\mathcal{B})$ 中的同构；即 $A^\bullet$ 计算 $RF$。
\end{lemma}

\begin{proof}
设 $A^\bullet$ 为由对 $F$ 右无环的对象组成的有界复形。
我们断言 $RF$ 在 $A^\bullet$ 处有定义，且
$F(A^\bullet) \to RF(A^\bullet)$ 是 $D^+(\mathcal{B})$ 中的同构。
事实上，由定义 \ref{derived-definition-derived-functor}，
这对至多含有一个非零且对 $F$ 右无环对象的复形成立。
再设 $n \not \in [a, b]$ 时 $A^n = 0$。
使用“笨”截断可得到逐项分裂的复形短正合序列
$$
0 \to \sigma_{\geq a + 1} A^\bullet \to A^\bullet \to
\sigma_{\leq a} A^\bullet \to 0
$$
见同调，章节 \ref{homology-section-truncations}。
从而得到特异三角
$(\sigma_{\geq a + 1} A^\bullet, A^\bullet, \sigma_{\leq a} A^\bullet)$。
由归纳假设，$RF$ 在外侧两个复形处有定义，且这两个复形都计算 $RF$。
由引理 \ref{derived-lemma-2-out-of-3-computes}，中间的复形亦然。

\medskip\noindent
设 $A^\bullet$ 是由无环对象组成的下有界复形，
且 $RF$ 在 $A^\bullet$ 处有定义。为证明
$F(A^\bullet) \to RF(A^\bullet)$ 是 $D^{+}(\mathcal{B})$ 中的同构，
只须证明对所有 $i$，
$H^i(F(A^\bullet)) \to H^i(RF(A^\bullet))$ 都是同构。
取 $i$。考虑逐项分裂的复形短正合序列
$$
0 \to \sigma_{\geq i + 2} A^\bullet \to A^\bullet \to
\sigma_{\leq i + 1} A^\bullet \to 0.
$$
注意，它诱导逐项分裂的短正合序列
$$
0 \to \sigma_{\geq i + 2} F(A^\bullet) \to F(A^\bullet) \to
\sigma_{\leq i + 1} F(A^\bullet) \to 0.
$$
因此得到特异三角
$$
(\sigma_{\geq i + 2} A^\bullet, A^\bullet,
\sigma_{\leq i + 1} A^\bullet)
\quad\text{且}\quad
(\sigma_{\geq i + 2} F(A^\bullet), F(A^\bullet),
\sigma_{\leq i + 1} F(A^\bullet))
$$
由于 $RF$ 在 $A^\bullet$ 处有定义（依假设），并且在
$\sigma_{\leq i + 1}A^\bullet$ 处有定义（依第一段），
故 $RF$ 在 $\sigma_{\geq i + 2}A^\bullet$ 处也有定义，
而且得到特异三角
$$
(RF(\sigma_{\geq i + 2} A^\bullet), RF(A^\bullet),
RF(\sigma_{\leq i + 1} A^\bullet))
$$
见引理 \ref{derived-lemma-2-out-of-3-defined}。
利用这些特异三角，可得到正合序列之间的映射
$$
\xymatrix{
H^i(\sigma_{\geq i + 2} F(A^\bullet)) \ar[r] \ar[d] &
H^i(F(A^\bullet)) \ar[r] \ar[d]^\alpha &
H^i(\sigma_{\leq i + 1} F(A^\bullet)) \ar[r] \ar[d]^\beta &
H^{i + 1}(\sigma_{\geq i + 2} F(A^\bullet)) \ar[d] \\
H^i(RF(\sigma_{\geq i + 2} A^\bullet)) \ar[r] &
H^i(RF(A^\bullet)) \ar[r] &
H^i(RF(\sigma_{\leq i + 1} A^\bullet)) \ar[r] &
H^{i + 1}(RF(\sigma_{\geq i + 2} A^\bullet))
}
$$
由第一段的结果，映射 $\beta$ 是同构。
直接看出左上角和右上角的对象为零。因此，为完成证明，只须说明
$H^i(RF(\sigma_{\geq i + 2} A^\bullet)) = 0$ 且
$H^{i + 1}(RF(\sigma_{\geq i + 2} A^\bullet)) = 0$。
这立即由引理 \ref{derived-lemma-negative-vanishing} 推出。
\end{proof}

\begin{proposition}
\label{derived-proposition-enough-acyclics}
\begin{slogan}
阿贝尔范畴之间的函子可借助由该函子的无环对象组成的复形作消解，
扩张到（下有界或上有界）导出范畴。
\end{slogan}
设 $F : \mathcal{A} \to \mathcal{B}$ 为阿贝尔范畴之间的加性函子。
\begin{enumerate}
\item 若 $\mathcal{A}$ 的每个对象都嵌入某个对 $RF$ 无环的对象，
则 $RF$ 在整个 $K^{+}(\mathcal{A})$ 上有定义，并得到正合函子
$$
RF : D^{+}(\mathcal{A}) \longrightarrow D^{+}(\mathcal{B})
$$
见 (\ref{derived-equation-everywhere})。此外，任意由对 $RF$ 无环的项
组成的下有界复形 $A^\bullet$ 都计算 $RF$。
\item 若 $\mathcal{A}$ 的每个对象都是某个对 $LF$ 无环对象的商，
% P02-E0035: frozen source omits the article in “is quotient of”.
则 $LF$ 在整个 $K^{-}(\mathcal{A})$ 上有定义，并得到正合函子
$$
LF : D^{-}(\mathcal{A}) \longrightarrow D^{-}(\mathcal{B})
$$
见 (\ref{derived-equation-everywhere})。此外，任意由对 $LF$ 无环的项
组成的上有界复形 $A^\bullet$ 都计算 $LF$。
\end{enumerate}
\end{proposition}

\begin{proof}
假设 $\mathcal{A}$ 的每个对象都嵌入某个对 $RF$ 无环的对象。
设 $\mathcal{I}$ 为对 $RF$ 无环的对象所成的集合。
设 $K^\bullet$ 为 $\mathcal{A}$ 中的下有界复形。
由引理 \ref{derived-lemma-subcategory-right-resolution}，
存在拟同构 $\alpha : K^\bullet \to I^\bullet$，其中
$I^\bullet$ 下有界且 $I^n \in \mathcal{I}$。
所以，为证明 (1)，只须说明：若
$s : I^\bullet \to (I')^\bullet$ 是由 $\mathcal{I}$ 中对象组成的
下有界复形之间的拟同构，则
$F(I^\bullet) \to F((I')^\bullet)$ 是拟同构；见引理
\ref{derived-lemma-find-existence-computes}。
注意，锥 $C(s)^\bullet$ 是无环下有界复形，且每一项都属于
$\mathcal{I}$。故只须说明：给定每一项都属于 $\mathcal{I}$ 的
无环下有界复形 $I^\bullet$，复形 $F(I^\bullet)$ 是无环的。

\medskip\noindent
设 $n < n_0$ 时 $I^n = 0$。令 $J^n = \Im(d^n)$，
便把 $I^\bullet$ 分解成短正合序列
$0 \to J^n \to I^{n + 1} \to J^{n + 1} \to 0$（$n \geq n_0$）。
由引理 \ref{derived-lemma-derived-canonical-delta-functor}，
这些序列在 $D^+(\mathcal{A})$ 中诱导特异三角
$(J^n, I^{n + 1}, J^{n + 1})$。
对每个 $k \in \mathbf{Z}$，用 $H_k$ 表示如下断言：
对所有 $n \leq k$，对象 $J^n$ 属于 $\mathcal{I}$。
当 $k < n_0$ 时，$H_k$ 显然成立。若 $H_n$ 成立，
则引理 \ref{derived-lemma-2-out-of-3-computes} 表明
$J^{n + 1}$ 属于 $\mathcal{I}$，从而 $H_{n + 1}$ 成立。
由命题 \ref{derived-proposition-derived-functor}，有特异三角
$(RF(J^n), RF(I^{n + 1}), RF(J^{n + 1}))$。
由于 $J^n, I^{n + 1}, J^{n + 1}$ 都属于 $\mathcal{I}$，
与该特异三角相伴的长正合上同调序列
(\ref{derived-equation-long-exact-cohomology-sequence-D}) 缩并为正合序列
$$
0 \to F(J^n) \to F(I^{n + 1}) \to F(J^{n + 1}) \to 0
$$
这又证明了 $F(I^\bullet)$ 正合。

\medskip\noindent
$LF$ 的情形之证明与此对偶。
\end{proof}

\begin{lemma}
\label{derived-lemma-right-derived-exact-functor}
设 $F : \mathcal{A} \to \mathcal{B}$ 为阿贝尔范畴之间的正合函子。
则：
\begin{enumerate}
\item $\mathcal{A}$ 的每个对象都对 $F$ 右无环；
\item $RF : D^{+}(\mathcal{A}) \to D^{+}(\mathcal{B})$ 处处有定义；
\item $RF : D(\mathcal{A}) \to D(\mathcal{B})$ 处处有定义；
\item 每个复形都计算 $RF$；换言之，对所有复形，典范映射
$F(K^\bullet) \to RF(K^\bullet)$ 都是同构；
\item 对 $i \not = 0$ 有 $R^iF = 0$。
\end{enumerate}
\end{lemma}

\begin{proof}
这是因为 $F$ 把无环复形送到无环复形，并把拟同构送到拟同构。
细节从略。
\end{proof}









\section{导出范畴的三角子范畴}
\label{derived-section-triangulated-sub}

\noindent
设 $\mathcal{A}$ 为阿贝尔范畴。本节考察某些严格满饱和三角子范畴
$\mathcal{D}' \subset D(\mathcal{A})$。

\medskip\noindent
设 $\mathcal{B} \subset \mathcal{A}$ 为弱 Serre 子范畴；见
同调，定义 \ref{homology-definition-serre-subcategory} 和
引理 \ref{homology-lemma-characterize-weak-serre-subcategory}。
% P02-E0036: frozen source omits “be” in “We let ... the full subcategory”.
令 $D_\mathcal{B}(\mathcal{A})$ 为 $D(\mathcal{A})$ 的满子范畴，
其对象为
$$
\Ob(D_\mathcal{B}(\mathcal{A}))
=
\{X \in \Ob(D(\mathcal{A})) \mid
H^n(X) \text{ 是下列范畴的对象： }\mathcal{B}\text{ 对所有 }n\}
$$
还定义
$D^{+}_\mathcal{B}(\mathcal{A}) =
D^{+}(\mathcal{A}) \cap D_\mathcal{B}(\mathcal{A})$，
其他有界版本同理。

\begin{lemma}
\label{derived-lemma-cohomology-in-serre-subcategory}
设 $\mathcal{A}$ 为阿贝尔范畴，
$\mathcal{B} \subset \mathcal{A}$ 为弱 Serre 子范畴。
则 $D_\mathcal{B}(\mathcal{A})$ 是 $D(\mathcal{A})$ 的
严格满饱和三角子范畴。各有界版本亦然。
\end{lemma}

\begin{proof}
显然，$D_\mathcal{B}(\mathcal{A})$ 是在平移函子下保持的加性子范畴。
若 $X \oplus Y$ 属于 $D_\mathcal{B}(\mathcal{A})$，则
% P02-E0037: frozen source says “kernels of maps between maps of objects”; one “maps” is extraneous.
$H^n(X)$ 和 $H^n(Y)$ 都是 $\mathcal{B}$ 中对象之间映射的核，
因为 $H^n(X \oplus Y) = H^n(X) \oplus H^n(Y)$。
所以 $X$ 和 $Y$ 都属于 $D_\mathcal{B}(\mathcal{A})$。由
引理 \ref{derived-lemma-triangulated-subcategory}，只须证明：
给定特异三角 $(X, Y, Z, f, g, h)$，若 $X$ 和 $Y$ 属于
$D_\mathcal{B}(\mathcal{A})$，则 $Z$ 是
$D_\mathcal{B}(\mathcal{A})$ 的对象。
长正合上同调序列 (\ref{derived-equation-long-exact-cohomology-sequence-D})
和弱 Serre 子范畴的定义（见同调，定义
\ref{homology-definition-serre-subcategory}）表明，
对所有 $n$，$H^n(Z)$ 都是 $\mathcal{B}$ 的对象。
故 $Z$ 是 $D_\mathcal{B}(\mathcal{A})$ 的对象。
\end{proof}

\noindent
继续假设 $\mathcal{B}$ 是阿贝尔范畴 $\mathcal{A}$ 的弱 Serre 子范畴。
那么 $\mathcal{B}$ 是阿贝尔范畴，且包含函子
$\mathcal{B} \to \mathcal{A}$ 正合。
所以得到导出函子 $D(\mathcal{B}) \to D(\mathcal{A})$；见引理
\ref{derived-lemma-right-derived-exact-functor}。显然，函子
$D(\mathcal{B}) \to D(\mathcal{A})$ 分解通过典范正合函子
\begin{equation}
\label{derived-equation-compare}
D(\mathcal{B}) \longrightarrow D_\mathcal{B}(\mathcal{A})
\end{equation}
毕竟，由 $\mathcal{B}$ 中对象组成的复形当然给出
$D_\mathcal{B}(\mathcal{A})$ 的对象；而
$D_\mathcal{B}(\mathcal{A})$ 中的特异三角，恰好是顶点都属于
$D_\mathcal{B}(\mathcal{A})$ 的 $D(\mathcal{A})$ 中的特异三角。
由于 $D(\mathcal{B}) \to D(\mathcal{A})$ 正合，故上述函子正合。
类似地，对各有界版本得到函子
$D^+(\mathcal{B}) \to D^+_\mathcal{B}(\mathcal{A})$、
$D^-(\mathcal{B}) \to D^-_\mathcal{B}(\mathcal{A})$ 和
$D^b(\mathcal{B}) \to D^b_\mathcal{B}(\mathcal{A})$。
很多情形中的一个关键问题是所展示的函子是否为等价。

\medskip\noindent
现在假设 $\mathcal{B}$ 是 $\mathcal{A}$ 的 Serre 子范畴。
此时有商函子 $\mathcal{A} \to \mathcal{A}/\mathcal{B}$；见
同调，引理 \ref{homology-lemma-serre-subcategory-is-kernel}。
此时 $D_\mathcal{B}(\mathcal{A})$ 是函子
$D(\mathcal{A}) \to D(\mathcal{A}/\mathcal{B})$ 的核。
因此由引理 \ref{derived-lemma-universal-property-quotient} 得到典范函子
$$
D(\mathcal{A})/D_\mathcal{B}(\mathcal{A})
\longrightarrow
D(\mathcal{A}/\mathcal{B})
$$
各有界版本同理。

\begin{lemma}
\label{derived-lemma-derived-of-quotient}
设 $\mathcal{A}$ 为阿贝尔范畴，
$\mathcal{B} \subset \mathcal{A}$ 为 Serre 子范畴。
则 $D(\mathcal{A}) \to D(\mathcal{A}/\mathcal{B})$ 本质满。
\end{lemma}

\begin{proof}
我们使用同调，引理 \ref{homology-lemma-serre-subcategory-is-kernel}
的证明中对范畴 $\mathcal{A}/\mathcal{B}$ 的描述。
设 $(X^\bullet, d^\bullet)$ 为 $\mathcal{A}/\mathcal{B}$ 的复形。
这意味着 $X^i$ 是 $\mathcal{A}$ 的对象，
$d^i : X^i \to X^{i + 1}$ 是 $\mathcal{A}/\mathcal{B}$ 中的态射，
且在 $\mathcal{A}/\mathcal{B}$ 中有 $d^i \circ d^{i - 1} = 0$。

\medskip\noindent
对 $i \geq 0$，可写 $d^i = (s^i, f^i)$，其中
$s^i : Y^i \to X^i$ 是 $\mathcal{A}$ 中的态射，
它的核和余核都属于 $\mathcal{B}$（等价地，$s^i$ 在商范畴中成为同构），
而 $f^i : Y^i \to X^{i + 1}$ 是 $\mathcal{A}$ 中的态射。
我们将归纳构造交换图
$$
\xymatrix{
& (X')^1 \ar@{..>}[r] & (X')^2 \ar@{..>}[r] & \ldots \\
X^0 \ar@{..>}[ru] &
X^1 \ar@{..>}[u] &
X^2 \ar@{..>}[u] &
\ldots \\
Y^0 \ar[u]_{s^0} \ar[ru]_{f^0} &
Y^1 \ar[u]_{s^1} \ar[ru]_{f^1} &
Y^2 \ar[u]_{s^2} \ar[ru]_{f^2} &
\ldots
}
$$
其中竖直箭头 $X^i \to (X')^i$ 在商范畴中成为同构。
具体地，先令
$(X')^1 = \Coker(Y^0 \to X^0 \oplus X^1)$
（更准确地说，是箭头 $s^0$ 和 $f^0$ 所成图的推出），
从而得到第一个交换图。接着取
$(X')^2 = \Coker(Y^1 \to (X')^1 \oplus X^2)$，依此类推。
再令 $n \leq 0$ 时 $(X')^n = X^n$，便可看出映射
$(X^\bullet, d^\bullet) \to ((X')^\bullet, (d')^\bullet)$
是 $\mathcal{A}/\mathcal{B}$ 中复形的同构。
所以可以假设：对 $n \geq 0$，
$d^n : X^n \to X^{n + 1}$ 由 $\mathcal{A}$ 中的映射
$X^n \to X^{n + 1}$ 给出。

\medskip\noindent
对偶地，对 $i < 0$ 可写 $d^i = (g^i, t^{i + 1})$，其中
$t^{i + 1} : X^{i + 1} \to Z^{i + 1}$ 是商范畴中的同构，
而 $g^i : X^i \to Z^{i + 1}$ 是态射。
我们将归纳构造交换图
% P02-E0040: frozen diagram uses subscripted t and g indices where the declared maps require superscripts.
$$
\xymatrix{
\ldots &
Z^{-2} &
Z^{-1} &
Z^0 \\
\ldots &
X^{-2} \ar[u]_{t_{-2}} \ar[ru]_{g_{-2}} &
X^{-1} \ar[u]_{t_{-1}} \ar[ru]_{g_{-1}} &
X^0 \ar[u]_{t^0} \\
\ldots &
(X')^{-2} \ar@{..>}[u] \ar@{..>}[r] &
(X')^{-1} \ar@{..>}[u] \ar@{..>}[ru]
}
$$
其中竖直箭头 $(X')^i \to X^i$ 在商范畴中成为同构。
具体地，取 $(X')^{-1} = X^{-1} \times_{Z^0} X^0$，
再取 $(X')^{-2} = X^{-2} \times_{Z^{-1}} (X')^{-1}$，依此类推。
再令 $n \geq 0$ 时 $(X')^n = X^n$，便可看出映射
$((X')^\bullet, (d')^\bullet) \to (X^\bullet, d^\bullet)$
是 $\mathcal{A}/\mathcal{B}$ 中复形的同构。
所以可以假设：对所有 $n \in \mathbf{Z}$，
$d^n : X^n \to X^{n + 1}$ 由 $\mathcal{A}$ 中的映射
$d^n : X^n \to X^{n + 1}$ 给出。

\medskip\noindent
在这种情形下，复合 $d^n \circ d^{n - 1}$ 在
$\mathcal{A}/\mathcal{B}$ 中为零。若对 $n > 0$，把 $X^n$ 替换为
$$
(X')^n = X^n/\sum\nolimits_{0 < k \leq n} \Im(\Im(X^{k - 2} \to X^k) \to X^n)
$$
则对 $n \geq 0$，复合 $d^n \circ d^{n - 1}$ 为零。
（与上面第二段类似，可得到复形的同构
$(X^\bullet, d^\bullet) \to ((X')^\bullet, (d')^\bullet)$。）
最后，对 $n < 0$，把 $X^n$ 替换为
$$
(X')^n = \bigcap\nolimits_{n \leq k < 0}
(X^n \to X^k)^{-1}\Ker(X^k \to X^{k + 2})
$$
并用同样的方法论证，便得到 $\mathcal{A}$ 中的复形，
其在 $\mathcal{A}/\mathcal{B}$ 中的像与给定复形同构。
\end{proof}

\begin{lemma}
\label{derived-lemma-quotient-by-serre-easy}
设 $\mathcal{A}$ 为阿贝尔范畴，
$\mathcal{B} \subset \mathcal{A}$ 为 Serre 子范畴。
假设函子 $v : \mathcal{A} \to \mathcal{A}/\mathcal{B}$
有左伴随 $u : \mathcal{A}/\mathcal{B} \to \mathcal{A}$，
且 $vu \cong \text{id}$。则
$$
D(\mathcal{A})/D_\mathcal{B}(\mathcal{A}) = D(\mathcal{A}/\mathcal{B})
$$
各有界版本亦然。
\end{lemma}

\begin{proof}
由引理 \ref{derived-lemma-derived-of-quotient}，函子
$D(v) : D(\mathcal{A}) \to D(\mathcal{A}/\mathcal{B})$ 本质满。
对 $D(\mathcal{A})$ 的对象 $X$，伴随映射
$c_X : uvX \to X$ 在 $D(\mathcal{A}/\mathcal{B})$ 中映为同构，
因为假设 $vu \cong \text{id}$ 意味着 $vuv \cong v$。
所以，在特异三角 $(uvX, X, Z, c_X, g, h)$ 中，
考察长正合上同调序列可知，对象 $Z$ 属于
$D_\mathcal{B}(\mathcal{A})$。故 $c_X$ 是用于定义商范畴
$D(\mathcal{A})/D_\mathcal{B}(\mathcal{A})$ 的乘法系的元素。
从而在 $D(\mathcal{A})/D_\mathcal{B}(\mathcal{A})$ 中有
$uvX \cong X$。
% P02-E0038: frozen source has an extra closing parenthesis in Ob(A)).
对 $X, Y \in \Ob(\mathcal{A}))$，映射
$$
\Hom_{D(\mathcal{A})/D_\mathcal{B}(\mathcal{A})}(X, Y)
\longrightarrow
\Hom_{D(\mathcal{A}/\mathcal{B})}(vX, vY)
$$
是双射，因为 $u$ 给出逆映射（由上述讨论）。
\end{proof}

\noindent
对某些 Serre 子范畴 $\mathcal{B} \subset \mathcal{A}$，
可以证明函子 $D(\mathcal{B}) \to D_\mathcal{B}(\mathcal{A})$ 满忠实。

\begin{lemma}
\label{derived-lemma-fully-faithful-embedding}
设 $\mathcal{A}$ 为阿贝尔范畴，
$\mathcal{B} \subset \mathcal{A}$ 为 Serre 子范畴。
假设对每个满射 $X \to Y$，其中
$X \in \Ob(\mathcal{A})$ 且 $Y \in \Ob(\mathcal{B})$，
存在 $X' \subset X$，$X' \in \Ob(\mathcal{B})$，且后者满射到 $Y$。
则 (\ref{derived-equation-compare}) 中的函子
$D^-(\mathcal{B}) \to D^-_\mathcal{B}(\mathcal{A})$ 是等价。
\end{lemma}

\begin{proof}
设 $X^\bullet$ 为 $\mathcal{A}$ 的上有界复形，且对所有
$i \in \mathbf{Z}$ 有 $H^i(X^\bullet) \in \Ob(\mathcal{B})$。
再假设对所有 $i \in \mathbf{Z}$ 给定
$B^i \subset X^i$，$B^i \in \Ob(\mathcal{B})$。
断言：存在子复形 $Y^\bullet \subset X^\bullet$，使得
\begin{enumerate}
\item $Y^\bullet \to X^\bullet$ 是拟同构；
\item 对所有 $i \in \mathbf{Z}$ 有 $Y^i \in \Ob(\mathcal{B})$；
\item 对所有 $i \in \mathbf{Z}$ 有 $B^i \subset Y^i$。
\end{enumerate}
为证明该断言，先利用引理的假设选取
$C^i \subset \Ker(d^i : X^i \to X^{i + 1})$，
$C^i \in \Ob(\mathcal{B})$，使其满射到 $H^i(X^\bullet)$。
令 $D^i = C^i + d^{i - 1}(B^{i - 1}) + B^i$，
可得到满足 (2) 和 (3) 的子复形 $D^\bullet$，
并且对所有 $i \in \mathbf{Z}$，
$H^i(D^\bullet) \to H^i(X^\bullet)$ 都为满射。
任取满足 $E^i \subset X^i$、$E^i \in \Ob(\mathcal{B})$ 且
$d^i(E^i) \subset D^{i + 1} + E^{i + 1}$ 的一族对象。
令 $Y^i = D^i + E^i$，便得到各项属于 $\mathcal{B}$ 的子复形，
它的上同调满射到 $X^\bullet$ 的上同调。显然，若
$d^i(E^i) = (D^{i + 1} + E^{i + 1}) \cap \Im(d^i)$，
则上同调上的映射也是单射。当 $n \gg 0$ 时可令 $E^n$ 等于 $0$。
按下降归纳法，可以对所有 $i$ 选取具有所需性质的 $E^i$。
具体地，给定 $E^{i + 1}, E^{i + 2}, \ldots$ 后，选取
$E^i \subset X^i$，使得
$d^i(E^i) = (D^{i + 1} + E^{i + 1}) \cap \Im(d^i)$。
这是可行的，因为引理的假设与如下事实共同保证了它：
$(D^{i + 1} + E^{i + 1}) \cap \Im(d^i)$ 属于 $\mathcal{B}$，
这是由于 $\mathcal{B}$ 是 $\mathcal{A}$ 的 Serre 子范畴。

\medskip\noindent
上述断言推出本引理。本质满性立即由该断言得到。
下面证明忠实性。设 $f : U^\bullet \to V^\bullet$ 是
$\mathcal{B}$ 的上有界复形之间的态射，且其在
$D(\mathcal{A})$ 中的像为零。则存在到 $\mathcal{A}$ 的某个
上有界复形中的拟同构 $s : V^\bullet \to X^\bullet$，
使得 $s \circ f$ 零伦。选取 $0$ 与 $s \circ f$ 之间的同伦
$h^i : U^i \to X^{i - 1}$。
对 $B^i = h^{i + 1}(U^{i + 1}) + s^i(V^i)$ 应用该断言。
所得映射 $s' : V^\bullet \to Y^\bullet$ 也是拟同构，且
$s' \circ f$ 零伦；这由 $h^i$ 分解通过 $Y^{i - 1}$ 显然可见。
这证明了忠实性。
% P02-E0039: frozen source says “Fully faithfulness”; the noun is “Full faithfulness”.
满忠实性也用完全相同的方法证明。
\end{proof}





\section{内射消解}
\label{derived-section-injective-resolutions}

\noindent
本节证明若干关于具有足够多内射对象的阿贝尔范畴中
内射消解存在性的引理。

\begin{definition}
\label{derived-definition-injective-resolution}
设 $\mathcal{A}$ 为阿贝尔范畴，$A \in \Ob(\mathcal{A})$。
$A$ 的一个{\it 内射消解}是复形 $I^\bullet$，
连同映射 $A \to I^0$，满足：
\begin{enumerate}
\item 对 $n < 0$ 有 $I^n = 0$。
\item 每个 $I^n$ 都是 $\mathcal{A}$ 的内射对象。
\item 映射 $A \to I^0$ 是到 $\Ker(d^0)$ 上的同构。
\item 对 $i > 0$ 有 $H^i(I^\bullet) = 0$。
\end{enumerate}
因此 $A[0] \to I^\bullet$ 是拟同构。
换言之，复形
$$
\ldots \to 0 \to A \to I^0 \to I^1 \to \ldots
$$
无环。设 $K^\bullet$ 为 $\mathcal{A}$ 中的复形。
$K^\bullet$ 的一个{\it 内射消解}是复形 $I^\bullet$，
连同复形映射 $\alpha : K^\bullet \to I^\bullet$，满足：
\begin{enumerate}
\item 对 $n \ll 0$ 有 $I^n = 0$；即 $I^\bullet$ 下有界。
\item 每个 $I^n$ 都是 $\mathcal{A}$ 的内射对象。
\item 映射 $\alpha : K^\bullet \to I^\bullet$ 是拟同构。
\end{enumerate}
\end{definition}

\noindent
换言之，内射消解 $K^\bullet \to I^\bullet$ 给出图
$$
\xymatrix{
\ldots \ar[r] & K^{n - 1} \ar[d] \ar[r] & K^n \ar[d] \ar[r] &
K^{n + 1} \ar[d] \ar[r] & \ldots \\
\ldots \ar[r] & I^{n - 1} \ar[r] & I^n \ar[r] & I^{n + 1} \ar[r] & \ldots
}
$$
它在每个次数的上同调对象上诱导同构。
$\mathcal{A}$ 的对象 $A$ 的内射消解，
与复形 $A[0]$ 的内射消解几乎是同一回事。

\begin{lemma}
\label{derived-lemma-cohomology-bounded-below}
设 $\mathcal{A}$ 为阿贝尔范畴，$K^\bullet$ 为 $\mathcal{A}$ 的复形。
\begin{enumerate}
\item 若 $K^\bullet$ 有内射消解，则当 $n \ll 0$ 时
$H^n(K^\bullet) = 0$。
\item 若对所有 $n \ll 0$ 有 $H^n(K^\bullet) = 0$，
则存在拟同构 $K^\bullet \to L^\bullet$，其中 $L^\bullet$ 下有界。
\end{enumerate}
\end{lemma}

\begin{proof}
从略。对第二个断言，取某个 $n \ll 0$，并使用
$L^\bullet = \tau_{\geq n}K^\bullet$。
截断 $\tau_{\geq n}$ 的定义见同调，章节
\ref{homology-section-truncations}。
\end{proof}

\begin{lemma}
\label{derived-lemma-injective-resolutions-exist}
设 $\mathcal{A}$ 为阿贝尔范畴。假设 $\mathcal{A}$ 有足够多内射对象。
\begin{enumerate}
\item $\mathcal{A}$ 的任意对象都有内射消解。
\item 若对所有 $n \ll 0$ 有 $H^n(K^\bullet) = 0$，
则 $K^\bullet$ 有内射消解。
\item 若 $K^\bullet$ 是满足 $n < a$ 时 $K^n = 0$ 的复形，
则存在内射消解 $\alpha : K^\bullet \to I^\bullet$，
使得 $n < a$ 时 $I^n = 0$，且每个
$\alpha^n : K^n \to I^n$ 都是单射。
\end{enumerate}
\end{lemma}

\begin{proof}
先证明 (1)。首先选取从 $A$ 到 $\mathcal{A}$ 的某个内射对象中的单射
$A \to I^0$。接着，选取到 $\mathcal{A}$ 的某个内射对象中的单射
% P02-E0041: frozen source uses I_0/A, inconsistent with the previously defined I^0.
$I_0/A \to I^1$。以 $d^0$ 表示所诱导的映射 $I^0 \to I^1$。
再选取到 $\mathcal{A}$ 的某个内射对象中的单射
$I^1/\Im(d^0) \to I^2$。以 $d^1$ 表示所诱导的映射
$I^1 \to I^2$，依此类推。
由引理 \ref{derived-lemma-cohomology-bounded-below}，(2) 由 (3) 推出。
(3) 是引理 \ref{derived-lemma-subcategory-right-resolution} 的特例。
\end{proof}

\begin{lemma}
\label{derived-lemma-acyclic-is-zero}
设 $\mathcal{A}$ 为阿贝尔范畴，$K^\bullet$ 为无环复形。
设 $I^\bullet$ 下有界并由内射对象组成。
则任意态射 $K^\bullet \to I^\bullet$ 都零伦。
\end{lemma}

\begin{proof}
设 $\alpha : K^\bullet \to I^\bullet$ 为复形态射。
注意，由于 $I^\bullet$ 下有界，当 $j \ll 0$ 时 $\alpha^j = 0$。
特别地，可以找到 $n$，使得对 $j \leq n$ 存在
$h^j :  K^j \to I^{j - 1}$，并且对 $j < n$ 有
$\alpha^j = d^{j - 1} \circ h^j + h^{j + 1} \circ d^j$。
我们将说明存在态射 $h^{n + 1} : K^{n + 1} \to I^n$，使得
$\alpha^n = d^{n - 1} \circ h^n + h^{n + 1} \circ d^n$。
注意
% P02-E0042: the frozen calculation has an alpha-index error and several h-index errors; see the ledger.
\begin{align*}
(\alpha^n - d^{n - 1} \circ h^{n - 1}) \circ d^{n - 1}
& =
\alpha^{n - 1} \circ d^{n - 1} - d^{n - 1} \circ h^{n - 1} \circ d^{n - 1} \\
& =
d^{n - 1} \circ \alpha^{n - 1} - d^{n - 1} \circ h^{n - 1} \circ d^{n - 1} \\
& =
d^{n - 1} \circ (d^{n - 2} \circ h^{n - 1} + h^n \circ d^{n - 1})
- d^{n - 1} \circ h^{n - 1} \circ d^{n - 1} \\
& = 0
\end{align*}
由于 $K^\bullet$ 无环，有
$d^{n - 1}(K^{n - 1}) = \Ker(K^n \to K^{n + 1})$。
所以可把 $\alpha^n -  d^{n - 1} \circ h^{n - 1}$ 看作
定义在 $K^{n + 1}$ 的子对象
$\Im(K^n \to K^{n + 1})$ 上、取值于 $I^n$ 的映射。
利用对象 $I^n$ 的内射性，可把它扩张成映射
$h^{n + 1} : K^{n + 1} \to I^n$。
读者可验证，在这一选择下确有
$\alpha^n = d^{n - 1} \circ h^n + h^{n + 1} \circ d^n$。

\medskip\noindent
对 $n$ 归纳可知，能够找到
$h  = (h^j)_{j \in \mathbf{Z}}$，它给出 $\alpha$ 与 $0$ 之间的同伦，
正合所需。
\end{proof}

\begin{remark}
\label{derived-remark-easier-proofs}
设 $\mathcal{A}$ 为阿贝尔范畴。
利用 $K(\mathcal{A})$ 是三角范畴这一事实，可以用引理
\ref{derived-lemma-acyclic-is-zero} 证明下面一些通常通过追图证明的引理。
具体而言，假设 $\alpha : K^\bullet \to L^\bullet$ 是复形的拟同构。
则
$$
(K^\bullet, L^\bullet, C(\alpha)^\bullet, \alpha, i, -p)
$$
是 $K(\mathcal{A})$ 中的特异三角
（引理 \ref{derived-lemma-the-same-up-to-isomorphisms}），而
$C(\alpha)^\bullet$ 是无环复形（引理 \ref{derived-lemma-acyclic}）。
再设 $I^\bullet$ 为由内射对象组成的下有界复形。则
$$
\xymatrix{
\Hom_{K(\mathcal{A})}(C(\alpha)^\bullet, I^\bullet) \ar[r] &
\Hom_{K(\mathcal{A})}(L^\bullet, I^\bullet) \ar[r] &
\Hom_{K(\mathcal{A})}(K^\bullet, I^\bullet) \ar[lld] \\
\Hom_{K(\mathcal{A})}(C(\alpha)^\bullet[-1], I^\bullet)
}
$$
是阿贝尔群的正合序列；见引理 \ref{derived-lemma-representable-homological}。
此时，引理 \ref{derived-lemma-acyclic-is-zero} 保证外侧两个群为零，故
$\Hom_{K(\mathcal{A})}(L^\bullet, I^\bullet) =
\Hom_{K(\mathcal{A})}(K^\bullet, I^\bullet)$。
\end{remark}

\begin{lemma}
\label{derived-lemma-morphisms-lift}
设 $\mathcal{A}$ 为阿贝尔范畴。考虑实线图
$$
\xymatrix{
K^\bullet \ar[r]_\alpha \ar[d]_\gamma & L^\bullet \ar@{-->}[dl]^\beta \\
I^\bullet
}
$$
其中 $I^\bullet$ 下有界且由内射对象组成，而 $\alpha$ 是拟同构。
\begin{enumerate}
\item 存在复形映射 $\beta$，使图交换至同伦。
\item 若 $\alpha$ 在每个次数上都是单射，则可找到使图交换的 $\beta$。
\end{enumerate}
\end{lemma}

\begin{proof}
(1) 的“正确”证明见注记 \ref{derived-remark-easier-proofs}。
这里也给出一个直接证明。

\medskip\noindent
先说明 (2) 推出 (1)。设
% P02-E0043: frozen source writes K rather than K^\bullet in the complex map.
$\tilde \alpha : K \to \tilde L^\bullet$、$\pi$、$s$
如引理 \ref{derived-lemma-make-injective} 中所述。
由于 $\tilde \alpha$ 为单射，由 (2) 存在态射
$\tilde \beta : \tilde L^\bullet \to I^\bullet$，使得
$\gamma = \tilde \beta \circ \tilde \alpha$。令
$\beta = \tilde \beta \circ s$。则
$$
\beta \circ \alpha
=
\tilde \beta \circ s \circ \pi \circ \tilde \alpha
\sim
\tilde \beta \circ \tilde \alpha
=
\gamma
$$
正合所需。

\medskip\noindent
假设 $\alpha : K^\bullet \to L^\bullet$ 为单射。
设已在所有次数 $\leq n - 1$ 上定义了与微分相容的 $\beta$，
且对所有 $j \leq n - 1$ 有
$\gamma^j = \beta^j \circ \alpha^j$。
考虑交换的实线图
$$
\xymatrix{
K^{n - 1} \ar[r] \ar@/_2pc/[dd]_\gamma \ar[d]^\alpha &
K^n \ar@/^2pc/[dd]^\gamma \ar[d]^\alpha \\
L^{n - 1} \ar[r] \ar[d]^\beta &
L^n \ar@{-->}[d] \\
I^{n - 1} \ar[r] &
I^n
}
$$
因此虚线箭头在子对象 $\alpha(K^n)$ 和
$d^{n - 1}(L^{n - 1})$ 上已有规定。
此外，这两个箭头在 $\alpha(d^{n - 1}(K^{n - 1}))$ 上相同。
所以，若
\begin{equation}
\label{derived-equation-qis}
\alpha(d^{n - 1}(K^{n - 1}))
=
\alpha(K^n) \cap d^{n - 1}(L^{n - 1})
\end{equation}
则这些态射可粘合成态射
$\alpha(K^n) + d^{n - 1}(L^{n - 1}) \to I^n$；
再利用 $I^n$ 的内射性，可把它扩张为从整个 $L^n$ 到 $I^n$ 的态射。
此后通过归纳，可同时对所有 $n$ 得到态射 $\beta$
（注意，因为 $I^\bullet$ 下有界，当 $n \ll 0$ 时可令
$\beta^n = 0$，从而启动归纳）。

\medskip\noindent
还须证明等式 (\ref{derived-equation-qis})。
建议读者用适当的追图自行论证；不过这里仍给出我们的论证。
包含关系
$\alpha(d^{n - 1}(K^{n - 1})) \subset \alpha(K^n) \cap d^{n - 1}(L^{n - 1})$
显然成立。取 $\mathcal{A}$ 的对象 $T$ 以及态射
$x : T \to L^n$，其像包含在子对象
$\alpha(K^n) \cap d^{n - 1}(L^{n - 1})$ 中。
由于 $\alpha$ 为单射，有 $x = \alpha \circ x'$，
其中 $x' : T \to K^n$。又因 $x$ 位于
$d^{n - 1}(L^{n - 1})$ 中，有 $d^n \circ x = 0$。
再次利用 $\alpha$ 的单射性，得到 $d^n \circ x' = 0$。
因此 $x'$ 给出态射 $[x'] : T \to H^n(K^\bullet)$。
另一方面，由 $x$ 诱导的相应映射
$[x] : T \to H^n(L^\bullet)$ 依假设为零。
由于 $\alpha$ 是拟同构，必有 $[x'] = 0$。
这恰好意味着 $x'$ 的像包含在
$d^{n - 1}(K^{n - 1})$ 中，证明完成。
\end{proof}

\begin{lemma}
\label{derived-lemma-morphisms-equal-up-to-homotopy}
设 $\mathcal{A}$ 为阿贝尔范畴。考虑实线图
$$
\xymatrix{
K^\bullet \ar[r]_\alpha \ar[d]_\gamma & L^\bullet \ar@{-->}[dl]^{\beta_i} \\
I^\bullet
}
$$
其中 $I^\bullet$ 下有界且由内射对象组成，而 $\alpha$ 是拟同构。
任意两个使图交换至同伦的态射 $\beta_1, \beta_2$ 都彼此同伦。
\end{lemma}

\begin{proof}
这由注记 \ref{derived-remark-easier-proofs} 推出。
这里也给出一个直接论证。

\medskip\noindent
% P02-E0043: this second frozen locus also writes K rather than K^\bullet.
设 $\tilde \alpha : K \to \tilde L^\bullet$、$\pi$、$s$
如引理 \ref{derived-lemma-make-injective} 中所述。
若能证明 $\beta_1 \circ\pi$ 与 $\beta_2 \circ \pi$ 同伦，
则由 $\pi \circ s$ 是恒等映射可推出
$\beta_1 \sim \beta_2$。因此可以假设，对所有 $n$，
$\alpha^n : K^n \to L^n$ 都是直和项的包含映射。
于是得到逐项分裂的复形短正合序列
$$
0 \to K^\bullet \to L^\bullet \to M^\bullet \to 0
$$
且 $M^\bullet$ 无环。选取分裂 $L^n = K^n \oplus M^n$，
于是有 $\beta_i^n : K^n \oplus M^n \to I^n$ 和
$\gamma^n : K^n \to I^n$。此时对 $\beta_i$ 的条件是：
存在态射 $h_i^n : K^n \to I^{n - 1}$，使得
$$
\gamma^n - \beta_i^n|_{K^n} = d \circ h_i^n + h_i^{n + 1} \circ d
$$
因此
$$
\beta_1^n|_{K^n} - \beta_2^n|_{K^n}
=
d \circ (h_1^n - h_2^n) + (h_1^{n + 1} - h_2^{n + 1}) \circ d
$$
考虑映射 $h^n : K^n \oplus M^n \to I^{n - 1}$，
它在第一个直和项上等于 $h_1^n - h_2^n$，在第二个直和项上为零。
于是可见
% P02-E0044: frozen formula has (d h^n + h^{n+1}) d rather than d h^n + h^{n+1} d.
$$
\beta_1^n - \beta_2^n
-
(d \circ h^n + h^{n + 1}) \circ d
$$
是复形态射 $L^\bullet \to I^\bullet$，
它在子复形 $K^\bullet$ 上恒为零。因此它分解为
$L^\bullet \to M^\bullet \to I^\bullet$。
所以本引理的结论由引理 \ref{derived-lemma-acyclic-is-zero} 推出。
\end{proof}

\begin{lemma}
\label{derived-lemma-morphisms-into-injective-complex}
设 $\mathcal{A}$ 为阿贝尔范畴。
设 $I^\bullet$ 为由内射对象组成的下有界复形，
$L^\bullet \in K(\mathcal{A})$。则
$$
\Mor_{K(\mathcal{A})}(L^\bullet, I^\bullet)
=
\Mor_{D(\mathcal{A})}(L^\bullet, I^\bullet).
$$
\end{lemma}

\begin{proof}
设 $a$ 为右边的元素。可以把它表示为
$a = \gamma\alpha^{-1}$，其中
$\alpha : K^\bullet \to L^\bullet$ 是拟同构，
$\gamma : K^\bullet \to I^\bullet$ 是复形映射。
由引理 \ref{derived-lemma-morphisms-lift}，可找到态射
$\beta : L^\bullet \to I^\bullet$，使得
$\beta \circ \alpha$ 与 $\gamma$ 同伦。这证明该映射为满射。
设 $b$ 为左边的元素，且其在右边的像为零。
那么 $b$ 是某个态射 $\beta : L^\bullet \to I^\bullet$ 的同伦类，
并存在拟同构 $\alpha : K^\bullet \to L^\bullet$，使得
$\beta  \circ \alpha$ 零伦。由引理
\ref{derived-lemma-morphisms-equal-up-to-homotopy}，
$\beta$ 本身也零伦；即 $b = 0$。
\end{proof}

\begin{lemma}
\label{derived-lemma-injective-resolution-ses}
设 $\mathcal{A}$ 为阿贝尔范畴。假设 $\mathcal{A}$ 有足够多内射对象。
对 $\text{Comp}^{+}(\mathcal{A})$ 的任意短正合序列
$0 \to A^\bullet \to B^\bullet \to C^\bullet \to 0$，
存在 $\text{Comp}^{+}(\mathcal{A})$ 中的交换图
$$
\xymatrix{
0 \ar[r] &
A^\bullet \ar[r] \ar[d] &
B^\bullet \ar[r] \ar[d] &
C^\bullet \ar[r] \ar[d] &
0 \\
0 \ar[r] &
I_1^\bullet \ar[r] &
I_2^\bullet \ar[r] &
I_3^\bullet \ar[r] &
0
}
$$
其中竖直箭头是内射消解，各行是复形的短正合序列。此外：
\begin{enumerate}
\item 给定任意内射消解 $A^\bullet \to I^\bullet$，
可以假设 $I_1^\bullet = I^\bullet$；
\item 若对 $n < 0$ 有 $A^n = B^n = C^n = 0$，
则可以假设对 $n < 0$ 有 $I_j^n = 0$；
\item 若还对 $n < 0$ 有 $I^n = 0$，则可合并 (1) 和 (2)；
% P02-E0045: frozen source contains the unfinished placeholder “add more here.”
\item 在此添加更多内容。
\end{enumerate}
\end{lemma}

\begin{proof}
步骤 1。选取内射消解 $A^\bullet \to I^\bullet$
（见引理 \ref{derived-lemma-injective-resolutions-exist}），或使用给定的消解。
回忆 $\text{Comp}^{+}(\mathcal{A})$ 是阿贝尔范畴；见同调，引理
\ref{homology-lemma-cat-cochain-abelian}。
因此可沿映射 $A^\bullet \to I^\bullet$ 作推出，得到
$$
\xymatrix{
0 \ar[r] &
A^\bullet \ar[r] \ar[d] &
B^\bullet \ar[r] \ar[d] &
C^\bullet \ar[r] \ar[d] &
0 \\
0 \ar[r] &
I^\bullet \ar[r] &
E^\bullet \ar[r] &
C^\bullet \ar[r] &
0
}
$$
由 $5$-引理和同调，引理
\ref{homology-lemma-long-exact-sequence-cochain} 的最后一个断言，
映射 $B^\bullet \to E^\bullet$ 是拟同构。
注意，下面的短正合序列逐项分裂；见同调，引理
\ref{homology-lemma-characterize-injectives}。
所以只须在
$0 \to A^\bullet \to B^\bullet \to C^\bullet \to 0$
逐项分裂时证明本引理。

\medskip\noindent
步骤 2。选取分裂。换言之，写成 $B^n = A^n \oplus C^n$。
以 $\delta : C^\bullet \to A^\bullet[1]$ 表示同调，引理
\ref{homology-lemma-ses-termwise-split-cochain} 中的态射。
选取内射消解 $f_1 : A^\bullet \to I_1^\bullet$ 和
$f_3 : C^\bullet \to I_3^\bullet$。
（若 $A^\bullet$ 是内射对象的复形，则使用
$I_1^\bullet = A^\bullet$。若对 $n < 0$ 有
$A^n = C^n = 0$，则由引理
\ref{derived-lemma-injective-resolutions-exist}，可选取消解使得
对 $n < 0$ 有 $I_1^n = I_3^n = 0$。）
可以假设 $f_3$ 在每个次数上都是单射。
由引理 \ref{derived-lemma-morphisms-lift}，可找到态射
$\delta' : I_3^\bullet \to I_1^\bullet[1]$，使得
$\delta' \circ f_3 = f_1[1] \circ \delta$
（复形态射的等式）。令 $I_2^n = I_1^n \oplus I_3^n$，定义
$$
d_{I_2}^n =
\left(
\begin{matrix}
d_{I_1}^n & (\delta')^n \\
0 & d_{I_3}^n
\end{matrix}
\right)
$$
并把映射 $B^n \to I_2^n$ 定义为映射
$A^n \to I_1^n$ 与 $C^n \to I_3^n$ 之和。
一切均显然。
\end{proof}








\section{投射消解}
\label{derived-section-projective-resolutions}

\noindent
本节与章节 \ref{derived-section-injective-resolutions} 对偶。
我们给出定义并陈述结果，但不再证明这些引理。

\begin{definition}
\label{derived-definition-projective-resolution}
设 $\mathcal{A}$ 为阿贝尔范畴，$A \in \Ob(\mathcal{A})$。
$A$ 的一个{\it 投射消解}是复形 $P^\bullet$，
连同映射 $P^0 \to A$，满足：
\begin{enumerate}
\item 对 $n > 0$ 有 $P^n = 0$。
\item 每个 $P^n$ 都是 $\mathcal{A}$ 的投射对象。
\item 映射 $P^0 \to A$ 诱导同构 $\Coker(d^{-1}) \to A$。
\item 对 $i < 0$ 有 $H^i(P^\bullet) = 0$。
\end{enumerate}
因此 $P^\bullet \to A[0]$ 是拟同构。
换言之，复形
$$
\ldots \to P^{-1} \to P^0 \to A \to 0 \to \ldots
$$
无环。设 $K^\bullet$ 为 $\mathcal{A}$ 中的复形。
$K^\bullet$ 的一个{\it 投射消解}是复形 $P^\bullet$，
连同复形映射 $\alpha : P^\bullet \to K^\bullet$，满足：
\begin{enumerate}
\item 对 $n \gg 0$ 有 $P^n = 0$；即 $P^\bullet$ 上有界。
\item 每个 $P^n$ 都是 $\mathcal{A}$ 的投射对象。
\item 映射 $\alpha : P^\bullet \to K^\bullet$ 是拟同构。
\end{enumerate}
\end{definition}

\begin{lemma}
\label{derived-lemma-cohomology-bounded-above}
设 $\mathcal{A}$ 为阿贝尔范畴，$K^\bullet$ 为 $\mathcal{A}$ 的复形。
\begin{enumerate}
\item 若 $K^\bullet$ 有投射消解，则当 $n \gg 0$ 时
$H^n(K^\bullet) = 0$。
\item 若当 $n \gg 0$ 时 $H^n(K^\bullet) = 0$，
则存在拟同构 $L^\bullet \to K^\bullet$，其中 $L^\bullet$ 上有界。
\end{enumerate}
\end{lemma}

\begin{proof}
与引理 \ref{derived-lemma-cohomology-bounded-below} 对偶。
\end{proof}

\begin{lemma}
\label{derived-lemma-projective-resolutions-exist}
设 $\mathcal{A}$ 为阿贝尔范畴。假设 $\mathcal{A}$ 有足够多投射对象。
\begin{enumerate}
\item $\mathcal{A}$ 的任意对象都有投射消解。
\item 若对所有 $n \gg 0$ 有 $H^n(K^\bullet) = 0$，
则 $K^\bullet$ 有投射消解。
\item 若 $K^\bullet$ 是满足 $n > a$ 时 $K^n = 0$ 的复形，
则存在投射消解 $\alpha : P^\bullet \to K^\bullet$，
使得 $n > a$ 时 $P^n = 0$，且每个
$\alpha^n : P^n \to K^n$ 都为满射。
\end{enumerate}
\end{lemma}

\begin{proof}
与引理 \ref{derived-lemma-injective-resolutions-exist} 对偶。
\end{proof}

\begin{lemma}
\label{derived-lemma-projective-into-acyclic-is-zero}
设 $\mathcal{A}$ 为阿贝尔范畴，$K^\bullet$ 为无环复形。
设 $P^\bullet$ 上有界并由投射对象组成。
则任意态射 $P^\bullet \to K^\bullet$ 都零伦。
\end{lemma}

\begin{proof}
与引理 \ref{derived-lemma-acyclic-is-zero} 对偶。
\end{proof}

\begin{remark}
\label{derived-remark-easier-projective}
设 $\mathcal{A}$ 为阿贝尔范畴。
假设 $\alpha : K^\bullet \to L^\bullet$ 是复形的拟同构，
$P^\bullet$ 是投射对象的上有界复形。则
$$
\Hom_{K(\mathcal{A})}(P^\bullet, K^\bullet)
\longrightarrow
\Hom_{K(\mathcal{A})}(P^\bullet, L^\bullet)
$$
是同构。这与注记 \ref{derived-remark-easier-proofs} 对偶。
\end{remark}

\begin{lemma}
\label{derived-lemma-morphisms-lift-projective}
设 $\mathcal{A}$ 为阿贝尔范畴。考虑实线图
$$
\xymatrix{
K^\bullet & L^\bullet \ar[l]^\alpha \\
P^\bullet \ar[u] \ar@{-->}[ru]_\beta
}
$$
其中 $P^\bullet$ 上有界且由投射对象组成，而 $\alpha$ 是拟同构。
\begin{enumerate}
\item 存在复形映射 $\beta$，使图交换至同伦。
\item 若 $\alpha$ 在每个次数上都是满射，则可找到使图交换的 $\beta$。
\end{enumerate}
\end{lemma}

\begin{proof}
与引理 \ref{derived-lemma-morphisms-lift} 对偶。
\end{proof}

\begin{lemma}
\label{derived-lemma-morphisms-equal-up-to-homotopy-projective}
设 $\mathcal{A}$ 为阿贝尔范畴。考虑实线图
$$
\xymatrix{
K^\bullet & L^\bullet \ar[l]^\alpha \\
P^\bullet \ar[u] \ar@{-->}[ru]_{\beta_i}
}
$$
其中 $P^\bullet$ 上有界且由投射对象组成，而 $\alpha$ 是拟同构。
任意两个使图交换至同伦的态射 $\beta_1, \beta_2$ 都彼此同伦。
\end{lemma}

\begin{proof}
与引理 \ref{derived-lemma-morphisms-equal-up-to-homotopy} 对偶。
\end{proof}

\begin{lemma}
\label{derived-lemma-morphisms-from-projective-complex}
设 $\mathcal{A}$ 为阿贝尔范畴。
设 $P^\bullet$ 为由投射对象组成的上有界复形，
$L^\bullet \in K(\mathcal{A})$。则
$$
\Mor_{K(\mathcal{A})}(P^\bullet, L^\bullet)
=
\Mor_{D(\mathcal{A})}(P^\bullet, L^\bullet).
$$
\end{lemma}

\begin{proof}
与引理 \ref{derived-lemma-morphisms-into-injective-complex} 对偶。
\end{proof}

\begin{lemma}
\label{derived-lemma-projective-resolution-ses}
设 $\mathcal{A}$ 为阿贝尔范畴。假设 $\mathcal{A}$ 有足够多投射对象。
% P02-E0046: the projective dual is bounded above, but frozen source names Comp^+ rather than Comp^- here and below.
对 $\text{Comp}^{+}(\mathcal{A})$ 的任意短正合序列
$0 \to A^\bullet \to B^\bullet \to C^\bullet \to 0$，
存在 $\text{Comp}^{+}(\mathcal{A})$ 中的交换图
$$
\xymatrix{
0 \ar[r] &
P_1^\bullet \ar[r] \ar[d] &
P_2^\bullet \ar[r] \ar[d] &
P_3^\bullet \ar[r] \ar[d] &
0 \\
0 \ar[r] &
A^\bullet \ar[r] &
B^\bullet \ar[r] &
C^\bullet \ar[r] &
0
}
$$
其中竖直箭头为投射消解，各行是复形的短正合序列。
事实上，给定任意投射消解 $P^\bullet \to C^\bullet$，
可以假设 $P_3^\bullet = P^\bullet$。
\end{lemma}

\begin{proof}
与引理 \ref{derived-lemma-injective-resolution-ses} 对偶。
\end{proof}

\begin{lemma}
\label{derived-lemma-precise-vanishing}
设 $\mathcal{A}$ 为阿贝尔范畴，$P^\bullet$、$K^\bullet$ 为复形，
$n \in \mathbf{Z}$。假设：
\begin{enumerate}
\item $P^\bullet$ 是由投射对象组成的有界复形；
\item 对 $i < n$ 有 $P^i = 0$；
\item 对 $i \geq n$ 有 $H^i(K^\bullet) = 0$。
\end{enumerate}
则
$\Hom_{K(\mathcal{A})}(P^\bullet, K^\bullet) =
\Hom_{D(\mathcal{A})}(P^\bullet, K^\bullet) = 0$。
\end{lemma}

\begin{proof}
第一个等号由引理 \ref{derived-lemma-morphisms-from-projective-complex} 推出。
注意，由注记 \ref{derived-remark-truncation-distinguished-triangle}，有特异三角
$$
(\tau_{\leq n - 1}K^\bullet, K^\bullet, \tau_{\geq n}K^\bullet, f, g, h)
$$
因此，由引理 \ref{derived-lemma-representable-homological}，只须证明
$\Hom_{K(\mathcal{A})}(P^\bullet, \tau_{\leq n - 1}K^\bullet) = 0$ 和
$\Hom_{K(\mathcal{A})}(P^\bullet, \tau_{\geq n} K^\bullet) = 0$。
第一个消失结论显然，第二个就是引理
\ref{derived-lemma-projective-into-acyclic-is-zero}。
\end{proof}

\begin{lemma}
\label{derived-lemma-lift-map}
设 $\mathcal{A}$ 为阿贝尔范畴。设
$\beta : P^\bullet \to L^\bullet$ 和
$\alpha : E^\bullet \to L^\bullet$ 为复形映射，
$n \in \mathbf{Z}$。假设：
\begin{enumerate}
\item $P^\bullet$ 是投射对象的有界复形，且对 $i < n$ 有 $P^i = 0$；
\item 当 $i > n$ 时 $H^i(\alpha)$ 为同构，
当 $i = n$ 时为满射。
\end{enumerate}
则存在复形映射 $\gamma : P^\bullet \to E^\bullet$，
使得 $\alpha \circ \gamma$ 与 $\beta$ 同伦。
\end{lemma}

\begin{proof}
考虑带有映射 $i : L^\bullet \to C^\bullet$ 的锥
$C^\bullet = C(\alpha)^\bullet$。
由引理 \ref{derived-lemma-precise-vanishing}，$i \circ \beta$ 为零。
所以由引理 \ref{derived-lemma-representable-homological}，
可把 $\beta$ 提升到 $E^\bullet$。
\end{proof}













\section{右导出函子与内射消解}
\label{derived-section-right-derived-functor}

\noindent
现在可以利用上述材料，定义从具有足够多内射对象的阿贝尔范畴
到一般阿贝尔范畴的加性函子的右导出函子。

\begin{lemma}
\label{derived-lemma-injective-acyclic}
设 $\mathcal{A}$ 为阿贝尔范畴，
$I \in \Ob(\mathcal{A})$ 为内射对象，
$I^\bullet$ 为 $\mathcal{A}$ 中内射对象的下有界复形。
\begin{enumerate}
\item 对任意到任意三角范畴 $\mathcal{D}$ 的正合函子
$F : K^{+}(\mathcal{A}) \to \mathcal{D}$，
$I^\bullet$ 都相对于 $\text{Qis}^{+}(\mathcal{A})$ 计算 $RF$。
\item 对任意到任意阿贝尔范畴 $\mathcal{B}$ 的加性函子
$F : \mathcal{A} \to \mathcal{B}$，$I$ 都右无环。
\end{enumerate}
\end{lemma}

\begin{proof}
% P02-E0047: frozen source says “a direct consequences”.
(2) 是 (1) 和定义 \ref{derived-definition-derived-functor} 的直接推论。
为证明 (1)，设 $\alpha : I^\bullet \to K^\bullet$ 为到某复形中的拟同构。
由引理 \ref{derived-lemma-morphisms-lift}，$\alpha$ 有左逆。
所以范畴 $I^\bullet/\text{Qis}^{+}(\mathcal{A})$
本质常值，值为 $\text{id} : I^\bullet \to I^\bullet$。
因而 ind-对象
$$
I^\bullet/\text{Qis}^{+}(\mathcal{A}) \longrightarrow \mathcal{D}, \quad
(I^\bullet \to K^\bullet) \longmapsto F(K^\bullet)
$$
也本质常值，值为 $F(I^\bullet)$。这证明了 (1)；见定义
\ref{derived-definition-right-derived-functor-defined} 和
\ref{derived-definition-computes}。
\end{proof}

\begin{lemma}
\label{derived-lemma-enough-injectives-right-derived}
设 $\mathcal{A}$ 为具有足够多内射对象的阿贝尔范畴。
\begin{enumerate}
\item 对任意到三角范畴 $\mathcal{D}$ 的正合函子
$F : K^{+}(\mathcal{A}) \to \mathcal{D}$，右导出函子
$$
RF : D^{+}(\mathcal{A}) \longrightarrow \mathcal{D}
$$
处处有定义。
\item 对任意到阿贝尔范畴 $\mathcal{B}$ 的加性函子
$F : \mathcal{A} \to \mathcal{B}$，右导出函子
$$
RF : D^{+}(\mathcal{A}) \longrightarrow D^{+}(\mathcal{B})
$$
处处有定义。
\end{enumerate}
\end{lemma}

\begin{proof}
对第二个断言，结合引理 \ref{derived-lemma-injective-acyclic} 和
命题 \ref{derived-proposition-enough-acyclics}。
为得到第一个断言，结合引理 \ref{derived-lemma-injective-resolutions-exist}、
引理 \ref{derived-lemma-injective-acyclic}、
引理 \ref{derived-lemma-existence-computes} 和等式
(\ref{derived-equation-everywhere})。
\end{proof}

\begin{lemma}
\label{derived-lemma-right-derived-properties}
设 $\mathcal{A}$ 为具有足够多内射对象的阿贝尔范畴，
$F : \mathcal{A} \to \mathcal{B}$ 为加性函子。
\begin{enumerate}
\item 函子 $RF$ 是正合函子
$D^{+}(\mathcal{A}) \to D^{+}(\mathcal{B})$。
\item 函子 $RF$ 诱导正合函子
$K^{+}(\mathcal{A}) \to D^{+}(\mathcal{B})$。
\item 函子 $RF$ 诱导 $\delta$-函子
$\text{Comp}^{+}(\mathcal{A}) \to D^{+}(\mathcal{B})$。
\item 函子 $RF$ 诱导 $\delta$-函子
$\mathcal{A} \to D^{+}(\mathcal{B})$。
\end{enumerate}
\end{lemma}

\begin{proof}
本引理只是回顾目前得到的一些结果。注意，由引理
\ref{derived-lemma-enough-injectives-right-derived}，$RF$ 处处有定义。
下面给出一些参考：
\begin{enumerate}
\item 导出函子正合：这归结为引理 \ref{derived-lemma-2-out-of-3-defined}。
\item 这是因为 $K^{+}(\mathcal{A}) \to D^{+}(\mathcal{A})$
正合，而正合函子的复合正合。
\item 这是因为 $\text{Comp}^{+}(\mathcal{A}) \to D^{+}(\mathcal{A})$
是 $\delta$-函子；见引理 \ref{derived-lemma-derived-canonical-delta-functor}。
\item 这是因为 $\mathcal{A} \to \text{Comp}^{+}(\mathcal{A})$
正合，而以正合函子预复合 $\delta$-函子仍得到 $\delta$-函子。
\end{enumerate}
\end{proof}

\begin{lemma}
\label{derived-lemma-higher-derived-functors}
设 $\mathcal{A}$ 为具有足够多内射对象的阿贝尔范畴，
$F : \mathcal{A} \to \mathcal{B}$ 为左正合函子。
\begin{enumerate}
\item 对 $\text{Comp}^{+}(\mathcal{A})$ 中复形的任意短正合序列
$0 \to A^\bullet \to B^\bullet \to C^\bullet \to 0$，
有相伴的长正合序列
$$
\ldots \to
H^i(RF(A^\bullet)) \to
H^i(RF(B^\bullet)) \to
H^i(RF(C^\bullet)) \to
H^{i + 1}(RF(A^\bullet)) \to \ldots
$$
\item 函子 $R^iF : \mathcal{A} \to \mathcal{B}$ 在 $i < 0$ 时为零。
并且 $R^0F = F : \mathcal{A} \to \mathcal{B}$。
\item 若 $i > 0$ 且 $I$ 内射，则 $R^iF(I) = 0$。
\item 序列 $(R^iF, \delta)$ 构成从 $\mathcal{A}$ 到 $\mathcal{B}$ 的
泛 $\delta$-函子（见同调，定义
\ref{homology-definition-universal-delta-functor}）。
\end{enumerate}
\end{lemma}

\begin{proof}
本引理只是回顾目前得到的一些结果。注意，由引理
\ref{derived-lemma-enough-injectives-right-derived}，$RF$ 处处有定义。
下面给出一些参考：
\begin{enumerate}
\item 这由引理 \ref{derived-lemma-right-derived-properties} 的 (3) 与
$D^{+}(\mathcal{B})$ 的长正合上同调序列
(\ref{derived-equation-long-exact-cohomology-sequence-D}) 结合得到。
\item 这就是引理 \ref{derived-lemma-left-exact-higher-derived}。
\item 这就是内射对象无环这一事实。
\item 这就是引理 \ref{derived-lemma-right-derived-delta-functor}。
\end{enumerate}
\end{proof}





\section{Cartan--Eilenberg 消解}
\label{derived-section-cartan-eilenberg}

\noindent
本节还可以扩充。这里的材料可以推广并应用于更多情形。
消解不必使用内射对象，而且在某些情形下该方法也适用于无界情形。

\begin{definition}
\label{derived-definition-cartan-eilenberg}
设 $\mathcal{A}$ 为阿贝尔范畴，$K^\bullet$ 为下有界复形。
$K^\bullet$ 的一个{\it Cartan--Eilenberg 消解}由双复形
$I^{\bullet, \bullet}$ 和复形态射
$\epsilon : K^\bullet \to I^{\bullet, 0}$ 给出，满足以下性质：
\begin{enumerate}
% P02-E0048: frozen source has the article error “There exists a i”.
\item 存在 $i \ll 0$，使得对所有 $p < i$ 和所有 $q$，
$I^{p, q} = 0$。
\item 若 $q < 0$，则 $I^{p, q} = 0$。
\item 复形 $I^{p, \bullet}$ 是 $K^p$ 的内射消解。
\item 复形 $\Ker(d_1^{p, \bullet})$ 是 $\Ker(d_K^p)$ 的内射消解。
\item 复形 $\Im(d_1^{p, \bullet})$ 是 $\Im(d_K^p)$ 的内射消解。
\item 复形 $H^p_I(I^{\bullet, \bullet})$ 是
$H^p(K^\bullet)$ 的内射消解。
\end{enumerate}
\end{definition}

\begin{lemma}
\label{derived-lemma-cartan-eilenberg}
设 $\mathcal{A}$ 为具有足够多内射对象的阿贝尔范畴，
$K^\bullet$ 为下有界复形。则 $K^\bullet$ 有 Cartan--Eilenberg 消解。
\end{lemma}

\begin{proof}
假设 $p < n$ 时 $K^p = 0$。如下把 $K^\bullet$ 分解成短正合序列：
令 $Z^p = \Ker(d^p)$、$B^p = \Im(d^{p - 1})$、
$H^p = Z^p/B^p$，并考虑
$$
\begin{matrix}
0 \to Z^n \to K^n \to B^{n + 1} \to 0 \\
0 \to B^{n + 1} \to Z^{n + 1} \to H^{n + 1} \to 0 \\
0 \to Z^{n + 1} \to K^{n + 1} \to B^{n + 2} \to 0 \\
0 \to B^{n + 2} \to Z^{n + 2} \to H^{n + 2} \to 0 \\
\ldots
\end{matrix}
$$
对 $p < n$ 令 $I^{p, q} = 0$。按如下方式归纳选取内射消解：
\begin{enumerate}
\item 选取内射消解 $Z^n \to J_Z^{n, \bullet}$，
使得 $m < 0$ 时 $J_Z^{n, m} = 0$；见引理
\ref{derived-lemma-injective-resolutions-exist}。
\item 利用引理 \ref{derived-lemma-injective-resolution-ses}，
选取内射消解 $K^n \to I^{n, \bullet}$、
$B^{n + 1} \to J_B^{n + 1, \bullet}$，使得
$m < 0$ 时 $I^{n, m} = 0$，且 $m < 0$ 时 $J_B^{n + 1, m} = 0$；
并选取复形的正合序列
$0 \to J_Z^{n, \bullet} \to I^{n, \bullet} \to J_B^{n + 1, \bullet} \to 0$，
使之与短正合序列
$0 \to Z^n \to K^n \to B^{n + 1} \to 0$ 相容。
\item 利用引理 \ref{derived-lemma-injective-resolution-ses}，
选取内射消解 $Z^{n + 1} \to J_Z^{n + 1, \bullet}$、
$H^{n + 1} \to J_H^{n + 1, \bullet}$，使得
$m < 0$ 时 $J_Z^{n + 1, m} = 0$，且 $m < 0$ 时
$J_H^{n + 1, m} = 0$；并选取复形的正合序列
$0 \to J_B^{n + 1, \bullet} \to J_Z^{n + 1, \bullet}
\to J_H^{n + 1, \bullet} \to 0$，
使之与短正合序列
$0 \to B^{n + 1} \to Z^{n + 1} \to H^{n + 1} \to 0$ 相容。
\item 依此类推。
\end{enumerate}
把映射 $d_1^\bullet : I^{p, \bullet} \to I^{p + 1, \bullet}$
取为复合
$I^{p, \bullet} \to J_B^{p + 1, \bullet} \to
J_Z^{p + 1, \bullet} \to I^{p + 1, \bullet}$，则一切均显然。
\end{proof}

\begin{lemma}
\label{derived-lemma-two-ss-complex-functor}
设 $F : \mathcal{A} \to \mathcal{B}$ 为阿贝尔范畴之间的左正合函子。
设 $K^\bullet$ 为 $\mathcal{A}$ 的下有界复形，
$I^{\bullet, \bullet}$ 为 $K^\bullet$ 的 Cartan--Eilenberg 消解。
与双复形 $F(I^{\bullet, \bullet})$ 相伴的谱序列
$({}'E_r, {}'d_r)_{r \geq 0}$ 和 $({}''E_r, {}''d_r)_{r \geq 0}$
满足关系
$$
{}'E_1^{p, q} = R^qF(K^p)
\quad
\text{且}
\quad
{}''E_2^{p, q} = R^pF(H^q(K^\bullet))
$$
此外，这些谱序列有界，收敛到 $H^*(RF(K^\bullet))$，
且在 $H^n(RF(K^\bullet))$ 上相伴诱导的过滤是有限的。
\end{lemma}

\begin{proof}
下面这些注记将在不另加说明的情况下使用：
\begin{enumerate}
\item 由于 $I^{p, \bullet}$ 是 $K^p$ 的内射消解，
可知 $RF$ 在 $K^p[0]$ 处有定义，值为 $F(I^{p, \bullet})$。
\item 由于 $H^p_I(I^{\bullet, \bullet})$ 是
$H^p(K^\bullet)$ 的内射消解，导出函子 $RF$ 在
$H^p(K^\bullet)[0]$ 处有定义，值为
$F(H^p_I(I^{\bullet, \bullet}))$。
\item 由同调，引理 \ref{homology-lemma-double-complex-gives-resolution}，
全复形 $\text{Tot}(I^{\bullet, \bullet})$ 是 $K^\bullet$ 的内射消解。
所以 $RF$ 在 $K^\bullet$ 处有定义，值为
$F(\text{Tot}(I^{\bullet, \bullet}))$。
\end{enumerate}
考虑与双复形 $L^{\bullet, \bullet} = F(I^{\bullet, \bullet})$ 相伴的
两个谱序列；见同调，引理 \ref{homology-lemma-ss-double-complex}。
二者都有界，收敛到 $H^*(\text{Tot}(L^{\bullet, \bullet}))$，
并在 $H^n(\text{Tot}(L^{\bullet, \bullet}))$ 上诱导有限过滤；见
同调，引理 \ref{homology-lemma-first-quadrant-ss}。
由于
$\text{Tot}(L^{\bullet, \bullet}) =
\text{Tot}(F(I^{\bullet, \bullet})) =
F(\text{Tot}(I^{\bullet, \bullet}))$ 计算 $RF(K^\bullet)$，
便得到本引理的最后一个断言。

\medskip\noindent
计算第一个谱序列。由于
${}'E_1^{p, q} = H^q(L^{p, \bullet})$，换言之，
$$
{}'E_1^{p, q} = H^q(F(I^{p, \bullet})) = R^qF(K^p)
$$
正合所需。留意以后将用到：映射
${}'d_1^{p, q} : {}'E_1^{p, q} \to {}'E_1^{p + 1, q}$
就是由 $K^p \to K^{p + 1}$ 以及 $R^qF$ 是函子这一事实所诱导的映射
$R^qF(K^p) \to R^qF(K^{p + 1})$。

\medskip\noindent
计算第二个谱序列。有
${}''E_1^{p, q} = H^q(L^{\bullet, p}) = H^q(F(I^{\bullet, p}))$。
注意，复形 $I^{\bullet, p}$ 下有界，由内射对象组成，
而且各微分的每个核、像和上同调群都是 $\mathcal{A}$ 的内射对象。
所以可以分裂这些微分；即每个微分都是到某个直和项上的分裂满射。
应用 $F$ 后亦然。因此
${}''E_1^{p, q} = F(H^q(I^{\bullet, p})) = F(H^q_I(I^{\bullet, p}))$。
其上的微分是 $(-1)^q$ 乘以 $F$ 作用于复形
% P02-E0049: frozen source says H_I^p here; the E_2 target requires the resolution of H^q(K^\bullet).
$H^p_I(I^{\bullet, \bullet})$ 的微分，而该复形是
$H^p(K^\bullet)$ 的内射消解。由此得到对 $E_2$ 项的描述。
\end{proof}

\begin{remark}
\label{derived-remark-functorial-ss}
引理 \ref{derived-lemma-two-ss-complex-functor} 的谱序列关于复形
$K^\bullet$ 具有函子性。这由 Cartan--Eilenberg 消解的函子性性质推出。
另一方面，二者都是一种更一般的谱序列的例子；该谱序列可与
$\mathcal{A}$ 的过滤复形相伴。函子性将由其构造推出。
我们将在过滤导出范畴一节回到这一点；见注记
\ref{derived-remark-final-functorial}。
\end{remark}











\section{右导出函子的复合}
\label{derived-section-composition-right-derived-functors}

\noindent
有时可以计算复合的右导出函子。
% P02-E0050: frozen source says “Suppose ... be abelian categories”.
设 $\mathcal{A}, \mathcal{B}, \mathcal{C}$ 为阿贝尔范畴。
设 $F : \mathcal{A} \to \mathcal{B}$ 和
$G : \mathcal{B} \to \mathcal{C}$ 为左正合函子。
假设右导出函子
$RF : D^{+}(\mathcal{A}) \to D^{+}(\mathcal{B})$、
$RG : D^{+}(\mathcal{B}) \to D^{+}(\mathcal{C})$ 和
$R(G \circ F) : D^{+}(\mathcal{A}) \to D^{+}(\mathcal{C})$
处处有定义。那么存在从 $D^{+}(\mathcal{A})$ 到
$D^{+}(\mathcal{C})$ 的函子之间的典范变换
$$
t : R(G \circ F) \longrightarrow RG \circ RF
$$
见引理 \ref{derived-lemma-compose-derived-functors-general}。
此变换未必总是同构。

\begin{lemma}
\label{derived-lemma-compose-derived-functors}
设 $\mathcal{A}, \mathcal{B}, \mathcal{C}$ 为阿贝尔范畴。
设 $F : \mathcal{A} \to \mathcal{B}$ 和
$G : \mathcal{B} \to \mathcal{C}$ 为左正合函子。
假设 $\mathcal{A}$、$\mathcal{B}$ 有足够多内射对象。
下列条件等价：
\begin{enumerate}
\item 对 $\mathcal{A}$ 的每个内射对象 $I$，$F(I)$ 都对 $G$ 右无环；
\item 典范映射
$$
t : R(G \circ F) \longrightarrow RG \circ RF.
$$
% P02-E0051: frozen source has a premature display period and omits the article in “is isomorphism”.
是从 $D^{+}(\mathcal{A})$ 到 $D^{+}(\mathcal{C})$ 的函子同构。
\end{enumerate}
\end{lemma}

\begin{proof}
若 (2) 成立，则把同构 $t$ 在 $RF(I) = F(I)$ 处取值即可得到 (1)。
反之，假设 (1) 成立。设 $A^\bullet$ 为 $\mathcal{A}$ 的下有界复形，
选取内射消解 $A^\bullet \to I^\bullet$。
映射 $t$ 由如下映射给出（见引理
\ref{derived-lemma-compose-derived-functors-general} 的证明）：
% P02-E0052: frozen source has an extra closing parenthesis in G(F(I^\bullet))).
$$
R(G \circ F)(A^\bullet) =
(G \circ F)(I^\bullet) =
G(F(I^\bullet))) \to
RG(F(I^\bullet)) =
RG(RF(A^\bullet))
$$
其中箭头由引理 \ref{derived-lemma-leray-acyclicity} 是同构。
\end{proof}

\begin{lemma}[Grothendieck 谱序列]
\label{derived-lemma-grothendieck-spectral-sequence}
% P02-E0053: frozen opening is a sentence fragment (“With assumptions ... and assuming ...”).
在引理 \ref{derived-lemma-compose-derived-functors} 的假设下，
并假定其中等价条件 (1) 和 (2) 成立。
设 $X$ 为 $D^{+}(\mathcal{A})$ 的对象。
存在由 $\mathcal{C}$ 的双分次对象 $E_r$ 组成的谱序列
$(E_r, d_r)_{r \geq 0}$，其中 $d_r$ 的双次数为
$(r, - r + 1)$，并且
$$
E_2^{p, q} = R^pG(H^q(RF(X)))
$$
此外，该谱序列有界，收敛到 $H^*(R(G \circ F)(X))$，
并在每个 $H^n(R(G \circ F)(X))$ 上诱导有限过滤。
\end{lemma}

\noindent
对 $\mathcal{A}$ 的对象 $A$，得到
$E_2^{p, q} = R^pG(R^qF(A))$，它收敛到
$R^{p + q}(G \circ F)(A)$。

\begin{proof}
可以用下有界复形 $A^\bullet$ 表示 $X$。
选取内射消解 $A^\bullet \to I^\bullet$。
利用引理 \ref{derived-lemma-cartan-eilenberg}，选取 Cartan--Eilenberg 消解
$F(I^\bullet) \to I^{\bullet, \bullet}$。
应用引理 \ref{derived-lemma-two-ss-complex-functor} 的第二个谱序列。
\end{proof}






\section{消解函子}
\label{derived-section-derived-category}

\noindent
设 $\mathcal{A}$ 为具有足够多内射对象的阿贝尔范畴。
以 $\mathcal{I}$ 表示 $\mathcal{A}$ 的满加性子范畴，
其对象是 $\mathcal{A}$ 的内射对象。
在这种情形下，$K^{+}(\mathcal{I})$ 与 $D^{+}(\mathcal{A})$ 等价
（见命题 \ref{derived-proposition-derived-category}）。
所以，为了许多目的，可以把 $D^{+}(\mathcal{A})$ 看作
（较易理解的）范畴 $K^{+}(\mathcal{I})$。

\begin{proposition}
\label{derived-proposition-derived-category}
设 $\mathcal{A}$ 为阿贝尔范畴。假设 $\mathcal{A}$ 有足够多内射对象。
以 $\mathcal{I} \subset \mathcal{A}$ 表示严格满加性子范畴，
其对象是 $\mathcal{A}$ 的内射对象。函子
$$
K^{+}(\mathcal{I}) \longrightarrow D^{+}(\mathcal{A})
$$
正合、满忠实且本质满；即它是三角范畴的等价。
\end{proposition}

\begin{proof}
该函子显然正合。由引理 \ref{derived-lemma-injective-resolutions-exist}，
它本质满。满忠实性由引理
% P02-E0039: another frozen “Fully faithfulness” locus; the noun is “Full faithfulness”.
\ref{derived-lemma-morphisms-into-injective-complex} 推出。
\end{proof}

\noindent
命题 \ref{derived-proposition-derived-category} 表明可以找到消解函子。
事实上，即使在阿贝尔范畴 $\mathcal{A}$ 是“大”范畴
（即具有对象的类）的某些情形下，也能证明消解函子存在。

\begin{definition}
\label{derived-definition-localization-functor}
设 $\mathcal{A}$ 为具有足够多内射对象的阿贝尔范畴。
$\mathcal{A}$ 的一个{\it 消解函子}\footnote{这很可能不是标准术语。}
由以下数据给出：
\begin{enumerate}
\item 对所有 $K^\bullet \in \Ob(K^{+}(\mathcal{A}))$，
给定内射对象的下有界复形 $j(K^\bullet)$；
\item 对所有 $K^\bullet \in \Ob(K^{+}(\mathcal{A}))$，
给定拟同构 $i_{K^\bullet} : K^\bullet \to j(K^\bullet)$。
\end{enumerate}
\end{definition}

\begin{lemma}
\label{derived-lemma-resolution-functor}
设 $\mathcal{A}$ 为具有足够多内射对象的阿贝尔范畴。
给定消解函子 $(j, i)$，存在唯一方式把 $j$ 变成函子，
把 $i$ 变成 $2$-同构，从而产生 $2$-交换图
$$
\xymatrix{
K^{+}(\mathcal{A}) \ar[rd] \ar[rr]_j & & K^{+}(\mathcal{I}) \ar[ld] \\
& D^{+}(\mathcal{A})
}
$$
其中 $\mathcal{I}$ 是 $\mathcal{A}$ 的满加性子范畴，
由内射对象组成。
\end{lemma}

\begin{proof}
对 $K^{+}(\mathcal{A})$ 的每个态射
$\alpha : K^\bullet \to L^\bullet$，存在唯一的
$K^{+}(\mathcal{I})$ 中的态射
$j(\alpha) : j(K^\bullet) \to j(L^\bullet)$，使得图
$$
\xymatrix{
K^\bullet \ar[r]_\alpha \ar[d]_{i_{K^\bullet}} &
L^\bullet \ar[d]^{i_{L^\bullet}} \\
j(K^\bullet) \ar[r]^{j(\alpha)} & j(L^\bullet)
}
$$
在 $K^{+}(\mathcal{A})$ 中交换。为此，可以使用引理
\ref{derived-lemma-morphisms-lift} 和
\ref{derived-lemma-morphisms-equal-up-to-homotopy}，或等价的引理
\ref{derived-lemma-morphisms-into-injective-complex}。
唯一性意味着 $j$ 是函子，而图的交换性意味着 $i$ 给出
$2$-态射，它见证本引理所述范畴图的 $2$-交换性。
\end{proof}

\begin{lemma}
\label{derived-lemma-into-derived-category}
设 $\mathcal{A}$ 为阿贝尔范畴。假设 $\mathcal{A}$ 有足够多内射对象。
则消解函子 $j$ 存在，并且在函子的唯一同构意义下唯一。
\end{lemma}

\begin{proof}
考虑 $K^{+}(\mathcal{A})$ 的所有对象 $K^\bullet$ 所成的集合。
（回忆依我们的约定，除非另有说明，任意范畴都有对象的集合。）
由引理 \ref{derived-lemma-injective-resolutions-exist}，每个对象都有内射消解。
由选择公理，可以对每个 $K^\bullet$ 选取内射消解
$i_{K^\bullet} : K^\bullet \to j(K^\bullet)$。
\end{proof}

\begin{lemma}
\label{derived-lemma-j-is-exact}
设 $\mathcal{A}$ 为具有足够多内射对象的阿贝尔范畴。
任意消解函子
$j : K^{+}(\mathcal{A}) \to K^{+}(\mathcal{I})$
都是正合的。
\end{lemma}

\begin{proof}
以 $i_{K^\bullet} : K^\bullet \to j(K^\bullet)$ 表示定义
\ref{derived-definition-localization-functor} 的典范映射。
先讨论函子同构 $j(K^\bullet[1]) \to j(K^\bullet)[1]$ 的存在性。
考虑图
$$
\xymatrix{
K^\bullet[1] \ar[d]^{i_{K^\bullet[1]}} \ar@{=}[rr] & &
K^\bullet[1] \ar[d]^{i_{K^\bullet}[1]} \\
j(K^\bullet[1]) \ar@{..>}[rr]^{\xi_{K^\bullet}} & & j(K^\bullet)[1]
}
$$
由引理 \ref{derived-lemma-morphisms-lift} 和
\ref{derived-lemma-morphisms-equal-up-to-homotopy}，
在 $K^{+}(\mathcal{I})$ 中存在唯一虚线箭头
$\xi_{K^\bullet}$，使图在 $K^{+}(\mathcal{A})$ 中交换。
我们略去这给出函子同构的验证。
（提示：再次使用引理 \ref{derived-lemma-morphisms-equal-up-to-homotopy}。）

\medskip\noindent
设 $(K^\bullet, L^\bullet, M^\bullet, f, g, h)$ 为
$K^{+}(\mathcal{A})$ 的特异三角。须证明
$(j(K^\bullet), j(L^\bullet), j(M^\bullet), j(f), j(g),
\xi_{K^\bullet} \circ j(h))$ 是 $K^{+}(\mathcal{I})$ 的特异三角。
注意，在 $K^{+}(\mathcal{A})$ 中有交换图
$$
\xymatrix{
K^\bullet \ar[r]_f \ar[d] &
L^\bullet \ar[r]_g \ar[d] &
M^\bullet \ar[rr]_h \ar[d] & &
K^\bullet[1] \ar[d] \\
j(K^\bullet) \ar[r]^{j(f)} &
j(L^\bullet) \ar[r]^{j(g)} &
j(M^\bullet) \ar[rr]^{\xi_{K^\bullet} \circ j(h)} & &
j(K^\bullet)[1]
}
$$
其竖直箭头为拟同构 $i_K, i_L, i_M$。
所以 $(j(K^\bullet), j(L^\bullet), j(M^\bullet), j(f), j(g),
\xi_{K^\bullet} \circ j(h))$ 在 $D^{+}(\mathcal{A})$ 中的像
与某个特异三角同构，故由 TR1 它本身就是特异三角。
于是由引理 \ref{derived-lemma-exact-equivalence}，
$(j(K^\bullet), j(L^\bullet), j(M^\bullet), j(f), j(g),
\xi_{K^\bullet} \circ j(h))$ 是 $K^{+}(\mathcal{I})$ 中的特异三角。
\end{proof}

\begin{lemma}
\label{derived-lemma-resolution-functor-quasi-inverse}
设 $\mathcal{A}$ 为具有足够多内射对象的阿贝尔范畴，
$j$ 为消解函子。以
$Q : K^{+}(\mathcal{A}) \to D^{+}(\mathcal{A})$ 表示自然函子。
则 $j = j' \circ Q$，其中唯一的函子
$j' : D^{+}(\mathcal{A}) \to K^{+}(\mathcal{I})$
是典范函子 $K^{+}(\mathcal{I}) \to D^{+}(\mathcal{A})$ 的拟逆。
\end{lemma}

\begin{proof}
由引理 \ref{derived-lemma-bounded-derived}，函子 $Q$ 是局部化。
为证明 $j'$ 存在，只须说明 $\text{Qis}^{+}(\mathcal{A})$
的任意元素在函子 $j$ 下映为同构；见引理
\ref{derived-lemma-universal-property-localization}。
考虑引理 \ref{derived-lemma-j-is-exact} 的证明中的交换方块。
在该方块中，若 $\alpha$ 是拟同构，则
$i_{L^\bullet}\circ\alpha=j(\alpha)\circ i_{K^\bullet}$
也是拟同构，因而 $j(\alpha)$ 也是拟同构。
所以由命题 \ref{derived-proposition-derived-category}，态射
$j(\alpha)$ 在 $K^+(\mathcal{I})$ 中是同构。
我们略去 $j'$ 是 $K^{+}(\mathcal{I}) \to D^{+}(\mathcal{A})$
的拟逆这一验证。
\end{proof}

\begin{remark}
\label{derived-remark-big-localization}
假设 $\mathcal{A}$ 是具有足够多内射对象的“大”阿贝尔范畴，
例如阿贝尔群范畴。此时构造消解函子时必须稍加谨慎，
因为不能对量词遍历一个类的情形使用选择公理。
但注意，该引理的证明确实表明任意两个局部化函子典范同构。
具体地，给定从下有界复形 $K^\bullet$ 到内射对象的下有界复形中的
拟同构 $i : K^\bullet \to I^\bullet$ 和
$i' : K^\bullet \to J^\bullet$，在 $K^{+}(\mathcal{I})$ 中存在唯一（！）
态射 $a : I^\bullet \to J^\bullet$，使得作为
% P02-E0054: frozen source has the ill-typed composition i' = i \circ a; it should compose a after i.
$K^{+}(\mathcal{I})$ 中的态射，有 $i' = i \circ a$。
所以唯一的问题是存在性；下一节将说明如何处理它。
\end{remark}









\section{函子性内射嵌入与消解函子}
\label{derived-section-functorial-injective-resolutions}

\noindent
本节在范畴 $\mathcal{A}$ 有函子性内射嵌入的情形下，
重新构造消解函子
$K^{+}(\mathcal{A}) \to K^{+}(\mathcal{I})$。
这样做有两个理由：(1) 证明更容易；(2) 当 $\mathcal{A}$ 是“大”
阿贝尔范畴时，该构造也适用。见下面的注记
\ref{derived-remark-big-abelian-category}。

\medskip\noindent
设 $\mathcal{A}$ 为阿贝尔范畴。与前面一样，以 $\mathcal{I}$ 表示
$\mathcal{A}$ 的加性满子范畴，由内射对象组成。
考虑箭头 $\alpha : K^\bullet \to I^\bullet$ 所成的范畴
$\text{InjRes}(\mathcal{A})$，其中 $K^\bullet$ 是
$\mathcal{A}$ 的下有界复形，$I^\bullet$ 是
$\mathcal{A}$ 中内射对象的下有界复形，而 $\alpha$ 是拟同构。
换言之，$\alpha$ 是内射消解，且 $K^\bullet$ 下有界。
有显然的函子
$$
s : \text{InjRes}(\mathcal{A}) \longrightarrow \text{Comp}^{+}(\mathcal{A})
$$
由 $(\alpha : K^\bullet \to I^\bullet) \mapsto K^\bullet$ 定义。
还有函子
$$
t : \text{InjRes}(\mathcal{A}) \longrightarrow K^{+}(\mathcal{I})
$$
由 $(\alpha : K^\bullet \to I^\bullet) \mapsto I^\bullet$ 定义。

\begin{lemma}
\label{derived-lemma-functorial-injective-resolutions}
设 $\mathcal{A}$ 为阿贝尔范畴。假设 $\mathcal{A}$ 有函子性内射嵌入；
见同调，定义 \ref{homology-definition-functorial-injective-embedding}。
\begin{enumerate}
\item 存在函子
$inj : \text{Comp}^{+}(\mathcal{A}) \to \text{InjRes}(\mathcal{A})$，
使得 $s \circ inj = \text{id}$。
\item 对任意满足 $s \circ inj = \text{id}$ 的函子
$inj : \text{Comp}^{+}(\mathcal{A}) \to \text{InjRes}(\mathcal{A})$，
可得到一个消解函子；见定义 \ref{derived-definition-localization-functor}。
\end{enumerate}
\end{lemma}

\begin{proof}
设 $A \mapsto (A \to J(A))$ 为函子性内射嵌入；见同调，定义
\ref{homology-definition-functorial-injective-embedding}。
先注意，可以假设 $J(0) = 0$。事实上，若非如此，
则对任意对象 $A$，有 $0 \to A \to 0$，从而给出直和分解
$J(A) = J(0) \oplus \Ker(J(A) \to J(0))$。
注意，函子性态射 $A \to J(A)$ 必须映入第二个直和项。
因此必要时可把函子替换为
$J'(A) = \Ker(J(A) \to J(0))$。

\medskip\noindent
设 $K^\bullet$ 为 $\mathcal{A}$ 的下有界复形。
设 $p < B$ 时 $K^p = 0$。我们将构造内射对象的双复形
$I^{\bullet, \bullet}$，以及映射
$\alpha : K^\bullet \to I^{\bullet, 0}$，使得 $\alpha$ 诱导
$K^\bullet$ 到 $I^{\bullet, \bullet}$ 的相伴全复形的拟同构。
首先，只要 $q < 0$，就令 $I^{p, q} = 0$。
接着，令 $I^{p, 0} = J(K^p)$，并令
$\alpha^p : K^p \to I^{p, 0}$ 为函子性嵌入。
由于 $J$ 是函子，$I^{\bullet, 0}$ 是复形，而 $\alpha$ 是复形态射。
每个 $\alpha^p$ 都为单射。并且由于 $J(0) = 0$，
当 $p < B$ 时 $I^{p, 0} = 0$。
再令 $I^{p, 1} = J(\Coker(K^p \to I^{p, 0}))$。
再次由函子性，$I^{\bullet, 1}$ 是复形。
而且当 $p < B$ 时 $I^{p, 1} = 0$。
还显然有 $K^p$ 同构地映到
$\Ker(I^{p, 0} \to I^{p, 1})$ 上。
第三步，取 $I^{p, 2} = J(\Coker(I^{p, 0} \to I^{p, 1}))$，
依此类推。

\medskip\noindent
现在可以应用同调，引理
\ref{homology-lemma-double-complex-gives-resolution}，得到映射
$$
\alpha : K^\bullet \longrightarrow \text{Tot}(I^{\bullet, \bullet})
$$
是拟同构。
% P02-E0056: frozen source says “it rests to show”; “it remains to show” is intended.
为证明得到函子 $inj$，还须说明上述构造具有函子性。
略去这一验证。

\medskip\noindent
假设已有满足 $s \circ inj = \text{id}$ 的函子 $inj$。
对 $\text{Comp}^{+}(\mathcal{A})$ 的每个对象 $K^\bullet$，可写
$$
inj(K^\bullet) = (i_{K^\bullet} : K^\bullet \to j(K^\bullet))
$$
这就给出定义 \ref{derived-definition-localization-functor} 所述的消解函子。
\end{proof}

\begin{remark}
\label{derived-remark-match}
假设 $inj$ 是引理 \ref{derived-lemma-functorial-injective-resolutions}
的 (2) 中满足 $s \circ inj = \text{id}$ 的函子。
如该引理的证明中那样，写成
$inj(K^\bullet) = (i_{K^\bullet} : K^\bullet \to j(K^\bullet))$。
设 $\alpha : K^\bullet \to L^\bullet$ 是下有界复形的映射。
考虑范畴 $\text{InjRes}(\mathcal{A})$ 中的映射 $inj(\alpha)$。
它诱导复形态射的交换图
% P02-E0057: frozen diagram writes j(K)^\bullet and j(L)^\bullet, inconsistent with j(K^\bullet), j(L^\bullet).
$$
\xymatrix{
K^\bullet
\ar[rr]^-{\alpha}
\ar[d]_{i_K} & &
L^\bullet \ar[d]^{i_L} \\
j(K)^\bullet
\ar[rr]^-{inj(\alpha)}
& &
j(L)^\bullet
}
$$
因此，考察引理 \ref{derived-lemma-resolution-functor} 的证明，
可见函子 $j : K^{+}(\mathcal{A}) \to K^{+}(\mathcal{I})$ 由规则
$$
j(\alpha\text{ 在同伦意义下}) = inj(\alpha)\text{ 在同伦意义下}\in
\Hom_{K^{+}(\mathcal{I})}(j(K^\bullet), j(L^\bullet))
$$
给出。所以此时 $j$ 与 $t \circ inj$ 相符；即图
$$
\xymatrix{
\text{Comp}^{+}(\mathcal{A}) \ar[rr]_{t \circ inj} \ar[rd] & &
K^{+}(\mathcal{I}) \\
& K^{+}(\mathcal{A}) \ar[ru]_j
}
$$
交换。
\end{remark}

\begin{remark}
\label{derived-remark-big-abelian-category}
设 $\textit{Mod}(\mathcal{O}_X)$ 为环空间
$(X, \mathcal{O}_X)$（或更一般地，环化位点）上的
$\mathcal{O}_X$-模范畴。稍后将看到，
$\textit{Mod}(\mathcal{O}_X)$ 有足够多内射对象，
而且事实上有函子性内射嵌入；见内射对象，定理
\ref{injectives-theorem-sheaves-modules-injectives}。
注意，引理 \ref{derived-lemma-into-derived-category} 的证明不适用于
$\textit{Mod}(\mathcal{O}_X)$，但引理
\ref{derived-lemma-functorial-injective-resolutions} 的证明适用于
$\textit{Mod}(\mathcal{O}_X)$。
因此得到
$$
j : K^{+}(\textit{Mod}(\mathcal{O}_X))
\longrightarrow
K^{+}(\mathcal{I})
$$
这是消解函子，其中 $\mathcal{I}$ 是内射
$\mathcal{O}_X$-模所成的加性范畴。该论证还适用于以下情形：
\begin{enumerate}
\item 环 $R$ 上的 $R$-模范畴 $\text{Mod}_R$。
\item 带环预层的位点上 $\mathcal{O}$-模预层的范畴
$\textit{PMod}(\mathcal{O})$。
\item 环化位点上 $\mathcal{O}$-模层的范畴
$\textit{Mod}(\mathcal{O})$。
% P02-E0058: frozen source contains the unfinished placeholder “Add more here as needed.”
\item 需要时在此添加更多内容。
\end{enumerate}
\end{remark}










\section{通过消解函子得到右导出函子}
\label{derived-section-right-derived-functor-via-resolutions}

\noindent
下面引理的内容是：若给定消解函子 $j$，
便可以直接定义 $RF(K^\bullet) = F(j(K^\bullet))$。

\begin{lemma}
\label{derived-lemma-right-derived-functor}
% P02-E0059: frozen source lacks punctuation after “enough injectives”.
设 $\mathcal{A}$ 为具有足够多内射对象的阿贝尔范畴。
设 $F : \mathcal{A} \to \mathcal{B}$ 为到某阿贝尔范畴中的加性函子。
设 $(i, j)$ 为消解函子；见定义
\ref{derived-definition-localization-functor}。
$F$ 的右导出函子 $RF$ 嵌入如下 $2$-交换图：
$$
\xymatrix{
D^{+}(\mathcal{A}) \ar[rd]_{RF} \ar[rr]^{j'} & &
K^{+}(\mathcal{I}) \ar[ld]^F \\
& D^{+}(\mathcal{B})
}
$$
其中 $j'$ 是引理 \ref{derived-lemma-resolution-functor-quasi-inverse}
中的函子。
\end{lemma}

\begin{proof}
由引理 \ref{derived-lemma-injective-acyclic}，有
$RF(K^\bullet) = F(j(K^\bullet))$。
\end{proof}

\begin{remark}
\label{derived-remark-right-derived-functor}
在引理 \ref{derived-lemma-right-derived-functor} 的情形中，
实际上已把右导出函子提升为正合函子
$F \circ j' : D^{+}(\mathcal{A}) \to K^{+}(\mathcal{B})$。
有时使用这样的分解很有用。
\end{remark}








\section{过滤导出范畴与内射消解}
\label{derived-section-filtered-derived}

\noindent
设 $\mathcal{A}$ 为阿贝尔范畴。本节将说明：若 $\mathcal{A}$
有足够多内射对象，则在某种意义下范畴
$\text{Fil}^f(\mathcal{A})$ 亦然。
可以利用这一观察在 $\mathcal{A}$ 的过滤导出范畴中作计算。

\medskip\noindent
范畴 $\text{Fil}^f(\mathcal{A})$ 是正合范畴的一个例子；见
内射对象，注记 \ref{injectives-remark-embed-exact-category}。
严格态射扮演特殊角色；见同调，定义
\ref{homology-definition-strict}，即满足
$\Coim(f) = \Im(f)$ 的态射 $f$。
若 $\text{Fil}^f(\mathcal{A})$ 中的复形 $A \to B \to C$
满足序列 $\text{gr}(A) \to \text{gr}(B) \to \text{gr}(C)$
在 $\mathcal{A}$ 中正合，则称该复形{\it 正合}。
这意味着 $A \to B$ 和 $B \to C$ 是严格态射；见同调，引理
\ref{homology-lemma-filtered-acyclic}。

\begin{definition}
\label{derived-definition-filtered-complexes-notation}
设 $\mathcal{A}$ 为阿贝尔范畴。
若 $\text{Fil}^f(\mathcal{A})$ 的对象 $I$ 的每个
$\text{gr}^p(I)$ 都是 $\mathcal{A}$ 的内射对象，
则称该对象为{\it 过滤内射对象}。
\end{definition}

\begin{lemma}
\label{derived-lemma-filtered-injective}
设 $\mathcal{A}$ 为阿贝尔范畴。
$\text{Fil}^f(\mathcal{A})$ 的对象 $I$ 为过滤内射对象，
当且仅当存在 $a \leq b$、$\mathcal{A}$ 的内射对象
$I_n$（$a \leq n \leq b$）以及同构
$I \cong \bigoplus_{a \leq n \leq b} I_n$，使得
$F^pI = \bigoplus_{n \geq p} I_n$。
\end{lemma}

\begin{proof}
这由如下事实推出：若 $J$ 是内射对象，则 $\mathcal{A}$ 的任意单射
$J \to M$ 都分裂。细节从略。
\end{proof}

\begin{lemma}
\label{derived-lemma-split-strict-monomorphism}
设 $\mathcal{A}$ 为阿贝尔范畴。
$\text{Fil}^f(\mathcal{A})$ 的任意严格单态射
$u : I \to A$，若 $I$ 为过滤内射对象，则它是分裂单射。
\end{lemma}

\begin{proof}
设 $p$ 为满足 $F^pI \not = 0$ 的最大整数。
特别地，$\text{gr}^p(I) = F^pI$。
设 $I'$ 为 $\text{Fil}^f(\mathcal{A})$ 的对象，其
$\mathcal{A}$ 中的底层对象为 $F^pI$，其过滤由如下条件给出：
当 $n > p$ 时 $F^nI' = 0$，当 $n \leq p$ 时
$F^nI' = I' = F^pI$。注意，$I' \to I$ 也是严格单态射。
$u$ 是严格单态射这一事实意味着
$F^pI \to A/F^{p + 1}(A)$ 为单射；见同调，引理
\ref{homology-lemma-characterize-strict}。
在 $\mathcal{A}$ 中选取分裂 $s : A/F^{p + 1}A \to F^pI$。
所诱导的态射 $s' : A \to I'$ 是过滤对象的严格态射，
并分裂复合 $I' \to I \to A$。所以可写
$A = I' \oplus \Ker(s')$ 以及
$I = I' \oplus \Ker(s'|_I)$。注意，
$\Ker(s'|_I) \to \Ker(s')$ 是严格单态射，且
$\Ker(s'|_I)$ 是过滤内射对象。
对 $I$ 的过滤长度归纳，映射
$\Ker(s'|_I) \to \Ker(s')$ 是分裂单射。证毕。
\end{proof}

\begin{lemma}
\label{derived-lemma-injective-property-filtered-injective}
设 $\mathcal{A}$ 为阿贝尔范畴。设
$u : A \to B$ 为 $\text{Fil}^f(\mathcal{A})$ 的严格单态射，
$f : A \to I$ 为从 $A$ 到
$\text{Fil}^f(\mathcal{A})$ 的过滤内射对象中的态射。
则存在态射 $g : B \to I$，使得 $f = g \circ u$。
\end{lemma}

\begin{proof}
$f$ 沿 $u$ 的推出 $f' : I \to I \amalg_A B$ 是严格单态射；见
同调，引理 \ref{homology-lemma-pushout-filtered}。
所以结论形式地由引理 \ref{derived-lemma-split-strict-monomorphism} 推出。
\end{proof}

\begin{lemma}
\label{derived-lemma-strict-monomorphism-into-filtered-injective}
设 $\mathcal{A}$ 为具有足够多内射对象的阿贝尔范畴。
对 $\text{Fil}^f(\mathcal{A})$ 的任意对象 $A$，
存在严格单态射 $A \to I$，其中 $I$ 是过滤内射对象。
\end{lemma}

\begin{proof}
选取 $a \leq b$，使得除非
$p \in \{a, a + 1, \ldots, b\}$，否则
$\text{gr}^p(A) = 0$。对每个
$n \in \{a, a + 1, \ldots, b\}$，选取单射
$u_n : A/F^{n + 1}A \to I_n$，其中 $I_n$ 是内射对象。
令 $I = \bigoplus_{a \leq n \leq b} I_n$，其过滤为
$F^pI = \bigoplus_{n \geq p} I_n$；并令 $u : A \to I$
等于各映射 $u_n$ 的直和。
\end{proof}

\begin{lemma}
\label{derived-lemma-filtered-injective-right-resolution-single-object}
设 $\mathcal{A}$ 为具有足够多内射对象的阿贝尔范畴。
对 $\text{Fil}^f(\mathcal{A})$ 的任意对象 $A$，
存在过滤拟同构 $A[0] \to I^\bullet$，其中
$I^\bullet$ 是过滤内射对象的复形，且对 $n < 0$ 有 $I^n = 0$。
\end{lemma}

\begin{proof}
首先选取从 $A$ 到某个过滤内射对象中的严格单态射
$u_0 : A \to I^0$；见引理
\ref{derived-lemma-strict-monomorphism-into-filtered-injective}。
接着，选取到 $\mathcal{A}$ 的某个过滤内射对象中的严格单态射
% P02-E0060: filtered injectives are objects of Fil^f(A), not of A as frozen prose states here and below.
$u_1 : \Coker(u_0) \to I^1$。以 $d^0$ 表示所诱导的映射
$I^0 \to I^1$。再选取到 $\mathcal{A}$ 的某个过滤内射对象中的
严格单态射 $u_2 : \Coker(u_1) \to I^2$。
以 $d^1$ 表示所诱导的映射 $I^1 \to I^2$，依此类推。
这是可行的，因为每个序列
$$
0 \to \Coker(u_n) \to I^{n + 1} \to \Coker(u_{n + 1}) \to 0
$$
都是短正合的，即应用 $\text{gr}$ 后诱导短正合序列。
为此使用同调，引理 \ref{homology-lemma-characterize-strict}。
\end{proof}

\begin{lemma}
\label{derived-lemma-filtered-injective-right-resolution-map}
设 $\mathcal{A}$ 为具有足够多内射对象的阿贝尔范畴。
设 $f : A \to B$ 为 $\text{Fil}^f(\mathcal{A})$ 的态射。
给定过滤拟同构 $A[0] \to I^\bullet$ 和
$B[0] \to J^\bullet$，其中 $I^\bullet, J^\bullet$ 是过滤内射对象的
复形，且对 $n < 0$ 有 $I^n = J^n = 0$，则存在交换图
$$
\xymatrix{
A[0] \ar[r] \ar[d] &
B[0] \ar[d] \\
I^\bullet \ar[r] &
J^\bullet
}
$$
\end{lemma}

\begin{proof}
% P02-E0061: the second object is B[0] in the lemma, but the frozen proof says C[0].
由于 $A[0] \to I^\bullet$ 和 $C[0] \to J^\bullet$ 是过滤拟同构，
可知 $a : A \to I^0$、$b : B \to J^0$ 以及所有态射
$d_I^n$、$d_J^n$ 都是严格态射；见同调，引理
\ref{homology-lemma-filtered-acyclic}。
我们将在下列交换图中归纳构造映射 $f^n$：
$$
\xymatrix{
A \ar[r]_a \ar[d]_f &
I^0 \ar[r] \ar[d]^{f^0} &
I^1 \ar[r] \ar[d]^{f^1} &
I^2 \ar[r] \ar[d]^{f^2} &
\ldots \\
B \ar[r]^b &
J^0 \ar[r] &
J^1 \ar[r] &
J^2 \ar[r] &
\ldots
}
$$
由于 $A \to I^0$ 是严格单态射，且 $J^0$ 是过滤内射对象，
由引理 \ref{derived-lemma-injective-property-filtered-injective}，可找到态射
$f^0 : I^0 \to J^0$，使得 $f^0 \circ a = b \circ f$。
复合 $d_J^0 \circ b \circ f$ 为零，故
$d_J^0 \circ f^0 \circ a = 0$，从而 $d_J^0 \circ f^0$
唯一地经由如下态射分解：
$$
\Coker(a) = \Coim(d_I^0) = \Im(d_I^0) \longrightarrow J^1.
$$
由于 $\Im(d_I^0) \to I^1$ 是严格单态射，再次应用引理
\ref{derived-lemma-injective-property-filtered-injective}，可将所示箭头延拓为态射
$f^1 : I^1 \to J^1$。依此类推。
\end{proof}

\begin{lemma}
\label{derived-lemma-filtered-injective-right-resolution-ses}
设 $\mathcal{A}$ 为具有足够多内射对象的阿贝尔范畴。
设 $0 \to A \to B \to C \to 0$ 为
$\text{Fil}^f(\mathcal{A})$ 中的短正合序列。
给定过滤拟同构 $A[0] \to I^\bullet$ 和
$C[0] \to J^\bullet$，其中 $I^\bullet, J^\bullet$ 是过滤内射对象的
复形，且对 $n < 0$ 有 $I^n = J^n = 0$，则存在交换图
$$
\xymatrix{
0 \ar[r] &
A[0] \ar[r] \ar[d] &
B[0] \ar[r] \ar[d] &
C[0] \ar[r] \ar[d] &
0 \\
0 \ar[r] &
I^\bullet \ar[r] &
M^\bullet \ar[r] &
J^\bullet \ar[r] &
0
}
$$
其中下行是逐项分裂的复形序列。
\end{lemma}

\begin{proof}
由于 $A[0] \to I^\bullet$ 和 $C[0] \to J^\bullet$ 是过滤拟同构，
可知 $a : A \to I^0$、$c : C \to J^0$ 以及所有态射
$d_I^n$、$d_J^n$ 都是严格态射；见同调，引理
% P02-E0062: the frozen cross-reference omits the homology- namespace used by the actual label.
\ref{derived-lemma-filtered-acyclic}。
我们将逐步构造下列交换图中朝向东南和朝南的箭头：
$$
\xymatrix{
B \ar[r]_\beta \ar[rd]^b &
C \ar[r]_c \ar[rd]^{\overline{b}} &
J^0 \ar[d]^{\delta^0} \ar[r] &
J^1 \ar[d]^{\delta^1} \ar[r] & \ldots \\
A \ar[u]^\alpha \ar[r]^a &
I^0 \ar[r] &
I^1 \ar[r] &
I^2 \ar[r] & \ldots
}
$$
由于 $A \to B$ 是严格单态射，由引理
\ref{derived-lemma-injective-property-filtered-injective}，可找到态射
$b : B \to I^0$，使得 $b \circ \alpha = a$。
由于 $A$ 是严格态射 $I^0 \to I^1$ 的核，且
$\beta = \Coker(\alpha)$，得到唯一的态射
$\overline{b} : C \to I^1$，使其嵌入图中。
由于 $c$ 是严格单态射且 $I^1$ 是过滤内射对象，
由引理 \ref{derived-lemma-injective-property-filtered-injective}，可找到
$\delta^0 : J^0 \to I^1$。
因为 $B \to C$ 是严格满态射，且
$B \to I^0 \to I^1 \to I^2$ 为零，所以
$C \to I^1 \to I^2$ 为零。因此 $d_I^1 \circ \delta^0$
在 $C \cong \Im(c)$ 上为零。
故 $d_I^1 \circ \delta^0$ 唯一地经由如下态射分解：
$$
\Coker(c) = \Coim(d_J^0) = \Im(d_J^0) \longrightarrow I^2.
$$
由于 $I^2$ 是过滤内射对象，且 $\Im(d_J^0) \to J^1$ 是
严格单态射，由引理 \ref{derived-lemma-injective-property-filtered-injective}，
可将所示态射延拓为态射 $\delta^1 : J^1 \to I^2$。
依此类推。令 $M^\bullet = I^\bullet \oplus J^\bullet$，其微分为
$$
d_M^n =
\left(
\begin{matrix}
d_I^n & (-1)^{n + 1}\delta^n \\
0 & d_J^n
\end{matrix}
\right)
$$
最后，映射 $B[0] \to M^\bullet$ 由
% P02-E0063: the frozen formula gives domain M although the map just defined has domain B.
$b \oplus c \circ \beta : M \to I^0 \oplus J^0$ 给出。
\end{proof}

\begin{lemma}
\label{derived-lemma-right-resolution-by-filtered-injectives}
设 $\mathcal{A}$ 为具有足够多内射对象的阿贝尔范畴。
对每个 $K^\bullet \in K^{+}(\text{Fil}^f(\mathcal{A}))$，
存在过滤拟同构 $K^\bullet \to I^\bullet$，其中 $I^\bullet$ 有下界，
每个 $I^n$ 都是过滤内射对象，且每个
$K^n \to I^n$ 都是严格单态射。
\end{lemma}

\begin{proof}
把 $K^\bullet$ 替换为某个平移后（这对证明无碍），可假设对
$n < 0$ 有 $K^n = 0$。考虑正合范畴
$\text{Fil}^f(\mathcal{A})$ 中的短正合序列
$$
\begin{matrix}
0 \to \Ker(d_K^0) \to K^0 \to \Coim(d_K^0) \to 0 \\
0 \to \Ker(d_K^1) \to K^1 \to \Coim(d_K^1) \to 0 \\
0 \to \Ker(d_K^2) \to K^2 \to \Coim(d_K^2) \to 0 \\
\ldots
\end{matrix}
$$
以及映射 $u_i : \Coim(d_K^i) \to \Ker(d_K^{i + 1})$。
对每个 $i \geq 0$，可以选取过滤拟同构
$$
\begin{matrix}
\Ker(d_K^i)[0] \to I_{ker, i}^\bullet \\
\Coim(d_K^i)[0] \to I_{coim, i}^\bullet
\end{matrix}
$$
其中 $I_{ker, i}^n, I_{coim, i}^n$ 是过滤内射对象，并且对
$n < 0$ 为零；见引理
\ref{derived-lemma-filtered-injective-right-resolution-single-object}。
由引理 \ref{derived-lemma-filtered-injective-right-resolution-map}，
可将 $u_i$ 提升为复形态射
$u_i^\bullet : I_{coim, i}^\bullet \to I_{ker, i + 1}^\bullet$。
最后，对每个 $i \geq 0$，由引理
\ref{derived-lemma-filtered-injective-right-resolution-ses}，可补全各图
$$
\xymatrix{
0 \ar[r] &
\Ker(d_K^i)[0] \ar[r] \ar[d] &
K^i[0] \ar[r] \ar[d] &
\Coim(d_K^i)[0] \ar[r] \ar[d] &
0 \\
0 \ar[r] &
I_{ker, i}^\bullet \ar[r]^{\alpha_i} &
I_i^\bullet \ar[r]^{\beta_i} &
I_{coim, i}^\bullet \ar[r] &
0
}
$$
使得下行是逐项分裂的正合序列。
对 $i \geq 0$，令 $d_i : I_i^\bullet \to I_{i + 1}^\bullet$
等于 $d_i =  \alpha_{i + 1} \circ u_i^\bullet \circ \beta_i$。
注意，$d_i \circ d_{i - 1} = 0$，因为
$\beta_i \circ \alpha_i = 0$。因此我们构造了交换图
$$
\xymatrix{
I_0^\bullet \ar[r] &
I_1^\bullet \ar[r] &
I_2^\bullet \ar[r] & \ldots \\
K^0[0] \ar[r] \ar[u] &
K^1[0] \ar[r] \ar[u] &
K^2[0] \ar[r] \ar[u] &
\ldots
}
$$
这里竖直箭头是过滤拟同构。
上行是复形的复形，且每个复形由过滤内射对象组成，
在次数 $< 0$ 处没有非零对象。
因此令 $I^{a, b} = I_a^b$，并用
$$
d_1^{a, b} : I^{a, b} = I_a^b \to I_{a + 1}^b = I^{a + 1, b}
$$
表示映射 $d_a^b$，又用
$$
d_2^{a, b} : I^{a, b} = I_a^b \to I_a^{b + 1} = I^{a, b + 1}
$$
表示映射 $d_{I_a}^b$，便得到双复形。
以 $\text{Tot}(I^{\bullet, \bullet})$ 表示与此双复形关联的全复形；见
同调，定义 \ref{homology-definition-associated-simple-complex}。
注意，映射 $K^n[0] \to I_n^\bullet$ 来自映射
$K^n \to I^{n, 0}$，后者诱导复形态射
$$
K^\bullet \longrightarrow \text{Tot}(I^{\bullet, \bullet})
$$
我们断言这是过滤拟同构。
由于 $\text{gr}(-)$ 是加性函子，可见
$\text{gr}(\text{Tot}(I^{\bullet, \bullet})) =
\text{Tot}(\text{gr}(I^{\bullet, \bullet}))$。
因此可以应用同调，引理
\ref{homology-lemma-double-complex-gives-resolution}，得出
$\text{gr}(K^\bullet) \to \text{gr}(\text{Tot}(I^{\bullet, \bullet}))$
如所需是拟同构。
\end{proof}

\begin{lemma}
\label{derived-lemma-filtered-acyclic-is-zero}
设 $\mathcal{A}$ 为阿贝尔范畴。
设 $K^\bullet, I^\bullet \in K(\text{Fil}^f(\mathcal{A}))$。
假设 $K^\bullet$ 过滤无环，并且 $I^\bullet$ 有下界且由过滤内射对象组成。
任意态射 $K^\bullet \to I^\bullet$ 都零同伦：
$\Hom_{K(\text{Fil}^f(\mathcal{A}))}(K^\bullet, I^\bullet) = 0$。
\end{lemma}

\begin{proof}
设 $\alpha : K^\bullet \to I^\bullet$ 为复形态射。
假设对 $j < n$ 有 $\alpha^j = 0$。
我们将证明存在态射 $h : K^{n + 1} \to I^n$，使得
$\alpha^n = h \circ d$。于是 $\alpha$ 将同伦于由下式定义的复形态射
$\beta$：
$$
\beta^j =
\left\{
\begin{matrix}
0 & \text{若} & j \leq n \\
\alpha^{n + 1} - d \circ h & \text{若} & j = n + 1 \\
\alpha^j & \text{若} & j > n + 1
\end{matrix}
\right.
$$
这显然（通过归纳）证明了引理。
为证明 $h$ 的存在性，注意
$\alpha^n \circ d_K^{n - 1} = 0$，因为
$\alpha^{n - 1} = 0$。由于 $K^\bullet$ 过滤无环，
$d_K^{n - 1}$ 和 $d_K^n$ 都是严格态射，并且
$$
0 \to \Im(d_K^{n - 1}) \to K^n \to \Im(d_K^n) \to 0
$$
是正合范畴 $\text{Fil}^f(\mathcal{A})$ 中的正合序列；见
同调，引理 \ref{homology-lemma-filtered-acyclic}。
因此可把 $\alpha^n$ 看作定义在 $\Im(d_K^n)$ 上、取值于
$I^n$ 的映射。
利用 $\Im(d_K^n) \to K^{n + 1}$ 是严格单态射且 $I^n$
为过滤内射对象，由引理
\ref{derived-lemma-injective-property-filtered-injective}，可如所需把此映射提升为
$h : K^{n + 1} \to I^n$。
\end{proof}

\begin{lemma}
\label{derived-lemma-morphisms-into-filtered-injective-complex}
设 $\mathcal{A}$ 为阿贝尔范畴。
设 $I^\bullet \in K(\text{Fil}^f(\mathcal{A}))$ 是由过滤内射对象组成的
有下界复形。
\begin{enumerate}
\item 设 $\alpha : K^\bullet \to L^\bullet$ 是
$K(\text{Fil}^f(\mathcal{A}))$ 中的过滤拟同构。则映射
$$
\Hom_{K(\text{Fil}^f(\mathcal{A}))}(L^\bullet, I^\bullet)
\to
\Hom_{K(\text{Fil}^f(\mathcal{A}))}(K^\bullet, I^\bullet)
$$
是双射。
\item 设 $L^\bullet \in K(\text{Fil}^f(\mathcal{A}))$。则
$$
\Hom_{K(\text{Fil}^f(\mathcal{A}))}(L^\bullet, I^\bullet)
=
\Hom_{DF(\mathcal{A})}(L^\bullet, I^\bullet).
$$
\end{enumerate}
\end{lemma}

\begin{proof}
证明 (1)。注意
$$
(K^\bullet, L^\bullet, C(\alpha)^\bullet, \alpha, i, -p)
$$
是 $K(\text{Fil}^f(\mathcal{A}))$ 中的特异三角
（引理 \ref{derived-lemma-the-same-up-to-isomorphisms}），并且
$C(\alpha)^\bullet$ 是过滤无环复形
（引理 \ref{derived-lemma-filtered-acyclic}）。于是
$$
\xymatrix{
\Hom_{K(\text{Fil}^f(\mathcal{A}))}(C(\alpha)^\bullet, I^\bullet) \ar[r] &
\Hom_{K(\text{Fil}^f(\mathcal{A}))}(L^\bullet, I^\bullet) \ar[r] &
\Hom_{K(\text{Fil}^f(\mathcal{A}))}(K^\bullet, I^\bullet) \ar[lld] \\
\Hom_{K(\text{Fil}^f(\mathcal{A}))}(C(\alpha)^\bullet[-1], I^\bullet)
}
$$
是阿贝尔群的正合序列；见引理
\ref{derived-lemma-representable-homological}。
此时引理 \ref{derived-lemma-filtered-acyclic-is-zero} 保证外侧两个群为零，
因而
% P02-E0064: the frozen conclusion drops Fil^f(A) from both homotopy categories.
$\Hom_{K(\mathcal{A})}(L^\bullet, I^\bullet) =
\Hom_{K(\mathcal{A})}(K^\bullet, I^\bullet)$。

\medskip\noindent
证明 (2)。
设 $a$ 为右侧的一个元素。
可将其表示成 $a = \gamma\alpha^{-1}$，其中
$\alpha : K^\bullet \to L^\bullet$ 是过滤拟同构，而
$\gamma : K^\bullet \to I^\bullet$ 是复形映射。由 (1)，
可找到态射 $\beta : L^\bullet \to I^\bullet$，使得
$\beta \circ \alpha$ 同伦于 $\gamma$。这证明该映射是满射。
设 $b$ 为左侧的元素，且在右侧中的像为零。则 $b$ 是某个态射
$\beta : L^\bullet \to I^\bullet$ 的同伦类，并且存在过滤拟同构
$\alpha : K^\bullet \to L^\bullet$，使得
$\beta  \circ \alpha$ 零同伦。于是 (1) 表明
$\beta$ 也零同伦，即 $b = 0$。
\end{proof}

\begin{lemma}
\label{derived-lemma-filtered-localization-functor}
设 $\mathcal{A}$ 为具有足够多内射对象的阿贝尔范畴。
以 $\mathcal{I}^f \subset \text{Fil}^f(\mathcal{A})$ 表示对象为
过滤内射对象的严格满加性子范畴。典范函子
$$
K^{+}(\mathcal{I}^f)
\longrightarrow
DF^{+}(\mathcal{A})
$$
正合、全忠实且本质满，即它是三角范畴的等价。此外各图
$$
\xymatrix{
K^{+}(\mathcal{I}^f) \ar[d]_{\text{gr}^p} \ar[r] &
DF^{+}(\mathcal{A}) \ar[d]_{\text{gr}^p} \\
K^{+}(\mathcal{I}) \ar[r] &
D^{+}(\mathcal{A})
}
\quad
\xymatrix{
K^{+}(\mathcal{I}^f) \ar[d]^{\text{遗忘 }F} \ar[r] &
DF^{+}(\mathcal{A}) \ar[d]^{\text{遗忘 }F} \\
K^{+}(\mathcal{I}) \ar[r] &
D^{+}(\mathcal{A})
}
$$
交换，其中 $\mathcal{I} \subset \mathcal{A}$ 是对象为内射对象的
严格满加性子范畴。
\end{lemma}

\begin{proof}
由引理 \ref{derived-lemma-right-resolution-by-filtered-injectives}，函子
$K^{+}(\mathcal{I}^f) \to DF^{+}(\mathcal{A})$ 本质满。
由引理 \ref{derived-lemma-morphisms-into-filtered-injective-complex}，它全忠实。
由我们关于特异三角的定义，它是正合函子。
各方块的交换性是立即的。
\end{proof}

\begin{remark}
\label{derived-remark-filtered-localization-big}
只有当 $\mathcal{A}$ 是我们意义下的范畴，即它具有对象的集合时，
我们才能逆转引理中的箭头。不过，假设给定一个具有足够多内射对象的
大阿贝尔范畴 $\mathcal{A}$，例如 $\textit{Mod}(\mathcal{O}_X)$。
那么对任意给定的对象集合 $\{A_i\}_{i\in I}$，存在包含所有这些对象且
具有足够多内射对象的阿贝尔子范畴
$\mathcal{A}' \subset \mathcal{A}$；见集合，引理
\ref{sets-lemma-abelian-injectives}。
因此可以对 $\mathcal{A}'$ 使用上述引理。
这实质上意味着，若我们只使用一集合之多的图等等，
那么使用该引理绝不会遇到麻烦。
\end{remark}

\noindent
设 $\mathcal{A}, \mathcal{B}$ 为阿贝尔范畴。
设 $T : \mathcal{A} \to \mathcal{B}$ 为左正合函子。
（我们不能用字母 $F$ 表示该函子，因为这会与我们用字母 $F$
表示过滤严重冲突。）注意，$T$ 按规则
$T(A, F) = (T(A), F)$ 诱导加性函子
$$
T : \text{Fil}^f(\mathcal{A}) \to \text{Fil}^f(\mathcal{B})
$$
其中 $F^pT(A) = T(F^pA)$；这有意义，因为 $T$ 左正合。
（警告：未必有 $\text{gr}(T(A)) = T(\text{gr}(A))$。）
这诱导三角范畴的函子
\begin{equation}
\label{derived-equation-induced-T-filtered}
T :
K^{+}(\text{Fil}^f(\mathcal{A}))
\longrightarrow
K^{+}(\text{Fil}^f(\mathcal{B})).
\end{equation}
$T$ 的过滤右导出函子，是把定义
\ref{derived-definition-right-derived-functor-defined} 应用于这个正合函子和乘法系
$\text{FQis}^{+}(\mathcal{A})$ 所得的右导出函子，再与正合函子
$K^{+}(\text{Fil}^f(\mathcal{B})) \to DF^{+}(\mathcal{B})$ 复合。
假设 $\mathcal{A}$ 具有足够多内射对象。此时可以重复
第 \ref{derived-section-right-derived-functor} 节的讨论，定义函子 $T$ 的
{\it 过滤右导出函子}
\begin{equation}
\label{derived-equation-filtered-derived-functor}
RT : DF^{+}(\mathcal{A}) \longrightarrow DF^{+}(\mathcal{B})
\end{equation}

\medskip\noindent
不过，我们将改为像第
\ref{derived-section-right-derived-functor-via-resolutions} 节那样进行；结果表明，
即使 $T$ 只是加性函子，也能定义 $RT$。
具体地，首先选取引理 \ref{derived-lemma-filtered-localization-functor} 中等价的
拟逆
$j' : DF^{+}(\mathcal{A}) \to K^{+}(\mathcal{I}^f)$。
由引理 \ref{derived-lemma-exact-equivalence}，可知 $j'$ 是三角范畴的正合函子。
接着注意，对过滤内射对象 $I$，由引理
\ref{derived-lemma-filtered-injective}，我们有（非典范）分解
\begin{equation}
\label{derived-equation-decompose}
I \cong \bigoplus\nolimits_{p \in \mathbf{Z}} I_p,
\quad\text{其中}\quad
F^pI = \bigoplus\nolimits_{q \geq p} I_q
\end{equation}
因此，若 $T$ 是任意加性函子 $T : \mathcal{A} \to \mathcal{B}$，
则令 $T_{ext}(I) = \bigoplus T(I_p)$ 且
$F^pT_{ext}(I) = \bigoplus_{q \geq p} T(I_q)$，便得到加性函子
\begin{equation}
\label{derived-equation-extend-T}
T_{ext} : \mathcal{I}^f \to \text{Fil}^f(\mathcal{B})
\end{equation}
注意，由构造有性质 $\text{gr}(T_{ext}(I)) = T(\text{gr}(I))$。
于是得到与 $\text{gr}$ 交换的函子
\begin{equation}
\label{derived-equation-extend-T-complexes}
T_{ext} : K^{+}(\mathcal{I}^f) \to K^{+}(\text{Fil}^f(\mathcal{B}))
\end{equation}
然后把 (\ref{derived-equation-filtered-derived-functor}) 定义为复合
\begin{equation}
\label{derived-equation-definition-filtered-derived-functor}
RT = T_{ext} \circ j'.
\end{equation}
由于 $RT : D^{+}(\mathcal{A}) \to D^{+}(\mathcal{B})$
也由内射消解计算；见引理 \ref{derived-lemma-injective-acyclic}，
$T$ 与 $\text{gr}$ 的交换性以及引理
\ref{derived-lemma-filtered-localization-functor} 中的交换图蕴含
\begin{equation}
\label{derived-equation-commute-gr}
\text{gr}^p \circ RT \cong RT \circ \text{gr}^p
\end{equation}
以及
\begin{equation}
\label{derived-equation-commute-forget}
(\text{遗忘 }F) \circ RT \cong RT \circ (\text{遗忘 }F)
\end{equation}
是函子 $DF^{+}(\mathcal{A}) \to D^{+}(\mathcal{B})$ 的同构。

\medskip\noindent
过滤导出函子 $RT$（\ref{derived-equation-filtered-derived-functor}）诱导函子
$$
\begin{matrix}
RT : \text{Fil}^f(\mathcal{A}) \to DF^{+}(\mathcal{B}), \\
RT : \text{Comp}^{+}(\text{Fil}^f(\mathcal{A})) \to DF^{+}(\mathcal{B}), \\
RT : KF^{+}(\mathcal{A}) \to DF^{+}(\mathcal{B}).
\end{matrix}
$$
注意，由于 $\text{Fil}^f(\mathcal{A})$ 和
$\text{Comp}^{+}(\text{Fil}^f(\mathcal{A}))$ 不再是阿贝尔范畴，
说 $RT$ 在它们上限制为 $\delta$-函子没有意义。
（把这些范畴视为正合范畴，并表述从正合范畴到三角范畴的
$\delta$-函子概念，可以修复这一点。）
但是，说 $RT$ 是三角范畴 $KF^{+}(\mathcal{A})$ 上的正合函子是有意义的，
并且由构造确实如此。

\begin{lemma}
\label{derived-lemma-ss-filtered-derived}
设 $\mathcal{A}, \mathcal{B}$ 为阿贝尔范畴。设
$T : \mathcal{A} \to \mathcal{B}$ 为左正合函子。
假设 $\mathcal{A}$ 具有足够多内射对象。
设 $(K^\bullet, F)$ 为
$\text{Comp}^{+}(\text{Fil}^f(\mathcal{A}))$ 的对象。
存在谱序列 $(E_r, d_r)_{r\geq 0}$，它由
$\mathcal{B}$ 的双分次对象 $E_r$ 和双次数为 $(r, - r + 1)$ 的
$d_r$ 组成，并且
$$
E_1^{p, q} = R^{p + q}T(\text{gr}^p(K^\bullet))
$$
此外，该谱序列有界，收敛到 $R^*T(K^\bullet)$，并在每个
$R^nT(K^\bullet)$ 上诱导有限过滤。该谱序列的构造关于
$\text{Comp}^{+}(\text{Fil}^f(\mathcal{A}))$ 的对象
$K^\bullet$ 具有函子性，而且对 $r \geq 1$，各项
$(E_r, d_r)$ 不依赖任何选择。
\end{lemma}

\begin{proof}
选取过滤拟同构 $K^\bullet \to I^\bullet$，其中 $I^\bullet$
是过滤内射对象的有下界复形；见引理
\ref{derived-lemma-right-resolution-by-filtered-injectives}。
考虑复形 $RT(K^\bullet) = T_{ext}(I^\bullet)$；见
(\ref{derived-equation-definition-filtered-derived-functor})。
于是可考虑把它视为 $\mathcal{B}$ 中的过滤复形而关联的谱序列
$(E_r, d_r)_{r \geq 0}$；见同调，第
\ref{homology-section-filtered-complex} 节。
由同调，引理
\ref{homology-lemma-spectral-sequence-filtered-complex}，有
$E_1^{p, q} = H^{p + q}(\text{gr}^p(T(I^\bullet)))$。
由式 (\ref{derived-equation-decompose})，有
$E_1^{p, q} = H^{p + q}(T(\text{gr}^p(I^\bullet)))$；而由过滤内射消解的
定义，映射 $\text{gr}^p(K^\bullet) \to \text{gr}^p(I^\bullet)$
是内射消解。因此
$E_1^{p, q} = R^{p + q}T(\text{gr}^p(K^\bullet))$。

\medskip\noindent
另一方面，每个 $I^n$ 都有有限过滤，因而每个 $T(I^n)$ 也有有限过滤。
所以可以应用同调，引理
\ref{homology-lemma-biregular-ss-converges}，得出该谱序列有界且收敛到
$H^n(T(I^\bullet)) = R^nT(K^\bullet)$，并且还在每一项上诱导有限过滤。

\medskip\noindent
假设 $K^\bullet \to L^\bullet$ 是
$\text{Comp}^{+}(\text{Fil}^f(\mathcal{A}))$ 的态射。
选取过滤拟同构 $L^\bullet \to J^\bullet$，其中 $J^\bullet$
是过滤内射对象的有下界复形；见引理
\ref{derived-lemma-right-resolution-by-filtered-injectives}。
由上述结果，例如引理
\ref{derived-lemma-morphisms-into-filtered-injective-complex}，存在图
$$
\xymatrix{
K^\bullet \ar[r] \ar[d] & L^\bullet \ar[d] \\
I^\bullet \ar[r] & J^\bullet
}
$$
它在同伦意义下交换。因此得到过滤复形的态射
$T(I^\bullet) \to T(J^\bullet)$，它诱导谱序列的态射；见同调，引理
\ref{homology-lemma-spectral-sequence-filtered-complex-functorial}。
最后一个陈述由此得出。
\end{proof}

\begin{remark}
\label{derived-remark-final-functorial}
正如评注 \ref{derived-remark-functorial-ss} 中所承诺的，我们讨论上述引理与
使用 Cartan--Eilenberg 消解的构造之间的联系。
设 $T : \mathcal{A} \to \mathcal{B}$ 为阿贝尔范畴之间的左正合函子，
假设 $\mathcal{A}$ 具有足够多内射对象，并设 $K^\bullet$ 为
$\mathcal{A}$ 的有下界复形。我们给出引理
\ref{derived-lemma-two-ss-complex-functor} 中谱序列 ${}'E$ 和 ${}''E$
的另一种构造。

\medskip\noindent
第一个谱序列。考虑 $K^\bullet$ 上的“愚蠢”过滤，其定义为
$F^p(K^\bullet) = \sigma_{\geq p}(K^\bullet)$；见同调，第
\ref{homology-section-truncations} 节。
% P02-E0065: frozen prose reads “Note that this stupid” and omits “is”.
注意，这在 $d(F^p(K^\bullet)) \subset F^{p + 1}(K^\bullet)$ 的意义下
是愚蠢的；比较同调，引理
\ref{homology-lemma-spectral-sequence-filtered-complex-d1}。
注意，在此过滤下 $\text{gr}^p(K^\bullet) = K^p[-p]$。
由引理 \ref{derived-lemma-ss-filtered-derived}，存在谱序列，其 $E_1$ 项为
$$
E_1^{p, q} = R^{p + q}T(K^p[-p]) = R^qT(K^p)
$$
正如谱序列 ${}'E_r$ 中那样。此外，注意微分
$E_1^{p, q} \to E_1^{p + 1, q}$ 与 $'{}E_1$ 中的微分相同；见
同调，引理 \ref{homology-lemma-spectral-sequence-filtered-complex-d1}
的 (2) 以及引理 \ref{derived-lemma-two-ss-complex-functor} 证明中对
${}'d_1$ 的描述。

\medskip\noindent
第二个谱序列。考虑 $K^\bullet$ 上由
$F^p(K^\bullet) = \tau_{\leq -p}(K^\bullet)$ 定义的过滤；见同调，第
\ref{homology-section-truncations} 节。
负号是得到递减过滤所必需的。注意，在此过滤下
$\text{gr}^p(K^\bullet)$ 拟同构于 $H^{-p}(K^\bullet)[p]$。
由引理 \ref{derived-lemma-ss-filtered-derived}，存在谱序列，其 $E_1$ 项为
$$
E_1^{p, q} = R^{p + q}T(H^{-p}(K^\bullet)[p])
= R^{2p + q}T(H^{-p}(K^\bullet)) = {}''E_2^{i, j}
$$
其中 $i = 2p + q$ 且 $j = -p$。（这看起来不自然，但注意，我们也
完全可以用递增过滤来发展整个过滤复形理论；结果便是这里看起来自然，
而另一个却不自然。）留给读者验证各微分相符。

\medskip\noindent
事实上，给定 Cartan--Eilenberg 消解
$K^\bullet \to I^{\bullet, \bullet}$，诱导的态射
$K^\bullet \to \text{Tot}(I^{\bullet, \bullet})$ 取值于相应的全复形；
对 $\text{Tot}(I^{\bullet, \bullet})$ 使用适当过滤后，
它对任一过滤都是过滤内射消解。
这可用于精确匹配这些谱序列。
\end{remark}






\section{Ext 群}
\label{derived-section-ext}

\noindent
本节开始描述阿贝尔范畴中对象的 Ext 群。首先给出如下非常一般的定义。

\begin{definition}
\label{derived-definition-ext}
设 $\mathcal{A}$ 为阿贝尔范畴。设 $i \in \mathbf{Z}$。设
$X, Y$ 为 $D(\mathcal{A})$ 的对象。$X$ 被 $Y$ 扩张的
{\it 第 $i$ 个 Ext 群}是群
$$
\Ext^i_\mathcal{A}(X, Y) =
\Hom_{D(\mathcal{A})}(X, Y[i]) =
\Hom_{D(\mathcal{A})}(X[-i], Y).
$$
若 $A, B \in \Ob(\mathcal{A})$，则令
$\Ext^i_\mathcal{A}(A, B) = \text{Ext}^i_\mathcal{A}(A[0], B[0])$。
\end{definition}

\noindent
由于 $\Hom_{D(\mathcal{A})}(X, -)$ 是同调函子，而
$\Hom_{D(\mathcal{A})}(-, Y)$ 是上同调函子；见引理
\ref{derived-lemma-representable-homological}，所以特异三角
$(Y, Y', Y'')$ 分别诱导长正合序列
$$
\ldots \to
\Ext^i_\mathcal{A}(X, Y) \to
\Ext^i_\mathcal{A}(X, Y') \to
\Ext^i_\mathcal{A}(X, Y'') \to
\Ext^{i + 1}_\mathcal{A}(X, Y) \to \ldots
$$
而特异三角 $(X, X', X'')$ 分别诱导
$$
\ldots \to
\Ext^i_\mathcal{A}(X'', Y) \to
\Ext^i_\mathcal{A}(X', Y) \to
\Ext^i_\mathcal{A}(X, Y) \to
\Ext^{i + 1}_\mathcal{A}(X'', Y) \to \ldots
$$
注意，由于 $D^+(\mathcal{A})$、$D^-(\mathcal{A})$、$D^b(\mathcal{A})$
是满子范畴，只要 $X$、$Y$ 属于它们，就可以用这些范畴中的 Hom 群
计算 Ext 群。

\medskip\noindent
若范畴 $\mathcal{A}$ 具有足够多内射对象或足够多投射对象，
则可分别用内射消解或投射消解计算 Ext 群。为避免混淆，回忆：
具有内射（相应地，投射）消解蕴含所有低（相应地，高）次数的同调消失；
见引理 \ref{derived-lemma-cohomology-bounded-below} 和
\ref{derived-lemma-cohomology-bounded-above}。

\begin{lemma}
\label{derived-lemma-compute-ext-resolutions}
设 $\mathcal{A}$ 为阿贝尔范畴。
设 $X^\bullet, Y^\bullet \in \Ob(K(\mathcal{A}))$。
\begin{enumerate}
\item 设 $Y^\bullet \to I^\bullet$ 为内射消解
（定义 \ref{derived-definition-injective-resolution}）。则
$$
\Ext^i_\mathcal{A}(X^\bullet, Y^\bullet) =
\Hom_{K(\mathcal{A})}(X^\bullet, I^\bullet[i]).
$$
\item 设 $P^\bullet \to X^\bullet$ 为投射消解
（定义 \ref{derived-definition-projective-resolution}）。则
$$
\Ext^i_\mathcal{A}(X^\bullet, Y^\bullet) =
\Hom_{K(\mathcal{A})}(P^\bullet[-i], Y^\bullet).
$$
\end{enumerate}
\end{lemma}

\begin{proof}
立即由引理 \ref{derived-lemma-morphisms-into-injective-complex} 和
引理 \ref{derived-lemma-morphisms-from-projective-complex} 得出。
\end{proof}

\noindent
在本节其余部分，我们讨论阿贝尔范畴本身的对象的扩张。
首先注意如下事实。

\begin{lemma}
\label{derived-lemma-negative-exts}
设 $\mathcal{A}$ 为阿贝尔范畴。
\begin{enumerate}
\item 设 $X$、$Y$ 为 $D(\mathcal{A})$ 的对象。给定
$a, b \in \mathbf{Z}$，使得当 $i > a$ 时 $H^i(X) = 0$，
而当 $j < b$ 时 $H^j(Y) = 0$，则当 $n < b - a$ 时
$\Ext^n_\mathcal{A}(X, Y) = 0$，并且
$$
\Ext^{b - a}_\mathcal{A}(X, Y) = \Hom_\mathcal{A}(H^a(X), H^b(Y))
$$
\item 设 $A, B \in \Ob(\mathcal{A})$。
当 $i < 0$ 时有 $\Ext^i_\mathcal{A}(B, A) = 0$。
并有 $\Ext^0_\mathcal{A}(B, A) = \Hom_\mathcal{A}(B, A)$。
\end{enumerate}
\end{lemma}

\begin{proof}
选取表示 $X$ 和 $Y$ 的复形 $X^\bullet$ 和 $Y^\bullet$。
由于 $Y^\bullet \to \tau_{\geq b}Y^\bullet$ 是拟同构，
可假设当 $j < b$ 时 $Y^j = 0$。
设 $L^\bullet \to X^\bullet$ 为任意拟同构。
则 $\tau_{\leq a}L^\bullet \to X^\bullet$ 是拟同构。
因此，$D(\mathcal{A})$ 中的态射 $X \to Y[n]$ 可表示成
$fs^{-1}$，其中 $s : L^\bullet \to X^\bullet$ 是拟同构，
$f : L^\bullet \to Y^\bullet[n]$ 是态射，并且当 $i < a$ 时
% P02-E0067: after tau_{<=a}, the frozen proof should say L^i=0 for i>a, not i<a.
$L^i = 0$。注意，$f$ 把 $L^i$ 映到 $Y^{i + n}$。
故若 $n < b - a$，则 $f = 0$，因为 $L^i$ 与 $Y^{i + n}$
总有一个为零。若 $n = b - a$，则 $f$ 恰好对应于态射
$H^a(X) \to H^b(Y)$。(2) 是 (1) 的特例。
\end{proof}

\noindent
设 $\mathcal{A}$ 为阿贝尔范畴。
假设 $0 \to A \to A' \to A'' \to 0$ 是
$\mathcal{A}$ 中对象的短正合序列。则
$0 \to A[0] \to A'[0] \to A''[0] \to 0$ 在
$D(\mathcal{A})$ 中诱导特异三角（见引理
\ref{derived-lemma-derived-canonical-delta-functor}），因而诱导 Ext 群的长正合序列
$$
0 \to \Ext^0_\mathcal{A}(B, A) \to
\Ext^0_\mathcal{A}(B, A') \to
\Ext^0_\mathcal{A}(B, A'') \to
\Ext^1_\mathcal{A}(B, A) \to \ldots
$$
类似地，给定短正合序列 $0 \to B \to B' \to B'' \to 0$，
得到 Ext 群的长正合序列
$$
0 \to \Ext^0_\mathcal{A}(B'', A) \to
\Ext^0_\mathcal{A}(B', A) \to
\Ext^0_\mathcal{A}(B, A) \to
\Ext^1_\mathcal{A}(B'', A) \to \ldots
$$
可以把这些 Ext 群视为导出范畴构造的一种应用。它表明，无须存在足够多
内射对象或投射对象，也能定义 Ext 群并构造 Ext 群的长正合序列。
Yoneda 给出了避免使用导出范畴的另一种 Ext 群构造；见
\cite{Yoneda}。

\begin{definition}
\label{derived-definition-yoneda-extension}
设 $\mathcal{A}$ 为阿贝尔范畴。设 $A, B \in \Ob(\mathcal{A})$。
当 $i \geq 1$ 时，$B$ 被 $A$ 扩张的次数为 $i$ 的
{\it Yoneda 扩张}是 $\mathcal{A}$ 中的正合序列
$$
E : 0 \to A \to Z_{i - 1} \to Z_{i - 2} \to \ldots \to Z_0 \to B \to 0
$$
若存在交换图
$$
\xymatrix{
0 \ar[r] & A \ar[r] & Z_{i - 1} \ar[r] & \ldots \ar[r] &
Z_0 \ar[r] & B \ar[r] & 0 \\
0 \ar[r] &
A \ar[r] \ar[u]^{\text{id}} \ar[d]_{\text{id}} &
Z''_{i - 1} \ar[r] \ar[u] \ar[d] &
\ldots \ar[r] &
Z''_0 \ar[r] \ar[u] \ar[d] &
B \ar[r] \ar[u]_{\text{id}} \ar[d]^{\text{id}} & 0 \\
0 \ar[r] & A \ar[r] & Z'_{i - 1} \ar[r] & \ldots \ar[r] &
Z'_0 \ar[r] & B \ar[r] & 0
}
$$
且中间一行也是 Yoneda 扩张，则称同次数的两个 Yoneda 扩张
$E$ 和 $E'$ {\it 等价}。
\end{definition}

\noindent
定义中的等价性是否为等价关系并非立即清楚。直接证明这一点颇有启发性，
不过它也将由下面的引理 \ref{derived-lemma-yoneda-extension} 得出。

\medskip\noindent
设 $\mathcal{A}$ 为具有对象 $A$、$B$ 的阿贝尔范畴。
给定 Yoneda 扩张
$E : 0 \to A \to Z_{i - 1} \to Z_{i - 2} \to \ldots \to Z_0 \to B \to 0$，
我们把相应的元素 $\delta(E) \in \Ext^i_\mathcal{A}(B, A)$ 定义为态射
$\delta(E) = fs^{-1} : B[0] \to A[i]$，其中 $s$ 是拟同构
$$
(\ldots \to 0 \to A \to Z_{i - 1} \to \ldots \to Z_0 \to 0 \to \ldots)
\longrightarrow
B[0]
$$
而 $f$ 是复形态射
$$
(\ldots \to 0 \to A \to Z_{i - 1} \to \ldots \to Z_0 \to 0 \to \ldots)
\longrightarrow
A[i]
$$
称 $\delta(E) = fs^{-1}$ 为该 Yoneda 扩张的{\it 类}。
结果表明，这个类刻画 Yoneda 扩张的等价类。

\begin{lemma}
\label{derived-lemma-yoneda-extension}
设 $\mathcal{A}$ 为具有对象 $A$、$B$ 的阿贝尔范畴，并设
$i \geq 1$。$\Ext^i_\mathcal{A}(B, A)$ 中任意元素都等于某个
$B$ 被 $A$ 扩张的次数为 $i$ 的 Yoneda 扩张的 $\delta(E)$。
给定两个同次数的 Yoneda 扩张 $E$、$E'$，则 $E$ 与 $E'$ 等价，
当且仅当 $\delta(E) = \delta(E')$。
\end{lemma}

\begin{proof}
设 $\xi : B[0] \to A[i]$ 为 $\Ext^i_\mathcal{A}(B, A)$ 的元素。
可以把 $\xi = f s^{-1}$ 写成某个拟同构
$s : L^\bullet \to B[0]$ 和映射 $f : L^\bullet \to A[i]$。
把 $L^\bullet$ 替换为 $\tau_{\leq 0}L^\bullet$ 后，
可假设当 $j > 0$ 时 $L^j = 0$。图为
$$
\xymatrix{
L^{- i - 1} \ar[r] & L^{-i} \ar[r] \ar[d] & \ldots \ar[r] &
L^0 \ar[r] & B \ar[r] & 0 \\
& A
}
$$
令 $Z_{i - 1} = (L^{- i + 1} \oplus A)/L^{-i}$，并且对
$j = i - 2, \ldots, 0$ 令 $Z_j = L^{-j}$，便得到 $B$ 被 $A$
扩张的次数为 $i$ 的扩张 $E$，其类 $\delta(E)$ 等于 $\xi$。

\medskip\noindent
由定义立即可知，等价的 Yoneda 扩张具有相同的类。假设
$E : 0 \to A \to Z_{i - 1} \to Z_{i - 2} \to \ldots \to Z_0 \to B \to 0$
和
$E' : 0 \to A \to Z'_{i - 1} \to Z'_{i - 2} \to \ldots \to Z'_0 \to B \to 0$
是具有相同类的 Yoneda 扩张。
由把 $D(\mathcal{A})$ 构造成 $K(\mathcal{A})$ 关于拟同构集合的局部化，
这意味着存在复形 $L^\bullet$ 以及拟同构
$$
t : L^\bullet \to
(\ldots \to 0 \to A \to Z_{i - 1} \to \ldots \to Z_0 \to 0 \to \ldots)
$$
和
$$
t' : L^\bullet \to
(\ldots \to 0 \to A \to Z'_{i - 1} \to \ldots \to Z'_0 \to 0 \to \ldots)
$$
使得 $s \circ t = s' \circ t'$ 且 $f \circ t = f' \circ t'$；见
范畴，第 \ref{categories-section-localization} 节。
设 $E''$ 为从证明第一段中的一对映射
$L^\bullet \to B[0]$ 和 $L^\bullet \to A[i]$ 构造出的
$B$ 被 $A$ 扩张的次数为 $i$ 的扩张。读者容易看出，存在定义等价
Yoneda 扩张时所述的次数为 $i$ 的 Yoneda 扩张“态射”
$E'' \to E$ 和 $E'' \to E'$（细节略去）。这完成证明。
\end{proof}

\begin{lemma}
\label{derived-lemma-ext-1}
设 $\mathcal{A}$ 为阿贝尔范畴。设 $A$、$B$ 为 $\mathcal{A}$ 的对象。
则 $\Ext^1_\mathcal{A}(B, A)$ 是同调，定义
\ref{homology-definition-ext-group} 中构造的群
$\Ext_\mathcal{A}(B, A)$。
\end{lemma}

\begin{proof}
这是引理 \ref{derived-lemma-yoneda-extension} 中 $i = 1$ 的情形。
\end{proof}

\noindent
给定阿贝尔范畴 $\mathcal{A}$ 以及 $D(\mathcal{A})$ 的对象
$X, Y, Z$，存在双线性且结合的复合律
$$
\Ext^j_\mathcal{A}(Y, Z) \times \Ext^i_\mathcal{A}(X, Y)
\longrightarrow
\Ext^{i + j}_\mathcal{A}(X, Z),\quad
(\eta, \xi) \longmapsto \eta \circ \xi
$$
具体地，若 $\xi : X \to Y[i]$ 且 $\eta : Y \to Z[j]$，
则把 $\eta \circ \xi$ 定义为 $\xi$ 与
$\eta[i] : Y[i] \to Z[i + j]$ 的复合。
若 $A, B, C$ 是 $\mathcal{A}$ 的对象且 $i, j \geq 1$，
则这个复合律
$$
\Ext^j_\mathcal{A}(B, C) \times \Ext^i_\mathcal{A}(A, B)
\longrightarrow
\Ext^{i + j}_\mathcal{A}(A, C)
$$
可以用 Yoneda 扩张描述如下：扩张
$$
0 \to A \to Z_{i - 1} \to Z_{i - 2} \to \ldots \to Z_0 \to B \to 0
$$
与
$$
0 \to B \to Z'_{j - 1} \to Z'_{j - 2} \to \ldots \to Z'_0 \to C \to 0
$$
的复合是 Yoneda 扩张
$$
0 \to A \to Z_{i - 1} \to Z_{i - 2} \to \ldots \to Z_0 \to 
Z'_{j - 1} \to Z'_{j - 2} \to \ldots \to Z'_0 \to C \to 0
$$

\begin{lemma}
\label{derived-lemma-cup-ext-1-zero}
设 $\mathcal{A}$ 为阿贝尔范畴。设
$0 \to A \to Z \to B \to 0$ 和
$0 \to B \to Z' \to C \to 0$ 为 $\mathcal{A}$ 中的短正合序列。
以 $[Z] \in \Ext^1_\mathcal{A}(B, A)$ 和
$[Z'] \in \Ext^1_\mathcal{A}(C, B)$ 表示它们的类。
则 $[Z] \circ [Z'] \in \Ext^2_\mathcal{A}(C, A)$ 为 $0$，
当且仅当存在交换图
$$
\xymatrix{
&
&
0 \ar[d] &
0 \ar[d]
\\
0 \ar[r] &
A \ar[r] \ar[d]^1 &
Z \ar[r] \ar[d] &
B \ar[r] \ar[d] &
0 \\
0 \ar[r] &
A \ar[r]  &
W \ar[r] \ar[d] &
Z' \ar[r] \ar[d] &
0 \\
&
&
C \ar[r]^1  \ar[d]&
C \ar[d]\\
&
&
0 &
0 
}
$$
其各行各列在 $\mathcal{A}$ 中正合。
\end{lemma}

\begin{proof}
略。提示：可以应用引理 \ref{derived-lemma-yoneda-extension}，并具体分析
$2$-扩张的类为零意味着什么。
或者，可以利用如下事实：若
$[Z] \circ [Z'] \in \Ext^2_\mathcal{A}(C, A)$ 为零，则由
$\Ext$ 的上同调长正合序列，元素
$[Z] \in \Ext^1_\mathcal{A}(B, A)$ 是
$\Ext^1_\mathcal{A}(Z', A)$ 中某个元素的像。
\end{proof}

\begin{lemma}
\label{derived-lemma-higher-ext-zero}
设 $\mathcal{A}$ 为阿贝尔范畴，并设 $p \geq 0$。
若对 $\mathcal{A}$ 中任意一对对象 $A$、$B$，都有
$\Ext^p_\mathcal{A}(B, A) = 0$，则对 $i \geq p$ 以及
$\mathcal{A}$ 中任意一对对象 $A$、$B$，都有
$\Ext^i_\mathcal{A}(B, A) = 0$。
\end{lemma}

\begin{proof}
当 $i > p$ 时，把任意类 $\xi$ 写成 $\delta(E)$，其中 $E$ 是
Yoneda 扩张
$$
E : 0 \to A \to Z_{i - 1} \to Z_{i - 2} \to \ldots \to Z_0 \to B \to 0
$$
由引理 \ref{derived-lemma-yoneda-extension}，这是可行的。
令 $C = \Ker(Z_{p - 1} \to Z_{p - 2}) = \Im(Z_p \to Z_{p - 1})$。
则 $\delta(E)$ 是 $\delta(E')$ 与 $\delta(E'')$ 的复合，其中
$$
E' : 0 \to C \to Z_{p - 1} \to \ldots \to Z_0 \to B \to 0
$$
和
$$
E'' : 0 \to A \to Z_{i - 1} \to Z_{i - 2} \to \ldots \to Z_p \to C \to 0
$$
由于 $\delta(E') \in \Ext^p_\mathcal{A}(B, C) = 0$，结论得证。
\end{proof}

\begin{lemma}
\label{derived-lemma-ext-2-zero-pre}
设 $\mathcal{A}$ 为阿贝尔范畴，并设 $K$ 为 $D^b(\mathcal{A})$ 的对象，
使得对所有 $p \geq 2$ 和 $i > j$，都有
$\Ext^p_\mathcal{A}(H^i(K), H^j(K)) = 0$。
则 $K$ 同构于其上同调的直和：
$K \cong \bigoplus H^i(K)[-i]$。
\end{lemma}

\begin{proof}
选取 $a, b$，使得当 $i \not \in [a, b]$ 时 $H^i(K) = 0$。
我们对 $b - a$ 归纳证明该引理。若 $b - a \leq 0$，结论显然。
若 $b - a > 0$，则考察截断的特异三角
$$
\tau_{\leq b - 1}K \to K \to H^b(K)[-b] \to (\tau_{\leq b - 1}K)[1]
$$
见评注 \ref{derived-remark-truncation-distinguished-triangle}。
由引理 \ref{derived-lemma-split}，若最后一个箭头为零，则
$K \cong \tau_{\leq b - 1}K \oplus H^b(K)[-b]$，归纳即得结论。
再次使用归纳可见
$$
\Hom_{D(\mathcal{A})}(H^b(K)[-b], (\tau_{\leq b - 1}K)[1]) =
\bigoplus\nolimits_{i < b} \Ext_\mathcal{A}^{b - i + 1}(H^b(K), H^i(K))
$$
由假设，该直和为零，证明完成。
\end{proof}

\begin{lemma}
\label{derived-lemma-ext-2-zero}
设 $\mathcal{A}$ 为阿贝尔范畴。假设对 $\mathcal{A}$ 中任意一对对象
$A$、$B$，都有 $\Ext^2_\mathcal{A}(B, A) = 0$。
则 $D^b(\mathcal{A})$ 的任意对象 $K$ 都同构于其上同调的直和：
$K \cong \bigoplus H^i(K)[-i]$。
\end{lemma}

\begin{proof}
由引理 \ref{derived-lemma-higher-ext-zero}，该假设蕴含：对 $i \geq 2$ 以及
$\mathcal{A}$ 中任意一对对象 $A, B$，都有
$\Ext^i_\mathcal{A}(B, A) = 0$。因此本引理是引理
\ref{derived-lemma-ext-2-zero-pre} 的特例。
\end{proof}






\section{K-群}
\label{derived-section-K-groups}

\noindent
简要讨论三角范畴的 $K_0$。

\begin{definition}
\label{derived-definition-K-zero}
设 $\mathcal{D}$ 为三角范畴。以 $K_0(\mathcal{D})$ 表示
{\it $\mathcal{D}$ 的第零 $K$-群}。它是按如下方式构造的阿贝尔群：
取 $\mathcal{D}$ 的对象所生成的自由阿贝尔群，并对每个特异三角
$X \to Y \to Z$ 加上关系 $[Y] - [X] - [Z] = 0$。
% P02-E0068: frozen prose says “objects on D” rather than “objects of D”.
\end{definition}

\noindent
注意，这蕴含 $[X[n]] = (-1)^n[X]$，因为存在特异三角
$(X, 0, X[1], 0, 0, -\text{id}[1])$。

\begin{lemma}
\label{derived-lemma-K-bounded-derived}
设 $\mathcal{A}$ 为阿贝尔范畴。则第零 $K$-群之间有典范等同
$K_0(D^b(\mathcal{A})) = K_0(\mathcal{A})$。
\end{lemma}

\begin{proof}
给定 $\mathcal{A}$ 的对象 $A$，以 $A[0]$ 表示视为集中在次数 $0$
的复形的对象 $A$。
若 $0 \to A \to A' \to A'' \to 0$ 是短正合序列，
则得到特异三角 $A[0] \to A'[0] \to A''[0] \to A[1]$；见
第 \ref{derived-section-canonical-delta-functor} 节。
这表明，把 $[A]$ 送到 $[A[0]]$ 给出映射
$K_0(\mathcal{A}) \to K_0(D^b(\mathcal{A}))$；请原谅这个糟糕的记号。

\medskip\noindent
另一方面，给定 $D^b(\mathcal{A})$ 的对象 $X$，可考虑元素
$$
c(X) = \sum (-1)^i[H^i(X)] \in K_0(\mathcal{A})
$$
给定特异三角 $X \to Y \to Z$，上同调长正合序列
(\ref{derived-equation-long-exact-cohomology-sequence-D}) 以及
$K_0(\mathcal{A})$ 中的关系表明 $c(Y) = c(X) + c(Z)$。
所以 $c$ 经由映射
$c : K_0(D^b(\mathcal{A})) \to K_0(\mathcal{A})$ 分解。

\medskip\noindent
我们要证明上述两个映射互为逆。显然，复合
$K_0(\mathcal{A}) \to
K_0(D^b(\mathcal{A})) \to K_0(\mathcal{A})$ 是恒等映射。
假设 $X^\bullet$ 是 $\mathcal{A}$ 的有界复形。
“愚蠢截断”的特异三角（见同调，第
\ref{homology-section-truncations} 节）
$$
\sigma_{\geq n}X^\bullet \to \sigma_{\geq n - 1}X^\bullet \to
X^{n - 1}[-n + 1] \to (\sigma_{\geq n}X^\bullet)[1]
$$
的存在性和归纳表明，在 $K_0(D^b(\mathcal{A}))$ 中
$$
[X^\bullet] = \sum (-1)^i[X^i[0]]
$$
（再次请原谅这个记号）。于是复合
$K_0(\mathcal{A}) \to K_0(D^b(\mathcal{A}))$ 是满射，证明完成。
\end{proof}

\begin{lemma}
\label{derived-lemma-map-K}
设 $F : \mathcal{D} \to \mathcal{D}'$ 为三角范畴的正合函子。
则 $F$ 诱导群同态 $K_0(\mathcal{D}) \to K_0(\mathcal{D}')$。
\end{lemma}

\begin{proof}
略。
\end{proof}

\begin{lemma}
\label{derived-lemma-homological-map-K}
设 $H : \mathcal{D} \to \mathcal{A}$ 为从三角范畴到阿贝尔范畴的
同调函子。假设对 $\mathcal{D}$ 中任意 $X$，对象
$H(X[i])$ 中只有有限多个在 $\mathcal{A}$ 中非零。
则 $H$ 诱导群同态 $K_0(\mathcal{D}) \to K_0(\mathcal{A})$，
它把 $[X]$ 送到 $\sum (-1)^i[H(X[i])]$。
\end{lemma}

\begin{proof}
略。
\end{proof}

\begin{lemma}
\label{derived-lemma-DBA-map-K}
设 $\mathcal{B}$ 为阿贝尔范畴 $\mathcal{A}$ 的弱 Serre 子范畴。
存在典范同构
$$
K_0(\mathcal{B}) \longrightarrow
K_0(D^b_\mathcal{B}(\mathcal{A})),\quad
[B] \longmapsto [B[0]]
$$
其逆把 $X$ 的类 $[X]$ 送到元素 $\sum (-1)^i[H^i(X)]$。
\end{lemma}

\begin{proof}
我们略去验证逆映射的规则给出良定义映射
$K_0(D^b_\mathcal{B}(\mathcal{A})) \to K_0(\mathcal{B})$。
显然，复合
$K_0(\mathcal{B}) \to
K_0(D^b_\mathcal{B}(\mathcal{A})) \to
K_0(\mathcal{B})$ 是恒等映射。另一方面，利用评注
\ref{derived-remark-truncation-distinguished-triangle} 中的特异三角和归纳论证，
读者可证明引理陈述中所示箭头是满射（细节略去）。引理得证。
\end{proof}

\begin{lemma}
\label{derived-lemma-bilinear-map-K}
设 $\mathcal{D}$、$\mathcal{D}'$、$\mathcal{D}''$ 为三角范畴。
设
$$
\otimes : \mathcal{D} \times \mathcal{D}' \longrightarrow \mathcal{D}''
$$
为如下函子：固定 $\mathcal{D}$ 中的 $X$ 时，函子
$X \otimes - : \mathcal{D}' \to \mathcal{D}''$ 正合；固定
$\mathcal{D}'$ 中的 $X'$ 时，函子
$- \otimes X' : \mathcal{D} \to \mathcal{D}''$ 正合。
则 $\otimes$ 诱导双线性映射
$K_0(\mathcal{D}) \times K_0(\mathcal{D}') \to K_0(\mathcal{D}'')$，
它把 $([X], [X'])$ 送到 $[X \otimes X']$。
\end{lemma}

\begin{proof}
略。
\end{proof}










\section{无界复形}
\label{derived-section-unbounded}

\noindent
本节材料的一个参考文献是 \cite{Spaltenstein}。
下列引理可用于寻找无界复形的“良好”左消解。

\begin{lemma}
\label{derived-lemma-special-direct-system}
设 $\mathcal{A}$ 为阿贝尔范畴。设
$\mathcal{P} \subset \Ob(\mathcal{A})$ 为子集。
假设 $\mathcal{P}$ 包含 $0$、在（有限）直和下封闭，并且
$\mathcal{A}$ 的每个对象都是 $\mathcal{P}$ 中某个元素的商。
设 $K^\bullet$ 为复形。存在复形范畴中的交换图
$$
\xymatrix{
P_1^\bullet \ar[d] \ar[r] & P_2^\bullet \ar[d] \ar[r] & \ldots \\
\tau_{\leq 1}K^\bullet \ar[r] & \tau_{\leq 2}K^\bullet \ar[r] & \ldots
}
$$
使得
\begin{enumerate}
\item 竖直箭头是拟同构且逐项满；
\item $P_n^\bullet$ 是各项属于 $\mathcal{P}$ 的有上界复形；
\item 箭头 $P_n^\bullet \to P_{n + 1}^\bullet$ 是逐项分裂单射，
且每个余核 $P^i_{n + 1}/P^i_n$ 都是 $\mathcal{P}$ 的元素。
\end{enumerate}
\end{lemma}

\begin{proof}
我们将使用同伦范畴 $K(\mathcal{A})$ 是三角范畴这一事实；见命题
\ref{derived-proposition-homotopy-category-triangulated}。
由引理 \ref{derived-lemma-subcategory-left-resolution}，可找到逐项满的复形映射
$P_1^\bullet \to \tau_{\leq 1}K^\bullet$，它是拟同构，且
$P_1^\bullet$ 的各项属于 $\mathcal{P}$。
归纳地，只须在给定 $P_1^\bullet, \ldots, P_n^\bullet$ 后构造
$P_{n + 1}^\bullet$ 以及映射
$P_n^\bullet \to P_{n + 1}^\bullet$ 和
$P_{n + 1}^\bullet \to \tau_{\leq n + 1}K^\bullet$。

\medskip\noindent
在 $K(\mathcal{A})$ 中选取特异三角
$P_n^\bullet \to \tau_{\leq n + 1}K^\bullet \to C^\bullet \to P_n^\bullet[1]$。
应用引理 \ref{derived-lemma-subcategory-left-resolution}，选取拟同构的复形映射
$Q^\bullet \to C^\bullet$，使得 $Q^\bullet$ 的各项属于 $\mathcal{P}$。
由三角范畴的公理，可把复合
$Q^\bullet \to C^\bullet \to P_n^\bullet[1]$ 嵌入
$K(\mathcal{A})$ 中的特异三角
$P_n^\bullet \to P_{n + 1}^\bullet \to Q^\bullet \to P_n^\bullet[1]$。
由引理 \ref{derived-lemma-improve-distinguished-triangle-homotopy}，
可以并且确实假设
$0 \to P_n^\bullet \to P_{n + 1}^\bullet \to Q^\bullet \to 0$
是逐项分裂的短正合序列。这蕴含 $P_{n + 1}^\bullet$ 的各项属于
$\mathcal{P}$，并且 $P_n^\bullet \to P_{n + 1}^\bullet$ 是余核属于
$\mathcal{P}$ 的逐项分裂单射。
由三角范畴的公理，得到三角范畴 $K(\mathcal{A})$ 中的特异三角态射
$$
\xymatrix{
P_n^\bullet \ar[r] \ar[d] &
P_{n + 1}^\bullet \ar[r] \ar[d] &
Q^\bullet \ar[r] \ar[d] &
P_n^\bullet[1] \ar[d] \\
P_n^\bullet \ar[r] &
\tau_{\leq n + 1}K^\bullet \ar[r] &
C^\bullet \ar[r] &
P_n^\bullet[1]
}
$$
选取实际的复形态射
$f : P_{n + 1}^\bullet \to \tau_{\leq n + 1}K^\bullet$。
上图左方块交换至同伦；但由于
$P_n^\bullet \to P_{n + 1}^\bullet$ 是逐项分裂单射，
可以提升该同伦并修改对 $f$ 的选择，使该方块交换。
最后，$f$ 是拟同构，因为 $P_n^\bullet \to P_n^\bullet$ 和
$Q^\bullet \to C^\bullet$ 都是拟同构。

\medskip\noindent
此时已有所有所需性质，唯一尚不确定的是映射
$f : P_{n + 1}^\bullet \to \tau_{\leq n + 1}K^\bullet$
是否逐项满。由于存在复形的交换图
$$
\xymatrix{
P_n^\bullet \ar[d] \ar[r] & P_{n + 1}^\bullet \ar[d] \\
\tau_{\leq n}K^\bullet \ar[r] & \tau_{\leq n + 1}K^\bullet
}
$$
由归纳假设可知，除可能在次数 $n$ 和 $n + 1$ 之外，
$f$ 在所有次数的项上都是满射。选取对象
$P \in \mathcal{P}$ 以及满射 $q : P \to K^n$。
考虑映射
$$
g :
P^\bullet = (\ldots \to 0 \to P \xrightarrow{1} P \to 0 \to \ldots)
\longrightarrow
\tau_{\leq n + 1}K^\bullet
$$
其中第一个 $P$ 位于次数 $n$，次数 $n$ 的映射由 $q$ 给出，
次数 $n + 1$ 的映射由 $d_K \circ q$ 给出。
它在次数 $n$ 是满射，而次数 $n + 1$ 的余核是
$H^{n + 1}(\tau_{\leq n + 1}K^\bullet)$；要看到这一点，回忆
$\tau_{\leq n + 1}K^\bullet$ 在次数 $n + 1$ 的项为
$\Ker(d_K^{n + 1})$。
不过，由于 $f$ 是拟同构，我们知道 $H^{n + 1}(f)$ 是满射。
因此，把 $f : P_{n + 1}^\bullet \to \tau_{\leq n + 1}K^\bullet$
替换为
$f \oplus g : P_{n + 1}^\bullet \oplus P^\bullet \to \tau_{\leq n + 1}K^\bullet$
后即得结论。
\end{proof}

\noindent
在某些情形下，可以用上述引理证明左导出函子处处有定义。

\begin{proposition}
\label{derived-proposition-left-derived-exists}
设 $F : \mathcal{A} \to \mathcal{B}$ 为阿贝尔范畴之间的右正合函子。
设 $\mathcal{P} \subset \Ob(\mathcal{A})$ 为子集。假设
\begin{enumerate}
\item $\mathcal{P}$ 包含 $0$、在（有限）直和下封闭，并且
$\mathcal{A}$ 的每个对象都是 $\mathcal{P}$ 中某个元素的商；
\item 对 $\mathcal{A}$ 的任意有上界无环复形 $P^\bullet$，
若对所有 $n$ 有 $P^n \in \mathcal{P}$，则复形 $F(P^\bullet)$ 正合；
\item $\mathcal{A}$ 和 $\mathcal{B}$ 都有以 $\mathbf{N}$ 为指标的系统的余极限；
\item 以 $\mathbf{N}$ 为指标的余极限在 $\mathcal{A}$ 和
$\mathcal{B}$ 中都正合；并且
\item $F$ 与以 $\mathbf{N}$ 为指标的余极限交换。
\end{enumerate}
则 $LF$ 在整个 $D(\mathcal{A})$ 上有定义。
\end{proposition}

\begin{proof}
由 (1) 和引理 \ref{derived-lemma-subcategory-left-resolution}，对任意有上界复形
$K^\bullet$，存在拟同构 $P^\bullet \to K^\bullet$，其中
$P^\bullet$ 有上界，且对所有 $n$ 有 $P^n \in \mathcal{P}$。
假设 $s : P^\bullet \to (P')^\bullet$ 是由 $\mathcal{P}$ 中对象组成的
有上界复形之间的拟同构。则
$F(P^\bullet) \to F((P')^\bullet)$ 是拟同构，因为由假设 (2)，
$F(C(s)^\bullet)$ 无环。这已经表明 $LF$ 在 $D^{-}(\mathcal{A})$
上有定义，并且由 $\mathcal{P}$ 中对象组成的有上界复形计算 $LF$；见
引理 \ref{derived-lemma-find-existence-computes}。

\medskip\noindent
接着，设 $K^\bullet$ 为 $\mathcal{A}$ 的任意复形。
选取引理 \ref{derived-lemma-special-direct-system} 中的图
$$
\xymatrix{
P_1^\bullet \ar[d] \ar[r] & P_2^\bullet \ar[d] \ar[r] & \ldots \\
\tau_{\leq 1}K^\bullet \ar[r] & \tau_{\leq 2}K^\bullet \ar[r] & \ldots
}
$$
注意，映射 $\colim P_n^\bullet \to K^\bullet$ 是拟同构，
因为 $\mathcal{A}$ 中以 $\mathbf{N}$ 为指标的余极限正合，且当
$n > i$ 时 $H^i(P_n^\bullet) = H^i(K^\bullet)$。
我们断言
$$
F(\colim P_n^\bullet) = \colim F(P_n^\bullet)
$$
（逐项余极限）是 $LF(K^\bullet)$，即 $\colim P_n^\bullet$ 计算
$LF$。要看到这一点，由引理 \ref{derived-lemma-find-existence-computes}，
只须证明如下断言。假设
$$
\colim Q_n^\bullet = Q^\bullet
\xrightarrow{\ \alpha\ }
P^\bullet = \colim P_n^\bullet
$$
是复形的拟同构，使得每个 $P_n^\bullet$、$Q_n^\bullet$ 都是
各项属于 $\mathcal{P}$ 的有上界复形，并且映射
$P_n^\bullet \to \tau_{\leq n}P^\bullet$ 和
$Q_n^\bullet \to \tau_{\leq n}Q^\bullet$ 是拟同构。
断言：$F(\alpha)$ 是拟同构。

\medskip\noindent
困难在于，我们没有假设 $\alpha$ 由复形 $P_n^\bullet$ 和
$Q_n^\bullet$ 之间的映射取余极限而给出。不过，对每个 $n$，
我们知道图中的实线箭头
% P02-E0070: alpha is declared Q -> P, while the entire downstream construction consistently uses P -> Q.
$$
\xymatrix{
& R^\bullet \ar@{..>}[d] \\
P_n^\bullet \ar[d] &
L^\bullet \ar@{..>}[l] \ar@{..>}[r] &
Q_n^\bullet \ar[d] \\
\tau_{\leq n}P^\bullet \ar[rr]^{\tau_{\leq n}\alpha} & &
\tau_{\leq n}Q^\bullet
}
$$
都是拟同构。由于拟同构在 $K(\mathcal{A})$ 中构成乘法系
（见引理 \ref{derived-lemma-acyclic}），可以找到拟同构
$L^\bullet \to P_n^\bullet$ 和复形映射
$L^\bullet \to Q_n^\bullet$，使得上图交换至同伦。
于是 $\tau_{\leq n}L^\bullet \to L^\bullet$ 是拟同构。
因此（由证明的第一部分）可找到各项属于 $\mathcal{P}$ 的有上界复形
$R^\bullet$ 以及拟同构 $R^\bullet \to L^\bullet$（如图所示）。
使用证明第一段的结果，可知
$F(R^\bullet) \to F(P_n^\bullet)$ 和
$F(R^\bullet) \to F(Q_n^\bullet)$ 是拟同构。
因此得到同构
$H^i(F(P_n^\bullet)) \to H^i(F(Q_n^\bullet))$，它嵌入交换图
$$
\xymatrix{
H^i(F(P_n^\bullet)) \ar[r] \ar[d] &
H^i(F(Q_n^\bullet)) \ar[d] \\
H^i(F(P^\bullet)) \ar[r] &
H^i(F(Q^\bullet))
}
$$
完全相同的论证表明，随着 $n$ 变化，这些映射也彼此相容。
由于 (4) 和 (5) 给出
$$
H^i(F(P^\bullet)) =
H^i(F(\colim P_n^\bullet)) =
H^i(\colim F(P_n^\bullet)) = \colim H^i(F(P_n^\bullet))
$$
且对 $Q^\bullet$ 也类似，所以得出
% P02-E0071: the frozen displayed inline map has unmatched parentheses and omits F(alpha); direction follows P02-E0070.
$H^i(\alpha) : H^i(F(P^\bullet) \to H^i(F(Q^\bullet)$ 是同构，
断言随之得证。
\end{proof}

\begin{lemma}
\label{derived-lemma-special-inverse-system}
设 $\mathcal{A}$ 为阿贝尔范畴。设
$\mathcal{I} \subset \Ob(\mathcal{A})$ 为子集。
假设 $\mathcal{I}$ 包含 $0$、在（有限）积下封闭，并且
$\mathcal{A}$ 的每个对象都是 $\mathcal{I}$ 中某个元素的子对象。
设 $K^\bullet$ 为复形。存在复形范畴中的交换图
$$
\xymatrix{
\ldots \ar[r] &
\tau_{\geq -2}K^\bullet \ar[r] \ar[d] &
\tau_{\geq -1}K^\bullet \ar[d] \\
\ldots \ar[r] & I_2^\bullet \ar[r] & I_1^\bullet
}
$$
使得
\begin{enumerate}
\item 竖直箭头是拟同构且逐项单；
\item $I_n^\bullet$ 是各项属于 $\mathcal{I}$ 的有下界复形；
\item 箭头 $I_{n + 1}^\bullet \to I_n^\bullet$ 是逐项分裂满射，
且 $\Ker(I^i_{n + 1} \to I^i_n)$ 是 $\mathcal{I}$ 的元素。
\end{enumerate}
\end{lemma}

\begin{proof}
本引理对偶于引理 \ref{derived-lemma-special-direct-system}。
\end{proof}







\section{导出伴随函子}
\label{derived-section-deriving-adjoints}

\noindent
设 $F : \mathcal{D} \to \mathcal{D}'$ 和
$G : \mathcal{D}' \to \mathcal{D}$ 为三角范畴的正合函子。
设 $S$（相应地，$S'$）为 $\mathcal{D}$（相应地，$\mathcal{D}'$）
中与三角结构相容的乘法系。
以 $Q : \mathcal{D} \to S^{-1}\mathcal{D}$ 和
$Q' : \mathcal{D}' \to (S')^{-1}\mathcal{D}'$ 表示局部化函子。
在此情形下，滥用记号，常以 $RF$ 表示对应于
$Q' \circ F : \mathcal{D} \to (S')^{-1}\mathcal{D}'$ 和乘法系 $S$
的部分定义的右导出函子。类似地，以 $LG$ 表示对应于
$Q \circ G : \mathcal{D}' \to S^{-1}\mathcal{D}$ 和乘法系 $S'$
的部分定义的左导出函子。图为
$$
\vcenter{
\xymatrix{
\mathcal{D} \ar[r]_F \ar[d]_Q & \mathcal{D}' \ar[d]^{Q'} \\
S^{-1}\mathcal{D} \ar@{..>}[r]^{RF} &
(S')^{-1}\mathcal{D}'
}
}
\quad\text{且}\quad
\vcenter{
\xymatrix{
\mathcal{D}' \ar[r]_G \ar[d]_{Q'} & \mathcal{D} \ar[d]^Q \\
(S')^{-1}\mathcal{D}' \ar@{..>}[r]^{LG} &
S^{-1}\mathcal{D}
}
}
$$

\begin{lemma}
\label{derived-lemma-pre-derived-adjoint-functors-general}
在上述情形中，假设 $F$ 是 $G$ 的右伴随。设
$K \in \Ob(\mathcal{D})$ 且 $M \in \Ob(\mathcal{D}')$。
若 $RF$ 在 $K$ 处有定义，且 $LG$ 在 $M$ 处有定义，
则存在典范同构
$$
\Hom_{(S')^{-1}\mathcal{D}'}(M, RF(K)) =
\Hom_{S^{-1}\mathcal{D}}(LG(M), K)
$$
在 $RF$ 和 $LG$ 有定义的 $S^{-1}\mathcal{D}$ 与
$(S')^{-1}\mathcal{D}'$ 的三角子范畴上，该同构关于两个变量都具有函子性。
\end{lemma}

\begin{proof}
由于 $RF$ 在 $K$ 处有定义，把 $S$ 中的
$s : K \to I$ 送到对象 $F(I)$ 的规则，作为
$(S')^{-1}\mathcal{D}'$ 的 ind-对象，本质恒定且值为 $RF(K)$。
类似地，把 $S'$ 中的 $t : P \to M$ 送到对象 $G(P)$ 的规则，
作为 $S^{-1}\mathcal{D}$ 的 pro-对象，本质恒定且值为 $LG(M)$。
因此有
\begin{align*}
\Hom_{(S')^{-1}\mathcal{D}'}(M, RF(K))
& =
\colim_{s : K \to I} \Hom_{(S')^{-1}\mathcal{D}'}(M, F(I)) \\
& =
\colim_{s : K \to I} \colim_{t : P \to M} \Hom_{\mathcal{D}'}(P, F(I)) \\
& =
\colim_{t : P \to M} \colim_{s : K \to I} \Hom_{\mathcal{D}'}(P, F(I)) \\
& =
\colim_{t : P \to M} \colim_{s : K \to I} \Hom_{\mathcal{D}}(G(P), I) \\
& =
\colim_{t : P \to M} \Hom_{S^{-1}\mathcal{D}}(G(P), K) \\
& =
\Hom_{S^{-1}\mathcal{D}}(LG(M), K)
\end{align*}
第一个等号由范畴，引理
\ref{categories-lemma-characterize-essentially-constant-ind} 成立。
% P02-E0072: the generalized proof refers to D(B), although no abelian category B is present here.
第二个等号由 $D(\mathcal{B})$ 中态射的定义成立；见范畴，评注
\ref{categories-remark-right-localization-morphisms-colimit}。
第三个等号由范畴，引理
\ref{categories-lemma-colimits-commute} 成立。
第四个等号成立是因为 $F$ 与 $G$ 伴随。
% P02-E0073: similarly, the generalized proof refers to D(A), although only S^{-1}D is defined.
第五个等号由 $D(\mathcal{A})$ 中态射的定义成立；见范畴，评注
\ref{categories-remark-left-localization-morphisms-colimit}。
第六个等号由范畴，引理
\ref{categories-lemma-characterize-essentially-constant-pro} 成立。
我们略去函子性的证明。
\end{proof}

\begin{lemma}
\label{derived-lemma-pre-derived-adjoint-functors}
设 $F : \mathcal{A} \to \mathcal{B}$ 和
$G : \mathcal{B} \to \mathcal{A}$ 为阿贝尔范畴的函子，
并且 $F$ 是 $G$ 的右伴随。设 $K^\bullet$ 为 $\mathcal{A}$ 的复形，
$M^\bullet$ 为 $\mathcal{B}$ 的复形。若 $RF$ 在 $K^\bullet$ 处有定义，
且 $LG$ 在 $M^\bullet$ 处有定义，则存在典范同构
$$
\Hom_{D(\mathcal{B})}(M^\bullet, RF(K^\bullet)) =
\Hom_{D(\mathcal{A})}(LG(M^\bullet), K^\bullet)
$$
在 $RF$ 和 $LG$ 有定义的 $D(\mathcal{A})$ 与 $D(\mathcal{B})$
的三角子范畴上，该同构关于两个变量都具有函子性。
\end{lemma}

\begin{proof}
这是非常一般的引理
\ref{derived-lemma-pre-derived-adjoint-functors-general} 的特例。
\end{proof}

\noindent
下列引理说明了为何使用无界导出范畴更为容易：若没有无界导出函子，
甚至无法陈述该引理。

\begin{lemma}
\label{derived-lemma-derived-adjoint-functors}
设 $F : \mathcal{A} \to \mathcal{B}$ 和
$G : \mathcal{B} \to \mathcal{A}$ 为阿贝尔范畴的函子，
并且 $F$ 是 $G$ 的右伴随。若导出函子
$RF : D(\mathcal{A}) \to D(\mathcal{B})$ 和
$LG : D(\mathcal{B}) \to D(\mathcal{A})$ 存在，则
$RF$ 是 $LG$ 的右伴随。
\end{lemma}

\begin{proof}
立即由引理 \ref{derived-lemma-pre-derived-adjoint-functors} 得出。
\end{proof}





\section{K-内射复形}
\label{derived-section-K-injective}

\noindent
下列类型的复形可用于计算无界导出范畴上的右导出函子。

\begin{definition}
\label{derived-definition-K-injective}
设 $\mathcal{A}$ 为阿贝尔范畴。若对每个无环复形 $M^\bullet$ 都有
$\Hom_{K(\mathcal{A})}(M^\bullet, I^\bullet) = 0$，
则称复形 $I^\bullet$ 为{\it K-内射的}。
\end{definition}

\noindent
在该定义的情形中，事实上对所有 $i$ 都有
$\Hom_{K(\mathcal{A})}(M^\bullet[i], I^\bullet) = 0$，
因为无环复形的平移仍无环。

\begin{lemma}
\label{derived-lemma-K-injective}
设 $\mathcal{A}$ 为阿贝尔范畴。设 $I^\bullet$ 为复形。
下列条件等价：
\begin{enumerate}
\item $I^\bullet$ 是 K-内射的；
\item 对每个拟同构 $M^\bullet \to N^\bullet$，映射
$$
\Hom_{K(\mathcal{A})}(N^\bullet, I^\bullet)
\to \Hom_{K(\mathcal{A})}(M^\bullet, I^\bullet)
$$
是双射；并且
\item 对每个复形 $N^\bullet$，映射
$$
\Hom_{K(\mathcal{A})}(N^\bullet, I^\bullet)
\to \Hom_{D(\mathcal{A})}(N^\bullet, I^\bullet)
$$
是同构。
\end{enumerate}
\end{lemma}

\begin{proof}
假设 (1)。则 (2) 成立，因为函子
$\Hom_{K(\mathcal{A})}( - , I^\bullet)$ 是上同调函子，
而拟同构的锥无环。

\medskip\noindent
假设 (2)。$D(\mathcal{A})$ 中的态射
$N^\bullet \to I^\bullet$ 具有形式
$fs^{-1} : N^\bullet \to I^\bullet$，其中
$s : M^\bullet \to N^\bullet$ 是拟同构，而
$f : M^\bullet \to I^\bullet$ 是映射。由 (2)，它对应于
$K(\mathcal{A})$ 中唯一的态射 $N^\bullet \to I^\bullet$，
即 (3) 成立。

\medskip\noindent
假设 (3)。若 $M^\bullet$ 无环，则 $M^\bullet$ 在 $D(\mathcal{A})$
中同构于零复形，故
$\Hom_{D(\mathcal{A})}(M^\bullet, I^\bullet) = 0$，
从而由 (3)，
$\Hom_{K(\mathcal{A})}(M^\bullet, I^\bullet) = 0$，
即 (1) 成立。
\end{proof}

\begin{lemma}
\label{derived-lemma-triangle-K-injective}
设 $\mathcal{A}$ 为阿贝尔范畴。设 $(K, L, M, f, g, h)$ 为
$K(\mathcal{A})$ 的特异三角。若 $K$、$L$、$M$ 中有两个是
K-内射复形，则第三个也是。
\end{lemma}

\begin{proof}
由定义、引理 \ref{derived-lemma-representable-homological} 以及
$K(\mathcal{A})$ 是三角范畴这一事实得出（命题
\ref{derived-proposition-homotopy-category-triangulated}）。
\end{proof}

\begin{lemma}
\label{derived-lemma-bounded-below-injectives-K-injective}
设 $\mathcal{A}$ 为阿贝尔范畴。内射对象的有下界复形是 K-内射的。
\end{lemma}

\begin{proof}
由引理 \ref{derived-lemma-K-injective} 和
\ref{derived-lemma-morphisms-into-injective-complex} 得出。
\end{proof}

\begin{lemma}
\label{derived-lemma-product-K-injective}
设 $\mathcal{A}$ 为阿贝尔范畴。设 $T$ 为集合，并对每个
$t \in T$ 设 $I_t^\bullet$ 为 K-内射复形。
若对所有 $n$，$I^n = \prod_t I_t^n$ 都存在，则 $I^\bullet$
是 K-内射复形。此外，$I^\bullet$ 表示各对象 $I_t^\bullet$
在 $D(\mathcal{A})$ 中的积。
\end{lemma}

\begin{proof}
% P02-E0074: frozen prose reads “an complex” rather than “a complex”.
设 $K^\bullet$ 为一个复形。注意，复形
$$
C :
\prod\nolimits_b \Hom(K^{-b}, I^{b - 1}) \to
\prod\nolimits_b \Hom(K^{-b}, I^b) \to
\prod\nolimits_b \Hom(K^{-b}, I^{b + 1})
$$
在中间的上同调为
$\Hom_{K(\mathcal{A})}(K^\bullet, I^\bullet)$。
类似地，复形
$$
C_t :
\prod\nolimits_b \Hom(K^{-b}, I_t^{b - 1}) \to
\prod\nolimits_b \Hom(K^{-b}, I_t^b) \to
\prod\nolimits_b \Hom(K^{-b}, I_t^{b + 1})
$$
计算 $\Hom_{K(\mathcal{A})}(K^\bullet, I_t^\bullet)$。
接着注意，由对 $I$ 的选择，我们有阿贝尔群复形的等式
$$
C = \prod\nolimits_{t \in T} C_t
$$
取积是阿贝尔群范畴上的正合函子。因此，若 $K^\bullet$ 无环，
则 $\Hom_{K(\mathcal{A})}(K^\bullet, I_t^\bullet) = 0$，
故 $C_t$ 无环，故 $C$ 无环，进而
$\Hom_{K(\mathcal{A})}(K^\bullet, I^\bullet) = 0$。
于是 $I^\bullet$ 是 K-内射的。
由此，再用引理 \ref{derived-lemma-K-injective} 可得
$$
\Hom_{D(\mathcal{A})}(K^\bullet, I^\bullet)
=
\prod\nolimits_{t \in T} \Hom_{D(\mathcal{A})}(K^\bullet, I_t^\bullet)
$$
而且 $I^\bullet$ 的确表示导出范畴中的积。
\end{proof}

\begin{lemma}
\label{derived-lemma-K-injective-defined}
设 $\mathcal{A}$ 为阿贝尔范畴。
设 $F : K(\mathcal{A}) \to \mathcal{D}'$ 为三角范畴的正合函子。
则 $RF$ 在 $K(\mathcal{A})$ 中每个拟同构于 K-内射复形的复形处
都有定义。事实上，每个 K-内射复形都计算 $RF$。
\end{lemma}

\begin{proof}
由引理 \ref{derived-lemma-derived-inverts}，只须证明 $RF$ 在 K-内射复形处
有定义，即只须证明 K-内射复形 $I^\bullet$ 计算 $RF$。
考虑拟同构 $I^\bullet \to N^\bullet$。由引理
\ref{derived-lemma-K-injective}，它有左逆。
因此 $I^\bullet \to I^\bullet$ 是
$I^\bullet/\text{Qis}(\mathcal{A})$ 的终对象，结论得证。
\end{proof}

\begin{lemma}
\label{derived-lemma-enough-K-injectives-implies}
设 $\mathcal{A}$ 为阿贝尔范畴。
假设每个复形都有取值于某个 K-内射复形的拟同构。
则三角范畴的任意正合函子
$F : K(\mathcal{A}) \to \mathcal{D}'$ 都有右导出函子
$$
RF : D(\mathcal{A}) \longrightarrow \mathcal{D}'
$$
且对 K-内射复形 $I^\bullet$，有
$RF(I^\bullet) = F(I^\bullet)$。
\end{lemma}

\begin{proof}
为此，对 K-内射复形的类 $\mathcal{I}$ 应用引理
\ref{derived-lemma-find-existence-computes}。由于 (1) 由假设成立，只须证明：
若 $I^\bullet \to J^\bullet$ 是 K-内射复形之间的拟同构，
则 $F(I^\bullet) \to F(J^\bullet)$ 是同构。
这是显然的，因为由引理 \ref{derived-lemma-K-injective}，
$I^\bullet \to J^\bullet$ 是同伦等价，即 $K(\mathcal{A})$ 中的同构。
\end{proof}

\noindent
下列引理可推广到以更大序数为指标的极限。

\begin{lemma}
\label{derived-lemma-limit-K-injectives}
\begin{slogan}
K-内射复形的“分裂”塔的极限是 K-内射的。
\end{slogan}
设 $\mathcal{A}$ 为阿贝尔范畴。设
$$
\ldots \to I_3^\bullet \to I_2^\bullet \to I_1^\bullet
$$
为复形的逆系统。假设
\begin{enumerate}
\item 每个 $I_n^\bullet$ 都是 $K$-内射的；
\item 每个映射 $I_{n + 1}^m \to I_n^m$ 都是分裂满射；
\item 极限 $I^m = \lim I_n^m$ 存在。
\end{enumerate}
则复形 $I^\bullet$ 是 K-内射的。
\end{lemma}

\begin{proof}
我们强烈建议读者跳过本引理的证明。
设 $M^\bullet$ 为无环复形。简记
$H_n(a, b) = \Hom_\mathcal{A}(M^a, I_n^b)$。用此记号，
$\Hom_{K(\mathcal{A})}(M^\bullet, I^\bullet)$ 是复形
$$
\prod_m \lim\limits_n H_n(m, m - 2)
\to
\prod_m \lim\limits_n H_n(m, m - 1)
\to
\prod_m \lim\limits_n H_n(m, m)
\to
\prod_m \lim\limits_n H_n(m, m + 1)
$$
从左起第三处的上同调。
可以交换 $\prod$ 与 $\lim$ 的次序，而每个复形
$$
\prod_m H_n(m, m - 2)
\to
\prod_m H_n(m, m - 1)
\to
\prod_m H_n(m, m)
\to
\prod_m H_n(m, m + 1)
$$
都由假设 (1) 正合。由假设 (2)，系统中的映射
$$
\ldots \to
\prod_m H_3(m, m - 2) \to
\prod_m H_2(m, m - 2) \to
\prod_m H_1(m, m - 2)
$$
都是满射。因此本引理由同调，引理
\ref{homology-lemma-apply-Mittag-Leffler} 得出。
\end{proof}

\noindent
看来，引理 \ref{derived-lemma-special-inverse-system}、
\ref{derived-lemma-bounded-below-injectives-K-injective} 与
\ref{derived-lemma-limit-K-injectives} 的组合，会为任意具有足够多内射对象和
可数积的阿贝尔范畴产生“足够多 K-内射对象”。
事实上，这可能行不通！解释见引理
\ref{derived-lemma-difficulty-K-injectives}。

\begin{lemma}
\label{derived-lemma-adjoint-preserve-K-injectives}
设 $\mathcal{A}$ 和 $\mathcal{B}$ 为阿贝尔范畴。
设 $u : \mathcal{A} \to \mathcal{B}$ 和
$v : \mathcal{B} \to \mathcal{A}$ 为加性函子。假设
\begin{enumerate}
\item $u$ 是 $v$ 的右伴随；并且
\item $v$ 正合。
\end{enumerate}
则 $u$ 把 K-内射复形变为 K-内射复形。
\end{lemma}

\begin{proof}
设 $I^\bullet$ 为 $\mathcal{A}$ 的 K-内射复形。
% P02-E0075: frozen prose reads “a acyclic complex”.
设 $M^\bullet$ 为 $\mathcal{B}$ 的无环复形。
由于 $v$ 正合，$v(M^\bullet)$ 是无环复形。由伴随性得到
$$
0 = \Hom_{K(\mathcal{A})}(v(M^\bullet), I^\bullet) =
\Hom_{K(\mathcal{B})}(M^\bullet, u(I^\bullet))
$$
故引理得证。
\end{proof}





\section{有界上同调维数}
\label{derived-section-bounded}

\noindent
还有另一种无界导出函子存在的情形，即函子具有有界上同调维数时。

\begin{lemma}
\label{derived-lemma-replace-resolution}
设 $\mathcal{A}$ 为阿贝尔范畴。设
$d : \Ob(\mathcal{A}) \to \{0, 1, 2, \ldots, \infty\}$ 为函数。
假设
\begin{enumerate}
\item $\mathcal{A}$ 的每个对象都是某个满足 $d(A) = 0$ 的对象
$A$ 的子对象；
\item 对 $A, B \in \mathcal{A}$，有
$d(A \oplus B) \leq \max \{d(A), d(B)\}$；并且
\item 若 $0 \to A \to B \to C \to 0$ 短正合，则
$d(C) \leq \max\{d(A) - 1, d(B)\}$。
\end{enumerate}
设 $K^\bullet$ 为复形，使得当 $n \to -\infty$ 时，
$n + d(K^n)$ 趋于 $-\infty$。则存在拟同构
$K^\bullet \to L^\bullet$，使得对所有 $n \in \mathbf{Z}$，
都有 $d(L^n) = 0$。
\end{lemma}

\begin{proof}
由引理 \ref{derived-lemma-subcategory-right-resolution}，可找到拟同构
$\sigma_{\geq 0}K^\bullet \to M^\bullet$，其中当 $n < 0$ 时
$M^n = 0$，而当 $n \geq 0$ 时 $d(M^n) = 0$。于是 $K^\bullet$
拟同构于复形
$$
\ldots \to K^{-2} \to K^{-1} \to M^0 \to M^1 \to \ldots
$$
因此可假设当 $n \gg 0$ 时 $d(K^n) = 0$。
注意，这一替换不破坏当 $n \to -\infty$ 时
$n + d(K^n) \to -\infty$ 的条件。

\medskip\noindent
我们将通过（无限）一列基本替换来改进 $K^\bullet$。
所谓{\it 基本替换}如下。选取满足 $d(K^n) > 0$ 的指标 $n$。
选取单射 $K^n \to M$，其中 $d(M) = 0$。
令 $M' = \Coker(K^n \to M \oplus K^{n + 1})$。
考虑复形映射
$$
\xymatrix{
K^\bullet : \ar[d] &
K^{n - 1} \ar[d] \ar[r] &
K^n \ar[d] \ar[r] &
K^{n + 1} \ar[d] \ar[r] &
K^{n + 2} \ar[d] \\
(K')^\bullet : &
K^{n - 1} \ar[r] &
M \ar[r] &
M' \ar[r] &
K^{n + 2}
}
$$
显然 $K^\bullet \to (K')^\bullet$ 是拟同构。
此外，显然 $d((K')^n) = 0$，并且
$$
d((K')^{n + 1}) \leq \max\{d(K^n) - 1, d(M \oplus K^{n + 1})\} \leq
\max\{d(K^n) - 1, d(K^{n + 1})\}
$$
其他值不变。

\medskip\noindent
为完成证明，我们仔细选择进行基本替换的顺序，使得对每个整数 $m$，
复形 $\sigma_{\geq m}K^\bullet$ 只被改变有限多次。为此令
$$
\xi(K^\bullet) = \max \{n + d(K^n) \mid d(K^n) > 0\}
$$
并令
$$
I = \{n \in \mathbf{Z} \mid \xi(K^\bullet) = n + d(K^n)
\text{ 且 }
 d(K^n) > 0\}
$$
我们的假设：当 $n \to -\infty$ 时 $n + d(K^n)$ 趋于
$-\infty$，以及当 $n >> 0$ 时 $d(K^n) = 0$ 这一事实，
蕴含 $\xi(K^\bullet) < +\infty$ 且 $I$ 是有限集。
显然，对上述基本变换，有
$\xi((K')^\bullet) \leq \xi(K^\bullet)$。
一次基本变换在次数 $\leq \xi(K^\bullet) + 1$ 处改变复形。
因此，若能找到降低 $\xi(K^\bullet)$ 的有限列基本变换，结论即得。
% P02-E0076: frozen prose says “find finite sequence” and omits the article.
不过，注意，若从最小元素 $n \in I$ 开始进行基本变换，
则或者 $I$ 的大小减小，或者 $\min I$ 增大。
由于 $I$ 的每个元素都 $\leq \xi(K^\bullet)$，经过有限步便得结论。
\end{proof}

\begin{lemma}
\label{derived-lemma-unbounded-right-derived}
设 $F : \mathcal{A} \to \mathcal{B}$ 为阿贝尔范畴之间的左正合函子。
假设
\begin{enumerate}
\item $\mathcal{A}$ 的每个对象都是某个对 $F$ 右无环的对象的子对象；
\item 存在整数 $n \geq 0$，使得 $R^nF = 0$。
\end{enumerate}
则
\begin{enumerate}
\item $RF : D(\mathcal{A}) \to D(\mathcal{B})$ 存在；
\item 由对 $F$ 右无环的对象组成的任意复形都计算 $RF$；
\item 任意复形都是某个取值于由对 $F$ 右无环对象组成的复形中的
拟同构的源；
\item 对 $E \in D(\mathcal{A})$：
\begin{enumerate}
% P02-E0077: the frozen map omits a closing parenthesis after RF(tau_{<=a}E).
\item 当 $i \leq a$ 时，
$H^i(RF(\tau_{\leq a}E) \to H^i(RF(E))$ 是同构；
\item 当 $i \geq b$ 时，
$H^i(RF(E)) \to H^i(RF(\tau_{\geq b - n + 1}E))$ 是同构；
\item 若对某个 $-\infty \leq a \leq b \leq \infty$，
当 $i \not \in [a, b]$ 时 $H^i(E) = 0$，则当
$i \not \in [a, b + n - 1]$ 时 $H^i(RF(E)) = 0$。
\end{enumerate}
\end{enumerate}
\end{lemma}

\begin{proof}
注意，第一个假设蕴含
$RF : D^+(\mathcal{A}) \to D^+(\mathcal{B})$ 存在；见命题
\ref{derived-proposition-enough-acyclics}。
设 $A$ 为 $\mathcal{A}$ 的对象。选取单射 $A \to A'$，其中
$A'$ 无环。则由上同调长正合序列，
$R^{n + 1}F(A) = R^nF(A'/A) = 0$。
因此 $R^{n + 1}F = 0$。继续如此归纳，得到对所有
$m \geq n$ 都有 $R^mF = 0$。

\medskip\noindent
我们将对函数
$d : \Ob(\mathcal{A}) \to \{0, 1, 2, \ldots \}$ 应用引理
\ref{derived-lemma-replace-resolution}，其中
% P02-E0078: frozen d(A) is malformed as max {0} union {i | ...}, rather than max of the union.
$d(A) = \max \{0\} \cup \{i \mid R^iF(A) \not = 0\}$。
引理 \ref{derived-lemma-replace-resolution} 的第一个假设就是我们的假设 (1)。
引理 \ref{derived-lemma-replace-resolution} 的第二个假设由
$RF(A \oplus B) = RF(A) \oplus RF(B)$ 得出。
引理 \ref{derived-lemma-replace-resolution} 的第三个假设由上同调长正合序列得出。
因此，对每个复形 $K^\bullet$，都存在取值于由对 $F$ 右无环的对象组成的
复形中的拟同构 $K^\bullet \to L^\bullet$。这证明陈述 (3)。

\medskip\noindent
我们断言，若 $L^\bullet \to M^\bullet$ 是由对 $F$ 右无环的对象组成的
复形之间的拟同构，则
$F(L^\bullet) \to F(M^\bullet)$ 是拟同构。
若证明此断言，则由引理 \ref{derived-lemma-find-existence-computes}，
得到引理的陈述 (1) 和 (2)。
为证明断言，选取整数 $i \in \mathbf{Z}$。
考虑特异三角
$$
\sigma_{\geq i - n - 1}L^\bullet \to
\sigma_{\geq i - n - 1}M^\bullet \to Q^\bullet,
$$
即令 $Q^\bullet$ 为第一个映射的锥。
注意，$Q^\bullet$ 有下界，且除可能在
$j = i - n - 1$ 或 $j = i - n - 2$ 外，$H^j(Q^\bullet)$ 均为零。
可以对 $Q^\bullet$ 应用 $RF$。
利用引理 \ref{derived-lemma-two-ss-complex-functor} 的第二个谱序列和
所假设的上同调消失 (2)，得出除可能对
$j \in \{i - n - 2, \ldots, i - 1\}$ 外，
$H^j(RF(Q^\bullet))$ 均为零。因此
$RF(\sigma_{\geq i - n - 1}L^\bullet) \to RF(\sigma_{\geq i - n - 1}M^\bullet)$
在次数 $\geq i$ 的上同调对象上诱导同构。
由命题 \ref{derived-proposition-enough-acyclics}，我们知道
$RF(\sigma_{\geq i - n - 1}L^\bullet) = \sigma_{\geq i - n - 1}F(L^\bullet)$
以及
$RF(\sigma_{\geq i - n - 1}M^\bullet) = \sigma_{\geq i - n - 1}F(M^\bullet)$。
所以如所需，$F(L^\bullet) \to F(M^\bullet)$ 在次数 $i$ 是同构。

\medskip\noindent
(4)(a) 由引理 \ref{derived-lemma-negative-vanishing} 得出。

\medskip\noindent
为证明 (4)(b)，以由对 $F$ 右无环的对象组成的复形 $L^\bullet$
表示 $E$。由 (2)，$RF(E)$ 由复形 $F(L^\bullet)$ 表示，而
% P02-E0079: frozen proof introduces sigma_{>=c} although c is not defined in this argument.
$RF(\sigma_{\geq c}L^\bullet)$ 由
$\sigma_{\geq c}F(L^\bullet)$ 表示。考虑评注
\ref{derived-remark-truncation-distinguished-triangle} 的特异三角
$$
H^{b - n}(L^\bullet)[n - b] \to
\tau_{\geq b - n}L^\bullet \to
\tau_{\geq b - n + 1}L^\bullet
$$
上述消失性给出：当 $i \geq b$ 时，
$H^i(RF(\tau_{\geq b - n}L^\bullet))$ 与
$H^i(RF(\tau_{\geq b - n + 1}L^\bullet))$ 相同。
考虑复形的短正合序列
$$
0 \to
\Im(L^{b - n - 1} \to L^{b - n})[n - b] \to
\sigma_{\geq b - n}L^\bullet \to
\tau_{\geq b - n}L^\bullet \to 0
$$
利用与此关联的特异三角（见第
\ref{derived-section-canonical-delta-functor} 节）以及与此前相同的消失性，
得出当 $i \geq b$ 时，
$H^i(RF(\tau_{\geq b - n}L^\bullet))$ 与
$H^i(RF(\sigma_{\geq b - n}L^\bullet))$ 相同。
由于映射 $RF(\sigma_{\geq b - n}L^\bullet) \to RF(L^\bullet)$
由 $\sigma_{\geq b - n}F(L^\bullet) \to F(L^\bullet)$ 表示，
所以如所需，当 $i \geq b$ 时，它又与
$H^i(RF(L^\bullet))$ 相同。

\medskip\noindent
证明 (4)(c)。在关于 $E$ 的假设下，有
$\tau_{\leq a - 1}E = 0$，并由 (4)(a) 得到当
$i \leq a - 1$ 时 $H^i(RF(E))$ 消失。
类似地，有 $\tau_{\geq b + 1}E = 0$，因而由 (4)(b) 得到当
$i \geq b + n$ 时 $H^i(RF(E))$ 消失。
\end{proof}

\begin{lemma}
\label{derived-lemma-unbounded-left-derived}
设 $F : \mathcal{A} \to \mathcal{B}$ 为阿贝尔范畴之间的右正合函子。
若
\begin{enumerate}
\item $\mathcal{A}$ 的每个对象都是某个对 $F$ 左无环的对象的商；
\item 存在整数 $n \geq 0$，使得 $L^nF = 0$。
\end{enumerate}
则
\begin{enumerate}
\item $LF : D(\mathcal{A}) \to D(\mathcal{B})$ 存在；
\item 由对 $F$ 左无环的对象组成的任意复形都计算 $LF$；
\item 任意复形都是一个拟同构的靶，该拟同构的源是由对 $F$
左无环的对象组成的复形；
\item 对 $E \in D(\mathcal{A})$：
\begin{enumerate}
% P02-E0080: the frozen map omits a closing parenthesis after LF(tau_{<=a+n-1}E).
\item 当 $i \leq a$ 时，
$H^i(LF(\tau_{\leq a + n - 1}E) \to H^i(LF(E))$ 是同构；
\item 当 $i \geq b$ 时，
$H^i(LF(E)) \to H^i(LF(\tau_{\geq b}E))$ 是同构；
\item 若对某个 $-\infty \leq a \leq b \leq \infty$，
当 $i \not \in [a, b]$ 时 $H^i(E) = 0$，则当
$i \not \in [a - n + 1, b]$ 时 $H^i(LF(E)) = 0$。
\end{enumerate}
\end{enumerate}
\end{lemma}

\begin{proof}
这对偶于引理 \ref{derived-lemma-unbounded-right-derived}。
\end{proof}






\section{导出余极限}
\label{derived-section-derived-colimit}

\noindent
三角范畴中有导出余极限的概念。

\begin{definition}
\label{derived-definition-derived-colimit}
设 $\mathcal{D}$ 为三角范畴。
设 $(K_n, f_n)$ 为 $\mathcal{D}$ 中对象的系统。
若直和 $\bigoplus K_n$ 存在，并且存在特异三角
$$
\bigoplus K_n \to \bigoplus K_n \to K \to \bigoplus K_n[1]
$$
其中映射 $\bigoplus K_n \to \bigoplus K_n$ 在次数 $n$
由 $1 - f_n$ 给出，则称对象 $K$ 为系统 $(K_n)$ 的
{\it 导出余极限}，或{\it 同伦余极限}。
在此情形下，有时用记号 $K = \text{hocolim} K_n$ 表示。
\end{definition}

\noindent
由 TR3，若导出余极限存在，则它在（不唯一的）同构意义下唯一。
此外，由 TR1，只要 $\bigoplus K_n$ 存在，$K_n$ 的导出余极限就存在。
阿贝尔群范畴的导出范畴 $D(\textit{Ab})$ 是所有同伦余极限都存在的
三角范畴的一个例子。

\medskip\noindent
不唯一性使得难以确定导出余极限。在更多代数，引理
\ref{more-algebra-lemma-map-from-hocolim} 中，读者可找到正合序列
$$
0 \to R^1\lim \Hom(K_n, L[-1]) \to \Hom(\text{hocolim} K_n, L)
\to \lim \Hom(K_n, L) \to 0
$$
它用通常的 $\Hom$ 群描述从同伦余极限出发的 $\Hom$ 群。

\begin{remark}
\label{derived-remark-uniqueness-derived-colimit}
设 $\mathcal{D}$ 为三角范畴。
设 $(K_n, f_n)$ 为 $\mathcal{D}$ 中对象的系统。
可以把导出余极限看作 $\mathcal{D}$ 中的对象 $K$，连同态射
$i_n : K_n \to K$，使得 $i_{n + 1} \circ f_n = i_n$，
并且存在态射 $c : K \to \bigoplus K_n[1]$，满足
$$
\bigoplus K_n \xrightarrow{1 - f_n} \bigoplus K_n \xrightarrow{i_n}
K \xrightarrow{c} \bigoplus K_n[1]
$$
为特异三角。若 $(K', i'_n, c')$ 是第二个导出余极限，
则存在同构 $\varphi : K \to K'$，使得
$\varphi \circ i_n = i'_n$ 且 $c' \circ \varphi = c$。
$\varphi$ 的存在性是 TR3，而 $\varphi$ 是同构这一事实由引理
\ref{derived-lemma-third-isomorphism-triangle} 给出。
\end{remark}

\begin{remark}
\label{derived-remark-functoriality-derived-colimit}
设 $\mathcal{D}$ 为三角范畴。
设 $(a_n) : (K_n, f_n) \to (L_n, g_n)$ 为 $\mathcal{D}$ 中对象系统的态射。
设 $(K, i_n, c)$ 为第一个系统的导出余极限，
$(L, j_n, d)$ 为第二个系统的导出余极限，记号如评注
\ref{derived-remark-uniqueness-derived-colimit}。
则存在态射 $a : K \to L$，使得
$a \circ i_n = j_n$ 且 $d \circ a = (a_n[1]) \circ c$。
这由对定义这些对象的特异三角应用 TR3 得出。
\end{remark}

\begin{lemma}
\label{derived-lemma-hocolim-subsequence}
设 $\mathcal{D}$ 为三角范畴。
设 $(K_n, f_n)$ 为 $\mathcal{D}$ 中对象的系统。
设 $n_1 < n_2 < n_3 < \ldots$ 为整数序列。
假设 $\bigoplus K_n$ 和 $\bigoplus K_{n_i}$ 存在。
则存在同构
$\text{hocolim} K_{n_i} \to \text{hocolim} K_n$，使得对所有 $i$，图
$$
\xymatrix{
K_{n_i} \ar[r] \ar[d]_{\text{id}} & \text{hocolim} K_{n_i} \ar[d] \\
K_{n_i} \ar[r] & \text{hocolim} K_n
}
$$
交换。
\end{lemma}

\begin{proof}
设 $g_i : K_{n_i} \to K_{n_{i + 1}}$ 为复合
$f_{n_{i + 1} - 1} \circ \ldots \circ f_{n_i}$。
如下构造交换图
$$
\vcenter{
\xymatrix{
\bigoplus\nolimits_i K_{n_i} \ar[r]_{1 - g_i} \ar[d]_b &
\bigoplus\nolimits_i K_{n_i} \ar[d]^a \\
\bigoplus\nolimits_n K_n \ar[r]^{1 - f_n} &
\bigoplus\nolimits_n K_n
}
}
\quad\text{且}\quad
\vcenter{
\xymatrix{
\bigoplus\nolimits_n K_n \ar[r]_{1 - f_n} \ar[d]_d &
\bigoplus\nolimits_n K_n \ar[d]^c \\
\bigoplus\nolimits_i K_{n_i} \ar[r]^{1 - g_i} &
\bigoplus\nolimits_i K_{n_i}
}
}
$$
令 $a_i = a|_{K_{n_i}}$ 为 $K_{n_i}$ 到直和中的包含映射。
换言之，$a$ 是自然包含。令 $b_i = b|_{K_{n_i}}$ 为映射
$$
K_{n_i}
\xrightarrow{1,\ f_{n_i},\ f_{n_i + 1} \circ f_{n_i},
\ \ldots,\ f_{n_{i + 1} - 2} \circ \ldots \circ f_{n_i}}
K_{n_i} \oplus K_{n_i + 1} \oplus \ldots \oplus K_{n_{i + 1} - 1}
$$
若 $n_{i - 1} < j \leq n_i$，则令 $c_j = c|_{K_j}$ 为映射
$$
K_j \xrightarrow{f_{n_i - 1} \circ \ldots \circ f_j} K_{n_i}
$$
若对任意 $i$ 都有 $j \not = n_i$，则令 $d_j = d|_{K_j}$ 为零；
并令 $d_{n_i}$ 为 $K_{n_i}$ 到直和中的自然包含。
换言之，$d$ 是自然投影。
由 TR3，这些图定义态射
$$
\varphi : \text{hocolim} K_{n_i} \to \text{hocolim} K_n
\quad\text{且}\quad
\psi : \text{hocolim} K_n \to \text{hocolim} K_{n_i}
$$
由于 $c \circ a$ 和 $d \circ b$ 是恒等映射，由引理
\ref{derived-lemma-third-isomorphism-triangle}，可知
$\varphi \circ \psi$ 是同构。
反过来，得到态射 $a \circ c$ 和 $b \circ d$。
考虑由如下规则给出的态射
$h = (h_j) : \bigoplus K_n \to \bigoplus K_n$：
对 $n_{i - 1} < j < n_i$，令
$$
h_j : K_j
\xrightarrow{1,\ f_j,\ f_{j + 1} \circ f_j,
\ \ldots,\ f_{n_i - 1} \circ \ldots \circ f_j}
K_j \oplus \ldots \oplus K_{n_i}
$$
于是读者可验证
$(1 - f) \circ h = \text{id} - a \circ c$ 且
$h \circ (1 - f) = \text{id} - b \circ d$。
这意味着，由引理 \ref{derived-lemma-third-map-square-zero}，
$\text{id} - \psi \circ \varphi$ 的平方为零（略去一个小论证）。
换言之，$\psi \circ \varphi$ 与恒等映射之差为幂零自同态，
故它是同构。因此如所需，$\varphi$ 与 $\psi$ 都是同构。
\end{proof}

\begin{lemma}
\label{derived-lemma-direct-sums}
设 $\mathcal{A}$ 为阿贝尔范畴。
若 $\mathcal{A}$ 有正合的可数直和，则 $D(\mathcal{A})$ 有可数直和。
事实上，给定以可数指标集 $I$ 为指标的一族复形 $K_i^\bullet$，
逐项直和 $\bigoplus K_i^\bullet$ 是各 $K_i^\bullet$ 在
$D(\mathcal{A})$ 中的直和。
\end{lemma}

\begin{proof}
设 $L^\bullet$ 为复形。假设给定 $D(\mathcal{A})$ 中的映射
$\alpha_i : K_i^\bullet \to L^\bullet$。
这意味着存在复形的拟同构
$s_i : M_i^\bullet \to K_i^\bullet$ 和复形映射
$f_i : M_i^\bullet \to L^\bullet$，使得
$\alpha_i = f_is_i^{-1}$。由假设，复形映射
$$
s : \bigoplus M_i^\bullet \longrightarrow \bigoplus K_i^\bullet
$$
是拟同构。因此，令 $f = \bigoplus f_i$，可见
$\alpha = fs^{-1}$ 是 $D(\mathcal{A})$ 中的映射，
它与余投影 $K_i^\bullet \to \bigoplus K_i^\bullet$ 的复合为
$\alpha_i$。我们略去验证 $\alpha$ 的唯一性。
\end{proof}

\begin{lemma}
\label{derived-lemma-compute-colimit}
设 $\mathcal{A}$ 为阿贝尔范畴。假设以 $\mathbf{N}$ 为指标的余极限存在且
正合。则可数直和存在且正合。此外，若 $(A_n, f_n)$ 是以
$\mathbf{N}$ 为指标的系统，则存在短正合序列
$$
0 \to \bigoplus A_n \to \bigoplus A_n \to \colim A_n \to 0
$$
其中第一个映射在次数 $n$ 由 $1 - f_n$ 给出。
\end{lemma}

\begin{proof}
第一个陈述由
$\bigoplus A_n = \colim (A_1 \oplus \ldots \oplus A_n)$ 得出。
对于第二个陈述，注意对每个 $n$ 都有短正合序列
$$
0 \to
A_1 \oplus \ldots \oplus A_{n - 1} \to
A_1 \oplus \ldots \oplus A_n \to A_n \to 0
$$
其中第一个映射由映射 $1 - f_i$ 给出，第二个映射是转移映射之和。
取余极限便得到引理中的序列。
\end{proof}

\begin{lemma}
\label{derived-lemma-colim-hocolim}
设 $\mathcal{A}$ 为阿贝尔范畴。设 $L_n^\bullet$ 为
$\mathcal{A}$ 中复形的系统。假设 $\mathcal{A}$ 中以
$\mathbf{N}$ 为指标的余极限存在且正合。
则逐项余极限 $L^\bullet = \colim L_n^\bullet$ 是该系统在
$D(\mathcal{A})$ 中的同伦余极限。
\end{lemma}

\begin{proof}
由引理 \ref{derived-lemma-compute-colimit}，存在复形的正合序列
$$
0 \to \bigoplus L_n^\bullet \to \bigoplus L_n^\bullet \to L^\bullet \to 0
$$
由引理 \ref{derived-lemma-direct-sums}，这些直和是 $D(\mathcal{A})$ 中的直和。
因此，结论由定义 \ref{derived-definition-derived-colimit} 中导出余极限的定义，
以及复形的短正合序列给出特异三角这一事实得出
（引理 \ref{derived-lemma-derived-canonical-delta-functor}）。
\end{proof}

\begin{lemma}
\label{derived-lemma-cohomology-of-hocolim}
设 $\mathcal{D}$ 为具有可数直和的三角范畴。
设 $\mathcal{A}$ 为具有正合的、以 $\mathbf{N}$ 为指标的余极限的
阿贝尔范畴。设 $H : \mathcal{D} \to \mathcal{A}$ 为与可数直和交换的
同调函子。则对 $\mathcal{D}$ 中任意对象系统，都有
$H(\text{hocolim} K_n) = \colim H(K_n)$。
\end{lemma}

\begin{proof}
写 $K = \text{hocolim} K_n$。对定义它的特异三角应用 $H$，得到
$$
\bigoplus H(K_n) \to \bigoplus H(K_n)
\to H(K) \to
\bigoplus H(K_n[1]) \to \bigoplus H(K_n[1])
$$
其中第一个映射由 $1 - H(f_n)$ 给出，最后一个映射由
$1 - H(f_n[1])$ 给出。应用引理
\ref{derived-lemma-compute-colimit} 即证明本引理。
\end{proof}

\noindent
下列引理告诉我们，从紧对象（稍后定义）出发取映射与导出余极限交换。

\begin{lemma}
\label{derived-lemma-commutes-with-countable-sums}
设 $\mathcal{D}$ 为具有可数直和的三角范畴。
设 $K \in \mathcal{D}$ 为如下对象：对任意可数对象集
$E_n \in \mathcal{D}$，典范映射
$$
\bigoplus \Hom_\mathcal{D}(K, E_n)
\longrightarrow
\Hom_\mathcal{D}(K, \bigoplus E_n)
$$
都是双射。则给定 $\mathcal{D}$ 中任意以 $\mathbf{N}$ 为指标的系统
$L_n$，若其导出余极限 $L = \text{hocolim} L_n$ 存在，则映射
$$
\colim \Hom_\mathcal{D}(K, L_n) \longrightarrow \Hom_\mathcal{D}(K, L)
$$
是双射。
\end{lemma}

\begin{proof}
由引理 \ref{derived-lemma-representable-homological}，这是引理
\ref{derived-lemma-cohomology-of-hocolim} 的特例。
\end{proof}







\section{导出极限}
\label{derived-section-derived-limit}

\noindent
在三角范畴中有导出极限这一概念。

\begin{definition}
\label{derived-definition-derived-limit}
设 $\mathcal{D}$ 为三角范畴。
设 $(K_n, f_n)$ 为 $\mathcal{D}$ 中对象的逆系统。
若直积 $\prod K_n$ 存在，且有特异三角
$$
K \to \prod K_n \to \prod K_n \to K[1]
$$
其中映射 $\prod K_n \to \prod K_n$ 由
$(k_n) \mapsto (k_n - f_{n + 1}(k_{n + 1}))$ 给出，
则称对象 $K$ 为系统 $(K_n)$ 的一个{\it 导出极限}，
或一个{\it 同伦极限}。在这种情况下，有时记作 $K = R\lim K_n$。
\end{definition}

\noindent
由 TR3，导出极限若存在，则在（非唯一的）同构意义下唯一。
此外，由 TR1，只要 $\prod K_n$ 存在，导出极限 $R\lim K_n$ 就存在。
阿贝尔群范畴的导出范畴 $D(\textit{Ab})$ 是所有导出极限均存在的
三角范畴的一个例子。

\medskip\noindent
这种非唯一性使得难以精确刻画导出极限。在 More on Algebra，
引理 \ref{more-algebra-lemma-map-into-Rlim} 中，读者可找到正合列
$$
0 \to R^1\lim \Hom(L, K_n[-1]) \to \Hom(L, R\lim K_n)
\to \lim \Hom(L, K_n) \to 0
$$
它用通常的 $\Hom$ 描述了到导出极限的 $\Hom$。

\begin{lemma}
\label{derived-lemma-products}
设 $\mathcal{A}$ 为具有正合可数直积的阿贝尔范畴。则
\begin{enumerate}
\item $D(\mathcal{A})$ 有可数直积；
\item $D(\mathcal{A})$ 中的可数直积 $\prod K_i$，可通过对代表
$K_i$ 的任意复形逐项取直积而得到；
\item $H^p(\prod K_i) = \prod H^p(K_i)$。
\end{enumerate}
\end{lemma}

\begin{proof}
设 $K_i^\bullet$ 为 $D(\mathcal{A})$ 中代表 $K_i$ 的复形。
设 $L^\bullet$ 为复形。假设给定 $D(\mathcal{A})$ 中的映射
$\alpha_i : L^\bullet \to K_i^\bullet$。
这意味着存在复形的拟同构 $s_i : K_i^\bullet \to M_i^\bullet$
和复形映射 $f_i : L^\bullet \to M_i^\bullet$，使得
$\alpha_i = s_i^{-1}f_i$。由假设，复形映射
$$
s : \prod K_i^\bullet \longrightarrow \prod M_i^\bullet
$$
是拟同构。因此，令 $f = \prod f_i$，可见
$\alpha = s^{-1}f$ 是 $D(\mathcal{A})$ 中的映射，
它与投影 $\prod K_i^\bullet \to K_i^\bullet$ 的复合为 $\alpha_i$。
我们略去验证 $\alpha$ 的唯一性。
\end{proof}

\noindent
这里应当陈述并证明引理 \ref{derived-lemma-compute-colimit}、
\ref{derived-lemma-colim-hocolim} 和
\ref{derived-lemma-commutes-with-countable-sums} 的对偶形式。
不过目前我们不知道这些引理有什么应用。

\begin{lemma}
\label{derived-lemma-inverse-limit-bounded-below}
设 $\mathcal{A}$ 为具有可数直积和足够多内射对象的阿贝尔范畴。
设 $(K_n)$ 为 $D^+(\mathcal{A})$ 中的逆系统。
则 $R\lim K_n$ 存在。
\end{lemma}

\begin{proof}
只需证明 $\prod K_n$ 在 $D(\mathcal{A})$ 中存在。
对每个 $n$，可用内射对象组成的下有界复形 $I_n^\bullet$ 代表
$K_n$（引理 \ref{derived-lemma-injective-resolutions-exist}）。
于是 $\prod K_n$ 由 $\prod I_n^\bullet$ 代表，见引理
\ref{derived-lemma-product-K-injective}。
\end{proof}

\begin{remark}
\label{derived-remark-functoriality-derived-limit}
设 $\mathcal{D}$ 为三角范畴。设 $(K_n)$ 和 $(L_n)$ 为
$\mathcal{D}$ 中具有导出极限 $K$ 和 $L$ 的对象逆系统。
假设有 pro-对象的态射 $a : (K_n) \to (L_n)$，见 Categories，
例 \ref{categories-example-pro-morphism-inverse-systems}。
这意味着给定映射 $a_n : K_{m(n)} \to L_n$。不妨假设
$m(1) < m(2) < m(3) < \ldots$，且各映射 $a_n$ 与系统中的
转移映射相容。于是可以考虑映射
$$
a', a'' : \prod K_n \longrightarrow \prod L_n
$$
它们由规则 $a'((k_n)) = (a_n(k_{m(n)}))$ 和
$$
a''((k_n)) =
(a_n(k_{m(n)} + f(k_{m(n) + 1}) + \ldots + f(k_{m(n + 1) - 1})))
$$
定义，
% P02-E0081: frozen source misspells “occurrence” as “occurence”.
其中每次出现的 $f$ 都表示逆系统 $(K_n)$ 的一个适当转移映射。
于是实线图
$$
\xymatrix{
K \ar[r] \ar@{..>}[d] & \prod K_n \ar[r] \ar[d]^{a'} &
\prod K_n \ar[d]^{a''} \\
L \ar[r] & \prod L_n \ar[r] & \prod L_n
}
$$
交换，并且由 TR3 得到虚线箭头，从而产生特异三角的态射。
我们提醒读者，映射 $K \to L$ 不唯一。稍后会看到，若 $a$
是 pro-同构，则 $K \to L$ 是同构，见 More on Algebra，
引理 \ref{more-algebra-lemma-pro-isomorphism-Rlim}。
\end{remark}

\begin{remark}
\label{derived-remark-map-into-derived-limit-truncations}
设 $\mathcal{A}$ 为阿贝尔范畴。设 $K^\bullet$ 为
$\mathcal{A}$ 中的复形。则 $\tau_{\geq -n}K^\bullet$ 构成一个复形逆系统，
% P02-E0082: frozen source says “is an inverse system of complexes and which in particular determines”.
因而特别给出 $D(\mathcal{A})$ 中的逆系统。假设
$R\lim \tau_{\geq -n}K^\bullet$ 存在。则典范映射
$c_n : K^\bullet \to \tau_{\geq -n}K^\bullet$ 与逆系统的转移映射相容。
由定义 \ref{derived-definition-derived-limit} 中的定义特异三角以及引理
\ref{derived-lemma-representable-homological}，可知存在态射
$$
c : K^\bullet \longrightarrow R\lim \tau_{\geq -n}K^\bullet
$$
在 $D(\mathcal{A})$ 中，使得 $c$ 与投影
$R\lim \tau_{\geq -n}K^\bullet \to \tau_{\geq -m}K^\bullet$ 的复合等于 $c_m$。
态射 $c$ 可能不唯一，但我们断言，$c$ 是否为同构与 $c$ 的选择
（以及同伦极限的选择）无关。事实上，对 $i \in \mathbf{Z}$ 和
$m > -i$，复合
$$
H^i(K^\bullet) \xrightarrow{H^i(c)} H^i(R\lim \tau_{\geq -n}K^\bullet)
\to H^i(\tau_{\geq -m}K^\bullet) = H^i(K^\bullet)
$$
是恒等映射。因此，$H^i(c)$ 为同构当且仅当第二个映射为同构。
这与 $c$ 无关，也与同伦极限的选择无关（因为任意两个选择都同构）。
\end{remark}

\begin{lemma}
\label{derived-lemma-difficulty-K-injectives}
设 $\mathcal{A}$ 为具有可数直积和足够多内射对象的阿贝尔范畴。
设 $K^\bullet$ 为复形。设 $I_n^\bullet$ 为引理
\ref{derived-lemma-special-inverse-system} 所构造的由内射对象组成的下有界复形逆系统。
则 $I^\bullet = \lim I_n^\bullet$ 存在，是 K-内射的，并且在
$D(\mathcal{A})$ 中代表 $R\lim \tau_{\geq -n}K^\bullet$；此外，下列条件等价：
\begin{enumerate}
\item 映射 $K^\bullet \to I^\bullet$（见证明）是拟同构；
\item 评注 \ref{derived-remark-map-into-derived-limit-truncations} 的映射
$K^\bullet \to R\lim \tau_{\geq -n}K^\bullet$ 是 $D(\mathcal{A})$ 中的同构。
\end{enumerate}
\end{lemma}

\begin{proof}
由于 $R\lim \tau_{\geq -n}K^\bullet$ 由引理
\ref{derived-lemma-inverse-limit-bounded-below} 存在，本引理的陈述是有意义的。
由引理 \ref{derived-lemma-bounded-below-injectives-K-injective}，每个复形
$I_n^\bullet$ 都是 K-内射的。对所有 $n \geq 1$，选择直和分解
$I_{n + 1}^p = C_{n + 1}^p \oplus I_n^p$。令 $C_1^p = I_1^p$。
复形 $I^\bullet = \lim I_n^\bullet$ 存在，因为可以取
$I^p = \prod_{n \geq 1} C_n^p$。固定 $p \in \mathbf{Z}$。
我们断言存在 $\mathcal{A}$ 中对象的分裂短正合列
$$
0 \to I^p \to \prod I_n^p \to \prod I_n^p \to 0
$$
这里第一个映射由投影映射 $I^p \to I_n^p$ 给出，第二个映射由
$(x_n) \mapsto (x_n - f^p_{n + 1}(x_{n + 1}))$ 给出，其中
$f^p_n : I_n^p \to I_{n - 1}^p$ 是转移映射。分裂来自映射
$\prod I_n^p \to \prod C_n^p = I^p$。由此得到逐项分裂的复形短正合列
$$
0 \to I^\bullet \to \prod I_n^\bullet \to \prod I_n^\bullet \to 0
$$
从而得到 $K(\mathcal{A})$ 和 $D(\mathcal{A})$ 中相应的特异三角。
由引理 \ref{derived-lemma-product-K-injective}，这些直积是 K-内射的，
并代表 $D(\mathcal{A})$ 中相应的直积。
因此 $I^\bullet$ 代表 $R\lim I_n^\bullet$
（定义 \ref{derived-definition-derived-limit}）。由于导出极限是在导出范畴层面定义的，
$R\lim I_n^\bullet \cong R\lim \tau_{\geq -n}K^\bullet$，
故 $I^\bullet$ 代表 $R\lim \tau_{\geq -n}K^\bullet$。
此外，由引理 \ref{derived-lemma-triangle-K-injective}，复形 $I^\bullet$ 是 K-内射的。
由引理 \ref{derived-lemma-special-inverse-system} 的交换图，并且由于当
$n \gg 0$ 时 $K^i = (\tau_{\geq -n}K^\bullet)^i$，可知得到唯一映射
$\gamma : K^\bullet \to I^\bullet$，使得各图
$$
\xymatrix{
K^\bullet \ar[r] \ar[d]_\gamma & \tau_{\geq -n} K^\bullet \ar[d] \\
I^\bullet \ar[r] & I_n^\bullet
}
$$
交换。因此，$\gamma$ 是复形映射，并且在 $D(\mathcal{A})$ 中代表评注
\ref{derived-remark-map-into-derived-limit-truncations} 的映射
$c : K^\bullet \to R\lim \tau_{\geq -n}K^\bullet$。
换言之，图
$$
\xymatrix{
K^\bullet \ar[r]_-c \ar[d]_\gamma & R\lim \tau_{\geq -n} K^\bullet
\ar[d]^{\cong} \\
I^\bullet \ar[r]^-{\cong} & R\lim I_n^\bullet
}
$$
在 $D(\mathcal{A})$ 中交换。引理得证。
\end{proof}

\begin{lemma}
\label{derived-lemma-enough-K-injectives-Ab4-star}
设 $\mathcal{A}$ 为具有足够多内射对象和正合可数直积的阿贝尔范畴。
则每个复形都有一个到 K-内射复形的拟同构。
\end{lemma}

\begin{proof}
由引理 \ref{derived-lemma-difficulty-K-injectives}，只需证明对
$D(\mathcal{A})$ 中的每个 $K$，$K \to R\lim\tau_{\geq -n}K$ 都是同构。
考虑定义特异三角
$$
R\lim\tau_{\geq -n}K \to
\prod \tau_{\geq -n}K \to
\prod \tau_{\geq -n}K \to
(R\lim\tau_{\geq -n}K)[1]
$$
由引理 \ref{derived-lemma-products}，有
$$
H^p(\prod \tau_{\geq -n}K) = \prod\nolimits_{p \geq -n} H^p(K)
$$
由所示特异三角的上同调长正合列，直接可得
$H^p(R\lim \tau_{\geq -n}K) = H^p(K)$。
\end{proof}




\section{全子范畴上的运算}
\label{derived-section-operate-on-full}

\noindent
设 $\mathcal{T}$ 为三角范畴。我们把 $\mathcal{T}$ 的全子范畴
与 $\Ob(\mathcal{T})$ 的子集视同。给定全子范畴
$\mathcal{A}, \mathcal{B}, \ldots$，作如下约定：
\begin{enumerate}
\item 对 $-\infty \leq a \leq b \leq \infty$，令
$\mathcal{A}[a, b]$ 为 $\mathcal{T}$ 的全子范畴，其对象是所有形如
$A[-i]$ 的对象，其中 $i \in [a, b] \cap \mathbf{Z}$ 且
$A \in \Ob(\mathcal{A})$（注意负号！）；
\item 令 $smd(\mathcal{A})$ 为 $\mathcal{T}$ 的全子范畴，其对象是
与 $\mathcal{A}$ 中对象的直和因子同构的所有对象；
\item 令 $add(\mathcal{A})$ 为 $\mathcal{T}$ 的全子范畴，其对象是
与 $\mathcal{A}$ 中对象的有限直和同构的所有对象；
\item 令 $\mathcal{A} \star \mathcal{B}$ 为 $\mathcal{T}$ 的全子范畴，
其对象是 $\mathcal{T}$ 中所有能嵌入特异三角 $A \to X \to B$ 的对象
$X$，其中 $A \in \Ob(\mathcal{A})$ 且 $B \in \Ob(\mathcal{B})$；
\item 令 $\mathcal{A}^{\star n} = \mathcal{A} \star \ldots \star \mathcal{A}$，
其中有 $n \geq 1$ 个因子（下文将看到 $\star$ 满足结合律）；
\item 令 $smd(add(\mathcal{A})^{\star n}) =
smd(add(\mathcal{A}) \star \ldots \star add(\mathcal{A}))$，
其中有 $n \geq 1$ 个因子。
\end{enumerate}
若 $E$ 是 $\mathcal{T}$ 的对象，有时也把 $E$ 看作
$\mathcal{T}$ 的仅含对象 $E$ 的全子范畴。于是可以考虑
$add(E[-1, 2])$ 等对象。我们提醒读者，这套记号并未得到普遍采用。

\begin{lemma}
\label{derived-lemma-associativity-star}
设 $\mathcal{T}$ 为三角范畴。给定全子范畴
$\mathcal{A}$、$\mathcal{B}$、$\mathcal{C}$，有
$(\mathcal{A} \star \mathcal{B}) \star \mathcal{C} =
\mathcal{A} \star (\mathcal{B} \star \mathcal{C})$。
\end{lemma}

\begin{proof}
若有特异三角 $A \to X \to B$ 和 $X \to Y \to C$，
则由公理 TR4，有特异三角 $A \to Y \to Z$ 和 $B \to Z \to C$。
\end{proof}

\begin{lemma}
\label{derived-lemma-smd-star}
设 $\mathcal{T}$ 为三角范畴。给定全子范畴
$\mathcal{A}$、$\mathcal{B}$，有
$smd(\mathcal{A}) \star smd(\mathcal{B}) \subset
smd(\mathcal{A} \star \mathcal{B})$，以及
$smd(smd(\mathcal{A}) \star smd(\mathcal{B})) =
smd(\mathcal{A} \star \mathcal{B})$。
\end{lemma}

\begin{proof}
假设有特异三角 $A_1 \to X \to B_1$，其中
$A_1 \oplus A_2 \in \Ob(\mathcal{A})$ 且
$B_1 \oplus B_2 \in \Ob(\mathcal{B})$。于是得到特异三角
$A_1 \oplus A_2 \to A_2 \oplus X \oplus B_2 \to B_1 \oplus B_2$，
这证明 $X$ 属于 $smd(\mathcal{A} \star \mathcal{B})$。
由此得到包含关系，而等式显然由该包含关系推出。
\end{proof}

\begin{lemma}
\label{derived-lemma-add-star}
设 $\mathcal{T}$ 为三角范畴。给定全子范畴
$\mathcal{A}$、$\mathcal{B}$，全子范畴
$add(\mathcal{A}) \star add(\mathcal{B})$ 和
$smd(add(\mathcal{A}))$ 对直和封闭。
\end{lemma}

\begin{proof}
事实上，若 $A \to X \to B$ 和 $A' \to X' \to B'$ 是特异三角，
$A, A' \in add(\mathcal{A})$ 且 $B, B' \in add(\mathcal{B})$，则
$A \oplus A' \to X \oplus X' \to B \oplus B'$ 是特异三角，且
$A \oplus A' \in add(\mathcal{A})$、$B \oplus B' \in add(\mathcal{B})$。
对 $smd(add(\mathcal{A}))$ 的结论是显然的。
\end{proof}

\begin{lemma}
\label{derived-lemma-cone-n}
设 $\mathcal{T}$ 为三角范畴。给定全子范畴 $\mathcal{A}$，对
$n \geq 1$，上面定义的子范畴
$$
\mathcal{C}_n = smd(add(\mathcal{A})^{\star n}) =
smd(add(\mathcal{A}) \star \ldots \star add(\mathcal{A}))
$$
是 $\mathcal{T}$ 的严格全子范畴，并且对直和及直和因子封闭；而且对所有
$n, m \geq 1$，都有
$\mathcal{C}_{n + m} = smd(\mathcal{C}_n \star \mathcal{C}_m)$。
\end{lemma}

\begin{proof}
立即由引理 \ref{derived-lemma-associativity-star}、\ref{derived-lemma-smd-star} 和
\ref{derived-lemma-add-star} 得出。
\end{proof}

\begin{remark}
\label{derived-remark-operations-functor}
设 $F : \mathcal{T} \to \mathcal{T}'$ 为三角范畴的正合函子。
给定 $\mathcal{T}$ 的全子范畴 $\mathcal{A}$，记
$F(\mathcal{A})$ 为 $\mathcal{T}'$ 的全子范畴，其对象是所有满足
$A \in \Ob(\mathcal{A})$ 的对象 $F(A)$。
% P02-E0083: frozen source has subject--verb disagreement in “whose objects consists”.
我们有
$$
F(\mathcal{A}[a, b]) = F(\mathcal{A})[a, b]
$$
$$
F(smd(\mathcal{A})) \subset smd(F(\mathcal{A})),
$$
$$
F(add(\mathcal{A})) \subset add(F(\mathcal{A})),
$$
$$
F(\mathcal{A} \star \mathcal{B}) \subset F(\mathcal{A}) \star F(\mathcal{B}),
$$
$$
F(\mathcal{A}^{\star n}) \subset F(\mathcal{A})^{\star n}.
$$
我们略去这些简单验证。
\end{remark}

\begin{remark}
\label{derived-remark-operations-unions}
设 $\mathcal{T}$ 为三角范畴。给定 $\mathcal{T}$ 的全子范畴
$\mathcal{A}_1 \subset \mathcal{A}_2 \subset \mathcal{A}_3 \subset \ldots$
和 $\mathcal{B}$，有
$$
\left(\bigcup \mathcal{A}_i\right)[a, b] = \bigcup \mathcal{A}_i[a, b]
$$
$$
smd\left(\bigcup \mathcal{A}_i\right) = \bigcup smd(\mathcal{A}_i),
$$
$$
add\left(\bigcup \mathcal{A}_i\right) = \bigcup add(\mathcal{A}_i),
$$
$$
\left(\bigcup \mathcal{A}_i\right) \star \mathcal{B} =
\bigcup \mathcal{A}_i \star \mathcal{B},
$$
$$
\mathcal{B} \star \left(\bigcup \mathcal{A}_i\right) =
\bigcup \mathcal{B} \star \mathcal{A}_i,
$$
$$
\left(\bigcup \mathcal{A}_i\right)^{\star n} =
\bigcup \mathcal{A}_i^{\star n}.
$$
我们略去这些简单验证。
\end{remark}

\begin{lemma}
\label{derived-lemma-in-cone-n}
设 $\mathcal{A}$ 为阿贝尔范畴，令 $\mathcal{D} = D(\mathcal{A})$。
设 $\mathcal{E} \subset \Ob(\mathcal{A})$ 为一个子集，并也把它视作
$\Ob(\mathcal{D})$ 的子集。设 $K$ 为 $\mathcal{D}$ 的对象。
\begin{enumerate}
\item 设 $b \geq a$，并假设当 $i \not \in [a, b]$ 时 $H^i(K)$ 为零，
而当 $i \in [a, b]$ 时 $H^i(K) \in \mathcal{E}$。则 $K$ 属于
$smd(add(\mathcal{E}[a, b])^{\star (b - a + 1)})$。
\item 设 $b \geq a$，并假设当 $i \not \in [a, b]$ 时 $H^i(K)$ 为零，
而当 $i \in [a, b]$ 时 $H^i(K) \in smd(add(\mathcal{E}))$。则 $K$ 属于
$smd(add(\mathcal{E}[a, b])^{\star (b - a + 1)})$。
\item 设 $b \geq a$，并假设 $K$ 可由复形 $K^\bullet$ 代表，且当
$i \not \in [a, b]$ 时 $K^i = 0$，当 $i \in [a, b]$ 时
$K^i \in \mathcal{E}$。则 $K$ 属于
$smd(add(\mathcal{E}[a, b])^{\star (b - a + 1)})$。
\item 设 $b \geq a$，并假设 $K$ 可由复形 $K^\bullet$ 代表，且当
$i \not \in [a, b]$ 时 $K^i = 0$，当 $i \in [a, b]$ 时
$K^i \in smd(add(\mathcal{E}))$。则 $K$ 属于
$smd(add(\mathcal{E}[a, b])^{\star (b - a + 1)})$。
\end{enumerate}
\end{lemma}

\begin{proof}
以下不再特别说明地使用引理 \ref{derived-lemma-cone-n}。
我们证明 (2)，它显然推出 (1)。对 $b - a$ 作归纳。
若 $b - a = 0$，则 $K$ 在 $\mathcal{D}$ 中同构于
$H^i(K)[-a]$，结论立即成立。
% P02-E0086: frozen source uses H^i(K) although the one surviving degree is a.
若 $b - a > 0$，则考虑特异三角
$$
\tau_{\leq b - 1}K^\bullet \to K^\bullet \to K^b[-b]
$$
并对 $b - a$ 使用归纳即可。我们略去 (3) 和 (4) 的证明。
\end{proof}

\begin{lemma}
\label{derived-lemma-forward-cone-n}
设 $\mathcal{T}$ 为三角范畴。设
$H : \mathcal{T} \to \mathcal{A}$ 为到阿贝尔范畴 $\mathcal{A}$ 的同调函子。
设 $a \leq b$，且 $\mathcal{E} \subset \Ob(\mathcal{T})$ 为如下子集：
当 $E \in \mathcal{E}$ 且 $i \not \in [a, b]$ 时，$H^i(E) = 0$。
则对 $X \in smd(add(\mathcal{E}[-m, m])^{\star n})$，当
$i \not \in [-m + na, m + nb]$ 时有 $H^i(X) = 0$。
\end{lemma}

\begin{proof}
略。这是关于这些定义的一道有趣练习。
\end{proof}





\section{三角范畴的生成元}
\label{derived-section-generators}

\noindent
本节简要介绍三角范畴生成元的几种不同概念。我们的术语取自
\cite{BvdB}（但他们所称的 ``\'epaisse''，我们称为“饱和”，
见定义 \ref{derived-definition-saturated}；此外，我们对 $add(\mathcal{A})$ 的定义不同）。

\medskip\noindent
设 $\mathcal{D}$ 为三角范畴，$E$ 为 $\mathcal{D}$ 的对象。
记 $\langle E \rangle_1$ 为 $\mathcal{D}$ 的严格全子范畴，
其对象是 $\mathcal{D}$ 中与 $E$ 的平移的有限直和
$$
\bigoplus\nolimits_{i = 1, \ldots, r} E[n_i]
$$
的直和因子同构的对象。显然，采用节 \ref{derived-section-operate-on-full} 的记号，有
$$
\langle E \rangle_1 = smd(add(E[-\infty, \infty]))
$$
对 $n > 1$，记 $\langle E \rangle_n$ 为 $\mathcal{D}$ 的全子范畴，
其对象是与能嵌入特异三角
$$
A \to X \to B \to A[1]
$$
的对象 $X$ 的直和因子同构的 $\mathcal{D}$ 中对象，其中 $A$ 是
$\langle E \rangle_1$ 的对象，$B$ 是 $\langle E \rangle_{n - 1}$ 的对象。
采用节 \ref{derived-section-operate-on-full} 的记号，有
$$
\langle E \rangle_n = smd(\langle E \rangle_1 \star \langle E \rangle_{n - 1})
$$
每个范畴 $\langle E \rangle_n$ 都是 $\mathcal{D}$ 的严格全加性子范畴
（由引理 \ref{derived-lemma-add-star}），并且在平移和取直和因子下保持不变。
然而，$\langle E \rangle_n$ 不一定对“取锥”或“扩张”封闭，
因而不一定是三角子范畴。如下子范畴则具有这一性质：
$$
\langle E \rangle = \bigcup\nolimits_n \langle E \rangle_n
$$
下面的引理将证明这一点。

\begin{lemma}
\label{derived-lemma-generated-by-E-explicit}
设 $\mathcal{T}$ 为三角范畴，$E$ 为 $\mathcal{T}$ 的对象。
对 $n \geq 1$，有
$$
\langle E \rangle_n =
smd(\langle E \rangle_1 \star \ldots \star \langle E \rangle_1) =
smd({\langle E \rangle_1}^{\star n}) =
\bigcup\nolimits_{m \geq 1} smd(add(E[-m, m])^{\star n})
$$
对 $n, n' \geq 1$，有 $\langle E \rangle_{n + n'} =
smd(\langle E \rangle_n \star \langle E \rangle_{n'})$。
\end{lemma}

\begin{proof}
所示公式的左边等式由引理 \ref{derived-lemma-associativity-star}、
\ref{derived-lemma-smd-star} 和归纳法得出。中间的等式只是记号问题。
由于
$\langle E \rangle_1 = smd(add(E[-\infty, \infty])])$
% P02-E0084: frozen source has an extra closing bracket in this formula.
并且 $E[-\infty, \infty] = \bigcup_{m \geq 1} E[-m, m]$，
由评注 \ref{derived-remark-operations-unions} 和引理 \ref{derived-lemma-smd-star}，
得到右边的等式。最后一个陈述则由该评注及引理
\ref{derived-lemma-cone-n} 的相应陈述得出。
\end{proof}

\begin{lemma}
\label{derived-lemma-find-smallest-containing-E}
设 $\mathcal{D}$ 为三角范畴，$E$ 为 $\mathcal{D}$ 的对象。子范畴
$$
\langle E \rangle = \bigcup\nolimits_n \langle E \rangle_n
= \bigcup\nolimits_{n, m \geq 1} smd(add(E[-m, m])^{\star n})
$$
是 $\mathcal{D}$ 的严格全、饱和三角子范畴，并且是包含对象 $E$ 的
$\mathcal{D}$ 的最小此类子范畴。
\end{lemma}

\begin{proof}
右边的等式由引理 \ref{derived-lemma-generated-by-E-explicit} 得出。
显然，$\langle E \rangle = \bigcup \langle E \rangle_n$ 包含 $E$，
并且在平移、直和及直和因子下保持不变。若
$A \in \langle E \rangle_a$、$B \in \langle E \rangle_b$，且
$A \to X \to B \to A[1]$ 是特异三角，则由引理
\ref{derived-lemma-generated-by-E-explicit}，有
$X \in \langle E \rangle_{a + b}$。因此
$\bigcup \langle E \rangle_n$ 也对扩张封闭，从而是三角子范畴。

\medskip\noindent
最后，设 $\mathcal{D}' \subset \mathcal{D}$ 为包含 $E$ 的
$\mathcal{D}$ 的严格全、饱和三角子范畴。则
$\mathcal{D}'[-\infty, \infty] \subset \mathcal{D}'$、
$add(\mathcal{D}) \subset \mathcal{D}'$、
% P02-E0085: frozen source uses add(D), although preservation here requires add(D').
$smd(\mathcal{D}') \subset \mathcal{D}'$，且
$\mathcal{D}' \star \mathcal{D}' \subset \mathcal{D}'$。
换言之，从 $E$ 构造 $\langle E \rangle$ 所用的全部运算都保持
$\mathcal{D}'$。因此 $\langle E \rangle \subset \mathcal{D}'$，证明完毕。
\end{proof}

\begin{definition}
\label{derived-definition-generators}
设 $\mathcal{D}$ 为三角范畴，$E$ 为 $\mathcal{D}$ 的对象。
\begin{enumerate}
\item 若 $\mathcal{D}$ 中包含 $E$ 的最小严格全、饱和三角子范畴
等于 $\mathcal{D}$，换言之，若 $\langle E \rangle = \mathcal{D}$，
则称 $E$ 为 $\mathcal{D}$ 的{\it 经典生成元}。
\item 若对某个 $n \geq 1$ 有 $\langle E \rangle_n = \mathcal{D}$，
则称 $E$ 为 $\mathcal{D}$ 的{\it 强生成元}。
\item 若对 $\mathcal{D}$ 的任意非零对象 $K$，都存在整数 $n$
和非零映射 $E \to K[n]$，则称 $E$ 为 $\mathcal{D}$ 的
{\it 弱生成元}，或简称{\it 生成元}。
\end{enumerate}
\end{definition}

\noindent
这一定义可以推广到对象族的情形。

\begin{lemma}
\label{derived-lemma-right-orthogonal}
设 $\mathcal{D}$ 为三角范畴，$E, K$ 为 $\mathcal{D}$ 的对象。
下列条件等价：
\begin{enumerate}
\item 对所有 $i \in \mathbf{Z}$，有 $\Hom(E, K[i]) = 0$；
\item 对所有 $E' \in \langle E \rangle$，有 $\Hom(E', K) = 0$。
\end{enumerate}
\end{lemma}

\begin{proof}
(2) $\Rightarrow$ (1) 立即成立。反过来，假设 (1)。
则对 $\langle E \rangle_1$ 中所有 $X$，都有 $\Hom(X, K) = 0$。
对 $n$ 归纳并使用引理 \ref{derived-lemma-representable-homological}，
可见对 $\langle E \rangle_n$ 中所有 $X$，都有 $\Hom(X, K) = 0$。
\end{proof}

\begin{lemma}
\label{derived-lemma-classical-generator-generator}
设 $\mathcal{D}$ 为三角范畴，$E$ 为 $\mathcal{D}$ 的对象。
若 $E$ 是 $\mathcal{D}$ 的经典生成元，则 $E$ 是生成元。
\end{lemma}

\begin{proof}
假设 $E$ 是经典生成元。设 $K$ 为 $\mathcal{D}$ 的对象，且对所有
$i \in \mathbf{Z}$ 都有 $\Hom(E, K[i]) = 0$。由引理
\ref{derived-lemma-right-orthogonal}，对 $\langle E \rangle$ 中所有 $E'$，都有
$\Hom(E', K) = 0$。然而，由于 $\mathcal{D} = \langle E \rangle$，
可知 $\text{id}_K = 0$，即 $K = 0$。
\end{proof}

\begin{lemma}
\label{derived-lemma-classical-generator-strong-generator}
设 $\mathcal{D}$ 为具有强生成元的三角范畴，$E$ 为 $\mathcal{D}$ 的对象。
若 $E$ 是 $\mathcal{D}$ 的经典生成元，则 $E$ 是强生成元。
\end{lemma}

\begin{proof}
设 $E'$ 为 $\mathcal{D}$ 的对象，且
$\mathcal{D} = \langle E' \rangle_n$。由于
$\mathcal{D} = \langle E \rangle$，由引理
\ref{derived-lemma-find-smallest-containing-E}，对某个 $m \geq 1$ 有
$E' \in \langle E \rangle_m$。于是
$\langle E' \rangle_1 \subset \langle E \rangle_m$，故
$$
\mathcal{D} =
\langle E' \rangle_n = smd(
\langle E' \rangle_1 \star \ldots \star \langle E' \rangle_1)
\subset
smd(
\langle E \rangle_m \star \ldots \star \langle E \rangle_m)
=
\langle E \rangle_{nm}
$$
正如所需。这里使用了引理 \ref{derived-lemma-generated-by-E-explicit}。
\end{proof}

\begin{remark}
\label{derived-remark-check-on-generator}
设 $\mathcal{D}$ 为三角范畴，$E$ 为 $\mathcal{D}$ 的对象。
设 $T$ 为 $\mathcal{D}$ 中对象的一个性质。假设：
\begin{enumerate}
\item 若 $K_i \in \mathcal{D}$，$i = 1, \ldots, r$，且对
$i = 1, \ldots, r$ 都有 $T(K_i)$，则 $T(\bigoplus K_i)$ 成立；
\item 若 $K \to L \to M \to K[1]$ 是特异三角，且 $T$ 对其中两个对象成立，
则 $T$ 对第三个对象也成立；
\item 若 $T(K \oplus L)$ 成立，则 $T(K)$ 和 $T(L)$ 成立；
\item 对所有 $n$，$T(E[n])$ 成立。
\end{enumerate}
则 $T$ 对 $\langle E \rangle$ 的所有对象成立。
\end{remark}












\section{紧对象}
\label{derived-section-compact}

\noindent
定义如下。

\begin{definition}
\label{derived-definition-compact-object}
设 $\mathcal{D}$ 为具有任意直和的加性范畴。$\mathcal{D}$ 的
{\it 紧对象}是满足如下条件的对象 $K$：对任意集合 $I$ 以及由
$i \in I$ 参数化的对象 $E_i \in \Ob(\mathcal{D})$，映射
$$
\bigoplus\nolimits_{i \in I} \Hom_{\mathcal{D}}(K, E_i)
\longrightarrow
\Hom_{\mathcal{D}}(K, \bigoplus\nolimits_{i \in I} E_i)
$$
是双射。
\end{definition}

\noindent
这一概念在代数几何中非常有用。它是对象所满足的一个内在条件，
迫使这些对象确实具有某种“紧性”。

\begin{lemma}
\label{derived-lemma-compact-objects-subcategory}
设 $\mathcal{D}$ 为具有直和的（预）三角范畴。则 $\mathcal{D}$ 的紧对象
构成 $\mathcal{D}$ 的一个 Karoubi、饱和、严格全（预）三角子范畴
$\mathcal{D}_c$ 的全部对象。
\end{lemma}

\begin{proof}
设 $(X, Y, Z, f, g, h)$ 为 $\mathcal{D}$ 的特异三角，且
$X$ 和 $Y$ 都是紧的。则由引理 \ref{derived-lemma-representable-homological}
和五引理（Homology，引理 \ref{homology-lemma-five-lemma}），
$Z$ 也是紧对象。显然，若 $X \oplus Y$ 是紧的，则 $X$、$Y$
也都是紧对象。因此 $\mathcal{D}_c$ 是饱和三角子范畴。
由引理 \ref{derived-lemma-projectors-have-images-triangulated}，$\mathcal{D}$
是 Karoubi 范畴，故 $\mathcal{D}_c$ 亦然。
\end{proof}

\begin{lemma}
\label{derived-lemma-write-as-colimit}
设 $\mathcal{D}$ 为具有直和的三角范畴。设 $E_i$，$i \in I$，
为 $\mathcal{D}$ 的一族紧对象，且 $\bigoplus E_i$ 生成 $\mathcal{D}$。
则 $\mathcal{D}$ 的每个对象 $X$ 都可写成
$$
X = \text{hocolim} X_n
$$
其中 $X_1$ 是 $E_i$ 的平移的直和，并且每个转移态射都嵌入特异三角
$Y_n \to X_n \to X_{n + 1} \to Y_n[1]$，其中 $Y_n$ 是
$E_i$ 的平移的直和。
\end{lemma}

\begin{proof}
令 $X_1 = \bigoplus_{(i, m, \varphi)} E_i[m]$，其中直和遍历所有满足
$i \in I$、$m \in \mathbf{Z}$ 及 $\varphi : E_i[m] \to X$ 的三元组
$(i, m, \varphi)$。于是 $X_1$ 带有典范态射 $X_1 \to X$。
给定 $X_n \to X$，令 $Y_n = \bigoplus_{(i, m, \varphi)} E_i[m]$，
其中直和遍历所有满足 $i \in I$、$m \in \mathbf{Z}$，且
$\varphi : E_i[m] \to X_n$ 是使得 $E_i[m] \to X_n \to X$ 为零的态射的
三元组 $(i, m, \varphi)$。选择特异三角
$Y_n \to X_n \to X_{n + 1} \to Y_n[1]$，并任取态射
$X_{n + 1} \to X$，使 $X_n \to X_{n + 1} \to X$ 等于给定态射；
由 $Y_n$ 的选择，这样的态射存在。由各映射 $X_n \to X$ 的构造，
得到态射 $\text{hocolim} X_n \to X$。选择特异三角
$$
C \to \text{hocolim} X_n \to X \to C[1]
$$
设 $E_i[m] \to C$ 为态射。由于 $E_i$ 是紧的，引理
\ref{derived-lemma-commutes-with-countable-sums} 蕴含复合
$E_i[m] \to C \to \text{hocolim} X_n$ 对某个 $n$ 经 $X_n$ 分解，
设相应态射为 $E_i[m] \to X_n$。于是由 $Y_n$ 的构造，复合
$E_i[m] \to X_n \to X_{n + 1}$ 为零。换言之，复合
$E_i[m] \to C \to \text{hocolim} X_n$ 为零。这意味着态射
$E_i[m] \to C$ 来自某个态射 $E_i[m] \to X[-1]$。
$X_1$ 的构造又表明，该态射提升到 $\text{hocolim} X_n$，因而态射
$E_i[m] \to C$ 为零。假设 $\bigoplus E_i$ 生成 $\mathcal{D}$，
遂蕴含 $C$ 为零，证明完毕。
\end{proof}

\begin{lemma}
\label{derived-lemma-factor-through}
沿用引理 \ref{derived-lemma-write-as-colimit} 的假设和记号。
% P02-E0087: frozen source sentence “With assumptions and notation ...” lacks a finite verb.
若 $C$ 是紧对象，且 $C \to X_n$ 是态射，则存在分解
$C \to E \to X_n$，其中对某些 $i_1, \ldots, i_t \in I$，
$E$ 是 $\langle E_{i_1} \oplus \ldots \oplus E_{i_t} \rangle$ 的对象。
\end{lemma}

\begin{proof}
对 $n$ 归纳。基础情形 $n = 1$ 是显然的。若 $n > 1$，
考虑复合 $C \to X_n \to Y_{n - 1}[1]$。它可经某个
$E'[1] \to Y_{n - 1}[1]$ 分解，其中 $E'$ 是 $E_i$ 的平移的有限直和。
令 $I' \subset I$ 为该直和中出现的指标所成的有限集。于是得到
$$
\xymatrix{
E' \ar[r] \ar[d] &
C' \ar[r] \ar[d] &
C \ar[r] \ar[d] &
E'[1] \ar[d] \\
Y_{n - 1} \ar[r] &
X_{n - 1} \ar[r] &
X_n \ar[r] &
Y_{n - 1}[1]
}
$$
由归纳假设，态射 $C' \to X_{n - 1}$ 经
$E'' \to X_{n - 1}$ 分解，其中对某个有限子集 $I'' \subset I$，
$E''$ 是 $\langle \bigoplus_{i \in I''} E_i \rangle$ 的对象。
选择特异三角
$$
E' \to E'' \to E \to E'[1]
$$
则 $E$ 是 $\langle \bigoplus_{i \in I' \cup I''} E_i \rangle$ 的对象。
由构造及三角范畴公理，可选择态射 $C \to E$ 和 $E \to X_n$，
使之嵌入三角的态射 $(E', C', C) \to (E', E'', E)$ 及
$(E', E'', E) \to (Y_{n - 1}, X_{n - 1}, X_n)$。
复合 $C \to E \to X_n$ 可能不等于给定态射 $C \to X_n$，
但它们到 $Y_{n - 1}$ 的复合相等。设 $C \to X_{n - 1}$ 为提升其差的态射。
由归纳假设，它可经态射 $E''' \to X_{n - 1}$ 分解，其中
$E''$ 是 $\langle \bigoplus_{i \in I'''} E_i \rangle$ 的对象，
% P02-E0088: frozen source says E'' here, although the newly introduced object is E'''.
而 $I' \subset I$ 是某个有限子集。
% P02-E0089: frozen source repeats I' here, although the new index set is I'''.
因此，考虑 $E \oplus E''' \to X_n$ 即得到所求解，因为
$E \oplus E'''$ 是
$\langle \bigoplus_{i \in I' \cup I'' \cup I'''} E_i \rangle$ 的对象。
\end{proof}

\begin{definition}
\label{derived-definition-compactly-generated}
设 $\mathcal{D}$ 为具有任意直和的三角范畴。若存在以 $i \in I$
为指标的一族紧对象 $E_i$，使得 $\bigoplus E_i$ 生成 $\mathcal{D}$，
则称 $\mathcal{D}$ 是{\it 紧生成的}。
\end{definition}

\noindent
下面的命题阐明经典生成元与弱生成元之间的关系。

\begin{proposition}
\label{derived-proposition-generator-versus-classical-generator}
设 $\mathcal{D}$ 为具有直和的三角范畴，$E$ 为 $\mathcal{D}$ 的紧对象。
下列条件等价：
\begin{enumerate}
\item $E$ 是 $\mathcal{D}_c$ 的经典生成元，且 $\mathcal{D}$ 是紧生成的；
\item $E$ 是 $\mathcal{D}$ 的生成元。
\end{enumerate}
\end{proposition}

\begin{proof}
若 $E$ 是 $\mathcal{D}_c$ 的经典生成元，则
$\mathcal{D}_c = \langle E \rangle$。由 $\mathcal{D}$ 紧生成的假设和引理
\ref{derived-lemma-right-orthogonal}，形式地推出 $E$ 是 $\mathcal{D}$ 的生成元。

\medskip\noindent
反向更有意思。假设 $E$ 是 $\mathcal{D}$ 的生成元。
设 $X$ 为 $\mathcal{D}$ 的紧对象。应用引理
\ref{derived-lemma-write-as-colimit}，取 $I = \{1\}$ 和 $E_1 = E$，按该引理写成
$$
X = \text{hocolim} X_n
$$
由于 $X$ 是紧的，对某个 $n$，映射
$X \to \text{hocolim} X_n$ 经 $X_n$ 分解（引理
\ref{derived-lemma-commutes-with-countable-sums}）。因此 $X$ 是 $X_n$ 的直和因子。
由引理 \ref{derived-lemma-factor-through}，可见 $X$ 是
$\langle E \rangle$ 的对象，命题得证。
% P02-E0091: frozen source closes this proposition proof with “the lemma is proven”.
\end{proof}











\section{Brown 可表性}
\label{derived-section-brown}

\noindent
本节内容的参考文献是 \cite{Neeman-Grothendieck}。

\begin{lemma}
\label{derived-lemma-brown}
\begin{reference}
\cite[Theorem 3.1]{Neeman-Grothendieck}.
\end{reference}
设 $\mathcal{D}$ 为具有直和的紧生成三角范畴。设
$H : \mathcal{D} \to \textit{Ab}$ 为把直和变为直积的反变上同调函子。
则 $H$ 可表。
\end{lemma}

\begin{proof}
设 $E_i$，$i \in I$，为一族紧对象，且
$\bigoplus_{i \in I} E_i$ 生成 $\mathcal{D}$。不妨并且确实假设
对象集 $\{E_i\}$ 在平移下保持不变。考虑序对 $(i, a)$，其中
$i \in I$ 且 $a \in H(E_i)$，并令
$$
X_1 = \bigoplus\nolimits_{(i, a)} E_i
$$
由于 $H(X_1) = \prod_{(i, a)} H(E_i)$，可见 $(a)_{(i, a)}$
定义元素 $a_1 \in H(X_1)$。令 $H_1 = \Hom_\mathcal{D}(- , X_1)$。
由 Yoneda 引理（Categories，引理 \ref{categories-lemma-yoneda}），
元素 $a_1$ 定义自然变换 $H_1 \to H$。

\medskip\noindent
我们将归纳构造 $X_n$ 和变换 $a_n : H_n \to H$，其中
$H_n = \Hom_\mathcal{D}(-, X_n)$。具体地，对函子
$\Ker(H_n \to H)$ 应用上面的步骤，得到对象
$$
K_{n + 1} = \bigoplus\nolimits_{(i, k),\ k \in \Ker(H_n(E_i) \to H(E_i))} E_i
$$
以及变换 $\Hom_\mathcal{D}(-, K_{n + 1}) \to \Ker(H_n \to H)$。
由 Yoneda 引理，复合 $\Hom_\mathcal{D}(-, K_{n + 1}) \to H_n$
给出态射 $K_{n + 1} \to X_n$。在 $\mathcal{D}$ 中选择特异三角
$$
K_{n + 1} \to X_n \to X_{n + 1} \to K_{n + 1}[1]
$$
由构造，元素 $a_n \in H(X_n)$ 在 $H(K_{n + 1})$ 中的像为零。
由于 $H$ 是上同调函子，可将它提升为元素
$a_{n + 1} \in H(X_{n + 1})$。

\medskip\noindent
我们断言 $X = \text{hocolim} X_n$ 表示 $H$。对定义特异三角
$$
\bigoplus X_n \to
\bigoplus X_n \to X \to
\bigoplus X_n[1]
$$
应用 $H$，得到正合列
$$
\prod H(X_n) \leftarrow 
\prod H(X_n) \leftarrow 
H(X)
$$
因此存在元素 $a \in H(X)$，它在 $\prod H(X_n)$ 中的像为 $(a_n)$。
从而有自然变换 $\Hom_\mathcal{D}(- , X) \to H$，使得
% P02-E0092: frozen source has the sentence fragment “Hence a natural transformation ...”.
$$
\Hom_\mathcal{D}(-, X_1) \to
\Hom_\mathcal{D}(-, X_2) \to
\Hom_\mathcal{D}(-, X_3) \to \ldots \to
\Hom_\mathcal{D}(-, X) \to H
$$
交换。对每个 $i$，映射 $\Hom_\mathcal{D}(E_i, X) \to H(E_i)$
由 $X_1$ 的构造是满射。另一方面，由 $X_n \to X_{n + 1}$ 的构造，
$\Hom_\mathcal{D}(E_i, X_n) \to H(E_i)$ 的核被映射
$\Hom_\mathcal{D}(E_i, X_n) \to \Hom_\mathcal{D}(E_i, X_{n + 1})$ 消去。
由于
$$
\Hom_\mathcal{D}(E_i, X) = \colim \Hom_\mathcal{D}(E_i, X_n)
$$
由引理 \ref{derived-lemma-commutes-with-countable-sums}，可见
$\Hom_\mathcal{D}(E_i, X) \to H(E_i)$ 是单射。

\medskip\noindent
为完成证明，考虑子范畴
$$
\mathcal{D}' =
\{Y \in \Ob(\mathcal{D}) \mid \Hom_\mathcal{D}(Y[n], X) \to H(Y[n])
\text{ 对所有 }n\text{ 均为同构}\}
$$
由于 $\Hom_\mathcal{D}(-, X) \to H$ 是上同调函子之间的变换，
子范畴 $\mathcal{D}'$ 是 $\mathcal{D}$ 的严格全、饱和三角子范畴
（略去细节；见引理 \ref{derived-lemma-homological-functor-kernel} 的证明）。
此外，由于 $H$ 和 $\Hom_\mathcal{D}(-, X)$ 都把直和变为直积，
$\mathcal{D}'$ 中对象的直和仍在 $\mathcal{D}'$ 中。因此
$\mathcal{D}'$ 中对象的导出余极限仍在 $\mathcal{D}'$ 中。
由于 $\{E_i\}$ 在平移下保持不变，可知对所有 $i$，$E_i$
都是 $\mathcal{D}'$ 的对象。由引理 \ref{derived-lemma-write-as-colimit}，
$\mathcal{D}' = \mathcal{D}$，证明完毕。
\end{proof}

\begin{proposition}
\label{derived-proposition-brown}
\begin{reference}
\cite[Theorem 4.1]{Neeman-Grothendieck}.
\end{reference}
设 $\mathcal{D}$ 为具有直和的紧生成三角范畴。设
$F : \mathcal{D} \to \mathcal{D}'$ 为把直和变为直和的三角范畴正合函子。
则 $F$ 有正合右伴随。
\end{proposition}

\begin{proof}
对 $\mathcal{D}'$ 的对象 $Y$，考虑反变函子
$$
\mathcal{D} \to \textit{Ab},\quad W \mapsto \Hom_{\mathcal{D}'}(F(W), Y)
$$
由于 $F$ 正合，该函子是上同调函子；又因为 $F$ 把直和变为直和，
该函子把直和变为直积。因此，由引理 \ref{derived-lemma-brown}，存在
$\mathcal{D}$ 的对象 $X$，使得
$\Hom_\mathcal{D}(W, X) = \Hom_{\mathcal{D}'}(F(W), Y)$。
伴随的存在性由 Categories，引理 \ref{categories-lemma-adjoint-exists} 得出，
正合性由引理 \ref{derived-lemma-adjoint-is-exact} 得出。
\end{proof}









\section{Brown 可表性之二}
\label{derived-section-brown-bis}

\noindent
本节说明一类具有适当生成元集的三角范畴的 Brown 可表性版本；
其他版本见 \cite{Franke}、\cite{Neeman} 和 \cite{Krause}。

\begin{lemma}
\label{derived-lemma-brown-bis}
\begin{reference}
\cite[Theorem A]{Krause} {\stackszhgreek 的弱化版本。}
\end{reference}
设 $\mathcal{D}$ 为具有直和的三角范畴。假设给定 $\mathcal{D}$
的对象集 $\mathcal{E}$，满足：
\begin{enumerate}
\item 若 $X$ 是 $\mathcal{D}$ 的非零对象，则存在
$E \in \mathcal{E}$ 和非零映射 $E \to X$；
\item 给定 $\mathcal{D}$ 的对象 $X_n$，$n \in \mathbf{N}$，
$E \in \mathcal{E}$，以及 $\alpha : E \to \bigoplus X_n$，
存在 $E_n \in \mathcal{E}$、$\beta_n : E_n \to X_n$ 和态射
$\gamma : E \to \bigoplus E_n$，使得
$\alpha = (\bigoplus \beta_n) \circ \gamma$。
\end{enumerate}
设 $H : \mathcal{D} \to \textit{Ab}$ 为把直和变为直积的反变上同调函子。
则 $H$ 可表。
\end{lemma}

\begin{proof}
本证明与引理 \ref{derived-lemma-brown} 的证明非常相似。可将 $\mathcal{E}$
替换为 $\bigcup_{i \in \mathbf{Z}} \mathcal{E}[i]$，并假设
$\mathcal{E}$ 在平移下保持不变。考虑序对 $(E, a)$，其中
$E \in \mathcal{E}$ 且 $a \in H(E)$，并令
$$
X_1 = \bigoplus\nolimits_{(E, a)} E
$$
由于 $H(X_1) = \prod_{(E, a)} H(E)$，可见 $(a)_{(E, a)}$
定义元素 $a_1 \in H(X_1)$。令 $H_1 = \Hom_\mathcal{D}(- , X_1)$。
由 Yoneda 引理（Categories，引理 \ref{categories-lemma-yoneda}），
元素 $a_1$ 定义自然变换 $H_1 \to H$。

\medskip\noindent
我们将归纳构造 $X_n$ 和变换 $a_n : H_n \to H$，其中
$H_n = \Hom_\mathcal{D}(-, X_n)$。具体地，对函子
$\Ker(H_n \to H)$ 应用上面的步骤，得到对象
$$
K_{n + 1} = \bigoplus\nolimits_{(E, k),\ k \in \Ker(H_n(E) \to H(E))} E
$$
以及变换 $\Hom_\mathcal{D}(-, K_{n + 1}) \to \Ker(H_n \to H)$。
由 Yoneda 引理，复合 $\Hom_\mathcal{D}(-, K_{n + 1}) \to H_n$
给出态射 $K_{n + 1} \to X_n$。在 $\mathcal{D}$ 中选择特异三角
$$
K_{n + 1} \to X_n \to X_{n + 1} \to K_{n + 1}[1]
$$
由构造，元素 $a_n \in H(X_n)$ 在 $H(K_{n + 1})$ 中的像为零。
由于 $H$ 是上同调函子，可将它提升为元素
$a_{n + 1} \in H(X_{n + 1})$。

\medskip\noindent
令 $X = \text{hocolim} X_n$。对定义特异三角
$$
\bigoplus X_n \to
\bigoplus X_n \to X \to
\bigoplus X_n[1]
$$
应用 $H$，得到正合列
$$
\prod H(X_n) \leftarrow
\prod H(X_n) \leftarrow
H(X)
$$
因此存在元素 $a \in H(X)$，它在 $\prod H(X_n)$ 中的像为 $(a_n)$。
从而有自然变换 $\Hom_\mathcal{D}(- , X) \to H$，使得
% P02-E0093: frozen source repeats the sentence fragment “Hence a natural transformation ...”.
$$
\Hom_\mathcal{D}(-, X_1) \to
\Hom_\mathcal{D}(-, X_2) \to
\Hom_\mathcal{D}(-, X_3) \to \ldots \to
\Hom_\mathcal{D}(-, X) \to H
$$
交换。我们断言 $\Hom_\mathcal{D}(-, X) \to H(-)$ 是同构。

\medskip\noindent
设 $E \in \mathcal{E}$。先证明
$$
\Hom_\mathcal{D}(E, \bigoplus X_n) \to  \Hom_\mathcal{D}(E, \bigoplus X_n)
$$
是单射。事实上，设 $\alpha : E \to \bigoplus X_n$。
则由假设 (2)，得到分解 $\alpha = (\bigoplus \beta_n) \circ \gamma$。
由于由构造 $E_n \to X_n \to X_{n + 1}$ 为零，可见复合
$\bigoplus E_n \to \bigoplus X_n \to \bigoplus X_n$
等于 $\bigoplus \beta_n$。因而复合
$E \to \bigoplus X_n \to \bigoplus X_n$ 也等于 $\alpha$。
这证明所述单射性，并且也证明
$$
\Hom_\mathcal{D}(E, \bigoplus X_n[1]) \to \Hom_\mathcal{D}(E, \bigoplus X_n[1])
$$
是单射。因此，对所有 $E \in \mathcal{E}$，都有正合列
$$
\Hom_\mathcal{D}(E, \bigoplus X_n) \to
\Hom_\mathcal{D}(E, \bigoplus X_n) \to
\Hom_\mathcal{D}(E, X) \to 0
$$

\medskip\noindent
设 $E \in \mathcal{E}$，并设 $f : E \to X$ 为态射。
由上一段，可选择提升 $f$ 的 $\alpha : E \to \bigoplus X_n$。
再由假设 (2)，得到分解 $\alpha = (\bigoplus \beta_n) \circ \gamma$。
对每个 $n$，存在态射 $\delta_n : E_n \to X_1$，使得
$\delta_n$ 和 $\beta_n$ 在 $H(E_n)$ 中的像相同。于是由
$X_n \to X_{n + 1}$ 的构造，复合
$$
E_n \to X_n \to X_{n + 1}
\quad\text{且}\quad
E_n \to X_1 \to X_{n + 1}
$$
相等。因此复合
$$
\bigoplus E_n \to \bigoplus X_n \to X
\quad\text{且}\quad
\bigoplus E_n \to \bigoplus X_1 \to X
$$
也相同。注意到 $\bigoplus X_1 \to X$ 分解为
$\bigoplus X_1 \to X_1 \to X$，可知
$$
\Hom_\mathcal{D}(E, X_1) \to \Hom_\mathcal{D}(E, X)
$$
是满射。由构造，映射 $\Hom_\mathcal{D}(E, X_1) \to H(E)$ 是满射，
且该映射的核被 $\Hom_\mathcal{D}(E, X_1) \to \Hom_\mathcal{D}(E, X)$
消去，故对所有 $E \in \mathcal{E}$，
$\Hom_\mathcal{D}(E, X) \to H(E)$ 都是双射。

\medskip\noindent
为完成证明，考虑子范畴
$$
\mathcal{D}' =
\{Y \in \Ob(\mathcal{D}) \mid \Hom_\mathcal{D}(Y[n], X) \to H(Y[n])
\text{ 对所有 }n\text{ 均为同构}\}
$$
由于 $\Hom_\mathcal{D}(-, X) \to H$ 是上同调函子之间的变换，
子范畴 $\mathcal{D}'$ 是 $\mathcal{D}$ 的严格全、饱和三角子范畴
（略去细节；见引理 \ref{derived-lemma-homological-functor-kernel} 的证明）。
此外，由于 $H$ 和 $\Hom_\mathcal{D}(-, X)$ 都把直和变为直积，
$\mathcal{D}'$ 中对象的直和仍在 $\mathcal{D}'$ 中。因此
$\mathcal{D}'$ 中对象的导出余极限仍在 $\mathcal{D}'$ 中。
由于 $\mathcal{E}$ 在平移下保持不变，由上一段的结果可知
$\mathcal{E} \subset \Ob(\mathcal{D}')$。为完成证明，还需证明
$\mathcal{D}' = \mathcal{D}$。

\medskip\noindent
设 $Y$ 为 $\mathcal{D}$ 的对象，并令 $H(-) = \Hom_\mathcal{D}(-, Y)$。
则 $H$ 是把直和变为直积的上同调函子。由证明第一部分的构造，
得到态射 $\colim X_n = X \to Y$，使得对所有
% P02-E0094: frozen source writes colim X_n although X was defined as hocolim X_n.
$E \in \mathcal{E}$，$\Hom_\mathcal{D}(E, X) \to \Hom_\mathcal{D}(E, Y)$
都是双射。假设 (1) 遂表明 $X \to Y$ 是同构！另一方面，由构造，
$X_1, X_2, \ldots$ 都属于 $\mathcal{D}'$，因而 $X$ 也属于其中。
于是 $Y \in \mathcal{D}'$，证明完毕。
\end{proof}

\begin{proposition}
\label{derived-proposition-brown-bis}
设 $\mathcal{D}$ 为具有直和的三角范畴。假设存在 $\mathcal{D}$ 的对象集
$\mathcal{E}$，满足引理 \ref{derived-lemma-brown-bis} 的条件 (1) 和 (2)。
设 $F : \mathcal{D} \to \mathcal{D}'$ 为把直和变为直和的三角范畴正合函子。
则 $F$ 有正合右伴随。
\end{proposition}

\begin{proof}
对 $\mathcal{D}'$ 的对象 $Y$，考虑反变函子
$$
\mathcal{D} \to \textit{Ab},\quad W \mapsto \Hom_{\mathcal{D}'}(F(W), Y)
$$
由于 $F$ 正合，该函子是上同调函子；又因为 $F$ 把直和变为直和，
该函子把直和变为直积。因此，由引理 \ref{derived-lemma-brown-bis}，存在
$\mathcal{D}$ 的对象 $X$，使得
$\Hom_\mathcal{D}(W, X) = \Hom_{\mathcal{D}'}(F(W), Y)$。
伴随的存在性由 Categories，引理 \ref{categories-lemma-adjoint-exists} 得出，
正合性由引理 \ref{derived-lemma-adjoint-is-exact} 得出。
\end{proof}








\section{容许子范畴}
\label{derived-section-admissible}

\noindent
本节的参考文献是 \cite[Section 1]{Bondal-Kapranov}。

\begin{definition}
\label{derived-definition-orthogonal}
设 $\mathcal{D}$ 为加性范畴，$\mathcal{A} \subset \mathcal{D}$ 为全子范畴。
$\mathcal{A}$ 的{\it 右正交} $\mathcal{A}^\perp$ 是如下对象组成的全子范畴：
$\mathcal{D}$ 的对象 $X$ 满足，对所有 $A \in \Ob(\mathcal{A})$，
$\Hom(A, X) = 0$。$\mathcal{A}$ 的{\it 左正交}
${}^\perp\mathcal{A}$ 是如下对象组成的全子范畴：$\mathcal{D}$ 的对象
$X$ 满足，对所有 $A \in \Ob(\mathcal{A})$，$\Hom(X, A) = 0$。
\end{definition}

\begin{lemma}
\label{derived-lemma-pre-prepare-adjoint}
设 $\mathcal{D}$ 为三角范畴，$\mathcal{A} \subset \mathcal{D}$
为在所有平移下不变的全子范畴。考虑 $\mathcal{D}$ 的特异三角
$$
X \to Y \to Z \to X[1]
$$
下列条件等价：
\begin{enumerate}
\item $Z$ 属于 $\mathcal{A}^\perp$；
\item 对所有 $A \in \Ob(\mathcal{A})$，有
$\Hom(A, X) = \Hom(A, Y)$。
\end{enumerate}
\end{lemma}

\begin{proof}
由引理 \ref{derived-lemma-representable-homological}，函子
$\Hom(A, -)$ 是同调函子，因而得到如
(\ref{derived-equation-long-exact-cohomology-sequence}) 中的长正合列。
假设 (1)，并令 $A \in \Ob(\mathcal{A})$。考虑正合列
$$
\Hom(A[1], Z) \to \Hom(A, X) \to \Hom(A, Y) \to \Hom(A, Z)
$$
由于 $A[1] \in \Ob(\mathcal{A})$，第一个和最后一个群均为零，
故得到 (2)。假设 (2)，并令 $A \in \Ob(\mathcal{A})$。
考虑正合列
$$
\Hom(A, X) \to \Hom(A, Y) \to \Hom(A, Z) \to \Hom(A[-1], X) \to \Hom(A[-1], Y)
$$
遂如所需得到 $\Hom(A, Z) = 0$。
\end{proof}

\begin{lemma}
\label{derived-lemma-pre-prepare-adjoint-dual}
设 $\mathcal{D}$ 为三角范畴，$\mathcal{B} \subset \mathcal{D}$
为在所有平移下不变的全子范畴。考虑 $\mathcal{D}$ 的特异三角
$$
X \to Y \to Z \to X[1]
$$
下列条件等价：
\begin{enumerate}
\item $X$ 属于 ${}^\perp\mathcal{B}$；
\item 对所有 $B \in \Ob(\mathcal{B})$，有
$\Hom(Y, B) = \Hom(Z, B)$。
\end{enumerate}
\end{lemma}

\begin{proof}
这是引理 \ref{derived-lemma-pre-prepare-adjoint} 的对偶。
\end{proof}

\begin{lemma}
\label{derived-lemma-orthogonal-triangulated}
设 $\mathcal{D}$ 为三角范畴，$\mathcal{A} \subset \mathcal{D}$
为在所有平移下不变的全子范畴。则 $\mathcal{A}$ 的右正交
$\mathcal{A}^\perp$ 和左正交 ${}^\perp\mathcal{A}$ 都是
$\mathcal{D}$ 的严格全、饱和\footnote{定义 \ref{derived-definition-saturated}。}
三角子范畴。
% P02-E0095: frozen source misspells “subcategories” as “subcagories”.
\end{lemma}

\begin{proof}
由定义立即可知，正交子范畴在平移、直和和直和因子下保持不变。
考虑 $\mathcal{D}$ 的特异三角
$$
X \to Y \to Z \to X[1]
$$
由引理 \ref{derived-lemma-triangulated-subcategory}，只需证明若 $X$ 和 $Y$
属于 $\mathcal{A}^\perp$，则 $Z$ 也属于 $\mathcal{A}^\perp$。
这立即由引理 \ref{derived-lemma-pre-prepare-adjoint} 得出。
\end{proof}

\begin{lemma}
\label{derived-lemma-prepare-adjoint}
设 $\mathcal{D}$ 为三角范畴，$\mathcal{A}$ 为 $\mathcal{D}$ 的全三角子范畴。
对 $\mathcal{D}$ 的对象 $X$，考虑性质 $P(X)$：在 $\mathcal{D}$ 中存在
特异三角 $A \to X \to B \to A[1]$，其中 $A$ 属于 $\mathcal{A}$，
$B$ 属于 $\mathcal{A}^\perp$。
\begin{enumerate}
\item 若 $X_1 \to X_2 \to X_3 \to X_1[1]$ 是特异三角，且 $P$
对三个对象中的两个成立，则它对第三个也成立。
\item 若 $P$ 对 $X_1$ 和 $X_2$ 成立，则它对 $X_1 \oplus X_2$ 也成立。
\end{enumerate}
\end{lemma}

\begin{proof}
设 $X_1 \to X_2 \to X_3 \to X_1[1]$ 为特异三角，并假设 $P$
对 $X_1$ 和 $X_2$ 成立。选择性质 $P$ 中的特异三角
$$
A_1 \to X_1 \to B_1 \to A_1[1]
\quad\text{且}\quad
A_2 \to X_2 \to B_2 \to A_2[1]
$$
由于由引理 \ref{derived-lemma-pre-prepare-adjoint} 有
$\Hom(A_1, A_2) = \Hom(A_1, X_2)$，存在唯一态射
$A_1 \to A_2$，使图
$$
\xymatrix{
A_1 \ar[d] \ar[r] & X_1 \ar[d] \\
A_2 \ar[r] & X_2
}
$$
交换。选择它到下图的一个扩张
$$
\xymatrix{
A_1 \ar[r] \ar[d] & X_1 \ar[r] \ar[d] & Q_1 \ar[r] \ar[d] & A_1[1] \ar[d] \\
A_2 \ar[r] \ar[d] & X_2 \ar[r] \ar[d] & Q_2 \ar[r] \ar[d] & A_2[1] \ar[d] \\
A_3 \ar[r] \ar[d] & X_3 \ar[r] \ar[d] & Q_3 \ar[r] \ar[d] & A_3[1] \ar[d] \\
A_1[1] \ar[r] & X_1[1] \ar[r] & Q_1[1] \ar[r] & A_1[2]
}
$$
如命题 \ref{derived-proposition-9} 中所述。由 TR3，$Q_1 \cong B_1$ 且
$Q_2 \cong B_2$，故 $Q_1, Q_2 \in \Ob(\mathcal{A}^\perp)$。
由于 $Q_1 \to Q_2 \to Q_3 \to Q_1[1]$ 是特异三角，
由引理 \ref{derived-lemma-orthogonal-triangulated}，有
$Q_3 \in \Ob(\mathcal{A}^\perp)$。由于 $\mathcal{A}$ 是全三角子范畴，
$A_3$ 与 $\mathcal{A}$ 的某个对象同构。因此 $X_3$ 满足 $P$。
(1) 的其他情形由平移归结为此情形。(2) 经引理 \ref{derived-lemma-split}
是 (1) 的特例。
\end{proof}

\begin{lemma}
\label{derived-lemma-prepare-adjoint-dual}
设 $\mathcal{D}$ 为三角范畴，$\mathcal{B}$ 为 $\mathcal{D}$ 的全三角子范畴。
对 $\mathcal{D}$ 的对象 $X$，考虑性质 $P(X)$：在 $\mathcal{D}$ 中存在
特异三角 $A \to X \to B \to A[1]$，其中 $B$ 属于 $\mathcal{B}$，
$A$ 属于 ${}^\perp\mathcal{B}$。
\begin{enumerate}
\item 若 $X_1 \to X_2 \to X_3 \to X_1[1]$ 是特异三角，且 $P$
对三个对象中的两个成立，则它对第三个也成立。
\item 若 $P$ 对 $X_1$ 和 $X_2$ 成立，则它对 $X_1 \oplus X_2$ 也成立。
\end{enumerate}
\end{lemma}

\begin{proof}
这是引理 \ref{derived-lemma-prepare-adjoint} 的对偶。
\end{proof}

\begin{lemma}
\label{derived-lemma-right-adjoint}
设 $\mathcal{D}$ 为三角范畴，$\mathcal{A} \subset \mathcal{D}$
为全三角子范畴。下列条件等价：
\begin{enumerate}
\item 嵌入函子 $\mathcal{A} \to \mathcal{D}$ 有右伴随；
\item 对 $\mathcal{D}$ 中每个 $X$，存在 $\mathcal{D}$ 中的特异三角
$$
A \to X \to B \to A[1]
$$
其中 $A \in \Ob(\mathcal{A})$ 且 $B \in \Ob(\mathcal{A}^\perp)$。
\end{enumerate}
若这些条件成立，则 $\mathcal{A}$ 是饱和的
（定义 \ref{derived-definition-saturated}）；若 $\mathcal{A}$ 是
$\mathcal{D}$ 的严格全子范畴，则
$\mathcal{A} = {}^\perp(\mathcal{A}^\perp)$。
\end{lemma}

\begin{proof}
假设 (1)，并以 $v : \mathcal{D} \to \mathcal{A}$ 表示右伴随。
设 $X \in \Ob(\mathcal{D})$，令 $A = v(X)$。可将伴随映射
$A \to X$ 扩张为特异三角 $A \to X \to B \to A[1]$。
由于对 $A' \in \Ob(\mathcal{A})$，有
$$
\Hom_\mathcal{A}(A', A) =
\Hom_\mathcal{A}(A', v(X)) =
\Hom_\mathcal{D}(A', X)
$$
由引理 \ref{derived-lemma-pre-prepare-adjoint}，可知
$B \in \Ob(\mathcal{A}^\perp)$。

\medskip\noindent
假设 (2)。我们显式构造伴随 $v$。设 $X \in \Ob(\mathcal{D})$，
选择 (2) 中的 $A \to X \to B \to A[1]$，并令 $v(X) = A$。
设 $f : X \to Y$ 为 $\mathcal{D}$ 中的态射。选择 (2) 中的
$A' \to Y \to B' \to A'[1]$。由于由引理
\ref{derived-lemma-pre-prepare-adjoint} 有 $\Hom(A, A') = \Hom(A, Y)$，
存在唯一态射 $f' : A \to A'$，使图
$$
\xymatrix{
A \ar[d]_{f'} \ar[r] & X \ar[d]^f \\
A' \ar[r] & Y
}
$$
交换。因此可令 $v(f) = f'$，从而得到函子。
再次使用引理 \ref{derived-lemma-pre-prepare-adjoint}，可见 $v$ 是嵌入的伴随。

\medskip\noindent
证明最后一个陈述。为证明 $\mathcal{A}$ 饱和，可以把 $\mathcal{A}$
替换为与 $\mathcal{A}$ 具有相同同构类的严格全子范畴；略去细节。假设 $\mathcal{A}$
严格全。若能证明 $\mathcal{A} = {}^\perp(\mathcal{A}^\perp)$，
则由引理 \ref{derived-lemma-orthogonal-triangulated}，$\mathcal{A}$ 将是饱和的。
包含关系 $\mathcal{A} \subset {}^\perp(\mathcal{A}^\perp)$ 是显然的，
% P02-E0096: frozen source misspells “inclusion” as “incusion”.
故只需证明另一个包含关系。设 $X$ 为
${}^\perp(\mathcal{A}^\perp)$ 的对象。选择 (2) 中的特异三角
$A \to X \to B \to A[1]$。由于由假设有 $\Hom(X, B) = 0$，
由引理 \ref{derived-lemma-split}，可见 $A \cong X \oplus B[-1]$。
由于 $B \in \mathcal{A}^\perp$，有 $\Hom(A, B[-1]) = 0$，
这蕴含 $B[-1] = 0$ 且 $A \cong X$，正如所需。
\end{proof}

\begin{lemma}
\label{derived-lemma-left-adjoint}
设 $\mathcal{D}$ 为三角范畴，$\mathcal{B} \subset \mathcal{D}$
为全三角子范畴。下列条件等价：
\begin{enumerate}
\item 嵌入函子 $\mathcal{B} \to \mathcal{D}$ 有左伴随；
\item 对 $\mathcal{D}$ 中每个 $X$，存在 $\mathcal{D}$ 中的特异三角
$$
A \to X \to B \to A[1]
$$
其中 $B \in \Ob(\mathcal{B})$ 且 $A \in \Ob({}^\perp\mathcal{B})$。
\end{enumerate}
若这些条件成立，则 $\mathcal{B}$ 是饱和的
（定义 \ref{derived-definition-saturated}）；若 $\mathcal{B}$ 是
$\mathcal{D}$ 的严格全子范畴，则
$\mathcal{B} = ({}^\perp\mathcal{B})^\perp$。
\end{lemma}

\begin{proof}
这是引理 \ref{derived-lemma-right-adjoint} 的对偶。
\end{proof}

\begin{definition}
\label{derived-definition-admissible}
设 $\mathcal{D}$ 为三角范畴。$\mathcal{D}$ 的{\it 右容许}子范畴，
是满足引理 \ref{derived-lemma-right-adjoint} 中等价条件的严格全三角子范畴。
$\mathcal{D}$ 的{\it 左容许}子范畴，是满足引理
\ref{derived-lemma-left-adjoint} 中等价条件的严格全三角子范畴。
同时右容许和左容许的子范畴称为{\it 双侧容许}子范畴。
\end{definition}

\noindent
设 $\mathcal{A}$ 为三角范畴 $\mathcal{D}$ 的右容许子范畴。
注意到，对 $X \in \mathcal{D}$，特异三角
$$
A \to X \to B \to A[1]
$$
其中 $A \in \mathcal{A}$ 且 $B \in \mathcal{A}^\perp$，在如下意义下是
典范的：对任意另一个满足 $A' \in \mathcal{A}$、
$B' \in \mathcal{A}^\perp$ 的特异三角 $A' \to X \to B' \to A'[1]$，
都存在三角的同构
$(\alpha, \text{id}_X, \beta) : (A, X, B) \to (A', X, B')$。
下面的命题总结以上所述。

\begin{proposition}
\label{derived-proposition-summarize-admissible}
设 $\mathcal{D}$ 为三角范畴，$\mathcal{A} \subset \mathcal{D}$ 和
$\mathcal{B} \subset \mathcal{D}$ 为子范畴。下列条件等价：
\begin{enumerate}
\item $\mathcal{A}$ 右容许，且 $\mathcal{B} = \mathcal{A}^\perp$；
\item $\mathcal{B}$ 左容许，且 $\mathcal{A} = {}^\perp\mathcal{B}$；
\item 对所有 $A \in \mathcal{A}$ 和 $B \in \mathcal{B}$，有
$\Hom(A, B) = 0$；并且对 $\mathcal{D}$ 中每个 $X$，存在
$\mathcal{D}$ 中的特异三角 $A \to X \to B \to A[1]$，其中
$A \in \mathcal{A}$ 且 $B \in \mathcal{B}$。
\end{enumerate}
若这些条件成立，则 $\mathcal{A} \to \mathcal{D}/\mathcal{B}$ 和
$\mathcal{B} \to \mathcal{D}/\mathcal{A}$ 是三角范畴的等价；
嵌入函子 $\mathcal{A} \to \mathcal{D}$ 的右伴随是
$\mathcal{D} \to \mathcal{D}/\mathcal{B} \to \mathcal{A}$；
嵌入函子 $\mathcal{B} \to \mathcal{D}$ 的左伴随是
$\mathcal{D} \to \mathcal{D}/\mathcal{A} \to \mathcal{B}$。
\end{proposition}

\begin{proof}
(1)、(2)、(3) 的等价性直接由引理 \ref{derived-lemma-right-adjoint} 和
\ref{derived-lemma-left-adjoint} 得出（略去一个小细节）。以
$v : \mathcal{D} \to \mathcal{A}$ 表示嵌入函子
$i : \mathcal{A} \to \mathcal{D}$ 的右伴随。立即有
$\Ker(v) = \mathcal{A}^\perp = \mathcal{B}$。因此，由商的泛性质，
$v$ 经函子 $\overline{v} : \mathcal{D}/\mathcal{B} \to \mathcal{A}$ 分解。
由于由 Categories，引理 \ref{categories-lemma-adjoint-fully-faithful} 有
$v \circ i = \text{id}_\mathcal{A}$，可见 $\overline{v}$ 是
$\overline{i} : \mathcal{A} \to \mathcal{D}/\mathcal{B}$ 的左拟逆。
我们还断言复合 $\overline{i} \circ \overline{v}$ 同构于
$\text{id}_{\mathcal{D}/\mathcal{B}}$。事实上，假设 $X$ 嵌入 (3) 中的
特异三角 $A \to X \to B \to A[1]$。如引理
\ref{derived-lemma-right-adjoint} 的证明中所见，$v(X) = A$。
把 $X$ 看作 $\mathcal{D}/\mathcal{B}$ 的对象，则有
$\overline{i}(\overline{v}(X)) = A$，并且在 $\mathcal{D}/\mathcal{B}$ 中
存在函子性同构 $\overline{i}(\overline{v}(X)) = A \to X$。
因此，确实有 $\overline{v} : \mathcal{D}/\mathcal{B} \to \mathcal{A}$
为等价。证明 $\mathcal{B} \to \mathcal{D}/\mathcal{A}$ 是等价，
且嵌入函子 $\mathcal{B} \to \mathcal{D}$ 的左伴随是
$\mathcal{D} \to \mathcal{D}/\mathcal{A} \to \mathcal{B}$，
与上述论证对偶。
\end{proof}




\section{Postnikov 系统}
\label{derived-section-postnikov}

\noindent
本节的参考文献是 \cite{Orlov-K3}。设 $\mathcal{D}$ 为三角范畴，设
$$
X_n \to X_{n - 1} \to \ldots \to X_0
$$
为 $\mathcal{D}$ 中的复形。本节考虑构造该复形的“全复形化”的问题。

\begin{definition}
\label{derived-definition-postnikov-system}
设 $\mathcal{D}$ 为三角范畴，设
$$
X_n \to X_{n - 1} \to \ldots \to X_0
$$
为 $\mathcal{D}$ 中的复形。{\it Postnikov 系统}归纳定义如下。
\begin{enumerate}
\item 若 $n = 0$，它是同构 $Y_0 \to X_0$。
\item 若 $n = 1$，它由同构 $Y_0 \to X_0$ 的一个选择以及特异三角
$$
Y_1 \to X_1 \to Y_0 \to Y_1[1]
$$
的一个选择组成，其中 $X_1 \to Y_0$ 与 $Y_0 \to X_0$ 的复合是
给定态射 $X_1 \to X_0$。
\item 若 $n > 1$，它由复形 $X_{n - 1} \to \ldots \to X_0$
的一个 Postnikov 系统以及特异三角
$$
Y_n \to X_n \to Y_{n - 1} \to Y_n[1]
$$
的一个选择组成，其中态射 $X_n \to Y_{n - 1}$ 与
$Y_{n - 1} \to X_{n - 1}$ 的复合是给定态射 $X_n \to X_{n - 1}$。
\end{enumerate}
给定 $\mathcal{D}$ 中两个等长复形之间的态射
\begin{equation}
\label{derived-equation-map-complexes}
\vcenter{
\xymatrix{
X_n \ar[r] \ar[d] &
X_{n - 1} \ar[r] \ar[d] &
\ldots \ar[r] &
X_0 \ar[d] \\
X'_n \ar[r] &
X'_{n - 1} \ar[r] &
\ldots \ar[r] &
X'_0
}
}
\end{equation}
显然可定义{\it Postnikov 系统的态射}。
\end{definition}

\noindent
下面是一个关键例子。

\begin{example}
\label{derived-example-key-postnikov}
设 $\mathcal{A}$ 为阿贝尔范畴，设 $\ldots \to A_2 \to A_1 \to A_0$
为 $\mathcal{A}$ 中的链复形。于是可以考虑 $D(\mathcal{A})$ 中的对象
$$
X_n = A_n
\quad\text{且}\quad
Y_n = (A_n \to A_{n - 1} \to \ldots \to A_0)[-n]
$$
连同显然的典范映射 $Y_n \to X_n$ 和
$Y_0 \to Y_1[1] \to Y_2[2] \to \ldots$，特异三角
$Y_n \to X_n \to Y_{n - 1} \to Y_n[1]$ 按定义
\ref{derived-definition-postnikov-system} 给出
$\ldots \to X_2 \to X_1 \to X_0$ 的 Postnikov 系统。
这里使用了 Postnikov 系统对 $D(\mathcal{A})$ 中无限复形的显然推广。
最后，若 $\mathcal{A}$ 中以 $\mathbf{N}$ 为指标的余极限存在且正合，则
$$
\text{hocolim} Y_n[n] = (\ldots \to A_2 \to A_1 \to A_0 \to 0 \to \ldots)
$$
在 $D(\mathcal{A})$ 中成立。这立即由引理 \ref{derived-lemma-colim-hocolim} 得出。
\end{example}

\noindent
给定复形 $X_n \to X_{n - 1} \to \ldots \to X_0$ 以及定义
\ref{derived-definition-postnikov-system} 中的 Postnikov 系统，可以考虑映射
$$
Y_0 \to Y_1[1] \to \ldots \to Y_n[n]
$$
这些映射共同嵌入某些特异三角，并与各 $X_i$ 之间的给定映射相容。
当 $n = 3$ 时，图形如下：
$$
\xymatrix{
Y_0 \ar[rr] & &
Y_1[1] \ar[dl] \ar[rr] & &
Y_2[2] \ar[dl] \ar[rr] & &
Y_3[3] \ar[dl] \\
& X_1[1] \ar[lu]_{+1} & &
X_2[2] \ar[ll]_{+1} \ar[lu]_{+1} & &
X_3[3] \ar[ll]_{+1} \ar[lu]_{+1}
}
$$
我们建议读者把 $Y_n[n]$ 看作由 $X_0, X_1[1], \ldots, X_n[n]$ 得到；
例如，若各映射 $X_i \to X_{i - 1}$ 为零，则可以取
$Y_n[n] = \bigoplus_{i = 0, \ldots, n} X_i[i]$。
Postnikov 系统并非总是存在。下面是关于较小 $n$ 的简单引理。

\begin{lemma}
\label{derived-lemma-postnikov-system-small-cases}
设 $\mathcal{D}$ 为三角范畴。考虑长度为 $n$ 的复形的 Postnikov 系统。
\begin{enumerate}
\item 当 $n = 0$ 时，Postnikov 系统总是存在，并且复形的任意态射
(\ref{derived-equation-map-complexes}) 唯一地扩张为 Postnikov 系统的态射。
\item 当 $n = 1$ 时，Postnikov 系统总是存在，并且复形的任意态射
(\ref{derived-equation-map-complexes}) 扩张为 Postnikov 系统的一个（非唯一）态射。
\item 当 $n = 2$ 时，Postnikov 系统总是存在，但一般而言，复形的态射
(\ref{derived-equation-map-complexes}) 不会扩张为 Postnikov 系统的态射。
\item 当 $n > 2$ 时，Postnikov 系统并非总是存在。
\end{enumerate}
\end{lemma}

\begin{proof}
$n = 0$ 的情形立即成立，因为同构可逆。$n = 1$ 的情形立即由
TR1（三角的存在性）和 TR3（把态射扩张为三角的态射）得出。
对 $n = 2$，论证如下。令 $Y_0 = X_0$。由 $n = 1$ 的情形，
可选择 Postnikov 系统
$$
Y_1 \to X_1 \to Y_0 \to Y_1[1]
$$
由于复合 $X_2 \to X_1 \to X_0$ 为零，由引理
\ref{derived-lemma-representable-homological}，可将 $X_2 \to X_1$
（非唯一地）分解为 $X_2 \to Y_1 \to X_1$。然后只需把态射
$X_2 \to Y_1$ 嵌入特异三角
$$
Y_2 \to X_2 \to Y_1 \to Y_2[1]
$$
便得到 $n = 2$ 的 Postnikov 系统。对 $n > 2$，不能类似论证，
因为不知道复合 $X_n \to X_{n - 1} \to Y_{n - 1}$ 在
$\mathcal{D}$ 中是否为零。
% P02-E0097: frozen source's last composition is ill-typed for the preceding Postnikov triangle.
\end{proof}

\begin{lemma}
\label{derived-lemma-maps-postnikov-systems-vanishing}
设 $\mathcal{D}$ 为三角范畴。给定态射
(\ref{derived-equation-map-complexes})，考虑条件
\begin{equation}
\label{derived-equation-P}
\Hom(X_i[i - j - 1], X'_j) = 0 \text{ 对 }i > j + 1
\end{equation}
则：
\begin{enumerate}
\item 若已有复形 $X'_n \to X'_{n - 1} \to \ldots \to X'_0$
的 Postnikov 系统，则性质 (\ref{derived-equation-P}) 蕴含
$$
\Hom(X_i[i - j - 1], Y'_j) = 0 \text{ 对 }i > j + 1
$$
\item 若两个复形都配备了 Postnikov 系统并且
(\ref{derived-equation-P}) 成立，则该态射扩张为 Postnikov 系统的一个（非唯一）态射。
\end{enumerate}
\end{lemma}

\begin{proof}
先对 $j$ 归纳证明 (1)。基础情形 $j = 0$ 无需证明，因为
$Y'_0 \to X'_0$ 是同构。假设结论对 $j - 1$ 成立。考虑特异三角
$$
Y'_j \to X'_j \to Y'_{j - 1} \to Y'_j[1]
$$
引理 \ref{derived-lemma-representable-homological} 的长正合列给出正合列
$$
\Hom(X_i[i - j - 1], Y'_{j - 1}[-1]) \to
\Hom(X_i[i - j - 1], Y'_j) \to
\Hom(X_i[i - j - 1], X'_j)
$$
由归纳假设和 (\ref{derived-equation-P})，两端的群均为零，故结论成立。

\medskip\noindent
证明 (2)。当 $n = 1$ 时，态射的存在性已由引理
\ref{derived-lemma-postnikov-system-small-cases} 证明。若 $n > 1$，
由归纳可假设已经给定长度 $n - 1$ 的 Postnikov 系统态射。
问题在于不知道图
$$
\xymatrix{
X_n \ar[r] \ar[d] & Y_{n - 1} \ar[d] \\
X'_n \ar[r] & Y'_{n - 1}
}
$$
是否交换。以 $\alpha : X_n \to Y'_{n - 1}$ 表示其差。
我们确实知道 $\alpha$ 与 $Y'_{n - 1} \to X'_{n - 1}$ 的复合为零
（这是长度 $n - 1$ 的 Postnikov 系统态射的含义）。由特异三角
$Y'_{n - 1} \to X'_{n - 1} \to Y'_{n - 2} \to Y'_{n - 1}[1]$，
这意味着 $\alpha$ 是 $Y'_{n - 2}[-1] \to Y'_{n - 1}$ 与某个映射
$\alpha' : X_n \to Y'_{n - 2}[-1]$ 的复合。由本引理 (1)，条件
(\ref{derived-equation-P}) 保证 $\alpha'$ 为零，故 $\alpha$ 为零。
为完成存在性的证明，交换性保证可由 TR3 选择下图中的虚线箭头：
$$
\xymatrix{
Y_{n - 1}[-1] \ar[d] \ar[r] &
Y_n \ar[r] \ar@{..>}[d] &
X_n \ar[r] \ar[d] &
Y_{n - 1} \ar[d] \\
Y'_{n - 1}[-1] \ar[r] &
Y'_n \ar[r] &
X'_n \ar[r] &
Y'_{n - 1}
}
$$
\end{proof}

\begin{lemma}
\label{derived-lemma-uniqueness-maps-postnikov-systems}
设 $\mathcal{D}$ 为三角范畴。给定态射
(\ref{derived-equation-map-complexes})，并假设两个复形都配备了 Postnikov 系统。
若下列任一条件成立：
\begin{enumerate}
\item 对 $i = 1, \ldots, n$，$\Hom(X_i[i], Y'_n[n]) = 0$；或
\item 对 $i = 1, \ldots, n$，$\Hom(Y_n[n], X'_{n - i}[n - i]) = 0$；或
\item 对 $j \geq i > 0$，有 $\Hom(X_{j - i}[-i + 1], X'_j) = 0$ 且
$\Hom(X_j, X'_{j - i}[-i]) = 0$，
\end{enumerate}
则这些 Postnikov 系统之间至多存在一个态射。
\end{lemma}

\begin{proof}
证明 (1)。考察下图：
$$
\xymatrix{
Y_0 \ar[r] \ar[d] &
Y_1[1] \ar[r] \ar[ld] &
Y_2[2] \ar[r] \ar[lld] &
\ldots \ar[r] &
Y_n[n] \ar[lllld] \\
Y'_n[n]
}
$$
各箭头是态射 $Y_n[n] \to Y'_n[n]$ 与态射
$Y_i[i] \to Y_n[n]$ 的复合。箭头 $Y_0 \to Y'_n[n]$ 是确定的，
因为它是复合 $Y_0 = X_0 \to X'_0 = Y'_0 \to Y'_n[n]$。
由于有特异三角 $Y_0 \to Y_1[1] \to X_1[1]$，
$\Hom(X_1[1], Y'_n[n]) = 0$ 保证第二条竖直箭头唯一。
由于有特异三角 $Y_1[1] \to Y_2[2] \to X_2[2]$，
$\Hom(X_2[2], Y'_n[n]) = 0$ 保证第三条竖直箭头唯一。依此类推。

\medskip\noindent
证明 (2)。复合 $Y_n[n] \to Y'_n[n] \to X_n[n]$ 与复合
% P02-E0098: frozen source ends this composition at X_n[n], although the primed system requires X'_n[n].
$Y_n[n] \to X_n[n] \to X'_n[n]$ 相同，因而唯一。
再用特异三角 $Y'_{n - 1}[n - 1] \to Y'_n[n] \to X'_n[n]$，
可见只需证明 $\Hom(Y_n[n], Y'_{n - 1}[n - 1]) = 0$。
利用特异三角
$$
Y'_{n - i - 1}[n - i - 1] \to Y'_{n - i}[n - i] \to X'_{n - i}[n - i]
$$
由假设得到该消失。略去小细节。

\medskip\noindent
证明 (3)。考察引理 \ref{derived-lemma-maps-postnikov-systems-vanishing} 的证明并
对 $n$ 归纳，只需证明三角态射
$$
\xymatrix{
Y_{n - 1}[-1] \ar[d] \ar[r] &
Y_n \ar[r] \ar@{..>}[d] &
X_n \ar[r] \ar[d] &
Y_{n - 1} \ar[d] \\
Y'_{n - 1}[-1] \ar[r] &
Y'_n \ar[r] &
X'_n \ar[r] &
Y'_{n - 1}
}
$$
中的虚线箭头唯一。由引理 \ref{derived-lemma-uniqueness-third-arrow} (5)，
只需证明 $\Hom(Y_{n - 1}, X'_n) = 0$ 和
$\Hom(X_n, Y'_{n - 1}[-1]) = 0$。
为证明第一个消失，对 $i > 0$ 使用特异三角
$Y_{n - i - 1}[-i] \to Y_{n - i}[-(i - 1)] \to X_{n - i}[-(i - 1)]$，
并对 $i$ 归纳，可见所假设的
$\Hom(X_{n - i}[-i + 1], X'_n)$ 的消失已经足够。
对第二个消失，类似地使用特异三角
$Y'_{n - i - 1}[-i - 1] \to Y'_{n - i}[-i] \to X'_{n - i}[-i]$，
可见所假设的 $\Hom(X_n, X'_{n - i}[-i])$ 的消失也足够。
\end{proof}

\begin{lemma}
\label{derived-lemma-existence-postnikov-system}
设 $\mathcal{D}$ 为三角范畴。设
$X_n \to X_{n - 1} \to \ldots \to X_0$ 为 $\mathcal{D}$ 中的复形。
若
$$
\Hom(X_i[i - j - 2], X_j) = 0 \text{ 对 }i > j + 2
$$
则存在 Postnikov 系统。若
$$
\Hom(X_i[i - j - 1], X_j) = 0 \text{ 对 }i > j + 1
$$
则任意两个 Postnikov 系统同构。
\end{lemma}

\begin{proof}
对 $n$ 归纳。$n = 0, 1, 2$ 的情形由引理
\ref{derived-lemma-postnikov-system-small-cases} 得出。假设 $n > 2$。
假设已给定复形 $X_{n - 1} \to X_{n - 2} \to \ldots \to X_0$
的 Postnikov 系统。把它扩张为长度 $n$ 的 Postnikov 系统的唯一障碍，
是必须找到态射 $X_n \to Y_{n - 1}$，使复合
$X_n \to Y_{n - 1} \to X_{n - 1}$ 等于给定映射
$X_n \to X_{n - 1}$。考虑特异三角
$$
Y_{n - 1} \to X_{n - 1} \to Y_{n - 2} \to Y_{n - 1}[1]
$$
以及由它和函子 $\Hom(X_n, -)$ 给出的长正合列
（见引理 \ref{derived-lemma-representable-homological}），可知只需证明复合
$X_n \to X_{n - 1} \to Y_{n - 2}$ 为零。
由于已知 $X_n \to X_{n - 1} \to X_{n - 2}$ 为零，
应用特异三角
$$
Y_{n - 2} \to X_{n - 2} \to Y_{n - 3} \to Y_{n - 2}[1]
$$
可见只要 $\Hom(X_n, Y_{n - 3}[-1]) = 0$ 即可。
完全仿照引理 \ref{derived-lemma-maps-postnikov-systems-vanishing} (1) 的证明，
读者容易看出这由本引理中的条件推出。

\medskip\noindent
关于同构的陈述，由引理 \ref{derived-lemma-maps-postnikov-systems-vanishing} (2)
所证明的、扩张复形恒等映射的 Postnikov 系统之间态射的存在性，以及引理
\ref{derived-lemma-third-isomorphism-triangle}（它表明所有映射均为同构）得出。
\end{proof}









\section{本质常值系统}
\label{derived-section-essentially-constant}

\noindent
先给出三角范畴中本质常值系统的一些预备引理。

\begin{lemma}
\label{derived-lemma-essentially-constant}
设 $\mathcal{D}$ 为三角范畴，$(A_i)$ 为 $\mathcal{D}$ 中的逆系统。
则 $(A_i)$ 是本质常值的（见 Categories，定义
\ref{categories-definition-essentially-constant-diagram}），当且仅当存在
$i$，并且对所有 $j \geq i$，存在直和分解 $A_j = A \oplus Z_j$，使得
(a) 各映射 $A_{j'} \to A_j$ 与直和分解相容，并且在 $A$ 上为恒等映射；
(b) 对所有 $j \geq i$，存在某个 $j' \geq j$，使 $Z_{j'} \to Z_j$ 为零。
\end{lemma}

\begin{proof}
假设 $(A_i)$ 是值为 $A$ 的本质常值系统。则 $A = \lim A_i$，
并且存在 $i$ 和态射 $A_i \to A$，使得 (1) 复合
$A \to A_i \to A$ 在 $A$ 上为恒等映射；(2) 对所有 $j \geq i$，
存在 $j' \geq j$，使 $A_{j'} \to A_j$ 分解为
$A_{j'} \to A_i \to A \to A_j$。
由 (1)，可知对 $j \geq i$，映射 $A \to A_j$ 与
$A_j \to A_i \to A$ 的复合在 $A$ 上为恒等映射。因此
$A_j \to A$ 有核 $Z_j$，并且映射 $A \oplus Z_j \to A_j$ 是同构，
见引理 \ref{derived-lemma-when-split} 和 \ref{derived-lemma-split}。
这些直和分解显然满足 (a)。由 (2)，对所有 $j$，存在
$j' \geq j$，使 $Z_{j'} \to Z_j$ 为零，故 (b) 成立。
略去反向的证明。
\end{proof}

\begin{lemma}
\label{derived-lemma-essentially-constant-2-out-of-3}
设 $\mathcal{D}$ 为三角范畴。设
$$
A_n \to B_n \to C_n \to A_n[1]
$$
为 $\mathcal{D}$ 中特异三角的逆系统。若 $(A_n)$ 和 $(C_n)$
本质常值，则 $(B_n)$ 也本质常值，并且它们的值嵌入特异三角
$A \to B \to C \to A[1]$；此外，对某个 $n \geq 1$，存在特异三角的态射
$$
\xymatrix{
A_n \ar[d] \ar[r] &
B_n \ar[d] \ar[r] &
C_n \ar[d] \ar[r] &
A_n[1] \ar[d] \\
A \ar[r] &
B \ar[r] &
C \ar[r] &
A[1]
}
$$
它诱导同构 $\lim_{n' \geq n} A_{n'} \to A$，对 $B$ 和 $C$ 亦然。
\end{lemma}

\begin{proof}
重新编号后，可假设 $A_n = A \oplus A'_n$ 且
$C_n = C \oplus C'_n$，其中逆系统 $(A'_n)$ 和 $(C'_n)$
本质为零，见引理 \ref{derived-lemma-essentially-constant}。特别地，对所有 $n$，
态射
$$
C \oplus C'_n \to (A \oplus A'_n)[1]
$$
把直和因子 $C$ 映入直和因子 $A[1]$，其上的映射
$\delta : C \to A[1]$ 与 $n$ 无关。选择特异三角
$$
A \to B \to C \xrightarrow{\delta} A[1]
$$
接着选择特异三角的态射
$$
(A_1 \to B_1 \to C_1 \to A_1[1]) \to
(A \to B \to C \to A[1])
$$
这由 TR3 可以做到。对 $\mathcal{D}$ 的任意对象 $D$，它诱导交换图
$$
\xymatrix{
\ldots \ar[r] &
\Hom_\mathcal{D}(C, D) \ar[r] \ar[d] &
\Hom_\mathcal{D}(B, D) \ar[r] \ar[d] &
\Hom_\mathcal{D}(A, D) \ar[r] \ar[d] &
\ldots \\
\ldots \ar[r] &
\colim \Hom_\mathcal{D}(C_n, D) \ar[r] &
\colim \Hom_\mathcal{D}(B_n, D) \ar[r] &
\colim \Hom_\mathcal{D}(A_n, D) \ar[r] &
\ldots
}
$$
左、右竖直箭头是同构，它们左边和右边的箭头也都是同构。
因此，由五引理，中间的箭头是同构。由 Categories，评注
\ref{categories-remark-pro-category-copresheaves} 中的讨论，
可知 $(B_n)$ 同构于值为 $B$ 的常值逆系统。
由 Categories，评注 \ref{categories-remark-pro-category}，
这等价于 $(B_n)$ 是值为 $B$ 的本质常值系统，证明完毕。
\end{proof}

\begin{lemma}
\label{derived-lemma-essentially-constant-cohomology}
设 $\mathcal{A}$ 为阿贝尔范畴，$A_n$ 为 $D(\mathcal{A})$ 中对象的逆系统。
假设：
\begin{enumerate}
\item 存在整数 $a \leq b$，使得当 $i \not \in [a, b]$ 时
$H^i(A_n) = 0$；
\item 对所有 $i \in \mathbf{Z}$，$\mathcal{A}$ 中的逆系统
$H^i(A_n)$ 都本质常值。
\end{enumerate}
则 $A_n$ 是 $D(\mathcal{A})$ 中的本质常值系统；其值 $A$ 满足：
对每个 $i \in \mathbf{Z}$，$H^i(A)$ 是常值系统 $H^i(A_n)$ 的值。
\end{lemma}

\begin{proof}
由评注 \ref{derived-remark-truncation-distinguished-triangle}，得到特异三角的逆系统
$$
\tau_{\leq a}A_n \to A_n \to \tau_{\geq a + 1}A_n \to (\tau_{\leq a}A_n)[1]
$$
当然，在 $D(\mathcal{A})$ 中有
$\tau_{\leq a}A_n = H^a(A_n)[-a]$。因此，由假设，它们构成本质常值系统。
对 $b - a$ 归纳，可知逆系统 $\tau_{\geq a + 1}A_n$ 本质常值，
设其值为 $A'$。由引理 \ref{derived-lemma-essentially-constant-2-out-of-3}，
$A_n$ 是本质常值系统。略去关于上同调的陈述的证明
（提示：使用引理 \ref{derived-lemma-essentially-constant-2-out-of-3} 的最后一部分）。
\end{proof}

\begin{lemma}
\label{derived-lemma-pro-isomorphism}
设 $\mathcal{D}$ 为三角范畴。设
$$
A_n \to B_n \to C_n \to A_n[1]
$$
为特异三角的逆系统。若系统 $C_n$ 是 pro-零的（即值为 $0$ 的本质常值系统），
则各映射 $A_n \to B_n$ 确定 pro-对象 $(A_n)$ 与 pro-对象 $(B_n)$
之间的 pro-同构。
\end{lemma}

\begin{proof}
对 $\mathcal{D}$ 的任意对象 $X$，考虑正合列
$$
\colim \Hom(C_n, X) \to
\colim \Hom(B_n, X) \to
\colim \Hom(A_n, X) \to
\colim \Hom(C_n[-1], X) \to
$$
正合性由引理 \ref{derived-lemma-representable-homological} 与 Algebra，引理
\ref{algebra-lemma-directed-colimit-exact} 联合得出。由假设，第一项和最后一项
为零。因此，对所有 $X$，映射
$\colim \Hom(B_n, X) \to \colim \Hom(A_n, X)$ 是同构。
本引理由此及 Categories，评注
\ref{categories-remark-pro-category-copresheaves} 得出。
\end{proof}

\begin{lemma}
\label{derived-lemma-pro-isomorphism-bis}
设 $\mathcal{A}$ 为阿贝尔范畴。
$$
A_n \to B_n
$$
为 $D(\mathcal{A})$ 中映射的逆系统。
% P02-E0100: frozen source omits “Let” before this displayed inverse system.
假设：
\begin{enumerate}
\item 存在整数 $a \leq b$，使得当 $i \not \in [a, b]$ 时，
$H^i(A_n) = 0$ 且 $H^i(B_n) = 0$；
\item 对所有 $i \in \mathbf{Z}$，$\mathcal{A}$ 中映射的逆系统
$H^i(A_n) \to H^i(B_n)$ 定义 $\mathcal{A}$ 中 pro-对象的同构。
% P02-E0101: frozen source has subject--verb disagreement in “the inverse system ... define”.
\end{enumerate}
则各映射 $A_n \to B_n$ 确定 pro-对象 $(A_n)$ 与 pro-对象 $(B_n)$
之间的 pro-同构。
\end{lemma}

\begin{proof}
由公理 TR3，可以归纳地把映射 $A_n \to B_n$ 扩张为特异三角的逆系统
$A_n \to B_n \to C_n \to A_n[1]$。由引理 \ref{derived-lemma-pro-isomorphism}，
只需证明 $C_n$ 是 pro-零的。由引理
\ref{derived-lemma-essentially-constant-cohomology}，只需证明对每个 $p$，
$H^p(C_n)$ 都是 pro-零的。这由假设 (2) 及长正合列
$$
H^p(A_n) \xrightarrow{\alpha_n} H^p(B_n)
\xrightarrow{\beta_n}
H^p(C_n) \xrightarrow{\delta_n} H^{p + 1}(A_n)
\xrightarrow{\epsilon_n}
H^{p + 1}(B_n)
$$
得出。事实上，对每个 $n$，可找到 $m > n$，使得
$\Im(\beta_m)$ 在 $H^p(C_n)$ 中的像为零，因为可选择 $m$，使
$H^p(B_m) \to H^p(B_n)$ 经
$\alpha_n : H^p(A_n) \to H^p(B_n)$ 分解。出于类似理由，
接着可以选择 $k > m$，使 $\Im(\delta_k)$ 在
$H^{p + 1}(A_m)$ 中的像为零。于是 $H^p(C_k) \to H^p(C_n)$ 为零，
因为 $H^p(C_k) \to H^p(C_m)$ 的像包含于 $\Ker(\delta_m)$，而
$H^p(C_m) \to H^p(C_n)$ 消去
$\Ker(\delta_m) = \Im(\beta_m)$。
\end{proof}
